\documentclass[a4paper,fleqn,colorlinks,unicode,lineno]{cas-dc}
\newcommand{\journal}{Applied Ocean Research}
\newcommand{\JOURNALTWOCOLUMN}{}
\newcommand{\whichjournal}{Applied Ocean Research}

\usepackage{amsmath}
\usepackage{amssymb}
\usepackage{soul}       % for highlighting text with \hl 
\usepackage{multirow}   % for multi row/col tables
\usepackage{multicol}
\usepackage{tikz}
\usepackage{array}      % required for text wrapping in tables
\usepackage{pdflscape}  % for landscape table
\usepackage{longtable}  % for table going multiple pages
\usepackage{rotating}   % for rotated table headers
\usepackage{color}      % also required for highlighting 
\usepackage{pdftexcmds} % for detecting if shell escape is enabled
\usepackage{enumitem}   % bullets
\usepackage{booktabs}   % for better looking tables

% save bibliography command before bibcop modifies it
\makeatletter
\let\bibcop@truebibliography\bibliography
\ifnum\pdf@shellescape=0
  % shell escape is disabled, so don't load bibcop
\else
  %\usepackage{bibcop}    % warnings for incorrect citations
\fi
\makeatother

\usepackage[title,titletoc]{appendix} % make it possible to do \ref of an appendix section without automatically getting the word "Appendix"

\usepackage{stmaryrd}  % element wise division symbol used in hessian scaling
\SetSymbolFont{stmry}{bold}{U}{stmry}{m}{n} % avoid warning about bold stmry font

\newcommand{\singleColMacro}[1]{%
    \ifdefined\JOURNALTWOCOLUMN
      \onecolumn%
    \fi
    {#1}%
    \ifdefined\JOURNALTWOCOLUMN
      \twocolumn%
    \fi
}%

% for 11x17 pages
\makeatletter
\@ifclassloaded{cas-sc}{}{%
  \@ifclassloaded{cas-dc}{}{%
    \usepackage[top=1in,
      bottom=1in,
      left=1in,
      right=1in]{geometry}
  }%
}%
\makeatother
\usepackage{caption}
\newenvironment{foldoutpage}{%
  \ifdefined\DISSERTATION
    \eject\pdfpageheight=11in\pdfpagewidth=17in
    \newgeometry{margin=1in}
    \textwidth=15in
    \linewidth=15in
    \textheight=9in
    \captionsetup{width=\textwidth, margin={0in,-8.5in}} 
  \fi
}{%
  \ifdefined\DISSERTATION
    \clearpage % otherwise it will float to another, non-resized page
    \eject\pdfpageheight=11in\pdfpagewidth=8.5in
    \restoregeometry
  \fi
}%

\usepackage{verbatim}  % word counts
\usepackage{textcomp}  % cent symbol

\makeatletter
\@ifclassloaded{cas-sc}{}{%
  \@ifclassloaded{cas-dc}{}{%
    \usepackage[unicode,hypertexnames=true,colorlinks=true,bookmarksnumbered=true]{hyperref}
  }%
}%
\makeatother
\usepackage{cleveref}

% glossary - must come after hyperref
% the shell escape conditional is because arxiv disables shell escape, so automake will fail
\ifnum\pdfshellescape>0
    \usepackage[nopostdot,automake,nogroupskip,nomain,acronym,toc=false,sanitizesort=false,section=subsection]{glossaries-extra}
    % patch to fix quoting of makeindex command from glossaries automake when auxdir is used
  \makeatletter
  \ExplSyntaxOn
  \cs_set:Npn \@gls@run@output@file #1
    {
      \glsquote
        { \l_glossaries_output_dir_tl / #1 }
    }
  \ExplSyntaxOff
  \makeatother
\else
    \usepackage[nopostdot         ,nogroupskip,nomain,acronym,toc=false,sanitizesort=false,section=subsection]{glossaries-extra}
\fi
\usepackage{morewrites}
\ifdefined\DISSERTATION
\else
  % glossary counter to avoid entries only used in the appendix
  \input{glossary/no-appendix-glossary.tex}
\fi

\newcommand\mynewglossary[4]{
  \newglossary{#1}{#2}{#3}{#4}
  \newglossary{alias-#1}{#2.alias}{#3.alias.gls}{#4}
}

\mynewglossary{sym-A}{A.sly}{A.gls}{A}
\mynewglossary{sym-B}{B.sly}{B.gls}{B}
\mynewglossary{sym-C}{C.sly}{C.gls}{C}
\mynewglossary{sym-D}{D.sly}{D.gls}{D}
\mynewglossary{sym-E}{E.sly}{E.gls}{E}
\mynewglossary{sym-F}{F.sly}{F.gls}{F}
\mynewglossary{sym-G}{G.sly}{G.gls}{G}
\mynewglossary{sym-H}{H.sly}{H.gls}{H}
\mynewglossary{sym-I}{I.sly}{I.gls}{I}
\mynewglossary{sym-J}{J.sly}{J.gls}{J}
\mynewglossary{sym-K}{K.sly}{K.gls}{K}
\mynewglossary{sym-L}{L.sly}{L.gls}{L}
\mynewglossary{sym-M}{M.sly}{M.gls}{M}
\mynewglossary{sym-N}{N.sly}{N.gls}{N}
\mynewglossary{sym-O}{O.sly}{O.gls}{O}
\mynewglossary{sym-P}{P.sly}{P.gls}{P}
\mynewglossary{sym-Q}{Q.sly}{Q.gls}{Q}
\mynewglossary{sym-R}{R.sly}{R.gls}{R}
\mynewglossary{sym-S}{S.sly}{S.gls}{S}
\mynewglossary{sym-T}{T.sly}{T.gls}{T}
\mynewglossary{sym-U}{U.sly}{U.gls}{U}
\mynewglossary{sym-V}{V.sly}{V.gls}{V}
\mynewglossary{sym-W}{W.sly}{W.gls}{W}
\mynewglossary{sym-X}{X.sly}{X.gls}{X}
\mynewglossary{sym-Y}{Y.sly}{Y.gls}{Y}
\mynewglossary{sym-Z}{Z.sly}{Z.gls}{Z}

\mynewglossary{sym-a}{a.sly}{a.gls}{a}
\mynewglossary{sym-b}{b.sly}{b.gls}{b}
\mynewglossary{sym-c}{c.sly}{c.gls}{c}
\mynewglossary{sym-d}{d.sly}{d.gls}{d}
\mynewglossary{sym-e}{e.sly}{e.gls}{e}
\mynewglossary{sym-f}{f.sly}{f.gls}{f}
\mynewglossary{sym-g}{g.sly}{g.gls}{g}
\mynewglossary{sym-h}{h.sly}{h.gls}{h}
\mynewglossary{sym-i}{i.sly}{i.gls}{i}
\mynewglossary{sym-j}{j.sly}{j.gls}{j}
\mynewglossary{sym-k}{k.sly}{k.gls}{k}
\mynewglossary{sym-l}{l.sly}{l.gls}{l}
\mynewglossary{sym-m}{m.sly}{m.gls}{m}
\mynewglossary{sym-n}{n.sly}{n.gls}{n}
\mynewglossary{sym-o}{o.sly}{o.gls}{o}
\mynewglossary{sym-p}{p.sly}{p.gls}{p}
\mynewglossary{sym-q}{q.sly}{q.gls}{q}
\mynewglossary{sym-r}{r.sly}{r.gls}{r}
\mynewglossary{sym-s}{s.sly}{s.gls}{s}
\mynewglossary{sym-t}{t.sly}{t.gls}{t}
\mynewglossary{sym-u}{u.sly}{u.gls}{u}
\mynewglossary{sym-v}{v.sly}{v.gls}{v}
\mynewglossary{sym-w}{w.sly}{w.gls}{w}
\mynewglossary{sym-x}{x.sly}{x.gls}{x}
\mynewglossary{sym-y}{y.sly}{y.gls}{y}
\mynewglossary{sym-z}{z.sly}{z.gls}{z}

\mynewglossary{sym-alpha}{alpha.sly}{alpha.gls}{$\boldsymbol\alpha$}
\mynewglossary{sym-beta}{beta.sly}{beta.gls}{$\boldsymbol\beta$}
\mynewglossary{sym-gamma}{gamma.sly}{gamma.gls}{$\boldsymbol\gamma$}
\mynewglossary{sym-delta}{delta.sly}{delta.gls}{$\boldsymbol\delta$}
\mynewglossary{sym-epsilon}{epsilon.sly}{epsilon.gls}{$\boldsymbol\epsilon$}
\mynewglossary{sym-zeta}{zeta.sly}{zeta.gls}{$\boldsymbol\zeta$}
\mynewglossary{sym-eta}{eta.sly}{eta.gls}{$\boldsymbol\eta$}
\mynewglossary{sym-theta}{theta.sly}{theta.gls}{$\boldsymbol\theta$}
\mynewglossary{sym-iota}{iota.sly}{iota.gls}{$\boldsymbol\iota$}
\mynewglossary{sym-kappa}{kappa.sly}{kappa.gls}{$\boldsymbol\kappa$}
\mynewglossary{sym-lambda}{lambda.sly}{lambda.gls}{$\boldsymbol\lambda$}
\mynewglossary{sym-mu}{mu.sly}{mu.gls}{$\boldsymbol\mu$}
\mynewglossary{sym-nu}{nu.sly}{nu.gls}{$\boldsymbol\nu$}
\mynewglossary{sym-xi}{xi.sly}{xi.gls}{$\boldsymbol\xi$}
%\mynewglossary{sym-omicron}{sly}{gls}{$\omicron$}
\mynewglossary{sym-pi}{pi.sly}{pi.gls}{$\boldsymbol\pi$}
\mynewglossary{sym-rho}{rho.sly}{rho.gls}{$\boldsymbol\rho$}
\mynewglossary{sym-sigma}{sigma.sly}{sigma.gls}{$\boldsymbol\sigma$}
\mynewglossary{sym-tau}{tau.sly}{tau.gls}{$\boldsymbol\tau$}
\mynewglossary{sym-upsilon}{upsilon.sly}{upsilon.gls}{$\boldsymbol\upsilon$}
\mynewglossary{sym-phi}{phi.sly}{phi.gls}{$\boldsymbol\phi$}
\mynewglossary{sym-chi}{chi.sly}{chi.gls}{$\boldsymbol\chi$}
\mynewglossary{sym-psi}{psi.sly}{psi.gls}{$\boldsymbol\psi$}
\mynewglossary{sym-omega}{omega.sly}{omega.gls}{$\boldsymbol\omega$}

%\mynewglossary{sym-Alpha}{sly}{gls}{$\Alpha$}
%\mynewglossary{sym-Beta}{sly}{gls}{$\Beta$}
\mynewglossary{sym-Gamma}{Gamma.sly}{Gamma.gls}{$\boldsymbol\Gamma$}
\mynewglossary{sym-Delta}{Delta.sly}{Delta.gls}{$\boldsymbol\Delta$}
%\mynewglossary{sym-Epsilon}{sly}{gls}{$\Epsilon$}
%\mynewglossary{sym-Zeta}{sly}{gls}{$\Zeta$}
%\mynewglossary{sym-Eta}{sly}{gls}{$\Eta$}
\mynewglossary{sym-Theta}{Theta.sly}{Theta.gls}{$\boldsymbol\Theta$}
%\mynewglossary{sym-Iota}{sly}{gls}{$\Iota$}
%\mynewglossary{sym-Kappa}{sly}{gls}{$\Kappa$}
\mynewglossary{sym-Lambda}{Lambda.sly}{Lambda.gls}{$\boldsymbol\Lambda$}
%\mynewglossary{sym-Mu}{sly}{gls}{$\Mu$}
%\mynewglossary{sym-Nu}{sly}{gls}{$\Nu$}
\mynewglossary{sym-Xi}{Xi.sly}{Xi.gls}{$\boldsymbol\Xi$}
%\mynewglossary{sym-Omicron}{sly}{gls}{$\Omicron$}
\mynewglossary{sym-Pi}{Pi.sly}{Pi.gls}{$\boldsymbol\Pi$}
%\mynewglossary{sym-Rho}{sly}{gls}{$\Rho$}
\mynewglossary{sym-Sigma}{Sigma.sly}{Sigma.gls}{$\boldsymbol\Sigma$}
%\mynewglossary{sym-Tau}{sly}{gls}{$\Tau$}
\mynewglossary{sym-Upsilon}{Upsilon.sly}{Upsilon.gls}{$\boldsymbol\Upsilon$}
\mynewglossary{sym-Phi}{Phi.sly}{Phi.gls}{$\boldsymbol\Phi$}
%\mynewglossary{sym-Chi}{sly}{gls}{$\Chi$}
\mynewglossary{sym-Psi}{Psi.sly}{Psi.gls}{$\boldsymbol\Psi$}
\mynewglossary{sym-Omega}{Omega.sly}{Omega.gls}{$\boldsymbol\Omega$}

\makeatletter



% ============================================================
% Greek letters
% ============================================================

\def\symgreek@Alpha{01}
\def\symgreek@Beta{02}
\def\symgreek@Gamma{03}
\def\symgreek@Delta{04}
\def\symgreek@Epsilon{05}
\def\symgreek@Zeta{06}
\def\symgreek@Eta{07}
\def\symgreek@Theta{08}
\def\symgreek@Iota{09}
\def\symgreek@Kappa{10}
\def\symgreek@Lambda{11}
\def\symgreek@Mu{12}
\def\symgreek@Nu{13}
\def\symgreek@Xi{14}
\def\symgreek@Omicron{15}
\def\symgreek@Pi{16}
\def\symgreek@Rho{17}
\def\symgreek@Sigma{18}
\def\symgreek@Tau{19}
\def\symgreek@Upsilon{20}
\def\symgreek@Phi{21}
\def\symgreek@Chi{22}
\def\symgreek@Psi{23}
\def\symgreek@Omega{24}

\def\symgreek@alpha{01}
\def\symgreek@beta{02}
\def\symgreek@gamma{03}
\def\symgreek@delta{04}
\def\symgreek@epsilon{05}
\def\symgreek@zeta{06}
\def\symgreek@eta{07}
\def\symgreek@theta{08}
\def\symgreek@iota{09}
\def\symgreek@kappa{10}
\def\symgreek@lambda{11}
\def\symgreek@mu{12}
\def\symgreek@nu{13}
\def\symgreek@xi{14}
\def\symgreek@omicron{15}
\def\symgreek@pi{16}
\def\symgreek@rho{17}
\def\symgreek@sigma{18}
\def\symgreek@tau{19}
\def\symgreek@upsilon{20}
\def\symgreek@phi{21}
\def\symgreek@chi{22}
\def\symgreek@psi{23}
\def\symgreek@omega{24}


% ============================================================
% Math formatting commands to ignore
% ============================================================

\def\symfmt@vec{1}
\def\symfmt@mathbf{1}
\def\symfmt@boldsymbol{1}
\def\symfmt@mathbb{1}
\def\symfmt@hat{1}
\def\symfmt@dot{1}
\def\symfmt@ddot{1}
\def\symfmt@bar{1}
\def\symfmt@tilde{1}
\def\symfmt@mathcal{1}
\def\symfmt@mathrm{1}
\def\symfmt@overline{1}
\def\symfmt@angle{1}
\def\symfmt@dDelta{1}
\def\symfmt@underline{1}

% ============================================================
% Sort one component
% ============================================================

\newcommand{\sym@sortcomponent}[1]{%
  \ifcsname symfmt@#1\endcsname
    % formatting command: ignore it
  \else
    \ifcsname symgreek@#1\endcsname
      \csname symgreek@#1\endcsname
    \else
      00#1%
    \fi
  \fi
}


% ============================================================
% Sort entire key
% ============================================================

\def\sym@sortaux#1-#2\@nil{%
  \ifcsname symfmt@#1\endcsname
    % skip formatting component
  \else
    \sym@sortcomponent{#1}%
    \ifx\relax#2\relax
    \else
      -%
    \fi
  \fi
  \ifx\relax#2\relax
  \else
    \sym@sortaux#2\@nil
  \fi
}

\newcommand{\mysortkey}[1]{%
  \sym@sortaux#1-\@nil
}


% ============================================================
% First actual symbol
% ============================================================

\def\getFirstChar#1#2\@nil{#1}
% Wrapper macro to make syntax cleaner
\newcommand{\firstletter}[1]{\getFirstChar#1\@nil}

\def\sym@firstaux#1-#2\@nil{%
  \ifcsname symfmt@#1\endcsname
    \sym@firstaux#2\@nil
  \else
    \ifcsname symgreek@#1\endcsname
      #1%
    \else
      \firstletter{#1}%
    \fi
  \fi
}

\newcommand{\mysymgroup}[1]{%
  \sym@firstaux#1-\@nil
}

% ============================================================
% Glossary entry
% ============================================================

\renewcommand{\glsxtrsetaliasnoindex}{%
    \glsxtrindexaliased
    \setkeys{glslink}{noindex}%
}

\NewDocumentCommand{\newsym}{m m m o}{%
  \ifglsentryexists{sym-#1}{%
    \PackageError{glossary-setup}{Duplicate symbol key `#1' (label `sym-#1') in \string\newsym}{%
      Remove or rename one of the duplicate \string\newsym definitions.\MessageBreak
      The symbol label `sym-#1' has already been defined.
    }%
  }{}%
  \edef\newsym@sortkey{\mysortkey{#1}}%
  \IfNoValueTF{#4}{%
    \newglossaryentry{sym-#1}{
      name={\ensuremath{#2}},
      category={symbol},
      sort={\newsym@sortkey},
      type={sym-\mysymgroup{#1}},
      description={#3}}%
  }{%
    \newglossaryentry{sym-#1}{
      name={\ensuremath{#2}},
      category={symbol},
      sort={\newsym@sortkey},
      type={alias-sym-\mysymgroup{#1}},
      description={#3},
      alias={sym-#4}}%
  }%
}
\makeatother

\newglossary[alg]{journal-acronyms}{acn}{acm.gls}{Acronyms}

\newglossary[clg]{control}{cld}{cln.gls}{Control Acronyms}
% control related
\newabbreviation[type=control]{acr-dp}{DP}{Dynamic Programming}
\newabbreviation[type=control]{acr-df}{DF}{Describing Function}
\newabbreviation[type=control]{acr-im}{IM}{Impedance Matching}
\newabbreviation[type=control]{acr-lqr}{LQR}{Linear Quadratic Regulator}
\newabbreviation[type=control]{acr-lmi}{LMI}{Linear Matrix Inequality}
\newabbreviation[type=control]{acr-pid}{PID}{Proportional Integral Derivative}
\newabbreviation[type=control]{acr-pi}{PI}{Proportional Integral}
\newabbreviation[type=control]{acr-pmp}{PMP}{Pontryagin Maximum Principle}
\newabbreviation[type=control]{acr-oc}{OC}{Optimal Control}
\newabbreviation[type=control,type=journal-acronyms]{acr-mfd}{MFD}{Modified Frequency Domain}
\newabbreviation[type=control]{acr-mpc}{MPC}{Model Predictive Control}
\newabbreviation[type=control]{acr-nmpc}{NMPC}{Nonlinear Model Predictive Control}
\newabbreviation[type=control,type=journal-acronyms]{acr-qlps}{QLPS}{Quasi-Linear Pseudo-Spectral}
\newabbreviation[type=control]{acr-ps}{PS}{Pseudo-Spectral}

\newglossary{methods}{mdd}{mdn.gls}{Computational Methods Acronyms}
% computational methodologies
\newabbreviation[type=methods]{acr-ccd}{CCD}{Control Co-Design}
\newabbreviation[type=methods,type=journal-acronyms]{acr-bem}{BEM}{Boundary Element Method}
\newabbreviation[type=methods]{acr-cfd}{CFD}{Computational Fluid Dynamics}
\newabbreviation[type=methods]{acr-fea}{FEA}{Finite Element Analysis}
\newabbreviation[type=methods]{acr-fem}{FEM}{Finite Element Method}
\newabbreviation[type=methods]{acr-lca}{LCA}{Life Cycle Analysis}
\newabbreviation[type=methods,type=journal-acronyms]{acr-meem}{MEEM}{Matched Eigenfunction Expansion Method}
\newabbreviation[type=methods,type=journal-acronyms]{acr-nlp}{NLP}{Nonlinear Program}
\newabbreviation[type=methods]{acr-mdo}{MDO}{Multidisciplinary Design Optimization}

\newglossary{org}{ord}{orn.gls}{Organization, Conference, and Publication Acronyms}
% organizations, institutions, professional societies, meetings, and publications
\newabbreviation[type=org]{acr-asme}{ASME}{American Society of Mechanical Engineers}
\newabbreviation[type=org]{acr-sae}{SAE}{Society of Automotive Engineers}
\newabbreviation[type=org]{acr-iec}{IEC}{International Electrotechnical Commission}
\newabbreviation[type=org]{acr-ieee}{IEEE}{Institute of Electrical and Electronics Engineers}
\newabbreviation[type=org]{acr-idetccie}{IDETC-CIE}{International Design Engineering Technical Conferences-Computers and Information in Engineering}
\newabbreviation[type=org]{acr-isope}{ISOPE}{International Society of Offshore and Polar Engineers}
\newabbreviation[type=org]{acr-isopeiopec}{IOPEC}{International Ocean and Polar Engineering Conference}
\newabbreviation[type=org]{acr-jfm}{JFM}{Journal of Fluid Mechanics}
\newabbreviation[type=org,description={National Renewable Energy Laboratory, currently known as National Laboratory of the Rockies}]{acr-nrel}{NREL}{National Renewable Energy Laboratory}
\newabbreviation[type=org,description={Mechanical and Aerospace Engineering Department at Cornell University}]{acr-mae}{MAE}{Mechanical and Aerospace Engineering}
\newabbreviation[type=org]{acr-mit}{MIT}{Massachusetts Institute of Technology}
\newabbreviation[type=org]{acr-siam}{SIAM}{Society for Industrial and Applied Mathematics}
\newabbreviation[type=org]{acr-umerc}{UMERC}{University Marine Energy Research Community}
\newabbreviation[type=org,description={Technical Specification standard published by the \gls{acr-iec}}]{acr-ts}{TS}{Technical Specification}

\newglossary{struct}{std}{stn.gls}{Structural Acronyms}
% structural related
\newabbreviation[type=struct]{acr-dlc}{DLC}{Design Load Case}
\newabbreviation[type=struct]{acr-fbd}{FBD}{Free Body Diagram}
\newabbreviation[type=struct]{acr-fls}{FLS}{Fatigue Limit State}
\newabbreviation[type=struct]{acr-fos}{FOS}{Factor of Safety}
\newabbreviation[type=struct]{acr-uls}{ULS}{Ultimate Limit State}

\newglossary{econ}{ecd}{ecn.gls}{Economic Acronyms} % Economic and Power Grid Acronyms
% economic related
\newabbreviation[type=econ]{acr-capex}{CAPEX}{Capital Expenditure}
\newabbreviation[type=econ]{acr-cbs}{CBS}{Cost Breakdown Structure}
\newabbreviation[type=econ]{acr-opex}{OPEX}{Operational Expenditures}
\newabbreviation[type=econ,type=journal-acronyms]{acr-aep}{AEP}{Annual Energy Production}
\newabbreviation[type=econ,type=journal-acronyms]{acr-fcr}{FCR}{Fixed Charge Rate}
\newabbreviation[type=econ,type=journal-acronyms]{acr-lcoe}{LCOE}{Levelized Cost of Energy}
\newabbreviation[type=econ]{acr-usd}{USD}{United States Dollar}
\newabbreviation[type=econ]{acr-nvoe}{NVOE}{Net Value of Energy}
% grid related
\newabbreviation[type=econ]{acr-cem}{CEM}{Capacity Expansion Model}
\newabbreviation[type=econ,description={Independent System Operator for New England, the organization which operates the region's electrical grid}]{acr-isone}{ISONE}{Independent System Operator for New England}

\newglossary{dynamics}{dyd}{dyn.gls}{Dynamics Acronyms}
% dynamics related
\newabbreviation[type=dynamics,type=journal-acronyms]{acr-dof}{DOF}{Degree of Freedom}
\newabbreviation[type=dynamics]{acr-abcd}{ABCD}{Cascade matrix}
\newabbreviation[type=dynamics]{acr-cw}{CW}{Capture width}
\newabbreviation[type=dynamics]{acr-dc}{DC}{Direct current}
\newabbreviation[type=dynamics]{acr-gr}{GR}{Gear Ratio}
\newabbreviation[type=dynamics]{acr-bw}{BW}{Bandwidth}
\newabbreviation[type=dynamics]{acr-rms}{RMS}{Root Mean Square}
\newabbreviation[type=dynamics]{acr-rom}{ROM}{Reduced Order Model}
\newabbreviation[type=dynamics,type=journal-acronyms]{acr-pto}{PTO}{Power Takeoff}

\newglossary{optim}{opd}{opn.gls}{Optimization and Search Acronyms}
% optimization and search related
\newabbreviation[type=optim]{acr-doe}{DOE}{Design of Experiments}
% \newabbreviation[type=optim]{acr-dv}{DV}{Design Variable}
\newabbreviation[type=optim]{acr-ga}{GA}{Genetic Algorithm}
\newabbreviation[type=optim,description={Karush-Kuhn-Tucker conditions for optimality}]{acr-kkt}{KKT}{Karush-Kuhn-Tucker}
% \newabbreviation[type=optim]{acr-lb}{LB}{Lower Bound}
\newabbreviation[type=optim]{acr-fd}{FD}{Finite Difference}
\newabbreviation[type=optim,type=journal-acronyms,description={Multiple Discipline Feasible \gls{acr-mdo} architecture}]{acr-mdf}{MDF}{Multiple Discipline Feasible}
\newabbreviation[type=optim,type=journal-acronyms,description={Simultaneous Analysis and Design \gls{acr-mdo} architecture}]{acr-sand}{SAND}{Simultaneous Analysis and Design}
% \newabbreviation[type=optim]{acr-ub}{UB}{Upper Bound}
\newabbreviation[type=optim,type=journal-acronyms]{acr-qcqp}{QCQP}{Quadratically Constrained Quadratic Program}
\newabbreviation[type=optim]{acr-qp}{QP}{Quadratic Program}
\newabbreviation[type=optim]{acr-sqp}{SQP}{Sequential Quadratic Programming}
\newabbreviation[type=optim,description={One at a Time method for design of experiments}]{acr-oat}{OAT}{One at a Time}
\newabbreviation[type=optim,type=journal-acronyms]{acr-xdsm}{XDSM}{Extended Design Structure Matrix}
% \newabbreviation[type=optim]{acr-bf}{BF}{Brute Force}

\newglossary{software}{swd}{swn.gls}{Software Tool Acronyms}
% software programs
\newabbreviation[description={System Advisor Model, an open-source techno-economic software published by \gls{acr-nrel}}]{acr-sam}{SAM}{System Advisor Model}
\newabbreviation[description={Wave Analysis Massachusetts Institute of Technology, a closed-source BEM package developed by \gls{acr-mit}}]{acr-wamit}{WAMIT}{Wave Analysis Massachusetts Institute of Technology}
\newabbreviation[description={}]{acr-wdrt}{WDRT}{WEC Design Response Toolbox}
\newabbreviation[description={}]{acr-jax}{JAX}{JAX Python package for automatic differentiation}
\newabbreviation[description={Response Amplitudes of Floating Turbines, an open-source software package}]{acr-raft}{RAFT}{Response Amplitudes of Floating Turbines}
\newabbreviation[description={}]{acr-nemoh}{NEMOH}{An open-source BEM package developed by Ecole Centrale de Nantes}
\newabbreviation[description={}]{acr-primre}{PRIMRE}{Portal and Repository for Information on Marine Renewable Energy}
\newabbreviation[description={Altair Inspire, a commercial tool for multidisciplinary optimization that integrates with \gls{acr-fea}}]{acr-alt}{ALT}{Altair Inspire}

\newglossary{wave}{wad}{wan.gls}{Wave and Wave Energy Acronyms}
% waves and wave energy
\newabbreviation[type=wave,type=journal-acronyms]{acr-wec}{WEC}{Wave Energy Converter}
\newabbreviation[type=wave]{acr-wets}{WETS}{Wave Energy Test Site}
\newabbreviation[type=wave]{acr-pbe}{PBE}{Powering the Blue Economy}
\newabbreviation[type=wave]{acr-reg}{REG}{Regular Waves}
\newabbreviation[type=wave,type=journal-acronyms]{acr-rm3}{RM3}{Reference Model 3 WEC}
\newabbreviation[type=wave]{acr-irr}{IRR}{Irregular waves}
\newabbreviation[type=wave,type=journal-acronyms,description={Joint Probability Density of wave height and wave period}]{acr-jpd}{JPD}{Joint Probability Density}
\newabbreviation[type=wave]{acr-sto}{STO}{Storm}
% wec types
\newabbreviation[type=wave]{acr-osw}{OSW}{Oscillating Surge WEC}
\newabbreviation[type=wave]{acr-otp}{OTP}{Overtopping WEC}
\newabbreviation[type=wave]{acr-owc}{OWC}{Oscillating Water Column WEC}
\newabbreviation[type=wave]{acr-pa}{PA}{Point Absorber WEC}
\newabbreviation[type=wave]{acr-trm}{TRM}{Terminator WEC}
\newabbreviation[type=wave]{acr-atn}{ATN}{Attenuator WEC}

% general computing, mathematics, and engineering
\newabbreviation[description={}]{acr-lts}{LTS}{Long-Term Support}
\newabbreviation[description={}]{acr-ram}{RAM}{Random Access Memory}
\newabbreviation[description={}]{acr-cpu}{CPU}{Central Processing Unit}
\newabbreviation[description={}]{acr-rgb}{RGB}{Red Green Blue}
\newabbreviation[description={}]{acr-si}{SI}{International System of Units}
\newabbreviation[description={}]{acr-thd}{THD}{Total Harmonic Distortion}

%%%%%%%%%%%%%%%%%%%%%%%%%%%%%%%%%%%%%%%%%%%%%%%%%%%%%%%%%%%%%
% \newabbreviation{acr-loc}{LOC}{Local Optimization Algorithm}
\newabbreviation{acr-uns}{UNS}{Unstructured}
\newabbreviation{acr-nl}{NL}{Nonlinear}
\newabbreviation{acr-alg}{ALG}{Algebraic}
\newabbreviation{acr-an}{AN}{Analytical}
\newabbreviation{acr-fit}{FIT}{Fit of model to numerical or experimental data}

% states
\newglossary{states}{sad}{san.gls}{States}
\newabbreviation[type=states]{acr-ca}{CA}{California}
\newabbreviation[type=states]{acr-hi}{HI}{Hawaii}

% wec subsystems or subsystem considerations
% \newabbreviation[description={}]{acr-ctl}{CTL}{Control}
\newabbreviation[description={}]{acr-dyn}{DYN}{Dynamics}
% \newabbreviation[description={}]{acr-econ}{ECON}{Economics}
\newabbreviation[description={}]{acr-geo}{GEO}{Geometry}
%\newabbreviation[description={}]{acr-ef}{EF}{Efficiency}
\newabbreviation[description={}]{acr-eff}{EFF}{Efficiency}
\newabbreviation[description={}]{acr-str}{STR}{Structures}
\newabbreviation[description={}]{acr-moor}{MOOR}{Mooring}
% \newabbreviation[description={}]{acr-amp}{AMP}{Motion amplitude}
\newabbreviation[description={}]{acr-ld}{LD}{Linear Damping}
\newabbreviation[description={}]{acr-pwr}{PWR}{Power}
% \newabbreviation[description={}]{acr-shp}{SHP}{Shape}
\newabbreviation[description={}]{acr-de}{DE}{Mooring physical design}
% \newabbreviation[description={}]{acr-sta}{STA}{Stability}
\newabbreviation[description={}]{acr-lo}{LO}{Loads}
% \newabbreviation[description={}]{acr-arry}{ARRY}{Array}
% \newabbreviation[description={}]{acr-var}{VAR}{Variation}
\makeglossaries

% generic stuff
\newsym{k-generic-index}{k}{Generic index}
\newsym{hat-e-z}{\hat{e}_z}{Unit vector in the vertical direction}
\newsym{mathbb-R-C-parent}{\mathbb{R},\mathbb{C}}{Set of (real, complex) numbers}
    \newsym{mathbb-R}{\mathbb{R}}{Set of real numbers}[mathbb-R-C-parent]
    \newsym{mathbb-C}{\mathbb{C}}{Set of complex numbers}[mathbb-R-C-parent]
\newsym{mathbf-X-x-parent}{\gls{sym-mathbf-X},\gls{sym-mathbf-x}}{Generic matrix}
    \newsym{mathbf-x}{\mathbf{x}}{Generic matrix}[mathbf-X-x-parent]
    \newsym{mathbf-X}{\mathbf{X}}{Generic matrix}[mathbf-X-x-parent]
\newsym{pi}{\pi}{Mathematical constant 3.14159...}
\newsym{j}{j}{Generic index representing an interior fluid region index in MEEM}
\newsym{vec-x}{\vec{x}}{Generic vector}
\newsym{mathbf-I}{\mathbf{I}}{Identity matrix}

% dimensions
    \newsym{D-d}{D_d}{Diameter of damping plate}[D-parent]
    \newsym{D-f}{D_f}{Outer diameter of float}[D-parent]
    \newsym{D-s}{D_s}{Outer diameter of spar}[D-parent]
    \newsym{D-f-in}{D_{f,\mathrm{in}}}{Inner diameter of float}[D-parent]
\newsym{t-parent}{t_{(*)}}{Structural thickness of (*) $\in \{d, f,b, s,r\} = \{$damping plate, float bottom plate, spar radial wall$\}$}
    \newsym{t-d}{t_d}{Structural thickness of damping plate}[t-parent]
    \newsym{t-s-r}{t_{s,r}}{Structural wall thickness of spar, measured radially}[t-parent]
    \newsym{t-f-b}{t_{f,b}}{Structural thickness of the float bottom plate}[t-parent]
\newsym{T-f-1}{T_{f,1}}{Draft of float, measured at outer diameter}
\newsym{h-fs-clear}{h_{fs,\mathrm{clear}}}{Float-spar clearance height}

% coordinates
\newsym{t}{t}{Time}
\newsym{r}{r}{Coordinate in the radial direction}
\newsym{theta}{\theta}{Angular coordinate along the circumference in cylindrical coordinate system}
\newsym{z}{z}{Cartesian and cylindrical coordinate in the vertical direction}

% wave environment
    \newsym{rho-w}{\rho_w}{Density of water}[rho-parent]
\newsym{g}{g}{Acceleration due to gravity}
\newsym{h}{h}{Water depth}
\newsym{omega}{\omega}{Angular frequency of equivalent regular wave}

% optim stuff
\newsym{J}{J}{Scalar objective function}

% dynamics stuff
\newsym{mathbf-M}{\mathbf{M}}{Mass matrix}
\newsym{ddot-vec-xi}{\ddot{\vec{\xi}}}{Body acceleration vector}
\newsym{dot-vec-xi}{\dot{\vec{\xi}}}{Body velocity vector}
\ifdefined\DISSERTATION
\newsym{H-wave-amplitude}{\frac{H}{2}}{Regular wave amplitude}
\else
\newsym{H-wave-amplitude}{A}{Wave amplitude}
\fi
\newsym{angle-xi-tol}{\angle \xi_\mathrm{tol}}{Dynamics solver convergence tolerance for angle of displacement phasor $\xi$}
\newsym{xi-tol}{\xi_{\mathrm{tol}}}{}

% meem stuff
\newsym{mathbf-A}{\mathbf{A}}{Matrix in MEEM, contains information about homogeneous potentials}
\newsym{vec-c}{\vec{c}}{MEEM output vector to obtain hydrodynamic coefficients from eigencoefficients}
\newsym{phi-i1}{\phi^{i1}}{Radiation velocity potential in the interior fluid region under the spar in MEEM}
\newsym{R-r}{R(r)}{Radial eigenfunction in MEEM}
\newsym{Z}{Z}{Vertical eigenfunction in MEEM}
\newsym{delta}{\delta}{Kronecker delta}
\newsym{mathbf-A-h}{\mathbf{A}_h}{Added mass matrix}[mathbf-A-h-parent]
\newsym{mathbf-B-h}{\mathbf{B}_h}{Radiation damping matrix}[mathbf-B-h-parent]
\newsym{mathbf-K-h}{\mathbf{K}_h}{Hydrostatic stiffness matrix}[mathbf-K-h-parent]
\newsym{Z-0-e}{Z_0^e}{Vertical eigenfunction for the exterior fluid region e, mode 0, in MEEM}
\newsym{a-parent}{a_{(\#)}}{Radius of body \# in MEEM, \# $\in \{1,2\}=\{$spar, float$\}$}
    \newsym{a-1}{a_1}{Radius of spar used for MEEM}[a-parent]
    \newsym{a-2}{a_2}{Radius of float used for MEEM}[a-parent]
\newsym{d-parent}{d_{(\#)}}{Draft of body \# in MEEM, \# $\in \{1,2\}=\{$spar, float$\}$}
    \newsym{d-1}{d_1}{Draft of spar used in MEEM}[d-parent]
    \newsym{d-2}{d_2}{Draft of float used in MEEM}[d-parent]
\newsym{vec-xi}{\vec{\xi}}{Time domain vector of vertical positions of each body degree of freedom}
\newsym{phi-e}       {\phi^e}{Total (particular + homogeneous) radiation potential in the exterior region $e$ in \glsxtrshort{acr-meem}}
\newsym{vec-b}{\vec{b}}{Right-hand side vector in MEEM matrix solve, contains information about particular potentials}

\ifdefined\whichjournal
\ifthenelse{\equal{\whichjournal}{Applied Ocean Research}}
{
    \newsym{N-0}{N_0}{Exterior-region vertical eigenfunction normalization term \gls{sym-N-k} evaluated at mode \gls{sym-k-index}=0.}
    \newsym{phi}{\phi}{Complex-valued radiation potential phasor as a function of space.}
}{}
\else
    \newsym{N-0}{N_0}{Exterior-region vertical eigenfunction normalization term \gls{sym-N-n-e} evaluated at mode \gls{sym-n-e}=0.}
    \newsym{phi}{\phi}{Complex-valued radiation potential phasor as a function of space. Shorthand for \gls{sym-phi-mathrm-R}.}
\fi

% bessel functions
\newsym{mathrm-H-0-1}{\mathrm{H}_0^1}{Hankel function of the first kind, order zero}
\newsym{mathrm-K-1}{\mathrm{K}_1}{Modified Bessel function of the second kind of order one}
\newsym{mathrm-K-0}{\mathrm{K}_0}{Modified Bessel function of the second kind of order zero}
\newsym{mathrm-I-0}{\mathrm{I}_0}{Modified Bessel function of the first kind of order zero}
\newsym{mathrm-I-1}{\mathrm{I}_1}{Modified Bessel function of the first kind of order 1}

% moved from paper-specific glossaries to fix undefined glossary entries
    \newsym{T-e}{T_e}{Energy period of irregular waves}[T-e-p-parent]
    \newsym{p-prime-0}{p^{'}_0}{Normalized price coefficient for constant capital and operational costs that do not scale with design}[p-parent]
    \newsym{T-s}{T_s}{Draft of spar}[T-parent]
    \newsym{F-max}{F_\mathrm{max}}{Peak instantaneous force that the generator is capable of producing}[A-A-max-parent]
    \newsym{P-max}{P_\mathrm{max}}{Maximum allowable power at the generator electrical port across all sea states}[A-A-max-parent]
\newsym{overline-P-elec}{\overline{P}_{\mathrm{elec}}}{Annual average electrical power generation}
\newsym{p}{p}{Instantaneous time-domain power at any port}
\newsym{p-pk-VI}{p_{\mathrm{pk},VI}}{Peak power over a wave period at the generator electrical port in a given sea state}
\newsym{F-p}{F_p}{\glsxtrshort{acr-pto} force on a single hydrodynamic degree of freedom}
\newsym{F-heave-mult}{F_{\mathrm{heave},\mathrm{mult}}}{Tuning multiplier for heave forces}
\newsym{xi-f}{\xi_f}{Time domain vertical position of the float}
\newsym{xi-s}{\xi_s}{Time domain vertical position of the spar}
\newsym{xi-lin}{\xi_{\mathrm{lin}}}{Maximum acceptable body displacement to maintain linear hydrodynamic assumptions}
\newsym{xi-slam}{\xi_\mathrm{slam}}{General slamming amplitude limit}
\newsym{h-fs-up}{h_{fs,\mathrm{up}}}{Allowable displacement of the float to move upwards relative to the spar}
\newsym{h-fs-down}{h_{fs,\mathrm{down}}}{Allowable displacement of the float to move downwards relative to the spar}
\newsym{N-harmonics}{N_{\mathrm{harmonics}}}{Number of harmonics used in \glsxtrshort{acr-meem} for all fluid regions}
\newsym{f}{f}{Generic scalar function}
\newsym{D-d-min}{D_{d,\mathrm{min}}}{Minimum permissible damping plate diameter}
\newsym{T-f-2}{T_{f,2}}{Draft of float, measured at inner diameter}
\newsym{h-f}{h_f}{Float height, including both above and below water}
\newsym{h-s}{h_s}{Spar height, including both above and below water}
\newsym{h-1-stiff-d}{h_{1,\mathrm{stiff},d}}{Height of the larger damping plate stiffener}
\newsym{h-d}{h_d}{Damping plate thickness}
\newsym{h-stiff-f}{h_{\mathrm{stiff},f}}{Height of float structural stiffeners}
\newsym{GM}{\overline{GM}}{Distance from center of gravity (G) to metacenter (M)}
\newsym{V-f-struct}{V_{f,\mathrm{struct}}}{Structural volume of the float}
\newsym{V-s-struct}{V_{s,\mathrm{struct}}}{Structural volume of the spar}
\newsym{T-period}{T}{Period of equivalent regular wave}
\newsym{T-struct}{T_\mathrm{struct}}{100-year (design storm) wave period used for structural load cases}
\newsym{BW}{BW}{Half-power bandwidth of the \glsxtrshort{acr-jpd} of a given location}
\newsym{H-s}{H_s}{Significant wave height of irregular waves}
\newsym{H-s-struct}{H_{s,\mathrm{struct}}}{Significant wave height used for structural checks in the storm loadcase}
\newsym{J-wave}{J_\mathrm{wave}}{Incident wave energy flux at a location}



\newcommand*\rot{\rotatebox{90}} % for rotated table headers
\newcommand{\phant}{\vphantom{\frac{-1}{\rho_M/\rho_w - 1}}} % to make inequality stacks the right height

\def\fillandplacepagenumber{% for landscape page numbers
 \par\thispagestyle{empty}%
 \vbox to 0pt{\vss}\vfill
 \vbox to 0pt{\baselineskip0pt
   \hbox to\linewidth{\hss}%
   \baselineskip\footskip
   \hbox to\linewidth{%
     \hfil\thepage\hfil}\vss}}

% allows bibcop to be used when bibliography includes multiple arguments
\makeatletter
% \AtBeginDocument{%
% \ifnum\pdf@shellescape=0
%   % shell escape is disabled, so don't patch bibcop
% \else
% \ifdefined\bibliography

%   \RequirePackage{etoolbox}

%   \renewcommand{\bibliography}[1]{%
%     \def\bibcop@list{#1}%
%     \@for\bibcop@file:=\bibcop@list\do{%
%       \edef\bibcop@file{\bibcop@file}%
%       \ifx\bibcop@file\@empty\else
%         \expandafter\bibcop@processbib\expandafter{\bibcop@file}%
%       \fi
%     }%
%     \bibcop@truebibliography{#1}%
%   }

%   \newcommand{\bibcop@processbib}[1]{%
%     % trim spaces
%     \edef\bibcop@file{\zap@space#1 \@empty}%

%     % skip empty entries
%     \ifx\bibcop@file\@empty
%       \relax
%     \else

%       % build .bib filename
%       \edef\bibcop@bibfile{\bibcop@file.bib}%

%       % use kpsewhich to resolve path using BIBINPUTS
%       \immediate\write18{kpsewhich \bibcop@bibfile > aux/\jobname.bibcop.tmp}%

%       % read result
%       \begingroup
%         \endlinechar=-1
%         \newread\bibcop@in
%         \openin\bibcop@in=aux/\jobname.bibcop.tmp
%         \read\bibcop@in to \bibcop@resolved
%         \closein\bibcop@in
%         \xdef\bibcop@resolved{\bibcop@resolved}%
%       \endgroup

%       % fallback if kpsewhich fails
%       \ifx\bibcop@resolved\@empty
%         \edef\bibcop@resolved{\bibcop@bibfile}%
%       \fi

%       % call original bibcop
%       \begingroup
%         \edef\x{\endgroup
%           \noexpand\bibcop@exec{\bibcop@resolved}%
%         }%
%       \x

%     \fi
%   }
% \fi
% \fi
% }
\makeatother

% references to appendix sections should be called "appendix" instead of "section"
% \AddToHook{cmd/appendices/before}{%
%   \crefalias{section}{suppsec}%
%   \crefalias{subsection}{suppsec}
%   \crefalias{subsubsection}{suppsec}
%   \crefalias{chapter}{suppsec}
% }
% \AddToHook{cmd/subappendices/before}{%
%   \crefalias{section}{suppsec}%
%   \crefalias{subsection}{suppsec}
%   \crefalias{subsubsection}{suppsec}
%   \crefalias{chapter}{suppsec}
% }
\crefname{suppsec}{appendix}{appendices}


\newcolumntype{P}[1]{>{\centering\arraybackslash}p{#1}}
\newcolumntype{M}[1]{>{\centering\arraybackslash}m{#1}}
\newcolumntype{L}[1]{>{\raggedright\arraybackslash}p{#1}}
\newcolumntype{?}{!{\vrule width 1.5pt}}

%% For including figures, graphicx.sty has been loaded in
%% elsarticle.cls. If you prefer to use the old commands
%% please give \usepackage{epsfig}
\usepackage{xurl}
\usepackage{microtype}
\usepackage{subcaption} % for subfigures
\usepackage{tikz}
\usepackage{fmtcount} % for \ordinalstringnum command in discussion.tex (imported from results)
%\usepackage[section]{placeins} % prevent figs from going in wrong section

\input{figure-widths.tex}
\input{table-column-specs.tex}


%%%%%%%%%%%%%%%%%%%%%%%%%%%%%%%%%%%%%%%%%% GENERAL %%%%%%%%%%%%%%%%%%%%%%%%%%%%%%%%%%%%%%%%%%
% generic variables
\newsym{a}{a}{Generic scalar constant}
\newsym{x}{x}{Generic scalar variable}
\newsym{tilde-x}{\widetilde{x}}{Generic complex time-harmonic variable}
\newsym{b}{b}{Generic scalar constant used for integration limits}

% mathematical constants

\newsym{epsilon}{\epsilon}{Small positive numerical tolerance}

% mathematical objects
\newsym{s}{s}{Laplace variable}
\newsym{vec-zero}{\vec{0}}{Vector of zeros}

% mathematical operators
\newsym{mathrm-diag}{\mathrm{diag}}{Operator to create a diagonal matrix from a vector}
\newsym{Sigma}{\Sigma}{Grand sum operator (sum of all elements in a matrix)}

% bessel functions
\newsym{mathrm-J-1}{\mathrm{J}_1}{Bessel function of the first kind of order one}

%%%%%%%%%%%%%%%%%%%%%%%%%%%%%%%%%%%%%%%%%% DIMENSIONS %%%%%%%%%%%%%%%%%%%%%%%%%%%%%%%%%%%%%%%%%%
\newsym{vec-x-shape}{\vec{x}_\mathrm{shape}}{Set of design input variables related to bulk geometry and hull shape}

% diameters
\newsym{D-f-nom}{D_{f,\mathrm{nom}}}{Outer diameter of float in nominal \gls{acr-rm3} design}
\newsym{D}{D}{Generic diameter}

% drafts
\newsym{T-draft}{T}{Draft}
\newsym{vec-T}{\vec{T}}{Vector of drafts for each hydrodynamic degree of freedom}
\newsym{T-s-nom}{T_{s,\mathrm{nom}}}{Draft of spar in nominal \gls{acr-rm3} design}

% non-draft vertically oriented bulk dimensions
\newsym{z-top}{z_{\mathrm{top}}}{Vertical coordinate of top of body for surge excitation force calculation}

% coords
\newsym{y}{y}{Cartesian coordinate in the direction of wave propagation}

% radii / radially oriented dims
\newsym{R-p}{R_p}{Radius of the damping plate}
\newsym{R-x}{R_x}{The smaller of the float radius and the damping plate radius, used for spar hydrodynamic coefficient calculation}
\newsym{h-0}{h_0}{Perpendicular height of the trapezoidal top and bottom float plates}
\newsym{alpha}{\alpha}{Nondimensional ratio of damping plate diameter to float diameter, \gls{sym-D-d}/\gls{sym-D-f}, used for spar excitation force scaling}
\newsym{b-inner-damping-plate}{b_\mathrm{plate}}{Inner radius of damping plate, \gls{sym-D-s}/2}
\newsym{a-outer-damping-plate}{a_\mathrm{plate}}{Outer radius of damping plate, \gls{sym-D-d}/2}
\newsym{beta-olaya}{\beta}{Nondimensional ratio of spar diameter to float diameter, \gls{sym-D-s}/\gls{sym-D-f}, used for spar excitation force scaling}
\newsym{sigma-ratio}{\varsigma}{Nondimensional ratio of spar diameter to damping plate diameter, \gls{sym-D-s}/\gls{sym-D-d}, used for spar added mass scaling}

% structural dimensions
\newsym{t-f-top-eq}{t_{f,\mathrm{top},eq}}{Equivalent thickness of the float top plate}

%%%%%%%%%%%%%%%%%%%%%%%%%%%%%%%%%%%%%%%%%% GEOMETRIC %%%%%%%%%%%%%%%%%%%%%%%%%%%%%%%%%%%%%%%%%%
% geometry relevant to pitch stability
\newsym{BM}{\overline{BM}}{Length of line segment connecting center of buoyancy (B) to metacenter (M)}
\newsym{KB}{\overline{KB}}{Length of line segment connecting keel (K) and center of buoyancy (B)}
\newsym{KG}{\overline{KG}}{Length of line segment connecting keel (K) and center of gravity (G)}
\newsym{G}{G}{Vertical coordinate of center of gravity}
\newsym{B}{B}{Vertical coordinate of center of buoyancy}
\newsym{K-keel}{K}{Vertical coordinate of keel}
\newsym{M-metacenter}{M}{Vertical coordinate of metacenter}

% areas
\newsym{A-w}{A_w}{Waterplane area}
\newsym{A-w-f}{A_{w,f}}{Waterplane area of float}
\newsym{A-w-s}{A_{w,s}}{Waterplane area of spar}
\newsym{A-wet}{A_{\mathrm{wet}}}{Wetted surface area}
\newsym{A}{A}{Structural cross sectional area of the spar for heave loads}

% volumes
\newsym{V-f-struct-max}{V_{f,\mathrm{struct},\mathrm{max}}}{}
\newsym{V-f-struct-min}{V_{f,\mathrm{struct},\mathrm{min}}}{}
\newsym{V-s-struct-max}{V_{s,\mathrm{struct},\mathrm{max}}}{}
\newsym{V-s-struct-min}{V_{s,\mathrm{struct},\mathrm{min}}}{}
\newsym{V-pto}{V_{\mathrm{pto}}}{Volume of power takeoff system}
\newsym{V-struct}{V_\mathrm{struct}}{Volume of structural material}
\newsym{V-sub}{V_\mathrm{sub}}{Submerged volume}
\newsym{V-surf}{V_{\mathrm{surf}}}{Surface volume (above the waterline)}

% masses
\newsym{m-bal}{m_{\mathrm{bal}}}{Mass of ballast}
\newsym{m-f}{m_f}{Mass of float}
\newsym{m-s}{m_s}{Mass of spar}
\newsym{M-op}{\mathbf{M}_{\mathrm{op}}}{Operational mass matrix}
\newsym{M-storm}{\mathbf{M}_{\mathrm{storm}}}{Storm-condition mass matrix}

%%%%%%%%%%%%%%%%%%%%%%%%%%%%%%%%%%%%%%%%%% DYNAMIC COEFFICIENTS %%%%%%%%%%%%%%%%%%%%%%%%%%%%%%%%%%%%%%%%%%
% added mass
\newsym{A-h}{A_h}{Diagonal element ($A_f$ or $A_s$) of the added mass matrix \gls{sym-mathbf-A-h-op}, computed via \gls{acr-meem}}
\newsym{mathbf-A-h-parent}{\mathbf{A}_h,\mathbf{A}_{h,\mathrm{op}},\mathbf{A}_{h,\mathrm{st}}}{Added mass matrix, in operational (op)/storm (st) conditions}
    \newsym{mathbf-A-h-op}{\mathbf{A}_{h,\mathrm{op}}}{Added mass matrix for operational conditions}[mathbf-A-h-parent]
    \newsym{mathbf-A-h-st}{\mathbf{A}_{h,\mathrm{st}}}{Added mass matrix for storm conditions}[mathbf-A-h-parent]
\newsym{A-parent}{A_{(\mathrm{body})}}{Added mass coefficient of (body)$\in\{f,s,m,c\}=\{$float,spar,merged,coupling$\}$}
    \newsym{A-m}{A_m}{Added mass coefficient of the merged float-spar system due to its own acceleration}[A-parent]
    \newsym{A-s}{A_s}{Added mass coefficient of the spar due to its own acceleration}[A-parent]
    \newsym{A-c}{A_c}{Added mass coupling coefficient, representing forces on the float due to the spar's acceleration and vice versa}[A-parent]
    \newsym{A-f}{A_f}{Added mass coefficient of the float due to its own acceleration}[A-parent]

% damping
\newsym{B-c}{B_c}{Radiation damping coefficient for float-spar coupling}
\newsym{mathbf-B-d}{\mathbf{B}_d}{Damping matrix due to drag}
\newsym{B-f}{B_f}{Radiation damping coefficient for float}
\newsym{B-h}{B_h}{Diagonal element (\gls{sym-B-f} or \gls{sym-B-s}) of the radiation damping matrix \gls{sym-mathbf-B-h-op}, computed via \gls{acr-meem}}
\newsym{mathbf-B-h-op}{\mathbf{B}_{h,\mathrm{op}}}{Operational radiation damping matrix}
\newsym{mathbf-B-h-st}{\mathbf{B}_{h,\mathrm{st}}}{Storm radiation damping matrix}
\newsym{B-l}{B_l}{Control damping coefficient at load}
\newsym{B-l-damp}{B_{l,\mathrm{damp}}}{Control damping coefficient at the load in damping control}
\newsym{B-l-reac}{B_{l,\mathrm{reac}}}{Control damping coefficient at the load in reactive control}
\newsym{B-m}{B_m}{Radiation damping coefficient for merged float-spar}
\newsym{B-s}{B_s}{Radiation damping coefficient for spar}
\newsym{mathbf-B-p}{\mathbf{B}_p}{\gls{acr-pto} damping matrix}
\newsym{B-h-ii}{B_{h,ii}}{}

% stiffnesses
\newsym{K-f}{K_f}{Hydrostatic stiffness of float}
\newsym{K-h}{K_h}{Diagonal element (\gls{sym-K-f} or \gls{sym-K-s}) of hydrostatic stiffness matrix \gls{sym-mathbf-K-h-op}}
\newsym{mathbf-K-h-op}{\mathbf{K}_{h,\mathrm{op}}}{Operational hydrostatic stiffness matrix}
\newsym{mathbf-K-h-st}{\mathbf{K}_{h,\mathrm{st}}}{Storm hydrostatic stiffness matrix}
\newsym{K-l}{K_l}{Control stiffness at load}
\newsym{K-l-reac}{K_{l,\mathrm{reac}}}{Control stiffness at load under reactive control}
\newsym{K-l-damp}{K_{l,\mathrm{damp}}}{Control stiffness at load under damping control}
\newsym{K-m}{K_m}{Hydrostatic stiffness of merged float-spar}
\newsym{mathbf-K-p}{\mathbf{K}_p}{\gls{acr-pto} stiffness matrix}
\newsym{K-s}{K_s}{Hydrostatic stiffness of spar}

% excitation coefficients
\newsym{gamma}{\gamma}{Hydrodynamic excitation coefficient}
\newsym{Upsilon}{\Upsilon}{Nondimensional helper function for spar excitation coefficient analysis}
\newsym{gamma-m}{\gamma_m}{Hydrodynamic excitation coefficient for merged float-spar}
\newsym{vec-gamma}{\vec{\gamma}}{Complex wave-excitation coefficient vector from hydrodynamic forcing}
\newsym{vec-gamma-d}{\vec{\gamma}_d}{Drag excitation coefficient vector}
\newsym{gamma-f}{\gamma_f}{Hydrodynamic excitation coefficient for float}
\newsym{gamma-s}{\gamma_s}{Hydrodynamic excitation coefficient for spar}
\newsym{gamma-i}{\gamma_i}{}

% impedances
\newsym{vec-Z-i-th}{\vec{Z}_{i,\mathrm{th}}}{Vector of Th\'{e}venin-equivalent intrinsic mechanical impedances}
\newsym{Z-i}{Z_i}{Intrinsic mechanical impedance scalar}
\newsym{mathbf-Z-i}{\mathbf{Z}_i}{Intrinsic impedance matrix}
\newsym{Z-i-c}{Z_{i,c}}{Coupling (off-diagonal) term of the intrinsic impedance matrix \gls{sym-mathbf-Z-i}, representing hydrodynamic coupling between float and spar}
\newsym{Z-i-f}{Z_{i,f}}{Mechanical intrinsic impedance of float}
\newsym{Z-i-s}{Z_{i,s}}{Mechanical intrinsic impedance of spar}
\newsym{Z-i-th}{Z_{i,\mathrm{th}}}{Thevenin-equivalent intrinsic mechanical impedance}
\newsym{Z-l}{Z_l}{Control electrical impedance at load}
\newsym{Z-l-reac}{Z_{l,\mathrm{reac}}}{Control electrical impedance at load in reactive control}
\newsym{mathbf-Z-p}{\mathbf{Z}_p}{Powertrain impedance matrix}
\newsym{Z-s-th}{Z_{s,\mathrm{th}}}{Thevenin-equivalent electrical source impedance}
\newsym{Z-s-th-beta}{Z_{s,\mathrm{th},\beta}}{Thevenin-equivalent electrical source impedance in the $\beta$th sea state}
\newsym{Z-shaft}{Z_{\mathrm{shaft}}}{Mechanical impedance of drivetrain shaft and transmission system}
\newsym{Z-term}{Z_{\mathrm{term}}}{Termination impedance at any port}
\newsym{Z-u}{Z_u}{Control mechanical impedance}
\newsym{Z-w}{Z_w}{Electrical impedance of generator winding}
\newsym{mathbf-Z-H-T-i-th}{\mathbf{Z}^{H,T}_{i,\mathrm{th}}}{Matrix of Thevenin-equivalent intrinsic mechanical impedance \gls{sym-Z-i-th} across all sea states \gls{sym-H}, \gls{sym-T-period}}
\newsym{mathbf-Z}{\mathbf{Z}}{Impedance matrix}

%%%%%%%%%%%%%%%%%%%%%%%%%%%%%%%%%%%%%%%%%% DYNAMIC SOLUTION PROCESS %%%%%%%%%%%%%%%%%%%%%%%%%%%%%%%%%%%%%%%%%%
% constraint/metric formulation
\newsym{mathbf-Q-P}{\mathbf{Q}_P}{Matrix to recover average power from complex port variables}
\newsym{mathbf-Q-hat-e-2}{\mathbf{Q}_{|\gls{sym-hat-e}|^2}}{Matrix to recover squared effort magnitude from complex port variables}
\newsym{mathbf-Q-hat-q-2}{\mathbf{Q}_{|\gls{sym-hat-q}|^2}}{Matrix to recover squared flow magnitude from complex port variables}
\newsym{mathbf-Q-M}{\mathbf{Q}_M}{Matrix to recover a generic quadratic metric from complex port variables}
\newsym{mathbf-Q-S}{\mathbf{Q}_S}{Matrix to recover complex power from complex port variables}
\newsym{underline-M}{\underline{M}}{Generic metric to be constrained or optimized, expressed as a quadratic product of complex port variables}

% cascade matrices
\newsym{mathbf-a}{\mathbf{a}}{Cascade matrix}
\newsym{mathbf-a-VI-leftarrow-F-dot-X}{\mathbf{a}_{VI \leftarrow F\dot{X}}}{Cascade matrix from PTO mechanical port to generator electrical port}
\newsym{mathbf-a-F-p-dot-xi-f-leftarrow-VI}{\mathbf{a}_{(F_p\dot{\xi})_f \leftarrow VI}}{Cascade matrix from generator electrical port to float mechanical port}
\newsym{mathbf-a-F-p-dot-xi-s-leftarrow-VI}{\mathbf{a}_{(F_p\dot{\xi})_s \leftarrow VI}}{Cascade matrix from generator electrical port to spar mechanical port}
\newsym{mathbf-a-F-dot-X-leftarrow-VI}{\mathbf{a}_{F\dot{X} \leftarrow VI}}{Cascade matrix from generator electrical port to PTO mechanical port}
\newsym{mathbf-a-F-dot-X-leftarrow-tau-Omega}{\mathbf{a}_{F\dot{X} \leftarrow \tau\Omega}}{Cascade matrix from generator mechanical port to PTO mechanical port}
\newsym{mathbf-a-tau-Omega-leftarrow-VI}{\mathbf{a}_{\tau\Omega \leftarrow VI}}{Cascade matrix from electrical port to generator mechanical variables}
\newsym{A-a}{A_\mathbf{a}}{First element of a cascade matrix \gls{sym-mathbf-a}}
\newsym{B-a}{B_\mathbf{a}}{Second element of a cascade matrix \gls{sym-mathbf-a}}
\newsym{C-a}{C_\mathbf{a}}{Third element of a cascade matrix \gls{sym-mathbf-a}}
\newsym{D-a}{D_\mathbf{a}}{Fourth element of a cascade matrix \gls{sym-mathbf-a}}

% slamming
\newsym{theta-surf}{\theta_\mathrm{surf}}{Worst-case phase angle of the body relative to the incident wave across the entire body surface where slamming and submersion should be prevented}

%%%%%%%%%%%%%%%%%%%%%%%%%%%%%%%%%%%%%%%%%% ECONOMICS %%%%%%%%%%%%%%%%%%%%%%%%%%%%%%%%%%%%%%%%%%
% cost
\newsym{C-1}{C_1}{Single-unit capital cost for economies of scale calculation}
\newsym{C}{C}{Capital cost (subscript shows cost category)}
\newsym{C-pto}{C_{\mathrm{pto}}}{Capital cost of the \gls{acr-pto}}
\newsym{C-pto-constant}{C_{\mathrm{pto,constant}}}{Constant term in the \gls{acr-pto} capital cost}
\newsym{C-struct}{C_\mathrm{struct}}{Capital cost of structural material}
\newsym{C-unit}{C_{\mathrm{unit}}}{Unit cost at a particular quantity, for economies of scale calculation}
\newsym{C-fixed}{C_\mathrm{fixed}}{Capital cost that is modeled as independent of the design}
\newsym{C-infty}{C_\infty}{Asymptotic unit cost as quantity approaches infinity}

% economies of scale
\newsym{N-WEC}{N_{\mathrm{WEC}}}{Number of \glspl{acr-wec} in array for economies-of-scale calculations (not dynamic interactions)}
\newsym{beta-expo}{\beta}{Exponent for economies of scale cost relationship}

% price
\newsym{p-prime-star}{p^{'}_{(*)}}{Normalized price coefficient for any quantity}
\newsym{p-s}{p_s}{Price coefficient for structures}
\newsym{p-star}{p^{'}_{(*)}}{Normalized price coefficient for any quantity}
\newsym{p-F}{p_F}{Price coefficient for generator force}
\newsym{p-P}{p_P}{Price coefficient for generator power}
\newsym{p-prime-F}{p^{'}_F}{Normalized price coefficient for generator force}
\newsym{p-prime-P}{p^{'}_P}{Normalized price coefficient for generator power}
\newsym{p-prime-s}{p^{'}_s}{Normalized price coefficient for structures}

%%%%%%%%%%%%%%%%%%%%%%%%%%%%%%%%%%%%%%%%%% WAVE ENVIRONMENT %%%%%%%%%%%%%%%%%%%%%%%%%%%%%%%%%%%%%%%%%%
% wave frequency/period
\newsym{T-p}{T_p}{Peak period of an irregular wave spectrum}
\newsym{k-wavenumber}{k_{\mathrm{wvn}}}{Wavenumber of incident wave}
\newsym{V-g}{V_g}{Group velocity of incident regular wave}

% wave height
\newsym{H}{H}{Wave height of equivalent regular wave}

% other environmental scalars

% representation of different sea states
\newsym{beta}{\beta}{Index of sea state in \gls{acr-jpd}}
\newsym{N-sea}{N_{\mathrm{sea}}}{Number of sea states in \gls{acr-jpd} with nonzero probability}
\newsym{mathbf-J-H-T-wave}{\mathbf{J}^{H,T}_\mathrm{wave}}{Matrix of incident energy density in each sea state}
\newsym{mathbf-CW-H-T-max}{\mathbf{CW}^{H,T}_{\mathrm{max}}}{Matrix of maximal capture widths in each sea state}
\newsym{w-beta}{w_\beta}{Weight assigned to sea state $\beta$ in annual average power calculations}

%%%%%%%%%%%%%%%%%%%%%%%%%%%%%%%%%%%%%%%%%% STRUCTURES %%%%%%%%%%%%%%%%%%%%%%%%%%%%%%%%%%%%%%%%%%
% structures - material properties
\newsym{rho-M}{\rho_M}{Density of structural material}
\newsym{E}{E}{Modulus of elasticity (Young's modulus) of structural material}
\newsym{P-r}{P_r}{Material constant for strength reduction in combined loading}
\newsym{nu-poisson}{\nu_M}{Poisson's ratio of structural material}

% structures - loads
\newsym{F-tube}{F_\mathrm{tube}}{Structural force in each of the four tubes connecting the damping plate to the spar}
\newsym{W}{W}{Vertical load applied to each trapezoidal float top plate, equal to the heave force divided by number of float sections}
\newsym{F-crit}{F_\mathrm{crit}}{Critical buckling load}

% structures - equivalent thickness
\newsym{t-eq}{t_{\mathrm{eq}}}{Equivalent thickness of stiffened plate}
\newsym{D-eq}{D_{\mathrm{eq}}}{Equivalent flexural rigidity of a stiffened plate}
\newsym{S-eq}{S_{\mathrm{eq}}}{Equivalent section modulus of stiffened plate}

% structural dimensions
\newsym{thick}{t}{Structural thickness}
\newsym{w}{w}{Effective structural width for plates (distance between stiffeners)}
\newsym{h-stiff}{h_{\mathrm{stiff}}}{Height of a structural stiffener for any component}
\newsym{b-1}{b_1}{Smaller base dimension of the float and bottom plate trapezoid}
\newsym{b-2}{b_2}{Larger base dimension of the float and bottom plate trapezoid}
\newsym{L}{L}{Beam length in a buckling calculation}
\newsym{r-0}{r_0}{Radius of load distribution on float top plate due to welded join with the spar attachment tube}
\newsym{I-zz}{I_{zz}}{Second moment of area of the spar column about the vertical axis}
\newsym{x-max}{x_\mathrm{max}}{Larger side length of a fixed-edge rectangular plate}
\newsym{x-min}{x_\mathrm{min}}{Smaller side length of a fixed-edge rectangular plate}
\newsym{z-max}{z_\mathrm{max}}{Largest fiber distance from the neutral axis of a stiffened plate}

% structures - nondim factors
\newsym{omega-s}{\omega_s}{Intermediate value in spar combined loading structural calculation}
\newsym{psi}{\psi}{Load distrubution factor for stiffened plate effective breadth ratio calculation}
\newsym{zeta-s}{\zeta_s}{Intermediate value in spar combined loading structural calculation}
\newsym{epsilon-slenderness}{\varepsilon}{Slenderness factor for stiffened plate effective breadth ratio calculation}
\newsym{K-end}{K_{\mathrm{end}}}{End fixity coefficient for spar column buckling}
\newsym{k-buckle}{k_\mathrm{buckle}}{Coefficient for plate buckling}
\newsym{rho-psi}{\rho_\psi}{Portion of the effective breadth ratio \gls{sym-rho} related to the load distribution factor \gls{sym-psi}}
\newsym{rho}{\rho}{Effective breadth ratio of a stiffened plate, used to measure the validity of using the equivalent-thickness method}
\newsym{rho-epsilon}{\rho_\varepsilon}{Component of the effective breadth ratio \gls{sym-rho} related to the slenderness factor \gls{sym-epsilon-slenderness}}
\newsym{beta-plate-0}{\beta_0}{Tabulated nondimensional coefficient used to find maximum bending stress under a uniformly distributed load as a function of plate aspect ratio}
\newsym{beta-plate-1}{\beta_1}{Tabulated nondimensional coefficient used to find bending stress at the plate's edge under a small circular distributed load as a function of plate aspect ratio}
\newsym{beta-plate-2}{\beta_2}{Tabulated nondimensional coefficient used to find bending stress at the plate's center under a small circular distributed load as a function of plate aspect ratio}
\newsym{m-trap}{m_\mathrm{trap}}{Slope of the trapezoidal float top and bottom plates, used to calculate the trapezoidal plate's equivalent rectangular plate dimensions for bending stress calculations}

% structural stresses
\newsym{f-1}{f_1}{Maximum compression stress along the width of a stiffened plate, for calculating the effective breadth ratio}
\newsym{f-2}{f_2}{Maximum tension stress along the width of a stiffened plate (or minimum compression stress, if there is no tension), for calculating the effective breadth ratio}
\newsym{sigma-s}{\sigma_s}{Structural stress of the spar column}
\newsym{sigma}{\sigma}{Structural stress}
\newsym{sigma-buckle}{\sigma_\mathrm{buckle}}{Critical buckling stress}
\newsym{sigma-cent}{\sigma_\mathrm{cent}}{Float plate stress at center of top plate}
\newsym{sigma-edge}{\sigma_\mathrm{edge}}{Float plate stress at edge of top plate}
\newsym{sigma-r}{\sigma_r}{Plate stress in radial direction}
\newsym{sigma-theta}{\sigma_\theta}{Plate stress in circumferential direction}
\newsym{q-pressure}{q}{Distributed pressure load applied to a structural component}
\newsym{sigma-Y}{\sigma_Y}{Structural yield stress}

% structures - stiffness, moments, deflections
\newsym{K-tube}{K_\mathrm{tube}}{Effective mechanical stiffness of tube connecting damping plate to spar}
\newsym{M-top}{M_\mathrm{top}}{Bending moment per unit width of the float top plate}
\newsym{M-bot}{M_\mathrm{bot}}{Bending moment per unit width of the float bottom plate}
\newsym{delta-plate}{\delta_\mathrm{plate}}{Deflection of damping plate}
\newsym{delta-tube}{\delta_\mathrm{tube}}{Deflection of damping plate connection tube}
\newsym{overline-M-r}{\overline{M}_r}{Nondimensional radial bending moment per unit width in damping plate}
\newsym{overline-delta-plate}{\overline{\delta}_\mathrm{plate}}{Nondimensional deflection of damping plate}

%%%%%%%%%%%%%%%%%%%%%%%%%%%%%%%%%%%%%%%%%% PTO %%%%%%%%%%%%%%%%%%%%%%%%%%%%%%%%%%%%%%%%%%
% PTO
\newsym{R}{R}{Equivalent radius for a static linear-rotational transmission, including any gear ratio}
\newsym{mathbf-h-kin}{\mathbf{h}_{\mathrm{kin}}}{Hybrid matrix of kinematics, square even for underactuated systems}
\newsym{R-gear}{R_{\mathrm{gear}}}{Equivalent radius of gear, pinion, or other linear-to-rotational drivetrain component}
\newsym{mathbf-h}{\mathbf{h}}{Hybrid matrix}
\newsym{k-tau}{k_\tau}{Torque coefficient of generator}
\newsym{mathbf-T-kin}{\mathbf{T}_{\mathrm{kin}}}{Rectangular kinematics matrix}
% aor - consider making aor notation match jfm

% meem dimensions
\newsym{a-3}{a_3}{Average radius of float and spar used for \glsxtrshort{acr-meem}}
\newsym{a-in}{a_{\mathrm{in}}}{Inner radius of body used for \glsxtrshort{acr-meem}}
\newsym{a-out}{a_{\mathrm{out}}}{Outer radius of body used for \glsxtrshort{acr-meem}}
\newsym{k-wvn-parent}{\gls{sym-k-wavenumber}, \gls{sym-m-0}}{Wavenumber of incident wave}
    \newsym{m-0}{m_0}{Wavenumber}[k-wvn-parent]

% matrices / outputs
\newsym{c}{c}{\glsxtrshort{acr-meem} output vector to obtain hydrodynamic coefficients from eigencoefficients}
\newsym{c-0}{c_0}{Constant term in the \glsxtrshort{acr-meem} hydrodynamic output expression}
\newsym{b-N-2M-1}{b_{N+2M+1}}{Element of \gls{sym-vec-b} at index $N+2M+1$ for which the high-frequency limit must be evaluated to avoid numerical overflow}
\newsym{Y}{Y}{Generic function used to demonstrate eigenfunction orthogonality properties}

% eigenmode counting
\newsym{N}{N}{Number of eigenmodes of the innermost region in \glsxtrshort{acr-meem}}
\newsym{n}{n}{Index of eigenmode in the innermost region in \glsxtrshort{acr-meem}}
\newsym{K}{K}{Number of eigenmodes of the exterior region in \glsxtrshort{acr-meem}}
\newsym{k}{k}{Index of eigenmode in the exterior region in \glsxtrshort{acr-meem}}
\newsym{M}{M}{Number of eigenmodes of the interior region in \glsxtrshort{acr-meem}}
\newsym{m}{m}{Index of eigenmode in the interior region in \glsxtrshort{acr-meem}}
\newsym{k-index}{k}{Index of eigenmode in the exterior fluid region in \glsxtrshort{acr-meem}}

% velocity potential
\newsym{phi-p}{\phi_p}{Particular potential in \glsxtrshort{acr-meem}}
\newsym{phi-h}{\phi_h}{Homogeneous potential in \glsxtrshort{acr-meem}}
\newsym{phi-i2}{\phi^{i2}}{Radiation velocity potential in the interior fluid region under the float in \glsxtrshort{acr-meem}}

% eigencoefficients
\newsym{C-1n-i1}{C_{1n}^{i1}}{Eigencoefficient for the $n$th mode of the $i1$st fluid region}
\newsym{C-1m-i2}{C_{1m}^{i2}}{First eigencoefficient for the $m$th mode of the $i2$th fluid region}
\newsym{C-2m-i2}{C_{2m}^{i2}}{Second eigencoefficient for the $m$th mode of the $i2$th fluid region}
\newsym{B-k-e}{B_k^e}{Eigencoefficient for the $\gls{sym-k-index}$-th coefficient of the exterior fluid region in \glsxtrshort{acr-meem}}
\newsym{vec-C-1m-i2}{\vec{C}_{1m}^{i2}}{Vector of \gls{sym-C-1m-i2} for m=0 to m=M-1 in \glsxtrshort{acr-meem}}
\newsym{alpha-generic}{\alpha}{Generic eigencoefficient to demonstrate properties of orthogonality used in \glsxtrshort{acr-meem}}
\newsym{vec-C-1n-i1}{\vec{C}_{1n}^{i1}}{Vector of \gls{sym-C-1n-i1} for n=0 to n=N-1 in \glsxtrshort{acr-meem}}
\newsym{vec-C-2m-i2}{\vec{C}_{2m}^{i2}}{Vector of \gls{sym-C-2m-i2} for m=0 to m=M-1 in \glsxtrshort{acr-meem}}
\newsym{vec-B-k-e}{\vec{B_k^e}}{Vector of exterior-region eigen-coefficients, for eigenmodes $\gls{sym-k-index}=0$ to $\gls{sym-k-index}=K-1$, in \glsxtrshort{acr-meem}}

% vertical eigenfunctions
\newsym{Z-m-i2}{Z_{m}^{i2}}{Vertical eigenfunction for the $m$th mode of the $i2$th fluid region}
\newsym{Z-n-i1}{Z_{n}^{i1}}{Vertical eigenfunction for the $n$th mode of the $i1$st fluid region}
\newsym{Z-k-e}{Z_k^e}{Vertical eigenfunction for the $\gls{sym-k-index}$th mode of the exterior ($e$) fluid region}

\newsym{N-k}{N_k}{Normalization term for the vertical eigenfunction in the exterior region in \glsxtrshort{acr-meem}}
\newsym{e-i}{e_i}{Generic eigenfunction of index $i$}
\newsym{e-j}{e_j}{Generic eigenfunction of index $\gls{sym-j}$}
\newsym{Z-j}{Z_j}{Generic vertical eigenfunction at index $\gls{sym-j}$, standing for $\gls{sym-n}$ or $\gls{sym-m}$ depending on region, used in the derivation of the \glsxtrshort{acr-meem} output vector $\gls{sym-vec-c}$}

% vectors of vertical eigenfunctions
\newsym{vec-Z-n-i1}{\vec{Z}_n^{i1}}{Vector of vertical eigenfunctions for modes $n=0$ to $n=N-1$ for the innermost fluid region ($i1$)}
\newsym{vec-Z-m-i2}{\vec{Z}_m^{i2}}{Vector of vertical eigenfunctions for modes $m=0$ to $m=M-1$ for the interior fluid region under the float ($i2$)}
\newsym{vec-Z-k-e}{\vec{Z}_k^e}{Vector of vertical eigenfunctions for modes $\gls{sym-k-index}=0$ to $\gls{sym-k-index}=K-1$ for the exterior fluid region ($e$)}

% radial eigenfunctions
\newsym{R-1n-i1}{R_{1n}^{i1}}{Radial eigenfunction for the $n$th mode of the $i1$st fluid region}
\newsym{R-1m-i2}{R_{1m}^{i2}}{First radial eigenfunction for the $m$th mode of the $i2$th fluid region}
\newsym{R-2m-i2}{R_{2m}^{i2}}{Second radial eigenfunction for the $m$th mode of the $i2$th fluid region}
\newsym{Lambda-k-e}{\Lambda_{k}^{e}}{Radial eigenfunction for the $\gls{sym-k-index}$th mode of the exterior fluid region}

% radial eigenfunction vectors
\newsym{vec-R-1m-i2}{\vec{R}_{1m}^{i2}}{Vector of the first radial eigenfunctions for modes $m=0$ to $m=M-1$ for the interior fluid region under the float ($i2$)}
\newsym{vec-R-1n-i1}{\vec{R}_{1n}^{i1}}{Vector of the first (and only) radial eigenfunctions for modes $n=0$ to $n=N-1$ for the innermost fluid region ($i1$)}
\newsym{vec-R-2m-i2}{\vec{R}_{2m}^{i2}}{Vector of the second radial eigenfunctions for modes $m=0$ to $m=M-1$ for the interior fluid region under the float ($i2$)}
\newsym{vec-Lambda-k-e}{\vec{\Lambda_k^e}}{Vector of the first (and only) radial eigenfunctions in the exterior region, for eigenmodes $\gls{sym-k-index}=0$ to $\gls{sym-k-index}=K-1$, in \glsxtrshort{acr-meem}}

% eigenvalues
\newsym{lambda-n-i1}{\lambda_{n}^{i1}}{Eigenvalue for the $n$th mode of the $i1$st fluid region}
\newsym{lambda-m-i2}{\lambda_{m}^{i2}}{Eigenvalue for the $m$th mode of the $i2$th fluid region}
\newsym{m-k}{m_k}{Eigenvalue of the $\gls{sym-k-index}$th eigenmode in the exterior region in \glsxtrshort{acr-meem}}
\newsym{lambda-m}{\lambda_m}{Eigenvalue of the $m$th eigenmode of the interior region in \glsxtrshort{acr-meem}}

% vectors of eigenvalues
\newsym{vec-lambda-m}{\vec{\lambda}_{m}^{i2}}{Vector of eigenvalues for all $m=1$ to $m=M-1$ modes of the $i2$th fluid region}
\newsym{vec-lambda-n}{\vec{\lambda}_{n}^{i1}}{Vector of eigenvalues for all $n=1$ to $n=N-1$ modes of the $i1$st fluid region}
\newsym{vec-m-k}{\vec{m}_{k}}{Vector of eigenvalues for all $\gls{sym-k-index}=1$ to $\gls{sym-k-index}=K-1$ modes of the exterior ($e$) fluid region}
\newsym{vec-lambda-j}{\vec{\lambda}_{j}}{Vector of eigenvalues for modes of the $j$th fluid region}

% meem numerics
\newsym{dDelta-z-f}{\Delta z_f}{Value of \gls{sym-dDelta-z-min} for the float}
\newsym{dDelta-z-min}{\Delta z_{\text{min}}}{Minimum fluid region height in \glsxtrshort{acr-meem}, below which the eigenfunctions experience numeric overflow unless exponentially scaled}
\newsym{dDelta-z-s}{\Delta z_s}{Value of \gls{sym-dDelta-z-min} for the spar}
\newsym{chi-max}{\chi_{max}}{Maximum admissible value of $\chi$ before numeric overflow occurs}
\newsym{chi-bessel}{\chi}{Generic Bessel function argument for \glsxtrshort{acr-meem} numeric overflow discussion}
\newsym{b-max}{b_\mathrm{max}}{}

% meem integrals
\newsym{boldsymbol-mathcal-R}{\boldsymbol{\mathcal{R}}}{Integrals of radial eigenfunctions used to calculate the \gls{sym-vec-c} output vector in \glsxtrshort{acr-meem}}
\newsym{boldsymbol-mathcal-Z}{\boldsymbol{\mathcal{Z}}}{Coupling integrals of vertical eigenfunctions used to calculate the \gls{sym-mathbf-A} matrix in \glsxtrshort{acr-meem}}
\newsym{boldsymbol-mathcal-Z-m-k-0}{\boldsymbol{\mathcal{Z}}_{m,k=0}}{}
\newsym{f-k}{f_k}{Helper term for computing the scale factor \gls{sym-N-k} in \glsxtrshort{acr-meem}}
\newsym{f-m}{f_{m}}{Helper term for computing the $m$th term in the high-frequency limit of vertical coupling integral \gls{sym-boldsymbol-mathcal-Z} in \glsxtrshort{acr-meem}}
\newsym{f-2m}{f_{2m}}{Helper term for computing the integral of radial eigenfunction $\gls{sym-boldsymbol-mathcal-R}_{2\gls{sym-m}}$ for $m=0$ in \glsxtrshort{acr-meem}}

%%%%%%%%%%%%%%%%%%%%%%%%%%%%%%%%%%%%%%%%%% Power/efficiency %%%%%%%%%%%%%%%%%%%%%%%%%%%%%%%%%%%%%%%%%%
% efficiencies
\newsym{eta}{\eta}{Static efficiency from generator electrical port to grid}
\newsym{eta-array}{\eta_{array}}{Static efficiency due to array uptime}
\newsym{boldsymbol-eta-H-T-parent}{\boldsymbol{\eta}^{H,T}_{(*)}}{Matrix of efficiency type (*) in each sea state, (*) $\in \{$d,elec,lim,rad$\}$ = $\{$drag, electrical, limits, radiation$\}$}
    \newsym{boldsymbol-eta-H-T}{\boldsymbol{\eta}^{H,T}}{Matrix of efficiencies in each sea state}[boldsymbol-eta-H-T-parent]
    \newsym{boldsymbol-eta-H-T-drag}{\boldsymbol{\eta}^{H,T}_\mathrm{d}}{Matrix of drag efficiency in each sea state}[boldsymbol-eta-H-T-parent]
    \newsym{boldsymbol-eta-H-T-drag-var}{\boldsymbol{\eta}^{H,T}_{\mathrm{d},\mathrm{var}}}{}[boldsymbol-eta-H-T-parent]
    \newsym{boldsymbol-eta-H-T-elec}{\boldsymbol{\eta}^{H,T}_{\mathrm{elec}}}{Matrix of electrical efficiency in each sea state}[boldsymbol-eta-H-T-parent]
    \newsym{boldsymbol-eta-H-T-limits}{\boldsymbol{\eta}^{H,T}_\mathrm{lim}}{Matrix of efficiency due to \glsxtrshort{acr-pto} limits in each sea state}[boldsymbol-eta-H-T-parent]
    \newsym{boldsymbol-eta-H-T-rad}{\boldsymbol{\eta}^{H,T}_{\mathrm{rad}}}{Matrix of radiation efficiency in each sea state}[boldsymbol-eta-H-T-parent]

% active, reactive, complex powers
\newsym{PQS-parent}{(P, Q, S)_{(\mathrm{port})}}{(Active, reactive, complex) power at any port}
    \newsym{P}{P}{Active power}[PQS-parent]
    \newsym{P-tau-Omega}{P_{\tau\Omega}}{Active power at the generator mechanical port}[PQS-parent]
    \newsym{P-VI}{P_{VI}}{Active power at the generator electrical port}[PQS-parent]
    \newsym{S}{S}{Complex power}[PQS-parent]
    \newsym{Q}{Q}{Reactive power}[PQS-parent]
\newsym{S-0}{S_0}{Impedance-matched apparent power in the extension of \glsxtrshort{acr-qcqp} to fourth-order apparent power constraints}

% instantaneous powers
\newsym{p-min}{p_{\mathrm{min}}}{Minimum power over a wave period at any port in a given sea state}

\newsym{p-avg-pk-parent}{p_{\mathrm{(*),(port)}}}{Power of type (*) over a wave period at any port in a given sea state, (*) $\in \{$avg, pk$\} = \{$average, peak$\}$}
    \newsym{p-avg}{p_{\mathrm{avg}}}{Average power at any port over a wave period in a given sea state}[p-avg-pk-parent]
    \newsym{p-avg-elec}{p_{\mathrm{avg},\mathrm{elec}}}{Average electrical power generation to the grid over a wave period in a given sea state, after accounting for static conversion efficiency}[p-avg-pk-parent]
    \newsym{p-avg-VI}{p_{\mathrm{avg},VI}}{Average power at the generator electrical port in a given sea state}[p-avg-pk-parent]
    \newsym{p-avg-VI-opt}{p_{\mathrm{avg},VI,\mathrm{opt}}}{Average power at the generator electrical port under optimal constrained control in a given sea state}[p-avg-pk-parent]
    \newsym{p-pk}{p_{\mathrm{pk}}}{Peak power over a wave period at any port in a given sea state}[p-avg-pk-parent]
    \newsym{p-pk-elec}{p_{\mathrm{pk},\mathrm{elec}}}{Peak power over a wave period at the grid in a given sea state}[p-avg-pk-parent]
    \newsym{p-pk-V-I-beta}{p_{\mathrm{pk},VI,\beta}}{Peak power over a wave period at the generator electrical port in the \gls{sym-beta}-th sea state}[p-avg-pk-parent]

% powers by sea state
\newsym{mathbf-P-H-T-elec}{\mathbf{P}^{H,T}_{\mathrm{elec}}}{Matrix of active electrical powers in each sea state}
\newsym{mathbf-P-H-T-F-dot-X}{\mathbf{P}^{H,T}_{F\dot{X}}}{Matrix of active powers at the \glsxtrshort{acr-pto} mechanical port in each sea state}
\newsym{mathbf-P-H-T-F-dot-X-ideal}{\mathbf{P}^{H,T}_{F\dot{X},\mathrm{ideal}}}{Matrix of active powers at the \glsxtrshort{acr-pto} mechanical port in each sea state under idealized (impedance-matched, drag-free) conditions}
\newsym{mathbf-P-H-T-IV}{\mathbf{P}^{H,T}_{IV}}{Matrix of active powers at the generator electrical port in each sea state}

%%%%%%%%%%%%%%%%%%%%%%%%%%%%%%%%%%%%%%%%%% Effort/flow %%%%%%%%%%%%%%%%%%%%%%%%%%%%%%%%%%%%%%%%%%
% generic flow phasor
\newsym{q-parent}{q, (\hat{q}, q(t))}{Flow variable at any port, in (phasor, time domain) form}
    \newsym{hat-q}{\hat{q}}{Flow phasor at any port}[q-parent]
    \newsym{q}{q}{Port flow variable}[q-parent]
    \newsym{q-t}{q(t)}{Time-domain port flow variable}[q-parent]
\newsym{vec-q}{\vec{q}}{Vector of flow phasors at any set of ports}

% generic effort phasor
\newsym{e-parent}{e, (\hat{e}, e(t))}{Effort variable at any port, in (phasor, time domain) form}
    \newsym{hat-e}{\hat{e}}{Complex phasor of a generic effort port variable}[e-parent]
    \newsym{e}{e}{Port effort variable}[e-parent]
    \newsym{e-t}{e(t)}{Time-domain port effort variable}[e-parent]
    \newsym{e-1}{e_1}{Effort at port 1}[e-parent]
\newsym{vec-e}{\vec{e}}{Vector of effort phasors at any set of ports}

%%%%%%%%%%%%%%%%%%%%%%%%%%%%%%%%%%%%%%%%%% DRAG INTEGRAL %%%%%%%%%%%%%%%%%%%%%%%%%%%%%%%%%%%%%%%%%%
\newsym{vec-hat-p-d}{\vec{\hat{p}}_d}{Vector of complex phasor of drag pressure on each hydrodynamic \glsxtrshort{acr-dof}}
\newsym{vec-p-d}{\vec{p}_d(\gls{sym-y},\gls{sym-t})}{Vector of time-domain pressure due to drag for each hydrodynamic \glsxtrshort{acr-dof}}
\newsym{vec-R}{\vec{R}}{Vector of outer radii of each body, used for drag integral}
\newsym{vec-r}{\vec{r}}{Amplitude ratio of the incident wave velocity and \glsxtrshort{acr-wec} velocity, for drag integral calculation}
\newsym{vec-R-in}{\vec{R}_{\mathrm{in}}}{Vector of inner radii of each body, used for drag integral}
\newsym{w-y}{w(y)}{Local strip width for a body, used for drag integral}
\newsym{vec-w-y}{\vec{w}(y)}{Vector of local strip widths for each body, used for drag integral}
\newsym{I-prime}{I'}{Alternative drag integral to avoid singularities}
\newsym{I-drag}{I}{Drag integral \gls{sym-I-G} or \gls{sym-I-B}}
\newsym{I-B}{I_B}{Real-valued nondimensional integral for drag damping}
\newsym{I-G}{I_G}{Complex-valued nondimensional integral for drag excitation}
\newsym{kappa}{\kappa}{Nondimensional wavelength for drag integral}

\newsym{C-d-parent}{C_{d,(\mathrm{body})}}{Drag coefficient: ratio between the pressure induced by drag and the dynamic pressure of the relative flow, for (body) $\in \{f,s\}=\{$float, spar$\}$}
    \newsym{C-d}{C_d}{Drag coefficient: ratio between the pressure induced by drag and the dynamic pressure of the relative flow}[C-d-parent]
    \newsym{C-d-f}{C_{d,f}}{Drag coefficient for float}[C-d-parent]
    \newsym{C-d-s}{C_{d,s}}{Drag coefficient for spar}[C-d-parent]
\newsym{vec-C-d}{\vec{C}_d}{Vector of drag coefficients for each hydrodynamic \glsxtrshort{acr-dof}}

%%%%%%%%%%%%%%%%%%%%%%%%%%%%%%%%%%%%%%%%%% FORCE %%%%%%%%%%%%%%%%%%%%%%%%%%%%%%%%%%%%%%%%%%
% force: general
\newsym{F-surge}{F_{\mathrm{surge}}}{Surge force acting on the hull of a body}
\newsym{F-p-op}{F_{p,\mathrm{op}}}{\glsxtrshort{acr-pto} force in operational waves}

% force: scalar time domain
\newsym{F-d-t}{F_d(t)}{Time-domain drag force on a generic hydrodynamic \glsxtrshort{acr-dof}}
\newsym{F-PTO}{F_{\mathrm{PTO}}}{Vector of time domain \glsxtrshort{acr-pto} forces for each \glsxtrshort{acr-pto} \glsxtrshort{acr-dof}}

% force: magnitudes (scalar and vector)
\newsym{F-heave}{F_\mathrm{heave}}{Total peak structural force in heave acting on the hull of a body}
\newsym{F-heave-f}{F_{\mathrm{heave},f}}{Total peak structural force in heave acting on the hull of the float}
\newsym{F-heave-s}{F_{\mathrm{heave},s}}{Total peak structural force in heave acting on the hull of the spar}

% force: scalar phasors
\newsym{hat-F-th}{\hat{F}_{\mathrm{th}}}{Thevenin-equivalent force}
\newsym{hat-F-p-guess}{\hat{F}_{p,\mathrm{guess}}}{Resultant powertrain force after running a dynamics iteration for a guessed reflection coefficient when using the implicit optimal control scheme}
\newsym{hat-F-p}{\hat{F}_p}{Complex phasor force of the \glsxtrshort{acr-pto} on a single body \glsxtrshort{acr-dof}}
\newsym{hat-F-PTO}{\hat{F}_{\mathrm{PTO}}}{Complex phasor \glsxtrshort{acr-pto} force for a single \glsxtrshort{acr-pto} \glsxtrshort{acr-dof}}
\newsym{hat-F-p-beta}{\hat{F}_{p,\beta}}{Complex phasor powertrain force in the $\beta$th sea state}
\newsym{hat-F-p-constr}{\hat{F}_{p,\mathrm{constr}}}{Constrained complex phasor \glsxtrshort{acr-pto} force}
\newsym{hat-F-p-unconstrained}{\hat{F}_{p,\mathrm{unconstrained}}}{Unconstrained complex phasor \glsxtrshort{acr-pto} force}
\newsym{hat-F-p-st}{\hat{F}_{p,\mathrm{st}}}{Complex phasor of \glsxtrshort{acr-pto} force on the floating body in the storm loadcase}

% force: vector/matrix
\newsym{vec-F-parent}{\vec{\hat{F}}_{(*)},\vec{F}_{(*)}(t)}{Vector of forces of type (*) on each hydrodynamic \glsxtrshort{acr-dof}, (*) $\in \{$d, e, p, rad, res$\} = \{$drag, excitation, powertrain, hydrodynamic radiation, hydrostatic restoring$\}$, in (complex phasor, time-domain) form}
    % force: vector phasors
    \newsym{vec-hat-F}{\vec{\hat{F}}}{Vector of complex phasor first-order hydrodynamic forces on each hydrodynamic \glsxtrshort{acr-dof} for interacting floating bodies, including wave excitation $\gls{sym-vec-hat-F-e}$, hydrodynamic radiation $\gls{sym-vec-hat-F-rad}$, and hydrostatic restoring $\gls{sym-vec-hat-F-res}$ contributions}[vec-F-parent]
    \newsym{vec-hat-F-p}{\vec{\hat{F}}_p}{Vector of complex phasor powertrain forces for each hydrodynamic \glsxtrshort{acr-dof}}[vec-F-parent]
    \newsym{vec-hat-F-d}{\vec{\hat{F}}_d}{Vector of complex phasor drag force on each hydrodynamic \glsxtrshort{acr-dof}}[vec-F-parent]
    \newsym{vec-hat-F-e}{\vec{\hat{F}}_e}{Vector of complex phasor excitation forces for each hydrodynamic \glsxtrshort{acr-dof}}[vec-F-parent]
    \newsym{vec-hat-F-rad}{\vec{\hat{F}}_{\mathrm{rad}}}{Hydrodynamic radiation force phasor vector}[vec-F-parent]
    \newsym{vec-hat-F-res}{\vec{\hat{F}}_{\mathrm{res}}}{Hydrostatic restoring force phasor}[vec-F-parent]
    % force: vector time domain
    \newsym{vec-F-d}{\vec{F}_d(t)}{Vector of time-domain drag force on each hydrodynamic \glsxtrshort{acr-dof}}[vec-F-parent]
    \newsym{vec-F-e}{\vec{F}_e(t)}{Vector of time-domain excitation forces for each hydrodynamic \glsxtrshort{acr-dof}}[vec-F-parent]
    \newsym{vec-F-p}{\vec{F}_p}{Vector of time-domain powertrain forces for each hydrodynamic \glsxtrshort{acr-dof}}[vec-F-parent]
    \newsym{vec-F-rad}{\vec{F}_{\mathrm{rad}}(t)}{Time-domain vector of radiation forces for each hydrodynamic \glsxtrshort{acr-dof}}[vec-F-parent]
    \newsym{vec-F-res}{\vec{F}_{\mathrm{res}}(t)}{Time-domain vector of hydrostatic restoring forces for each hydrodynamic \glsxtrshort{acr-dof}}[vec-F-parent]


\newsym{vec-hat-F-th}{\vec{\hat{F}}_{\mathrm{th}}}{Complex phasor Thevenin-equivalent force}
\newsym{mathbf-hat-F-H-T-th}{\mathbf{\hat{F}}^{H,T}_{\mathrm{th}}}{Matrix of complex phasor Thevenin-equivalent forces in each sea state}
\newsym{vec-hat-F-PTO}{\vec{\hat{F}}_{\mathrm{PTO}}}{Vector of complex phasor \glsxtrshort{acr-pto} forces for each \glsxtrshort{acr-pto} \glsxtrshort{acr-dof}}
%%%%%%%%%%%%%%%%%%%%%%%%%%%%%%%%%%%%%%%%%% BODY MOTION %%%%%%%%%%%%%%%%%%%%%%%%%%%%%%%%%%%%%%%%%%

% xi: scalar time-domain
\newsym{ddot-xi-s}{\ddot{\xi}_s}{Time domain vertical acceleration of the spar}
\newsym{dot-xi}{\dot{\xi}}{Time domain vertical velocity of a body}
\newsym{ddot-xi}{\ddot{\xi}}{Time domain vertical acceleration of a body}
\newsym{xi}{\xi}{Time domain vertical position of a body}
\newsym{xi-t}{\xi(t)}{Time-domain vertical body position} % dup
\newsym{dot-xi-Delta}{\dot{\xi}_\Delta}{Time domain change in vertical velocity of a body induced by the presence of the \glsxtrshort{acr-pto}}
\newsym{ddot-xi-f}{\ddot{\xi}_f}{Time domain vertical acceleration of the float}
\newsym{xi-Delta}{\xi_\Delta}{Time domain change in vertical position of a body induced by the presence of the \glsxtrshort{acr-pto}}

% xi: vector time-domain
\newsym{vec-dot-xi-0}{\vec{\dot{\xi}}_0}{Vector of complex phasor velocities of each body \glsxtrshort{acr-dof} due to wave excitation in the absence of power generation}
\newsym{vec-dot-xi}{\vec{\dot{\xi}}}{Vector of time-domain velocities of each body \glsxtrshort{acr-dof}}
\newsym{vec-dot-xi-Delta}{\vec{\dot{\xi}}_\Delta}{Vector of the change in oscillation velocities of each hydrodynamic \glsxtrshort{acr-dof} due to the presence of the \glsxtrshort{acr-pto}}

% xi: scalar phasors
\newsym{hat-xi-parent}{(\hat{\xi},\hat{\dot{\xi}}, \hat{\ddot{\xi}})_{(\mathrm{body})}}{Complex phasor (displacement, velocity, acceleration) of (body) $\in \{f,s,m\} = \{$float, spar, merged$\}$}
    \newsym{hat-dot-xi}{\hat{\dot{\xi}}}{Complex phasor velocity of a body}[hat-xi-parent]
    \newsym{hat-xi}{\hat{\xi}}{Complex phasor displacement of a body}[hat-xi-parent]
    \newsym{hat-xi-f}{\hat{\xi}_f}{Complex phasor displacement of the float}[hat-xi-parent]
    \newsym{hat-xi-f-beta}{\hat{\xi}_{f,\beta}}{Complex phasor displacement of the float for sea state \(\beta\)}[hat-xi-parent]
    \newsym{hat-xi-m}{\hat{\xi}_m}{Complex phasor displacement of the merged float-spar body}[hat-xi-parent]
    \newsym{hat-xi-s}{\hat{\xi}_s}{Complex phasor displacement of the spar}[hat-xi-parent]
    \newsym{hat-xi-s-beta}{\hat{\xi}_{s,\beta}}{Complex phasor displacement of the spar for sea state \(\beta\)}[hat-xi-parent]
\newsym{hat-xi-f-0}{\hat{\xi}_{f,0}}{Complex phasor displacement of the float due to wave excitation in the absence of power generation}
\newsym{hat-xi-s-0}{\hat{\xi}_{s,0}}{Complex phasor displacement of the spar due to wave excitation in the absence of power generation}

% xi: vectors/matrices of phasors
\newsym{vec-hat-xi-parent}{(\vec{\hat{\xi}},\vec{\hat{\dot{\xi}}}, \vec{\hat{\ddot{\xi}}})_{(\mathrm{lc})}}{Vector of complex phasor (displacements, velocities, accelerations) of each body \glsxtrshort{acr-dof} in a given loadcase (lc) $\in \{\mathrm{op}, \mathrm{st}\} = \{$operational, storm$\}$}
    \newsym{vec-hat-xi-st}{\vec{\hat{\xi}}_{\mathrm{st}}}{Vector of complex phasor displacements of each body \glsxtrshort{acr-dof} in the storm loadcase}[vec-hat-xi-parent]
    \newsym{vec-hat-ddot-xi}{\vec{\hat{\ddot{\xi}}}}{Vector of complex phasor accelerations of each body \glsxtrshort{acr-dof}}[vec-hat-xi-parent]
    \newsym{vec-hat-dot-xi}{\vec{\hat{\dot{\xi}}}}{Vector of complex phasor velocities of each body \glsxtrshort{acr-dof}}[vec-hat-xi-parent]
    \newsym{vec-hat-xi}{\vec{\hat{\xi}}}{Vector of complex phasor displacements of each body \glsxtrshort{acr-dof}}[vec-hat-xi-parent]
    \newsym{vec-hat-xi-op}{\vec{\hat{\xi}}_{\mathrm{op}}}{Vector of complex phasor displacements of each body \glsxtrshort{acr-dof} in an operational sea state}[vec-hat-xi-parent]
\newsym{vec-hat-dot-xi-0}{\vec{\hat{\dot{\xi}}}_0}{Vector of complex phasor velocities of each body \glsxtrshort{acr-dof} due to wave excitation in the absence of power generation}
\newsym{boldsymbol-hat-xi-H-T-f}{\boldsymbol{\hat{\xi}}^{H,T}_{f}}{Matrix of complex phasor displacements of the float in each sea state}


%%%%%%%%%%%%%%%%%%%%%%%%%%%%%%%%%%%%%%%%%% PTO STATES %%%%%%%%%%%%%%%%%%%%%%%%%%%%%%%%%%%%%%%%%%
% phasors - X PTO
\newsym{hat-X-PTO}{\hat{X}_{\mathrm{PTO}}}{Complex phasor displacement of the \glsxtrshort{acr-pto}. Shortened to \gls{sym-hat-X} in subscripts for readability.}
    \newsym{hat-X}{\hat{X}}{Compact notation for \gls{sym-hat-X-PTO} used in cascade matrix subscripts for readability}[hat-X-PTO]
\newsym{dot-X-PTO}{\dot{X}_{\mathrm{PTO}}}{\glsxtrshort{acr-pto} relative velocity}
\newsym{dot-hat-X-PTO}{\dot{\hat{X}}_{\mathrm{PTO}}}{Complex phasor of \glsxtrshort{acr-pto} relative velocity}
\newsym{vec-hat-X-PTO}{\vec{\hat{X}}_{\mathrm{PTO}}}{Vector of \glsxtrshort{acr-pto} displacement phasors}
\newsym{vec-hat-dot-X-PTO}{\vec{\hat{\dot{X}}}_{\mathrm{PTO}}}{Vector of \glsxtrshort{acr-pto} velocity phasors}

\newsym{dot-X}{\dot{X}}{Shorthand for \gls{sym-dot-X-PTO} used in cascade matrix subscripts for readability}
\newsym{vec-X}{\vec{X}}{}


% phasors - Omega
\newsym{Omega}{\Omega}{Generator rotational mechanical speed}
\newsym{hat-Omega}{\hat{\Omega}}{Generator rotational mechanical speed phasor}

% phasors - torque
\newsym{hat-tau}{\hat{\tau}}{Complex phasor of torque at the generator mechanical port}
\newsym{tau}{\tau}{Generator torque}

% phasors - voltage
\newsym{hat-V-s-th}{\hat{V}_{s,\mathrm{th}}}{Thevenin-equivalent source voltage phasor as viewed from the generator electrical port}
    \newsym{hat-V-s-th-beta}{\hat{V}_{s,\mathrm{th},\beta}}{Thevenin-equivalent source voltage phasor as viewed from the generator electrical port in the $\beta$th sea state}[hat-V-s-th]
\newsym{hat-V}{\hat{V}}{Generator quadrature voltage phasor}
    \newsym{V}{V}{Generator quadrature voltage}[hat-V]

% phasors - current
\newsym{hat-I}{\hat{I}}{Generator quadrature current phasor}
\newsym{I}{I}{Generator quadrature current}
\newsym{hat-I-p}{\hat{I}_p}{Complex phasor generator current in the particular solution of the system dynamics (chosen here to be the unconstrained impedance-matched solution)}

%%%%%%%%%%%%%%%%%%%%%%%%%%%%%%%%%%%%%%%%%% WAVE STATES %%%%%%%%%%%%%%%%%%%%%%%%%%%%%%%%%%%%%%%%%%

% relative vel
\newsym{vec-v-rel-parent}{\gls{sym-vec-v-rel}(t),\gls{sym-vec-hat-v-rel}}{Vector of wave-\glsxtrshort{acr-wec} relative velocities for each body, in (time domain, complex phasor) form}
\newsym{vec-v-rel}{\vec{v}_{\mathrm{rel}}}{Relative velocity vector between body and fluid for drag calculations}[vec-v-rel-parent]
\newsym{vec-hat-v-rel}{\vec{\hat{v}}_{\mathrm{rel}}}{Vector of the complex phasors of the relative velocity between the incident wave and the \glsxtrshort{acr-wec} for each body}[vec-v-rel-parent]

% wave vel
\newsym{vec-v-wave}{\vec{v}_{\mathrm{wave}}}{Time-domain vector of wave velocity for each body as a function of position $y$, used in drag integral}
\newsym{hat-V-wave}{\hat{V}_{\mathrm{wave}}}{Wave velocity phasor}
\newsym{vec-hat-V-wave}{\vec{\hat{V}}_{\mathrm{wave}}}{Wave-velocity phasor vector}

% fluid vel
\newsym{vec-v}{\vec{v}}{Fluid velocity vector field in \glsxtrshort{acr-meem}}
\newsym{vec-alpha-v-y}{\vec{\alpha}_v(y)}{Vector of amplitude ratios of the relative \glsxtrshort{acr-wec}-fluid velocity to the \glsxtrshort{acr-wec} velocity as a function of lateral coordinate \gls{sym-y}, used in drag integral calculations}

% wave elevation zeta
\newsym{hat-zeta}{\hat{\zeta}}{Wave elevation phasor}
\newsym{zeta}{\zeta(t)}{Time-domain wave elevation}
\newsym{zeta-y-t}{\zeta(y,t)}{Wave elevation at horizontal coordinate y and time t}
\newsym{tilde-zeta}{\tilde{\zeta}}{Complex time-domain wave elevation}

%%%%%%%%%%%%%%%%%%%%%%%%%%%%%%%%%%%%%%%%%% LIMITS %%%%%%%%%%%%%%%%%%%%%%%%%%%%%%%%%%%%%%%%%%
\newsym{A-A-max-parent}{(*)_\mathrm{max}}{Maximum permissible value of any quantity (*)}

% power limits
    \newsym{S-max-k}{S_{\mathrm{max},k}}{Maximum permissible apparent power in of the $\gls{sym-k-generic-index}$th apparent power constraint}[A-A-max-parent]
    \newsym{P-max-elec}{P_{\mathrm{max},\mathrm{elec}}}{Maximum allowable power at the grid electrical port across all sea states}[A-A-max-parent]
    \newsym{P-tau-Omega-max}{P_{\tau\Omega,\mathrm{max}}}{Maximum allowable mechanical port power at the generator shaft}[A-A-max-parent]

% torque limits
    \newsym{tau-max}{\tau_\mathrm{max}}{Maximum permissible generator torque}[A-A-max-parent]

% xi limits
    \newsym{xi-s-max}{\xi_{s,\mathrm{max}}}{Maximum permissible spar displacement amplitude}[A-A-max-parent]
    \newsym{xi-f-max}{\xi_{f,\mathrm{max}}}{Maximum permissible float displacement amplitude}[A-A-max-parent]
\newsym{xi-max-slam}{\xi_{\mathrm{max,slam}}}{Maximum permitted position from still water line, above which slamming occurs}
\newsym{xi-min-slam}{\xi_{\mathrm{min,slam}}}{Minimum permitted position from still water line, below which slamming occurs}
\newsym{xi-f-lin}{\xi_{f,\mathrm{lin}}}{Maximum permissible float displacement amplitude to maintain the assumptions of linear wave theory}
\newsym{xi-f-slam}{\xi_{f,\mathrm{slam}}}{Maximum permissible float displacement amplitude to prevent slamming}
\newsym{xi-s-lin}{\xi_{s,\mathrm{lin}}}{Maximum permissible spar displacement amplitude to maintain the assumptions of linear wave theory}
\newsym{xi-s-slam}{\xi_{s,\mathrm{slam}}}{Maximum permissible spar displacement amplitude to prevent slamming}

% z limits
\newsym{dDelta-z-slam}{\Delta z_{\mathrm{slam}}}{Distance between still water level and body bottom used in slamming limits}
\newsym{dDelta-z-sub}{\Delta z_{\mathrm{sub}}}{Distance between still water level and body top used in submersion limits}
\newsym{dDelta-z-surf}{\Delta z_{\mathrm{surf}}}{Relevant vertical body dimension used in surface-piercing (slamming or submersion) limits}

% X PTO limits
\newsym{X-PTO-max}{X_{\mathrm{PTO},\mathrm{max}}}{Maximum allowable \glsxtrshort{acr-pto} relative displacement}[A-A-max-parent]

%%%%%%%%%%%%%%%%%%%%%%%%%%%%%%%%%%%%%%%%%% QCQP %%%%%%%%%%%%%%%%%%%%%%%%%%%%%%%%%%%%%%%%%%

\newsym{vec-n}{\vec{n}}{Null space vector of dynamics}
\newsym{vec-x-opt}{\vec{x}_{\mathrm{opt}}}{Vector of dynamics state phasors under optimal control for a given sea state}
\newsym{vec-x-p}{\vec{x}_p}{Vector of particular solution for state phasors for a given sea state, chosen as the unconstrained impedance-matched solution}

\newsym{mu-nu-parent}{\mu,\nu}{Index of constraint in \glsxtrshort{acr-qcqp}}
    \newsym{mu}{\mu}{Index of constraint in \glsxtrshort{acr-qcqp}}[mu-nu-parent]
    \newsym{nu}{\nu}{Index of constraint in \glsxtrshort{acr-qcqp}}[mu-nu-parent]

\newsym{g-mu}{g_\mu}{Constraint function for the $\mu$th dynamic inequality constraint in \glsxtrshort{acr-qcqp}}
\newsym{lambda-mu}{\lambda_\mu}{Lagrange multiplier of the $\mu$th constraint in \glsxtrshort{acr-qcqp}}

% reflection coeff
\newsym{Gamma-int-mu-star-nu-star}{\Gamma_{\mathrm{int},\mu^*,\nu^*}}{Complex optimal reflection coefficient at the intersection of two active constraint circles in \glsxtrshort{acr-qcqp}}
\newsym{Gamma-mu-nu}{\Gamma_{\mu,\nu}}{Complex reflection coefficient at the intersection of the $\mu$th and $\nu$th constraint circles in \glsxtrshort{acr-qcqp}}
\newsym{Gamma-opt}{\Gamma_\mathrm{opt}}{Complex optimal reflection coefficient}
    \newsym{Gamma-beta-opt}{\Gamma_{\beta,\mathrm{opt}}}{Optimal reflection coefficient in the $\beta$th sea state}[Gamma-opt]
\newsym{Gamma}{\Gamma}{Complex reflection coefficient at the generator electrical port}
\newsym{Gamma-cl}{\Gamma_{\mathrm{cl}}}{Reflection coefficient at the closest point of a constraint circle in \glsxtrshort{acr-qcqp}}
\newsym{Gamma-int}{\Gamma_{\mathrm{int}}}{Reflection coefficient at the intersection of constraints in \glsxtrshort{acr-qcqp}}
\newsym{Gamma-cl-mu}{\Gamma_{\mathrm{cl},\mu}}{Complex reflection coefficient at the closest point of the $\mu$th constraint circle in \glsxtrshort{acr-qcqp}}
\newsym{Gamma-cl-mu-star}{\Gamma_{\mathrm{cl},\mu^*}}{Complex optimal reflection coefficient at the closest point of a single active constraint circle in \glsxtrshort{acr-qcqp}}

% circle centers and radii
\newsym{Gamma-c-mu}{\Gamma_{c,\mu}}{Complex center of the $\mu$th constraint circle in \glsxtrshort{acr-qcqp}}
\newsym{Gamma-c-mu-star}{\Gamma_{c,\mu^*}}{Complex center of the $\mu$th constraint circle in \glsxtrshort{acr-qcqp}, when $\mu$ is the only active constraint in the optimal solution}
\newsym{Gamma-c-nu}{\Gamma_{c,\nu}}{Complex center of the $\nu$th constraint circle in \glsxtrshort{acr-qcqp}}
\newsym{Gamma-c-nu-star}{\Gamma_{c,\nu^*}}{Complex center of the $\nu$th constraint circle in \glsxtrshort{acr-qcqp}, when $\nu$ is one of multiple active constraints in the optimal solution}
\newsym{r-mu}{r_\mu}{Radius of the $\mu$-th constraint circle, in \glsxtrshort{acr-qcqp}}
\newsym{r-nu}{r_\nu}{Radius of the $\nu$-th constraint circle in \glsxtrshort{acr-qcqp}}

% a - one index
\newsym{vec-a-mu}{\vec{a}_\mu}{Vector of linear coefficients for the $\mu$th constraint in \glsxtrshort{acr-qcqp}, in multi-\glsxtrshort{acr-dof} form}

% a - two indices
\newsym{a-mu-nu}{a_{\mu\nu}}{Distance from the center of the $\gls{sym-mu}$th circle to the midpoint of the two intersection points of the $\mu$th and $\nu$th circles in \glsxtrshort{acr-qcqp} after collapsing to single-\glsxtrshort{acr-dof} form}
\newsym{a-mu-star-nu-star}{a_{\mu^*\nu^*}}{Distance from the center of the $\gls{sym-mu}^*$th circle to the midpoint of the two intersection points of the $\mu^*$th and $\nu^*$th circles in \glsxtrshort{acr-qcqp} after collapsing to single-\glsxtrshort{acr-dof} form, where $\mu^*$ and $\nu^*$ are the two active constraints in the optimal solution}

% b
\newsym{b-mu-parent}{\gls{sym-b-mu},\gls{sym-b-nu}}{Scalar right-hand-side term for the $\mu$th or $\nu$th constraint limit in \glsxtrshort{acr-qcqp} multi-\glsxtrshort{acr-dof} form.}
    \newsym{b-mu}{b_\mu}{Scalar right-hand-side term for the $\mu$th constraint limit in \glsxtrshort{acr-qcqp} multi-\glsxtrshort{acr-dof} form. Also appears in \glsxtrshort{acr-qcqp} sensitivity analysis.}[b-mu-parent]
    \newsym{b-nu}{b_\nu}{See \gls{sym-b-mu} for the $\nu$th rather than $\mu$th constraint}[b-mu-parent]

\newsym{b-mu-star}{b_{\mu^*}}{\gls{sym-b-mu} for the active constraint $\mu^*$ in the \glsxtrshort{acr-qcqp} optimal solution}
\newsym{b-nu-star}{b_{\nu^*}}{\gls{sym-b-nu} for the active constraint $\nu^*$ in the \glsxtrshort{acr-qcqp} optimal solution}

% c
\newsym{c-mu}{c_\mu}{Helper term for the $\mu$th constraint circle independent of constraint limits used in the numerator of the \glsxtrshort{acr-qcqp} sensitivity analysis}
\newsym{c-mu-star}{c_{\mu^*}}{\gls{sym-c-mu} for the active constraint $\mu^*$ in the \glsxtrshort{acr-qcqp} optimal solution}

% d - two indices
\newsym{d-mu-nu}{d_{\mu\nu}}{Distance between the centers of the $\mu$th and $\nu$th constraint circles in \glsxtrshort{acr-qcqp}}
\newsym{d-mu-star-nu-star}{d_{\mu^*\nu^*}}{\gls{sym-d-mu-nu} for the two active constraints $\mu^*$ and $\nu^*$ in the \glsxtrshort{acr-qcqp} optimal solution}

% d - one index
\newsym{d-mu}{d_\mu}{Helper term for the $\mu$th constraint circle. It is independent of constraint limits and used in the denominator of the \glsxtrshort{acr-qcqp} sensitivity analysis}
\newsym{d-mu-star}{d_{\mu^*}}{\gls{sym-d-mu} for the active constraint $\mu^*$ in the \glsxtrshort{acr-qcqp} optimal solution}

% Capital letters
\newsym{mathbf-Q-mu}{\mathbf{Q}_\mu}{Quadratic coefficient for $\mu$th constraint in multi-\glsxtrshort{acr-dof} form}
\newsym{A-q-mu}{A_{q,\mu}}{Quadratic coefficient of the $\mu$th constraint in \glsxtrshort{acr-qcqp}, after collapsing to single-\glsxtrshort{acr-dof} form}
\newsym{B-q-mu}{B_{q,\mu}}{Linear coefficient of the $\mu$th constraint in \glsxtrshort{acr-qcqp}, after collapsing to single-\glsxtrshort{acr-dof} form}
\newsym{C-q-mu}{C_{q,\mu}}{Constant term in the $\mu$th constraint in \glsxtrshort{acr-qcqp}, after collapsing to single-\glsxtrshort{acr-dof} form}
\newsym{S-mu}{S_\mu}{Sign indicator for the $\mu$th constraint in \glsxtrshort{acr-qcqp}, used to determine which side of the constraint circle is feasible after collapsing to single-\glsxtrshort{acr-dof} form}
\newsym{L-mu-beta}{L_{\mu\beta}}{Constraint limit for the $\mu$th constraint in the $\beta$th sea state, for \glsxtrshort{acr-qcqp} sensitivity}
\newsym{L-mu}{L_\mu}{Constraint limit for the $\mu$th constraint, for \glsxtrshort{acr-qcqp} sensitivity}

% feasibility
\newsym{mathcal-F}{\mathcal{F}}{Set of feasible points for analytical \glsxtrshort{acr-qcqp}}
\newsym{mathbb-I-mu}{\mathbb{I}_\mu}{Indicator function for feasibility of a candidate reflection coefficient for the $\mu$th constraint in \glsxtrshort{acr-qcqp}}
\newsym{f-penalty}{f_\mathrm{pen}}{Penalty function to ensure feasibility of a candidate reflection coefficient in \glsxtrshort{acr-qcqp}}

% candidate points
\newsym{mathcal-C}{\mathcal{C}}{Set of candidate points for analytical \glsxtrshort{acr-qcqp}}
\newsym{mathcal-C-cl}{\mathcal{C}_{\mathrm{cl}}}{Set of candidate points for analytical \glsxtrshort{acr-qcqp} representing the closest point to the origin of each circle}
\newsym{mathcal-C-cl-oval}{\mathcal{C}_{\mathrm{cl},\mathrm{oval}}}{Set of candidate points for analytical \glsxtrshort{acr-qcqp} representing the closest point to the origin of each oval}
\newsym{mathcal-C-int}{\mathcal{C}_{\mathrm{int}}}{Set of candidate points for analytical \glsxtrshort{acr-qcqp} representing the intersections of each circle}
\newsym{mathcal-C-int-oval-circ}{\mathcal{C}_{\mathrm{int},\mathrm{oval-circ}}}{Set of \glsxtrshort{acr-qcqp} candidate points at oval-circle intersections}

% qcqp weights
\newsym{pi-1-mu-nu-beta}{\varpi_{1\mu\nu\beta}}{First transformed weighting term for the $\mu,\nu$ constraint pair in sea state $\beta$}
\newsym{pi-2-mu-nu-beta}{\varpi_{2\mu\nu\beta}}{Second transformed weighting term for the $\mu,\nu$ constraint pair in sea state $\beta$}

\ifdefined\DISSERTATION
\else
    \makeglossaries
\fi

% Source - https://tex.stackexchange.com/a/515729
% Posted by user194703, modified by community. See post 'Timeline' for change history
% Retrieved 2026-06-21, License - CC BY-SA 4.0
\usetikzlibrary{shapes.geometric}
\newcommand{\Stars}[2][fill=black,draw=black]{
    \begin{tikzpicture}[baseline=-0.35em,#1]
    \foreach \X in {1,...,3}
    {\pgfmathsetmacro{\xfill}{min(1,max(1+#2-\X,0))}
    \path (\X*1.1em,0) 
    node[star,draw,star point height=0.25em,minimum size=1em,inner sep=0pt,
    path picture={\fill (path picture bounding box.south west) 
    rectangle  ([xshift=\xfill*1em]path picture bounding box.north west);}]{};
    }
    \end{tikzpicture}
}


\newcommand{\assumption}[1]{\refstepcounter{asmpnumber}\label{aor:asmp:#1}A\arabic{asmpnumber}}
\newcommand{\asmpref}[1]{A\ref{aor:asmp:#1}}

%\usepackage{graphicx}
%\usepackage{float}
%\usepackage{placeins}

%% The amsthm package provides extended theorem environments
%% \usepackage{amsthm}

%% The lineno packages adds line numbers. Start line numbering with
%% \begin{linenumbers}, end it with \end{linenumbers}. Or switch it on
%% for the whole article with \linenumbers.
%% \usepackage{lineno}

\makeatletter%
\newcommand\resultsAOR[1][all]{%
  \ifnum\pdfstrcmp{#1}{all}=0%
    \def\resultsAOR@out{%
    \{%
      "simRuntimeOldSeconds": 0.04, %
      "forceScaleFactor": "0.28", %
      "powerScaleFactor": "0.49", %
      "massScaleFactor": "1.12", %
      "postpro\_time": 404.323235, %
      "analysis\_time": 3683.159218, %
      "nominalLCOE": "0.89 \$/kWh", %
      "surgeForceFloatNominal": "10 MN", %
      "surgeForceSparNominal": "3 MN", %
      "heaveForceNominal": "9 MN", %
      "capexDesignNominal": 2190114.3223833675, %
      "powerAvgNominal": 72822.42633091439, %
      "simRuntime": "152 ms", %
      "pctRuntimeDynamics": "20\%", %
      "pctRuntimeMEEM": "72\%", %
      "pctRuntimeOther": "8\%", %
      "percent\_nans": 18.0288, %
      "n\_combos": 93750, %
      "geom\_sweep\_analysis\_time": 2402.779923, %
      "ptoSweepMinLCOE\_F\_max": "6.00 MN", %
      "ptoSweepMinLCOE\_P\_max": "400 kW", %
      "ptoSweepMinLCOE\_power": "76 kW", %
      "ptoSweepMinLCOE\_cost": "2.25 \$M", %
      "ptoSweepMinLCOE\_LCOE": "0.86 \$/kWh", %
      "err\_structs": {[}%
        \{%
          "pct\_error\_baseline": {[}%
            0.8351547681041923, %
            {-}0.6753493890717724%
          {]}, %
          "pct\_error\_drag": {[}%
            10.790884267350643, %
            5.002294575560603%
          {]}, %
          "pct\_error\_meem": {[}%
            {-}13.649714349255992, %
            19.55578164189912%
          {]}, %
          "pct\_error\_total": {[}%
            {-}12.715966046141624, %
            5.090339861457517%
          {]}%
        \}, %
        \{%
          "pct\_error\_baseline": {[}%
            {-}0.1673331619860154, %
            {-}0.09472855056901647%
          {]}, %
          "pct\_error\_drag": {[}%
            5.756780832854308, %
            {-}4.8276990600241785%
          {]}, %
          "pct\_error\_meem": {[}%
            {-}0.28101071566998165, %
            5.870308306405258%
          {]}, %
          "pct\_error\_total": {[}%
            {-}2.0486004608632875, %
            {-}6.27926538734057%
          {]}%
        \}, %
        \{%
          "pct\_error\_baseline": {[}%
            39.22645866181762, %
            {-}27.657557186404567, %
            null, %
            null%
          {]}, %
          "pct\_error\_drag": {[}%
            65.47109436867527, %
            {-}11.439039626615426, %
            null, %
            null%
          {]}, %
          "pct\_error\_meem": {[}%
            39.22645866181762, %
            {-}32.7477682461545, %
            null, %
            null%
          {]}, %
          "pct\_error\_total": {[}%
            65.47109436867527, %
            {-}10.810656387995223, %
            null, %
            null%
          {]}%
        \}%
      {]}, %
      "wecsimAvgPowerErrorBestCase": "0.17\%", %
      "wecsimAvgPowerErrorWorstCase": "12.72\%", %
      "MEEMRuntime": "110 ms", %
      "meemSweepTime": "40 minutes", %
      "meemSweepTimePerCombo": "26 ms", %
      "simRuntimeOld": "40 ms", %
      "multiplierOldRuntime": "4", %
      "meemSweepSize": "93,750"%
    \}}%
  \else%
    \ifnum\pdfstrcmp{#1}{simRuntimeOldSeconds}=0%
      \def\resultsAOR@out{%
    0.04}%
    \else%
      \ifnum\pdfstrcmp{#1}{forceScaleFactor}=0%
      \def\resultsAOR@out{%
    0.28}%
    \else%
      \ifnum\pdfstrcmp{#1}{powerScaleFactor}=0%
      \def\resultsAOR@out{%
    0.49}%
    \else%
      \ifnum\pdfstrcmp{#1}{massScaleFactor}=0%
      \def\resultsAOR@out{%
    1.12}%
    \else%
      \ifnum\pdfstrcmp{#1}{postpro_time}=0%
      \def\resultsAOR@out{%
    404.323235}%
    \else%
      \ifnum\pdfstrcmp{#1}{analysis_time}=0%
      \def\resultsAOR@out{%
    3683.159218}%
    \else%
      \ifnum\pdfstrcmp{#1}{nominalLCOE}=0%
      \def\resultsAOR@out{%
    0.89 \$/kWh}%
    \else%
      \ifnum\pdfstrcmp{#1}{surgeForceFloatNominal}=0%
      \def\resultsAOR@out{%
    10 MN}%
    \else%
      \ifnum\pdfstrcmp{#1}{surgeForceSparNominal}=0%
      \def\resultsAOR@out{%
    3 MN}%
    \else%
      \ifnum\pdfstrcmp{#1}{heaveForceNominal}=0%
      \def\resultsAOR@out{%
    9 MN}%
    \else%
      \ifnum\pdfstrcmp{#1}{capexDesignNominal}=0%
      \def\resultsAOR@out{%
    2190114.3223833675}%
    \else%
      \ifnum\pdfstrcmp{#1}{powerAvgNominal}=0%
      \def\resultsAOR@out{%
    72822.42633091439}%
    \else%
      \ifnum\pdfstrcmp{#1}{simRuntime}=0%
      \def\resultsAOR@out{%
    152 ms}%
    \else%
      \ifnum\pdfstrcmp{#1}{pctRuntimeDynamics}=0%
      \def\resultsAOR@out{%
    20\%}%
    \else%
      \ifnum\pdfstrcmp{#1}{pctRuntimeMEEM}=0%
      \def\resultsAOR@out{%
    72\%}%
    \else%
      \ifnum\pdfstrcmp{#1}{pctRuntimeOther}=0%
      \def\resultsAOR@out{%
    8\%}%
    \else%
      \ifnum\pdfstrcmp{#1}{percent_nans}=0%
      \def\resultsAOR@out{%
    18.0288}%
    \else%
      \ifnum\pdfstrcmp{#1}{n_combos}=0%
      \def\resultsAOR@out{%
    93750}%
    \else%
      \ifnum\pdfstrcmp{#1}{geom_sweep_analysis_time}=0%
      \def\resultsAOR@out{%
    2402.779923}%
    \else%
      \ifnum\pdfstrcmp{#1}{ptoSweepMinLCOE_F_max}=0%
      \def\resultsAOR@out{%
    6.00 MN}%
    \else%
      \ifnum\pdfstrcmp{#1}{ptoSweepMinLCOE_P_max}=0%
      \def\resultsAOR@out{%
    400 kW}%
    \else%
      \ifnum\pdfstrcmp{#1}{ptoSweepMinLCOE_power}=0%
      \def\resultsAOR@out{%
    76 kW}%
    \else%
      \ifnum\pdfstrcmp{#1}{ptoSweepMinLCOE_cost}=0%
      \def\resultsAOR@out{%
    2.25 \$M}%
    \else%
      \ifnum\pdfstrcmp{#1}{ptoSweepMinLCOE_LCOE}=0%
      \def\resultsAOR@out{%
    0.86 \$/kWh}%
    \else%
      \ifnum\pdfstrcmp{#1}{err_structs}=0%
      \let\resultsAOR@out\resultsAOR@I%
    \else%
      \ifnum\pdfstrcmp{#1}{wecsimAvgPowerErrorBestCase}=0%
      \def\resultsAOR@out{%
    0.17\%}%
    \else%
      \ifnum\pdfstrcmp{#1}{wecsimAvgPowerErrorWorstCase}=0%
      \def\resultsAOR@out{%
    12.72\%}%
    \else%
      \ifnum\pdfstrcmp{#1}{MEEMRuntime}=0%
      \def\resultsAOR@out{%
    110 ms}%
    \else%
      \ifnum\pdfstrcmp{#1}{meemSweepTime}=0%
      \def\resultsAOR@out{%
    40 minutes}%
    \else%
      \ifnum\pdfstrcmp{#1}{meemSweepTimePerCombo}=0%
      \def\resultsAOR@out{%
    26 ms}%
    \else%
      \ifnum\pdfstrcmp{#1}{simRuntimeOld}=0%
      \def\resultsAOR@out{%
    40 ms}%
    \else%
      \ifnum\pdfstrcmp{#1}{multiplierOldRuntime}=0%
      \def\resultsAOR@out{%
    4}%
    \else%
      \ifnum\pdfstrcmp{#1}{meemSweepSize}=0%
      \def\resultsAOR@out{%
    93,750}%
    \else%
      \def\resultsAOR@out{%
    ??}%
    \fi\fi\fi\fi\fi\fi\fi\fi\fi\fi\fi\fi\fi\fi\fi\fi\fi\fi\fi\fi\fi\fi\fi\fi\fi\fi\fi\fi\fi\fi\fi\fi\fi
  \fi
  \resultsAOR@out
}\newcommand\resultsAOR@I[1][all]{%
  \ifnum\pdfstrcmp{#1}{all}=0%
    \def\resultsAOR@I@out{%
    {[}%
      \{%
        "pct\_error\_baseline": {[}%
          0.8351547681041923, %
          {-}0.6753493890717724%
        {]}, %
        "pct\_error\_drag": {[}%
          10.790884267350643, %
          5.002294575560603%
        {]}, %
        "pct\_error\_meem": {[}%
          {-}13.649714349255992, %
          19.55578164189912%
        {]}, %
        "pct\_error\_total": {[}%
          {-}12.715966046141624, %
          5.090339861457517%
        {]}%
      \}, %
      \{%
        "pct\_error\_baseline": {[}%
          {-}0.1673331619860154, %
          {-}0.09472855056901647%
        {]}, %
        "pct\_error\_drag": {[}%
          5.756780832854308, %
          {-}4.8276990600241785%
        {]}, %
        "pct\_error\_meem": {[}%
          {-}0.28101071566998165, %
          5.870308306405258%
        {]}, %
        "pct\_error\_total": {[}%
          {-}2.0486004608632875, %
          {-}6.27926538734057%
        {]}%
      \}, %
      \{%
        "pct\_error\_baseline": {[}%
          39.22645866181762, %
          {-}27.657557186404567, %
          null, %
          null%
        {]}, %
        "pct\_error\_drag": {[}%
          65.47109436867527, %
          {-}11.439039626615426, %
          null, %
          null%
        {]}, %
        "pct\_error\_meem": {[}%
          39.22645866181762, %
          {-}32.7477682461545, %
          null, %
          null%
        {]}, %
        "pct\_error\_total": {[}%
          65.47109436867527, %
          {-}10.810656387995223, %
          null, %
          null%
        {]}%
      \}%
    {]}}%
  \else%
    \ifnum\pdfstrcmp{#1}{0}=0%
      \let\resultsAOR@I@out\resultsAOR@II%
    \else%
      \ifnum\pdfstrcmp{#1}{1}=0%
      \let\resultsAOR@I@out\resultsAOR@III%
    \else%
      \ifnum\pdfstrcmp{#1}{2}=0%
      \let\resultsAOR@I@out\resultsAOR@IV%
    \else%
      \def\resultsAOR@I@out{%
    ??}%
    \fi\fi\fi
  \fi
  \resultsAOR@I@out
}\newcommand\resultsAOR@II[1][all]{%
  \ifnum\pdfstrcmp{#1}{all}=0%
    \def\resultsAOR@II@out{%
    \{%
      "pct\_error\_baseline": {[}%
        0.8351547681041923, %
        {-}0.6753493890717724%
      {]}, %
      "pct\_error\_drag": {[}%
        10.790884267350643, %
        5.002294575560603%
      {]}, %
      "pct\_error\_meem": {[}%
        {-}13.649714349255992, %
        19.55578164189912%
      {]}, %
      "pct\_error\_total": {[}%
        {-}12.715966046141624, %
        5.090339861457517%
      {]}%
    \}}%
  \else%
    \ifnum\pdfstrcmp{#1}{pct_error_baseline}=0%
      \let\resultsAOR@II@out\resultsAOR@V%
    \else%
      \ifnum\pdfstrcmp{#1}{pct_error_drag}=0%
      \let\resultsAOR@II@out\resultsAOR@VI%
    \else%
      \ifnum\pdfstrcmp{#1}{pct_error_meem}=0%
      \let\resultsAOR@II@out\resultsAOR@VII%
    \else%
      \ifnum\pdfstrcmp{#1}{pct_error_total}=0%
      \let\resultsAOR@II@out\resultsAOR@VIII%
    \else%
      \def\resultsAOR@II@out{%
    ??}%
    \fi\fi\fi\fi
  \fi
  \resultsAOR@II@out
}\newcommand\resultsAOR@III[1][all]{%
  \ifnum\pdfstrcmp{#1}{all}=0%
    \def\resultsAOR@III@out{%
    \{%
      "pct\_error\_baseline": {[}%
        {-}0.1673331619860154, %
        {-}0.09472855056901647%
      {]}, %
      "pct\_error\_drag": {[}%
        5.756780832854308, %
        {-}4.8276990600241785%
      {]}, %
      "pct\_error\_meem": {[}%
        {-}0.28101071566998165, %
        5.870308306405258%
      {]}, %
      "pct\_error\_total": {[}%
        {-}2.0486004608632875, %
        {-}6.27926538734057%
      {]}%
    \}}%
  \else%
    \ifnum\pdfstrcmp{#1}{pct_error_baseline}=0%
      \let\resultsAOR@III@out\resultsAOR@IX%
    \else%
      \ifnum\pdfstrcmp{#1}{pct_error_drag}=0%
      \let\resultsAOR@III@out\resultsAOR@X%
    \else%
      \ifnum\pdfstrcmp{#1}{pct_error_meem}=0%
      \let\resultsAOR@III@out\resultsAOR@XI%
    \else%
      \ifnum\pdfstrcmp{#1}{pct_error_total}=0%
      \let\resultsAOR@III@out\resultsAOR@XII%
    \else%
      \def\resultsAOR@III@out{%
    ??}%
    \fi\fi\fi\fi
  \fi
  \resultsAOR@III@out
}\newcommand\resultsAOR@IV[1][all]{%
  \ifnum\pdfstrcmp{#1}{all}=0%
    \def\resultsAOR@IV@out{%
    \{%
      "pct\_error\_baseline": {[}%
        39.22645866181762, %
        {-}27.657557186404567, %
        null, %
        null%
      {]}, %
      "pct\_error\_drag": {[}%
        65.47109436867527, %
        {-}11.439039626615426, %
        null, %
        null%
      {]}, %
      "pct\_error\_meem": {[}%
        39.22645866181762, %
        {-}32.7477682461545, %
        null, %
        null%
      {]}, %
      "pct\_error\_total": {[}%
        65.47109436867527, %
        {-}10.810656387995223, %
        null, %
        null%
      {]}%
    \}}%
  \else%
    \ifnum\pdfstrcmp{#1}{pct_error_baseline}=0%
      \let\resultsAOR@IV@out\resultsAOR@XIII%
    \else%
      \ifnum\pdfstrcmp{#1}{pct_error_drag}=0%
      \let\resultsAOR@IV@out\resultsAOR@XIV%
    \else%
      \ifnum\pdfstrcmp{#1}{pct_error_meem}=0%
      \let\resultsAOR@IV@out\resultsAOR@XV%
    \else%
      \ifnum\pdfstrcmp{#1}{pct_error_total}=0%
      \let\resultsAOR@IV@out\resultsAOR@XVI%
    \else%
      \def\resultsAOR@IV@out{%
    ??}%
    \fi\fi\fi\fi
  \fi
  \resultsAOR@IV@out
}\newcommand\resultsAOR@V[1][all]{%
  \ifnum\pdfstrcmp{#1}{all}=0%
    \def\resultsAOR@V@out{%
    {[}%
      0.8351547681041923, %
      {-}0.6753493890717724%
    {]}}%
  \else%
    \ifnum\pdfstrcmp{#1}{0}=0%
      \def\resultsAOR@V@out{%
    0.8351547681041923}%
    \else%
      \ifnum\pdfstrcmp{#1}{1}=0%
      \def\resultsAOR@V@out{%
    {-}0.6753493890717724}%
    \else%
      \def\resultsAOR@V@out{%
    ??}%
    \fi\fi
  \fi
  \resultsAOR@V@out
}\newcommand\resultsAOR@VI[1][all]{%
  \ifnum\pdfstrcmp{#1}{all}=0%
    \def\resultsAOR@VI@out{%
    {[}%
      10.790884267350643, %
      5.002294575560603%
    {]}}%
  \else%
    \ifnum\pdfstrcmp{#1}{0}=0%
      \def\resultsAOR@VI@out{%
    10.790884267350643}%
    \else%
      \ifnum\pdfstrcmp{#1}{1}=0%
      \def\resultsAOR@VI@out{%
    5.002294575560603}%
    \else%
      \def\resultsAOR@VI@out{%
    ??}%
    \fi\fi
  \fi
  \resultsAOR@VI@out
}\newcommand\resultsAOR@VII[1][all]{%
  \ifnum\pdfstrcmp{#1}{all}=0%
    \def\resultsAOR@VII@out{%
    {[}%
      {-}13.649714349255992, %
      19.55578164189912%
    {]}}%
  \else%
    \ifnum\pdfstrcmp{#1}{0}=0%
      \def\resultsAOR@VII@out{%
    {-}13.649714349255992}%
    \else%
      \ifnum\pdfstrcmp{#1}{1}=0%
      \def\resultsAOR@VII@out{%
    19.55578164189912}%
    \else%
      \def\resultsAOR@VII@out{%
    ??}%
    \fi\fi
  \fi
  \resultsAOR@VII@out
}\newcommand\resultsAOR@VIII[1][all]{%
  \ifnum\pdfstrcmp{#1}{all}=0%
    \def\resultsAOR@VIII@out{%
    {[}%
      {-}12.715966046141624, %
      5.090339861457517%
    {]}}%
  \else%
    \ifnum\pdfstrcmp{#1}{0}=0%
      \def\resultsAOR@VIII@out{%
    {-}12.715966046141624}%
    \else%
      \ifnum\pdfstrcmp{#1}{1}=0%
      \def\resultsAOR@VIII@out{%
    5.090339861457517}%
    \else%
      \def\resultsAOR@VIII@out{%
    ??}%
    \fi\fi
  \fi
  \resultsAOR@VIII@out
}\newcommand\resultsAOR@IX[1][all]{%
  \ifnum\pdfstrcmp{#1}{all}=0%
    \def\resultsAOR@IX@out{%
    {[}%
      {-}0.1673331619860154, %
      {-}0.09472855056901647%
    {]}}%
  \else%
    \ifnum\pdfstrcmp{#1}{0}=0%
      \def\resultsAOR@IX@out{%
    {-}0.1673331619860154}%
    \else%
      \ifnum\pdfstrcmp{#1}{1}=0%
      \def\resultsAOR@IX@out{%
    {-}0.09472855056901647}%
    \else%
      \def\resultsAOR@IX@out{%
    ??}%
    \fi\fi
  \fi
  \resultsAOR@IX@out
}\newcommand\resultsAOR@X[1][all]{%
  \ifnum\pdfstrcmp{#1}{all}=0%
    \def\resultsAOR@X@out{%
    {[}%
      5.756780832854308, %
      {-}4.8276990600241785%
    {]}}%
  \else%
    \ifnum\pdfstrcmp{#1}{0}=0%
      \def\resultsAOR@X@out{%
    5.756780832854308}%
    \else%
      \ifnum\pdfstrcmp{#1}{1}=0%
      \def\resultsAOR@X@out{%
    {-}4.8276990600241785}%
    \else%
      \def\resultsAOR@X@out{%
    ??}%
    \fi\fi
  \fi
  \resultsAOR@X@out
}\newcommand\resultsAOR@XI[1][all]{%
  \ifnum\pdfstrcmp{#1}{all}=0%
    \def\resultsAOR@XI@out{%
    {[}%
      {-}0.28101071566998165, %
      5.870308306405258%
    {]}}%
  \else%
    \ifnum\pdfstrcmp{#1}{0}=0%
      \def\resultsAOR@XI@out{%
    {-}0.28101071566998165}%
    \else%
      \ifnum\pdfstrcmp{#1}{1}=0%
      \def\resultsAOR@XI@out{%
    5.870308306405258}%
    \else%
      \def\resultsAOR@XI@out{%
    ??}%
    \fi\fi
  \fi
  \resultsAOR@XI@out
}\newcommand\resultsAOR@XII[1][all]{%
  \ifnum\pdfstrcmp{#1}{all}=0%
    \def\resultsAOR@XII@out{%
    {[}%
      {-}2.0486004608632875, %
      {-}6.27926538734057%
    {]}}%
  \else%
    \ifnum\pdfstrcmp{#1}{0}=0%
      \def\resultsAOR@XII@out{%
    {-}2.0486004608632875}%
    \else%
      \ifnum\pdfstrcmp{#1}{1}=0%
      \def\resultsAOR@XII@out{%
    {-}6.27926538734057}%
    \else%
      \def\resultsAOR@XII@out{%
    ??}%
    \fi\fi
  \fi
  \resultsAOR@XII@out
}\newcommand\resultsAOR@XIII[1][all]{%
  \ifnum\pdfstrcmp{#1}{all}=0%
    \def\resultsAOR@XIII@out{%
    {[}%
      39.22645866181762, %
      {-}27.657557186404567, %
      null, %
      null%
    {]}}%
  \else%
    \ifnum\pdfstrcmp{#1}{0}=0%
      \def\resultsAOR@XIII@out{%
    39.22645866181762}%
    \else%
      \ifnum\pdfstrcmp{#1}{1}=0%
      \def\resultsAOR@XIII@out{%
    {-}27.657557186404567}%
    \else%
      \ifnum\pdfstrcmp{#1}{2}=0%
      \def\resultsAOR@XIII@out{%
    None}%
    \else%
      \ifnum\pdfstrcmp{#1}{3}=0%
      \def\resultsAOR@XIII@out{%
    None}%
    \else%
      \def\resultsAOR@XIII@out{%
    ??}%
    \fi\fi\fi\fi
  \fi
  \resultsAOR@XIII@out
}\newcommand\resultsAOR@XIV[1][all]{%
  \ifnum\pdfstrcmp{#1}{all}=0%
    \def\resultsAOR@XIV@out{%
    {[}%
      65.47109436867527, %
      {-}11.439039626615426, %
      null, %
      null%
    {]}}%
  \else%
    \ifnum\pdfstrcmp{#1}{0}=0%
      \def\resultsAOR@XIV@out{%
    65.47109436867527}%
    \else%
      \ifnum\pdfstrcmp{#1}{1}=0%
      \def\resultsAOR@XIV@out{%
    {-}11.439039626615426}%
    \else%
      \ifnum\pdfstrcmp{#1}{2}=0%
      \def\resultsAOR@XIV@out{%
    None}%
    \else%
      \ifnum\pdfstrcmp{#1}{3}=0%
      \def\resultsAOR@XIV@out{%
    None}%
    \else%
      \def\resultsAOR@XIV@out{%
    ??}%
    \fi\fi\fi\fi
  \fi
  \resultsAOR@XIV@out
}\newcommand\resultsAOR@XV[1][all]{%
  \ifnum\pdfstrcmp{#1}{all}=0%
    \def\resultsAOR@XV@out{%
    {[}%
      39.22645866181762, %
      {-}32.7477682461545, %
      null, %
      null%
    {]}}%
  \else%
    \ifnum\pdfstrcmp{#1}{0}=0%
      \def\resultsAOR@XV@out{%
    39.22645866181762}%
    \else%
      \ifnum\pdfstrcmp{#1}{1}=0%
      \def\resultsAOR@XV@out{%
    {-}32.7477682461545}%
    \else%
      \ifnum\pdfstrcmp{#1}{2}=0%
      \def\resultsAOR@XV@out{%
    None}%
    \else%
      \ifnum\pdfstrcmp{#1}{3}=0%
      \def\resultsAOR@XV@out{%
    None}%
    \else%
      \def\resultsAOR@XV@out{%
    ??}%
    \fi\fi\fi\fi
  \fi
  \resultsAOR@XV@out
}\newcommand\resultsAOR@XVI[1][all]{%
  \ifnum\pdfstrcmp{#1}{all}=0%
    \def\resultsAOR@XVI@out{%
    {[}%
      65.47109436867527, %
      {-}10.810656387995223, %
      null, %
      null%
    {]}}%
  \else%
    \ifnum\pdfstrcmp{#1}{0}=0%
      \def\resultsAOR@XVI@out{%
    65.47109436867527}%
    \else%
      \ifnum\pdfstrcmp{#1}{1}=0%
      \def\resultsAOR@XVI@out{%
    {-}10.810656387995223}%
    \else%
      \ifnum\pdfstrcmp{#1}{2}=0%
      \def\resultsAOR@XVI@out{%
    None}%
    \else%
      \ifnum\pdfstrcmp{#1}{3}=0%
      \def\resultsAOR@XVI@out{%
    None}%
    \else%
      \def\resultsAOR@XVI@out{%
    ??}%
    \fi\fi\fi\fi
  \fi
  \resultsAOR@XVI@out
}%
\makeatother%

\usepackage[authoryear]{natbib}
\tracinglostchars=1 % get rid of nullfont warnings
\usepackage{lineno}
\linenumbers
\usepackage[section]{placeins} % float barrier between sections

%%Author macros
\def\tsc#1{\csdef{#1}{\textsc{\lowercase{#1}}\xspace}}
\tsc{WGM}
\tsc{QE}
%%%

\begin{document}

\let\WriteBookmarks\relax
\def\floatpagepagefraction{1}
\def\textpagefraction{.001}

\shorttitle{Development, Validation, and Benchmarking of a Multidisciplinary Semi-Analytical Model for Wave Energy Converters}
\shortauthors{R. McCabe et~al.}

\title[mode=title]{Development, Validation, and Benchmarking of a Multidisciplinary Semi-Analytical Model for Wave Energy Converters}

% First author
%
% Options: Use if required
% eg: \author[1,3]{Author Name}[type=editor,
%       style=chinese,
%       auid=000,
%       bioid=1,
%       prefix=Sir,
%       orcid=0000-0000-0000-0000,
%       facebook=<facebook id>,
%       twitter=<twitter id>,
%       linkedin=<linkedin id>,
%       gplus=<gplus id>]

\author[MAE]{Rebecca McCabe}[
orcid=0000-0001-5108-998X,
auid=57980095200,
linkedin=rebecca-mccabe]
\cormark[1]
\fnmark[1]
\ead{rgm222@cornell.edu}
\ead[url]{http://sea.mae.cornell.edu}

% Optional:
\credit{Conceptualization, Methodology, Software, Validation, Writing - Original draft}

\author[MAE]{Madison Dietrich}[
    orcid=0009-0000-3590-5334
]
\credit{Software, Visualization,  Writing - Review \& Editing}

\author[mich]{Maha Haji}[
    orcid=0000-0002-2953-7253
]

\credit{Supervision, Writing - Review \& Editing}

\affiliation[MAE]{
    organization={Sibley School of Mechanical and Aerospace Engineering, Cornell University},
    addressline={124 Hoy Rd.},
    city={Ithaca},
    postcode={14853},
    state={NY},
    country={USA}
}

\affiliation[mich]{
    organization={Department of Mechanical Engineering, University of Michigan},
    addressline={G.G. Brown Laboratory, 2350 Hayward},
    city={Ann Arbor},
    postcode={48109},
    state={MI},
    country={USA}
}

\cortext[1]{Corresponding author}

\begin{abstract}
Wave energy converters (WECs) require system-level techno-economic analysis to balance power production, cost, and survivability,
yet existing simulation tools are often either too computationally intensive for large-scale optimization
or too narrow in disciplinary scope to capture realistic viability criteria.
This work presents MDOcean, a novel modular WEC simulation framework developed for rapid early-stage design exploration, parametric analysis, and multidisciplinary optimization.
MDOcean integrates analytical and semi-analytical models of hydrodynamics, dynamics, structures, and economics within a computationally efficient architecture, substantially reducing runtime while maintaining high accuracy. 

The framework includes a mesh-free eigenfunction-based linear hydrodynamic solver, a quasi-linearized frequency-domain dynamics engine capable of modeling drag and saturation nonlinearities, a structural sizing module incorporating yield, fatigue, ultimate, and buckling design criteria, and a simple cost model for techno-economic assessment.
The \acrlong{acr-qlps} optimal control formulation receives particular emphasis, unifying describing functions with an analytical quadratically-constrained quadratic program.
This formulation preserves the efficiency and intuition of standard frequency-domain techniques while handling nonlinearities and dynamic constraints.

Validation and benchmarking demonstrate that MDOcean's \resultsAOR[simRuntime] runtime is orders of magnitude faster than leading WEC simulation tools while maintaining agreement with higher-fidelity baselines to within a few percent in most cases.
Beyond computational performance, the framework also provides insight into limiting behaviors, scaling laws, subsystem interactions, and key tradeoffs governing WEC design and techno-economic performance.
Limiting assumptions include regular waves and point absorber archetype.
The MDOcean software is released open-source to accelerate WEC research, design, and optimization toward the clean energy transition.
\end{abstract}

\begin{graphicalabstract}
\includegraphics[\FigureOptions{aor:figs/from-matlab/graphical_abstract_modeling.pdf}]{figs/from-matlab/graphical_abstract_modeling.pdf}
\end{graphicalabstract}

\begin{highlights}
\item Presents MDOcean, an open-source semi-analytical wave energy converter simulation.
\item Integrates hydrodynamics, powertrain, structures, \& economics for coupling intuition.

\ifdefined\DISSERTATION
    \item Quasi-linear frequency domain dynamics model captures nonlinearities such as generator force saturation and drag.
    \item Derives analytical constrained optimal control solution for underactuated wave-to-wire multi-port model respecting effort, flow, and power limits.
    \item \gls{acr-an} stiffened-plate structural model assesses stress and stability failure modes under fatigue and storm survival loadcases.
\else
    \item Derives analytical frequency-domain constrained optimal control with nonlinearities.
\fi

\item Achieves \resultsAOR[simRuntime] runtime, 10-1000x faster than existing tools.
\item Verifies <10\% power \& amplitude error for standard reference model 3 design.

\end{highlights}

\begin{keywords}
marine renewable energy \sep
semi-analytical hydrodynamics \sep
linearized pseudo-spectral optimal control \sep
structural survivability \sep
techno-economic modeling \sep
model validation \sep
computational benchmarking
\end{keywords}

\maketitle

% added for my convenience, should take out before submission
%\tableofcontents
%\listoffigures
%\listoftables
%\clearpage
%\summarytexcount{main}
%\detailtexcount{sections/model}{subsection}
%\detailtexcount{sections/other-appendices}{section}
%\newpage

\renewcommand{\glossarysection}[2][]{} % no glossary title

% glossary: symbols
% \printglossary[type=sym,title={General Symbols}]
\printglossary[type=sym-letter-A]
\printglossary[type=sym-letter-a]
\printglossary[type=sym-letter-B]
\printglossary[type=sym-letter-b]
\printglossary[type=sym-letter-C]
\printglossary[type=sym-letter-c]
\printglossary[type=sym-letter-D]
\printglossary[type=sym-letter-d]
% \printglossary[type=sym-letter-E]
\printglossary[type=sym-letter-e]
\printglossary[type=sym-letter-F]
% \printglossary[type=sym-letter-f]
\printglossary[type=sym-letter-G]
\printglossary[type=sym-letter-g]
\printglossary[type=sym-letter-H]
\printglossary[type=sym-letter-h]
\printglossary[type=sym-letter-I]
\printglossary[type=sym-letter-i]
\printglossary[type=sym-letter-J]
% \printglossary[type=sym-letter-j]
\printglossary[type=sym-letter-K]
\printglossary[type=sym-letter-k]
% \printglossary[type=sym-letter-L]
\printglossary[type=sym-letter-l]
\printglossary[type=sym-letter-M]
\printglossary[type=sym-letter-m]
\printglossary[type=sym-letter-N]
% \printglossary[type=sym-letter-n]
\printglossary[type=sym-letter-O]
\printglossary[type=sym-letter-o]
\printglossary[type=sym-letter-P]
\printglossary[type=sym-letter-p]
\printglossary[type=sym-letter-Q]
\printglossary[type=sym-letter-q]
\printglossary[type=sym-letter-R]
% \printglossary[type=sym-letter-r]
% \printglossary[type=sym-letter-S]
\printglossary[type=sym-letter-s]
\printglossary[type=sym-letter-T]
\printglossary[type=sym-letter-t]
\printglossary[type=sym-letter-U]
\printglossary[type=sym-letter-u]
\printglossary[type=sym-letter-V]
\printglossary[type=sym-letter-v]
% \printglossary[type=sym-letter-W]
\printglossary[type=sym-letter-w]
\printglossary[type=sym-letter-X]
\printglossary[type=sym-letter-x]
% \printglossary[type=sym-letter-Y]
\printglossary[type=sym-letter-y]
\printglossary[type=sym-letter-Z]
% \printglossary[type=sym-letter-z]

%\printglossary[type=sym-letter-Alpha]
% \printglossary[type=sym-letter-alpha]
%\printglossary[type=sym-letter-Beta]
\printglossary[type=sym-letter-beta]
\printglossary[type=sym-letter-Gamma]
\printglossary[type=sym-letter-gamma]
\printglossary[type=sym-letter-Delta]
% \printglossary[type=sym-letter-delta]
%\printglossary[type=sym-letter-Epsilon]
% \printglossary[type=sym-letter-epsilon]
%\printglossary[type=sym-letter-Zeta]
\printglossary[type=sym-letter-zeta]
%\printglossary[type=sym-letter-Eta]
\printglossary[type=sym-letter-eta]
\printglossary[type=sym-letter-Theta]
% \printglossary[type=sym-letter-theta]
%\printglossary[type=sym-letter-Iota]
\printglossary[type=sym-letter-iota]
%\printglossary[type=sym-letter-Kappa]
% \printglossary[type=sym-letter-kappa]
% \printglossary[type=sym-letter-Lambda]
% \printglossary[type=sym-letter-lambda]
%\printglossary[type=sym-letter-Mu]
\printglossary[type=sym-letter-mu]
%\printglossary[type=sym-letter-Nu]
% \printglossary[type=sym-letter-nu]
\printglossary[type=sym-letter-Xi]
\printglossary[type=sym-letter-xi]
%\printglossary[type=sym-letter-Omicron]
%\printglossary[type=sym-letter-omicron]
\printglossary[type=sym-letter-Pi]
\printglossary[type=sym-letter-pi]
%\printglossary[type=sym-letter-Rho]
\printglossary[type=sym-letter-rho]
\printglossary[type=sym-letter-Sigma]
% \printglossary[type=sym-letter-sigma]
\printglossary[type=sym-letter-Tau]
% \printglossary[type=sym-letter-tau]
% \printglossary[type=sym-letter-Upsilon]
\printglossary[type=sym-letter-upsilon]
\printglossary[type=sym-letter-Phi]
% \printglossary[type=sym-letter-phi]
%\printglossary[type=sym-letter-Chi]
% \printglossary[type=sym-letter-chi]
\printglossary[type=sym-letter-Psi]
% \printglossary[type=sym-letter-psi]
% \printglossary[type=sym-letter-Omega]
\printglossary[type=sym-letter-omega]
%
% glossary: acronyms
% \printglossary[type=methods]
% \printglossary[type=control]
% \printglossary[type=org]
% \printglossary[type=struct]
% \printglossary[type=econ]
% \printglossary[type=dynamics]
% \printglossary[type=optim]
% \printglossary[type=software]
% \printglossary[type=wave]
\singleColMacro{
    % the other glossaries are printed in 2-col layout and conveniently end at the end of the page, but this glossary is shorter so is printed with mcoltree style in 1-col layout to keep the left and right cols the same height
\printglossary[type=journal-acronyms,style=mcoltree,nonumberlist]
}
\section{Introduction}
\label{sec:intro}
\subsection{Wave Energy Overview}
%(Why WECs, in renewable energy context):

% stronger opening sentence
The global climate crisis requires a transition to carbon-free energy sources such as ocean wave energy.
Ocean waves have higher consistency, predictability, and energy density than other renewable energy sources such as wind and solar, while the temporal complementarity of waves with other resources can improve grid resilience and capacity adequacy, decrease energy prices, and decrease requirements for energy storage and balancing power \citep{akdemir_opportunities_2023,bhattacharya_timing_2021,pennock_temporal_2022}.
\glspl{acr-wec}, the devices that harness and convert this energy, could provide both electricity for the grid at large scale and power for smaller offshore technologies like aquaculture, desalination, and marine sensing \citep{livecchi_powering_2019}.
\ifdefined\DISSERTATION
Despite the potential advantages of \glspl{acr-wec}, a significant decrease in the cost per unit energy production is required before the technology can be deployed at large scales. %\hl{(Add a more specific number to show that the LCOE needs to come down an order of magnitude).} Meanwhile, a lack of design convergence (...)
Thus, a design process that emphasizes techno-economic viability from an early stage is necessary.
\else
    However, \glspl{acr-wec} require substantial cost reduction before large-scale deployment, motivating a design process that emphasizes techno-economic viability from an early stage.
\fi

%(Why multidisciplinary):
\ifdefined\DISSERTATION
    However, the presence of strong interdisciplinary coupling complicates the application of typical design and optimization techniques to \glspl{acr-wec}.
    The powertrain, controller, bulk dimensions, hydrodynamic shape, and structural thicknesses must be selected to balance the complex tradeoff between power production and cost while ensuring survivability and practical feasibility.
    Accurately capturing all major relationships and requirements that drive structural, control, hydrodynamic, and electrical design is technically challenging and computationally prohibitive.
    Nonetheless, leveraging the coupling between subsystems and across technical domains can produce large gains in techno-economic viability.
    Automating this process through optimization allows \gls{acr-wec} designers to more quickly and systematically explore tradeoffs of major design decisions, and could eventually standardize comparisons between different \gls{acr-wec} architectures.
    Thus, performing a system-level \gls{acr-wec} optimization that simultaneously considers the coupled techno-economic goals, design decisions, and requirements is highly advantageous.
\else
    The strong interdisciplinary coupling of between powertrain, controller, bulk dimensions, hydrodynamic shape, and structural thicknesses complicates conventional \gls{acr-wec} design techniques.
    Systematically exploring these coupled tradeoffs requires a fast multidisciplinary simulation that can be embedded in an optimization loop.
\fi
%\Cref{fig:disciplines} is a non-exhaustive outline of the most common disciplines in a wave energy converter, and the level and direction of coupling between them.

% what is RM3
\ifdefined\DISSERTATION
    The U.S.\ Department of Energy's 2010-2014 Reference Model Project \citep{RM3} represents an early attempt to consolidate a multidisciplinary \gls{acr-wec} design process, though without optimization.
    The project's report \citep{RM3} introduces six benchmark marine energy device designs (three \glspl{acr-wec} and three tidal/current turbines) and describes their design considerations and performance calculations.
    Its website releases further design artifacts.\footnote{\url{https://openei.org/wiki/PRIMRE/Signature_Projects/Reference_Model}} 
    The third reference model, known as \acrshort{acr-rm3}, is a two-body point absorber \gls{acr-wec} designed for a reference site located in Humboldt Bay, \gls{acr-ca}.
    The device consists of a surface float and a spar comprising a vertical column and a subsurface damping plate.
    \Cref{fig:dims} pictures components and dimensions.
    The float oscillates vertically with incident waves while the spar stays mostly stationary, and this relative heave motion produces electrical power.
    The present study builds on considerable scholarship surrounding \gls{acr-rm3} and presents a multidisciplinary techno-economic simulation for \gls{acr-rm3}-like \glspl{acr-wec}.
\else
    This paper builds on the U.S.\ Department of Energy's Reference Model Project, specifically the third reference model (\acrshort{acr-rm3}) \citep{RM3}: a two-body point-absorber \gls{acr-wec} designed for Humboldt Bay, \gls{acr-ca}, comprising a surface float and a spar with a subsurface damping plate (\Cref{fig:dims}).
    Electrical power is produced from the relative heave motion of the float against the spar.
\fi

The remainder of this introduction reviews \gls{acr-wec} modeling literature, identifies the gap motivating semi-analytical system models, and summarizes the paper's contributions.

\subsection{\gls{acr-wec} System Modeling: State of the Art and Gaps}
\label{sec:lit}
\ifdefined\DISSERTATION
    Existing models of \glspl{acr-wec} in the literature vary widely in their disciplinary scope, fidelity, and suitability for optimization.
    In the wave energy industry, multidisciplinary modeling is common due to the imperative of considering all design elements, although different disciplines are often modeled separately and sequentially rather than concurrently.
    \citet{trueworthy_wave_2020} provides a critical summary of the \gls{acr-wec} design process and shares survey results from 25 \gls{acr-wec} designers and developers.
    50\% of respondents design all subsystems concurrently.
    The others tend to start with the \gls{acr-wec} shape and \gls{acr-pto}, then move to the controller and moorings, and finish with the power transportation system.
    \citet{trueworthy_wave_2020} concludes that concurrent design has promise to take advantage of subsystem interactions and is an under-utilized technique worthy of further study, emphasizing the importance of incorporating all design requirements at an early stage.
\else
    Existing \gls{acr-wec} models vary widely in disciplinary scope, fidelity, and suitability for optimization.
In industry, multidisciplinary modeling is common but typically sequential rather than concurrent; a survey of 25 \gls{acr-wec} designers \citep{trueworthy_wave_2020} found that only 50\% design all subsystems concurrently, and identified concurrent design as an under-utilized technique worthy of further study.
\fi

The academic state-of-the-art typically focuses on modeling and optimizing one or a few disciplines at a time, with the most common being hydrodynamics, dynamics, and controls.
Major modeling choices include the representation (or lack thereof) of the \gls{acr-wec}'s hydrodynamics, dynamics, controls, mooring, powertrain, structures, and economic viability, as well as the sea states considered.

% figure of taxonomy
\Cref{fig:model-taxonomy} provides a simplified taxonomy of the various \gls{acr-wec} system modeling approaches used in the literature,
and \Cref{tab:lit} compares the most relevant studies along this taxonomy.


% table summary
\ifdefined\DISSERTATION
    After examining the overall trends, each table entry will be described in more detail.
    \Acrfull{acr-pa} seem to be the most common device architecture, though most major device types are represented, and several studies optimize arrays of multiple devices.
    Nearly all model hydrodynamics and power production in some way.
    Hydrodynamic modeling is typically performed with the \acrfull{acr-bem}, occasionally with semi-analytical models like \acrfull{acr-meem}, and rarely with high fidelity \acrfull{acr-cfd}, fitting a model to numerical or experimental data (\acrshort{acr-fit}), or \acrfull{acr-alg}.
    Nonlinear forces such as drag are most often linearized and included as an additional \acrfull{acr-ld}, though other methods exist including quadratic representation (Q), \acrfull{acr-df}, strip-theory integration of Morison's equation (M), and \acrfull{acr-cfd}.
    Dynamic modeling is frequently performed both in the frequency domain (F) and time domain (T), with a handful of studies using extensions of frequency domain techniques (F+) to handle constraints and nonlinearities that are typically only handled in the time domain, and a few using the newer \acrfull{acr-ps}.
    A range of controllers are considered, including proportional-integral / reactive (PI), proportional / pure damping (P), their extensions to incorporate constraints (P+/PI+), nonlinear structured (\acrshort{acr-nl}), and unstructured (\acrshort{acr-uns}).

    Optimizations that incorporate powertrain, structures, and economics are less ubiquitous but still common.
    When a discipline is included, its fidelity can vary widely.
    For instance, some studies model the \gls{acr-pto} only with an \acrfull{acr-eff}, while others model the \gls{acr-pto} dynamics.
    Structural models are split relatively evenly between analytical models for stress as a function of dimensions (\acrshort{acr-an}), calculating loads without considering dimension-dependent stress (\acrshort{acr-lo}), and \acrfull{acr-fea}.
    Cost models most often use a geometric cost proxy (\acrshort{acr-geo}) such as hull volume or area.
    Some estimate costs of the structural material (\acrshort{acr-str}), power take-off (\acrshort{acr-pto}), mooring (\acrshort{acr-moor}), or include a comprehensive cost model (FULL).
    The use of a mooring model is relatively rare, with four studies modeling the mooring system physical design (\acrshort{acr-de}) and one modeling its dynamic properties (\acrshort{acr-dyn}).
    It is most common to simulate irregular waves (\acrshort{acr-irr}), though several studies assume regular waves (\acrshort{acr-reg}).
    Others incorporate multiple sea states via a joint probability density matrix (\acrshort{acr-irr}+/\acrshort{acr-reg}+).
    Only a few consider storm sea states (\acrshort{acr-sto}).
\else
    The dominant trends across these studies are: point absorbers as the most common architecture;
    \acrfull{acr-bem} hydrodynamics with occasional semi-analytical (\gls{acr-meem}) treatments;
    drag handled via linear damping or describing functions;
    dynamics in either frequency or time domain (with extensions of either to handle constraints); 
    and limited inclusion of structures, mooring, and cost modeling. 
\fi

% Per-study paragraphs
\ifdefined\DISSERTATION
    \paragraph{Notable large-scale optimization studies}
    More specifically, a few studies \citep{gaudin_single_2021,khanal_multi-objective_2024,edwards_optimisation_2022,garcia-teruel_reliability-based_2021} perform \gls{acr-wec} optimization with several disciplines.
    These studies are notable for their large scales of 120 \glspl{acr-dof} and 14, 15, and 66 design variables respectively.
    \citet{gaudin_single_2021} optimizes an array of 20 submerged cylinders, holding the \gls{acr-wec} design constant but optimizing site selection, array layout, and mooring design while considering cable and pile cost as well as pile load capacity.
    \citet{khanal_multi-objective_2024} optimizes an array of four cylindrical point absorbers, with 4 modeled disciplines of array layout, hydrodynamics, controls, and economics.
    \citet{edwards_optimisation_2022} is a point absorber shape optimization study that minimizes surface area as a cost proxy, accounts for stability constraints, and considers only geometries with power extraction equal to theoretical radiation, amplitude, and steepness limits.
    \citet{garcia-teruel_reliability-based_2021}, along with the smaller scale study \citep{garcia-teruel_design_2022} by the same authors, apply multi-objective optimization to trade off power with fatigue load and device volume, respectively.
    \citet{cotten_multi-objective_2022} similarly optimizes fatigue load for an attenuator-style \gls{acr-wec} with many more \glspl{acr-dof} but fewer (8) design variables.
    A different shape optimization \citep{abdulkadir_control_2024} notably calculates hydrodynamic coefficients via the semi-analytical \gls{acr-meem}, rather than from boundary element method as is typical, but does not include disciplines beyond hydrodynamics and dynamics/controls.
    A team at \gls{acr-nrel} attempted local optimization for a terminator shape study but found they could only reasonably optimize 3 of the intended 13 design variables at a time due to the high computational cost of their time-domain simulation \citep{housner_numerical_2024}.
    The previous studies \citet{garcia-teruel_reliability-based_2021,garcia-teruel_design_2022,cotten_multi-objective_2022} avoided this issue by simulating primarily in the frequency domain and using the time domain only to address \gls{acr-pto} constraints.

    \paragraph{Structural modeling and optimization}
    Structural optimization of \glspl{acr-wec} faces a similar challenge in the high compute time of \gls{acr-fea}.
    References \citep{coe_survey_2018, ove_arup__partners_ltd_structural_2016,paduano_towards_2024,giannini_wave_2022} provide guidelines for \gls{acr-wec} survivability modeling and design, including the prediction of both extreme loads and structural failure.
    Fatigue modeling methods include spectral analysis, rainflow counting, damage equivalent loads, and Miner's law, augmented with standard stress concentration factors or \gls{acr-fea} for local hotspots \citep{ove_arup__partners_ltd_structural_2016}.
    Strength modeling methods consist primarily of various fidelities of \gls{acr-fea}, although hand and spreadsheet calculations permit sizing to a specified stress early in the design process \citep{ove_arup__partners_ltd_structural_2016}.
    \gls{acr-fea} of \glspl{acr-wec} is most often performed with manual design iterations, as in the studies \citep{RM3,mi_multi-scale_2025}.
    These two studies, while not optimization, are extremely multidisciplinary, encompassing mooring design and detailed economic evaluation in addition to the more common hydrodynamics, powertrain, controls, and structures.
    Automating the design process reported there remains a long-term vision for \gls{acr-wec} multidisciplinary design.

    Notably, one study \citep{an_optimal_2024} performs true optimization of \gls{acr-wec} structural \gls{acr-fea} using the commercial computer aided engineering software Altair Inspire.
    The study examines four design variables, although it keeps the power calculations entirely separate from the structural optimization, neglecting the interdisciplinary coupling.
    Another structural optimization does not consider power at all, but is able to perform a brute-force sweep of 3 design variables by using analytical equations for stress and other failure modes \citep{ambuhl_reliability-based_2014}.
    The new Project SEA Stack modeling suite adds hydrodynamics to Project Chrono, a multiphysics software with the ability to perform finite element analysis using beam, shell, or solid elements, although it appears that so far the structural modeling capabilities have only been used for offshore wind turbines rather than \glspl{acr-wec} \citep{ogden_hydrochrono_2023}.
    Finally, several studies consider structures indirectly by examining the load on the device rather than the stress, e.g.\ \citep{nguyen_theoretical_2024, ferri_balancing_2014}, or equivalently assuming that stress scales with load and does not depend on any other dimensional design variable, e.g.\ \citep{garcia-teruel_reliability-based_2021, cotten_multi-objective_2022}.
    Like the earlier \gls{acr-meem} study, \citet{nguyen_theoretical_2024} uses the equivalent analytical hydrodynamics in elliptical coordinates for an oscillating surge \gls{acr-wec}.

    \paragraph{Controls and powertrain co-design}
    The study \citep{ferri_balancing_2014} is notable because in addition to structural loads, it compares various control schemes and simulates second-order drivetrain dynamics.
    The \acrfull{acr-ccd} methodology \citep{garcia-sanz_control_2019} emphasizes the importance of considering drivetrain dynamics and generator efficiency due to their strong effect on electrical power generation \citep{coe_useful_2023}.
    A number of recent studies \citep{rosati_control_2023,son_performance_2016,anderson_re-imagining_2024,devin_high-dimensional_2024,grasberger_control_2024} simultaneously optimize the \gls{acr-pto} with the controller.
    \citet{rosati_control_2023} optimizes an \gls{acr-owc} and is notable for its comparatively detailed cost modeling and its optimization of both economic and power variability metrics.
    \citet{mccabe_multidisciplinary_2022} also optimized power variation, although that study's use of regular waves makes the variation formulation less useful.
    Meanwhile, one point absorber \gls{acr-pto} optimization \citep{son_performance_2016} is additionally notable for its use of the \gls{acr-meem} method.
    \citet{gaebele_tpl_2025} is relevant because it optimizes the dimensional scale and \gls{acr-pto} of the \gls{acr-rm3} \gls{acr-wec} directly using a \gls{acr-ccd} approach.
    It also efficiently incorporates many time-domain constraints (maximum force, speed, average power, \gls{acr-rms} torque, and stroke) while maintaining frequency-domain hydrodynamics via the \gls{acr-ps} modeling method, which \citep{devin_high-dimensional_2024,grasberger_control_2024} also use.

    Other \gls{acr-ccd} papers optimize \gls{acr-wec} dimensions without modeling the \gls{acr-pto}.
    \citet{herber_dynamic_2014} performs a direct transcription co-optimization of cylinder draft and radius with control.
    \citet{lin_fast_2025} introduces a new \gls{acr-ccd} formulation for constrained control that is even faster than \gls{acr-ps}.
    The method leverages the analytical \gls{acr-pmp} and performs a case study to optimize a single geometric variable.
    A different \gls{acr-pmp} formulation \citep{abdulkadir_optimal_2024} incorporates \gls{acr-pto}-constrained array dynamics, although it has not yet been used for design optimization.
\else
    Existing studies therefore tend to sacrifice disciplinary breadth and/or search space
    in order to maintain computational tractability.

    High-fidelity multidisciplinary models are generally too expensive for large-scale optimization.
    For example, \linelabel{lne:struct-refs-start}the cost of \gls{acr-fea} means that most structural studies analyze rather than optimize
    \citep{coe_survey_2018,ove_arup__partners_ltd_structural_2016,paduano_towards_2024,giannini_wave_2022}.

    Meanwhile, optimization-oriented models often omit structures, powertrain dynamics, or economic assessment.
    The few \gls{acr-wec} structural optimizations \citep{an_optimal_2024,ambuhl_reliability-based_2014} do not calculate power\linelabel{lne:struct-refs-end}.
    Large-scale optimizations invariably exclude either structures or powertrain include
    \citep{gaudin_single_2021,khanal_multi-objective_2024,edwards_optimisation_2022,garcia-teruel_reliability-based_2021,garcia-teruel_design_2022,cotten_multi-objective_2022}.
    \gls{acr-ccd} studies
    \citep{rosati_control_2023,son_performance_2016,gaebele_tpl_2025,devin_high-dimensional_2024,grasberger_control_2024,herber_dynamic_2014,lin_fast_2025,abdulkadir_optimal_2024,gaebele_tpl_2025}
    efficiently handle constraints via \gls{acr-ps} methods or analytical \gls{acr-pmp},
    but none integrate structures or full economic costs.
    \citep{gaebele_tpl_2025} is directly comparable since it optimizes \gls{acr-rm3} dimensional scale and \gls{acr-pto}.
\fi

\paragraph{Software tools and gaps}
\ifdefined\DISSERTATION
    Beyond models used in individual studies, it is also helpful to analyze the capabilities and gaps of available simulation tools.
    The popular time-domain hydrodynamics tool WEC-Sim \citep{ruehl_wec-simwec-sim_2024} can become multidisciplinary through the use of integrations like \gls{acr-nemoh}/Capytaine (hydrodynamic coefficients), \gls{acr-pto}-Sim (electric and hydraulic components), WEC-Sim Applications (optimal control and other useful tools), and MoorDyn (mooring).
    However, as \citep{housner_numerical_2024} reveals, it is too slow to use for large-scale optimization.
    Meanwhile, the newer pseudospectral tool WecOptTool \citep{coe_initial_2020} is faster and suitable for \gls{acr-ccd}, but its disciplines are currently limited to hydrodynamics (via Capytaine) and control, with the ability to easily add custom dynamics like \gls{acr-pto} and mooring as needed, and it still can take considerable time for large optimizations.

    A 2023 marine energy software assessment \citep{ruehl_next-generation_2023} identified that industry-wide needs broadly include using state-of-the-art computation; overcoming limitations of existing software; ensuring appropriate language, architecture, and license of new software; accessing data; and ensuring usability and accuracy of free open source software.
    Specifically, relevant gaps include tools for powertrain and controls modeling, co-design, and optimization; high-fidelity models with faster runtime and low/mid-fidelity models with higher accuracy; inter-operabilty of models across disciplines and fidelities; and data for verification and validation \citep{ruehl_next-generation_2023}.
    This lack of a multidisciplinary \gls{acr-wec} simulation platform that is sufficiently fast for optimization helps explain why no optimization study in \Cref{tab:lit} simultaneously optimizes the device geometry, structures, and powertrain.
\else
    The two leading open-source \gls{acr-wec} simulation tools illustrate the gap.
    WEC-Sim \citep{ruehl_wec-simwec-sim_2024} is a time-domain simulator extensible to multidisciplinary studies via integrations
    (\acrshort{acr-nemoh}/Capytaine for hydrodynamics, PTO-Sim for powertrain, MoorDyn for mooring),
    but is too slow for large-scale optimization \citep{housner_numerical_2024}.
    WecOptTool \citep{coe_initial_2020} is a faster \gls{acr-ps} tool suitable for \gls{acr-ccd} but currently limited to hydrodynamics and control.
    A 2023 marine energy software assessment \citep{ruehl_next-generation_2023} explicitly identifies the absence of a fast multidisciplinary optimization platform as an industry-wide gap, which explains why no study in \Cref{tab:lit} simultaneously optimizes geometry, structures, and powertrain.
    This tradeoff motivates the development of a computationally efficient framework capable of concurrently capturing the major techno-economic couplings governing \gls{acr-wec} performance.
\fi
\begin{landscape}
\begingroup
\begin{figure}[htbp]
\centering
\includegraphics[\FigureOptions{aor:figs/from-matlab/taxonomy.pdf}]{figs/from-matlab/taxonomy.pdf}
\caption[\gls{acr-wec} system modeling approaches taxonomy]{\gls{acr-wec} system modeling approaches, with MDOcean's capabilities outlined in black.}\label{fig:model-taxonomy}
\fillandplacepagenumber
\end{figure}
\endgroup
\end{landscape}
%Need to copy sources over from slide 27 and 35 of F24 reseach slides with original charts (with color)
{\renewcommand{\arraystretch}{0.98}
\singleColMacro{
\setlength{\tabcolsep}{1pt}
\expanded{
    \noexpand\longtable{\TableColumnSpecs{aor:lit-review}}
}
\rot{\textbf{Ref}} & \rot{\textbf{Type}} & \rot{\textbf{Hydro}}& \rot{\textbf{Drag}} & \rot{\textbf{Dyn}} & \rot{\textbf{Ctrl}} & \rot{\textbf{Moor}} & \rot{\textbf{\acrshort{acr-pto}}} & \rot{\textbf{Struct}} & \rot{\textbf{Econ}} & \rot{\textbf{Waves}} \\
\hline
{\setlength{\baselineskip}{0.5\baselineskip} This work\par}
 & 2-\acrshort{acr-pa} &
{\setlength{\baselineskip}{0.5\baselineskip} \acrshort{acr-meem}, \acrshort{acr-alg}, \acrshort{acr-fit}\par}
 & \acrshort{acr-df}, M & F+ & PI+ & - & DY & \acrshort{acr-an} & \acrshort{acr-str}, \acrshort{acr-pto} & \acrshort{acr-reg}+, \acrshort{acr-sto} \\

\cite{mccabe_multidisciplinary_2022} & 2-\acrshort{acr-pa} & \acrshort{acr-alg} & - & F+ & PI+ & - & EF & \acrshort{acr-an} & \acrshort{acr-str} & \acrshort{acr-reg}+, \acrshort{acr-sto} \\

\cite{khanal_multi-objective_2024} & 1-\acrshort{acr-pa}* & \acrshort{acr-bem} & - & F & PI+ & - & - & - & \acrshort{acr-geo} & \acrshort{acr-reg} \\

\cite{gaudin_single_2021} & 1-\acrshort{acr-pa}* & \acrshort{acr-meem} & - & F & PI & \acrshort{acr-de} & - & \acrshort{acr-an} & \acrshort{acr-moor} & \acrshort{acr-irr}+, \acrshort{acr-sto} \\

\cite{edwards_optimisation_2022} & 1-\acrshort{acr-pa} & \acrshort{acr-bem} & - & F & P & - & - & - & \acrshort{acr-geo} & \acrshort{acr-reg} \\

\cite{garcia-teruel_reliability-based_2021}& 1-\acrshort{acr-pa} & \acrshort{acr-bem} & \acrshort{acr-ld} & F+ & PI & - & - & \acrshort{acr-ld} & - & \acrshort{acr-irr}+ \\

\cite{garcia-teruel_design_2022}& 1-\acrshort{acr-pa} & \acrshort{acr-bem} & - & F+ & PI & - & - & - & - & \acrshort{acr-irr} \\

\cite{cotten_multi-objective_2022} & \acrshort{acr-atn} & \acrshort{acr-bem} & - & F+ & P & - & - & \acrshort{acr-ld} & - & \acrshort{acr-irr}+ \\

\cite{abdulkadir_control_2024} & 1-\acrshort{acr-pa} & \acrshort{acr-meem} & \acrshort{acr-ld} & T & \acrshort{acr-nl} P & - & - & - & - & \acrshort{acr-irr} \\

\cite{housner_numerical_2024} & \acrshort{acr-trm} & \acrshort{acr-bem} & - & T & PI & DY & - & - & - & \acrshort{acr-reg} \\

\cite{al_shami_parameter_2019} & 2-\acrshort{acr-pa} & \acrshort{acr-bem} & \acrshort{acr-df} & F & PI & - & - & - & - & \acrshort{acr-reg} \\

\cite{RM3} & 2-\acrshort{acr-pa} & \acrshort{acr-bem} & Q & T & P & \acrshort{acr-de} & - & \acrshort{acr-fea} & FULL & \acrshort{acr-irr}+, \acrshort{acr-sto} \\

\cite{mi_multi-scale_2025} & \acrshort{acr-osw} & \acrshort{acr-bem} & \footnotesize{Various} & T & PI & \acrshort{acr-de} & DY & \acrshort{acr-fea} & FULL & \acrshort{acr-irr}+, \acrshort{acr-sto} \\

\cite{an_optimal_2024} & \acrshort{acr-otp} & \acrshort{acr-cfd} & \acrshort{acr-cfd} & T & - & - & - & \acrshort{acr-fea} & - & \acrshort{acr-sto} \\

\cite{ambuhl_reliability-based_2014} & 1-\acrshort{acr-pa} & - & Q & - & - & \acrshort{acr-de} & - & \acrshort{acr-an} & \acrshort{acr-str} & \acrshort{acr-sto} \\

\cite{nguyen_theoretical_2024} & \acrshort{acr-osw} & \acrshort{acr-meem} & \acrshort{acr-ld} & F & PI & - & - & \acrshort{acr-ld} & - & \acrshort{acr-irr} \\

\cite{ferri_balancing_2014} & 1-\acrshort{acr-pa} & \acrshort{acr-bem} & \acrshort{acr-ld} & T & P, PI+, US & - & DY & \acrshort{acr-ld} & \acrshort{acr-str} & \acrshort{acr-irr}+ \\

\cite{rosati_control_2023} & \acrshort{acr-owc} & \acrshort{acr-bem} & - & T & \acrshort{acr-nl} P+ & - & EF & - & \acrshort{acr-pto} & \acrshort{acr-irr}+ \\

\cite{son_performance_2016} & 2-\acrshort{acr-pa} & \acrshort{acr-meem} & \acrshort{acr-ld} & F & P & - & EF & - & - & \acrshort{acr-reg} \\

\cite{gaebele_tpl_2025} & 2-\acrshort{acr-pa} & \acrshort{acr-bem} & - & \acrshort{acr-ps} & PI & - & DY & - & \acrshort{acr-pto}, \acrshort{acr-geo} & \acrshort{acr-irr} \\

\cite{devin_high-dimensional_2024} & 1-\acrshort{acr-pa} & \acrshort{acr-fit} & - & \acrshort{acr-ps} & US & - & DY & - & - & \acrshort{acr-reg} \\

\cite{grasberger_control_2024} & \acrshort{acr-osw} & \acrshort{acr-bem} & - & \acrshort{acr-ps} & PI & - & DY & - & \acrshort{acr-geo} & \acrshort{acr-irr} \\

\cite{herber_dynamic_2014} & 1-\acrshort{acr-pa} & \acrshort{acr-fit} & \acrshort{acr-ld} & T & US & - & - & - & - & \acrshort{acr-irr} \\

\cite{lin_fast_2025} & 1-\acrshort{acr-pa} & \acrshort{acr-bem} & \acrshort{acr-ld} & F+ & US & - & - & - & - & \acrshort{acr-irr}+ \\

\cite{abdulkadir_optimal_2024} & 1-\acrshort{acr-pa}* & \acrshort{acr-bem} & - & T & US & - & - & - & - & \acrshort{acr-irr} \\

%\cite{zhang_performance_2024} & 3-PA & MEEM & LD & 3 & F & P & DY & - & - & - & IRR
%& BF & PWR & - & CTL & 2 \\

%\cite{wen_shape_2018} & 1-PA & BEM & - & 1 & F & P & - & - & - & - & IRR+
%& DOE, LOC & PWR & STA & DIM, SHP & 3 \\

\hline 
\caption[Comparison of disciplinary scope and fidelity of \gls{acr-wec} models]{Comparison of the disciplinary scope and fidelity of optimization-relevant \gls{acr-wec} models.
See \Cref{fig:model-taxonomy} for abbreviations.
An asterisk * in type column means array effects from multiple interacting devices are modeled.}
\label{tab:lit}
%\fillandplacepagenumber
\endlongtable
} % end single col macro
} % end array stretch

\subsection{Relevance of Semi-Analytical Modeling}   
An optimization study can use any kind of simulation model to assess the design's performance, including simple algebraic models, theory-intensive semi-analytical models, and standard numerical models.
\Cref{tab:model-types} summarizes the tradeoffs.

\begin{table}[htbp]
    \centering
    \caption{Comparison of model types}
    \label{tab:model-types}
\begin{tabular}{cM{0.14\linewidth}cM{0.13\linewidth}}
         \textbf{Model type}&  \textbf{Expertise required}&  \textbf{Accuracy}&  \textbf{Computation cost}\\ \hline
         \textbf{Algebraic}&  Low&  Low&  Low\\
         \textbf{Semi-analytical}&  High&  Med&  Low\\
         \textbf{Numerical}&  Med&  Med/High&  High\\
    \end{tabular}
    \end{table}

For early-stage \gls{acr-wec} design optimization, where the model must be run many thousands of times, semi-analytical models offer the most beneficial combination of speed and accuracy.

\ifdefined\DISSERTATION
    Semi-analytical models can be orders of magnitude faster than numerical models, often with minimal sacrifices in accuracy.
    Additionally, semi-analytical models typically vary smoothly with their inputs and can make it easier to obtain accurate gradients, a feature that is often not the case for fully numerical models that may contain non-differentiable artifacts.
    Accurate gradients can also speed up convergence of the optimization and allow it to navigate high-dimension design spaces that would be intractable with heuristic algorithms.

    Unfortunately, semi-analytical models are less likely to be available commercially or open-source as numerical models are, and developing them from scratch requires substantial background in mathematics, numerical methods, and the technical domain of the model (for example, hydrodynamics or constrained optimal control).
    Often, semi-analytical models remain in the realm of theoretical academics, and ``application'' studies that use them to perform design/optimization usually consider just one discipline and have low practical realism.
    This is especially true if the models are closed-source, preventing others besides the model authors from using them.
    In the context of \glspl{acr-wec}, low practical realism usually means studies are limited to mechanical power production, without considering cost and realistic constraints, or perhaps using volume or surface area as a cost proxy.
    The design studies with high practical realism tend to use numerical models instead.
    As the literature review in \Cref{tab:lit} revealed, the high computational cost of the numerical models means that often a study is either multidisciplinary without optimization \citep{RM3,mi_multi-scale_2025}, or performs optimization for a single discipline, such as controls, powertrain, or structures, not all at once.

    Of the few studies that perform multidisciplinary optimization, those that use numerical models take many hours to run, especially when using heuristic optimizers, preventing designers from iterating quickly.
    For example, the \gls{acr-bem}-based 4-\gls{acr-wec} array optimization in the study \citep{khanal_multi-objective_2024} took 36 hours to run on a high-end workstation for only a single frequency.
    This type of model is too slow for use in realistic large-scale optimizations.
    On the other extreme, \citet{mccabe_multidisciplinary_2022} uses an algebraic hydrodynamic model that runs quickly but lacks sufficient accuracy to draw meaningful design conclusions.
    The present paper introduces a set of interoperable semi-analytical models that are designed to have the speed and robustness required for multidisciplinary optimization.
\else
    The directly comparable studies above either use slow numerical models that limit optimization scope (e.g.\ a \gls{acr-bem}-based 4-\gls{acr-wec} array optimization that took 36 hours per frequency on a high-end workstation \citep{khanal_multi-objective_2024}) or fast algebraic models that lack accuracy for meaningful design conclusions (e.g.\ \citep{mccabe_multidisciplinary_2022}).
    Semi-analytical models offer a useful intermediate that has not yet been combined with multidisciplinary scope in an open-source WEC tool, which is the gap MDOcean addresses.
    %\hl{Merge this section with the previous one, and just say that I can explain all the stuff in 1.2 by the type of model they use.}
\fi

\subsection{Paper Contributions and Roadmap}
\ifdefined\DISSERTATION
    To address the gaps just described, the authors of this paper create a fast, multidisciplinary, and open-source \gls{acr-wec} simulation software.
    The developed software, called MDOcean \citep{mccabe_mdocean_2024}, implements relevant semi-analytical models from various fields, integrates them in a scalable framework for \gls{acr-wec} simulation and concurrent optimization, and validates them in a \gls{acr-wec} system context.
    The modeling framework articulated here is scalable to more disciplines, and enables systematic \gls{acr-wec} design optimization.
\else
    This paper presents MDOcean \citep{mccabe_mdocean_2024}, an open-source semi-analytical \gls{acr-wec} simulation framework addressing the gap identified above.
    MDOcean is available at \url{github.com/symbiotic-engineering/MDOcean}.\linelabel{lne:github-link-intro}
    The work provides the following contributions:
    \begin{itemize}[leftmargin=*]
        \item \textbf{Integrates semi-analytical models} from hydrodynamics, dynamics, control, structures, and economics in a single \linelabel{lne:open-source}open-source framework suitable for optimization
        \item Formulates an underactuated optimal control problem with coupled nonlinear multi-port dynamics, and introduces the \textbf{\acrfull{acr-qlps} method} to approximate time-domain constraints and nonlinearities
        \item \textbf{Analytically solves} the \gls{acr-qlps} optimal control problem as a univariate \gls{acr-qcqp} with a geometric interpretation on the complex reflection coefficient plane
        \item \textbf{Validates} against WEC-Sim and the \gls{acr-rm3} reference design across power, structural, and economic outputs
        \item Demonstrates \textbf{computational performance} (\resultsAOR[simRuntime] simulation runtime) enabling optimization studies that are otherwise infeasible
        \item Illustrates the \textbf{mathematical structure and effect magnitude} of subsystem interactions and techno-economic scaling relationships, providing intuition to facilitate early-stage design tradeoffs
    \end{itemize}
\fi
\ifdefined\DISSERTATION
    In \Cref{sec:model-structure}, the research objective and design scope are first used to define a broad problem formulation and module decomposition.
    Next, \Cref{sec:modules} describes the construction of the simulation model.
    It develops the assumptions, analysis method, and implementation details of each module based on the appropriate balance of accuracy and speed.
    \Cref{sec:validation-benchmarking} discusses model validation and runtime benchmarking.
    Finally, \Cref{sec:discussion} showcases results from model sweeps and insights from the model structure.
\else
    \Cref{sec:model-structure} introduces MDOcean's modular structure and the rationale behind its analysis architecture.
    \Cref{sec:modules} documents the assumptions, analysis methods, and implementation choices of each module -- geometry, hydrodynamics, dynamics and control, structures, and economics.
    \Cref{sec:validation-benchmarking} compares MDOcean against WEC-Sim and the \gls{acr-rm3} reference design, demonstrates the runtime gain over established tools, and discusses where the model is reliable.
    \Cref{sec:discussion} presents results from variable sweeps and model-derived insights, along with limitations and future work.
    Extended derivations, validation studies, and implementation details are provided in the appendices to preserve readability of the primary narrative while maintaining reproducibility.
    A companion paper \citep{mccabe_leveraging_2026} applies MDOcean to a multidisciplinary techno-economic optimization of the \gls{acr-rm3}.

    \linelabel{lne:summary-objective-start}
    In summary, this study aims to develop a model that has sufficient speed and scope to support early-stage multidisciplinary design optimization and tradeoff analysis of wave energy converters.
    \linelabel{lne:summary-objective-end}
\fi
%\clearpage
\section{Model Structure}
\label{sec:model-structure}

The overall scope of \gls{acr-wec} design includes the selection of discrete architecture choices, such the type of device and \gls{acr-pto} (top row of \Cref{fig:design-scope});
continuous design variables in the physical domain, such as dimensions (second row of \Cref{fig:design-scope});
static analysis to evaluate the properties of these physical designs (blue boxes in third row of \Cref{fig:design-scope});
dynamic analysis to quantify the effect of these static properties on coupled dynamic properties like power production (blue dotted outline of \Cref{fig:design-scope});
and performance analysis to evaluate the viability of the design based on its static and dynamic properties (far right of \Cref{fig:design-scope}).
MDOcean performs static analysis for a single type of floating body (an \gls{acr-rm3}-like two-body point absorber)
as well as dynamic analysis and optimal control for the full coupled system,
assuming that the static properties for other components (e.g. drivetrain, generator) are known or swept inputs.

\ifdefined\DISSERTATION
    \begin{foldoutpage}
\fi
        \begin{figure}
        \centering
        \includegraphics[width=\textwidth]{figs/two_port_dimension_coeff_chain.pdf}
        \caption[Full \gls{acr-wec} design space and scope of MDOcean's analyses]{Flowchart showing the full \gls{acr-wec} design space and the scope of MDOcean's static, dynamic, and optimal control analyses.
        For visual clarity, some dependencies are omitted, such as that of the constraint limits on static analyses.
        These dependencies are included in the \gls{acr-xdsm} diagram in \Cref{fig:n2}.}
        \label{fig:design-scope}
        \end{figure}
\ifdefined\DISSERTATION
    \end{foldoutpage}
\fi

The \gls{acr-mdo} and \gls{acr-ccd} literature offers various structures to combine dynamic and static analysis, design optimization, and optimal control.
After considering eight possible analysis and control architectures, MDOcean adopts the explicit-\gls{acr-mdf} formulation, highlighted in \Cref{fig:control-arch},
which is the most computationally efficient given the coupling and computational cost of its various disciplines.
The analysis architecture determines whether it is a solver or an optimization constraint which converges the dynamic residuals, while the control architecture determines whether control variables are determined via simulation, residual solve, or simultaneous/nested optimization.
The implementation follows the module structure depicted in \Cref{fig:n2}.
\ifdefined\DISSERTATION
\else
    \citet{mccabe_dissertation_2026} provides more detailed architecture rationale, coupling analysis, the relation to \gls{acr-mdo} and \gls{acr-ccd} literature, and describes an earlier version of MDOcean using the implicit-\gls{acr-mdf} architecture.
\fi
% Point of this section: out of the 8 examples of CCD architecture and \gls{acr-mdo} architectures, I use the one in the box because it's fast.
% I implement it with the module structure of the XDSM diagram.

\begin{figure}[b!]
\centering
\includegraphics[\FigureOptions{aor:figs/from-matlab/control_analysis_flowcharts.pdf}]{figs/from-matlab/control_analysis_flowcharts.pdf}
\caption[Optimization architectures by analysis and control method]{Optimization architectures organized by analysis method (\gls{acr-mdf} and \gls{acr-sand}) and control method (explicit, implicit, simultaneous, and nested).}
\label{fig:control-arch}
\end{figure}

\subsection{Multiple-Discipline-Feasible Analysis Architecture}
\ifdefined\DISSERTATION
    All modules except the dynamics module are explicit, meaning they require no internal iteration to converge.
    The dynamics module requires iteration to incorporate nonlinearities into the quasi-linear frequency domain model, shown in the \gls{acr-xdsm} diagram as the orange box.
    The dynamics iteration occurs separately for each outer iteration of the optimizer.
    This structure is known as the multiple discipline feasible (\gls{acr-mdf}) architecture.
    In principle, it is possible to remove the dynamics iteration and instead incorporate the nonlinear dynamics as a residual equality constraint within the optimization, which is known as the simultaneous analysis and design (\gls{acr-sand}) architecture \citep{martins_multidisciplinary_2013}.
    \gls{acr-sand} generally decreases the runtime of the simulation (analysis) but requires more optimization iterations, since it is the optimization rather than the simulation that must converge residuals. 
    The WecOptTool control co-optimization software pursues the \gls{acr-sand} strategy by collocating the dynamics constraints with the pseudo-spectral method \citep{coe_initial_2020}.
    
    Readers familiar with trajectory optimization should note that the choice of \gls{acr-mdf} versus \gls{acr-sand} as an \gls{acr-mdo} architecture parallels that of shooting versus collocation as a transcription method, in the sense that one chooses to enforce governing equations with simulation (\gls{acr-mdf} and shooting) versus optimization (\gls{acr-sand} and collocation) \citep{tedrake_underactuated_2024}.
    In \gls{acr-mdo}, the appropriate decision depends on characteristics of all modules, not merely the module whose governing equations are in question.
    As \Cref{sec:sim-runtime} will demonstrate, the hydrodynamics module, not dynamics and controls, is the dominant computational cost of the simulation.
    This means that even a substantial dynamics speedup represents only a small speedup of the full simulation and is unlikely to outweigh the slowdown of more (hydrodynamics-dominated) simulation evaluations.
    For this reason, \gls{acr-sand} would likely increase the runtime of the full optimization, motivating MDOcean's selection of \gls{acr-mdf} as the more suitable architecture.
    Another advantage of \gls{acr-mdf} over \gls{acr-sand} is that the latter invites the possibility of inconsistent dynamics if the optimization terminates unsuccessfully \citep{martins_multidisciplinary_2013} and requires optimizing a dummy constant objective to perform simulation without optimization.
\else
    In \gls{acr-mdo}, when governing equations circularly depend on one another,
    % within a module (an ``implicit'' discipline) or between modules (an interdisciplinary ``feedback loop''),
    dependencies are often resolved with numerical iteration.
    Two common architectures for doing so include \acrfull{acr-mdf} and \acrfull{acr-sand}.
    They enforce governing equations differently: 
    \gls{acr-mdf} converges residuals to zero \emph{inside the simulation},
    while \gls{acr-sand} incorporates the residuals as equality constraints in an \emph{external optimization} \citep{martins_multidisciplinary_2013}.
    \gls{acr-sand} reduces the cost of each simulation but can increase the number of simulations required to optimize,
    so the preference is problem-dependent.
    Separately, \gls{acr-sand} has practical limitations:
    it invites the possibility of inconsistent dynamics if the optimization terminates unsuccessfully \citep{martins_multidisciplinary_2013}
    and requires optimizing a dummy constant objective to perform simulation without optimization.

    Only one of MDOcean's governing equations requires such iteration:
    the resolution of nonlinearities in the quasi-linear frequency-domain formulation.
    % governing equations contains no interdisciplinary feedback loops and the only implicit module is dynamics (to resolve nonlinearities in the quasi-linear frequency-domain formulation).
    This means that the ratio of dynamics simulation cost to total simulation cost and the expected number of iterations required for convergence determine the choice of \gls{acr-mdf} and \gls{acr-sand}.
    Hydrodynamics, not dynamics, will be shown to dominate computational cost (\Cref{sec:sim-runtime}), and the dynamics iteration converges more rapidly than the design optimization, so \gls{acr-mdf} is chosen.
\fi

\subsection{Explicit Optimal Control Architecture}
\ifdefined\DISSERTATION
    In \gls{acr-ccd}, the nested and simultaneous architectures are often contrasted \citep{sundarrajan_towards_2021,herber_nested_2018}.
    The nested architecture, known as ``asymmetric subspace optimization'' in \gls{acr-mdo} literature \citep{chittick_asymmetric_2009}, optimizes control variables in an inner loop for each iteration of the outer design optimization.
    This takes advantage of the fact that the inner optimization is often convex and can be solved quickly.
    On the other hand, the simultaneous architecture uses a single optimizer to optimize design and control variables concurrently, creating a larger (typically non-convex) optimization problem and avoiding the difficulty of inner loop infeasibility that may arise in the nested approach \citep{herber_nested_2018}. 
    We presently distinguish two additional types of \gls{acr-ccd} architectures in which the controller is not numerically optimized and instead obtained through an optimal control law.
    The optimal control law can be either explicitly derived and applied in the simulation (``explicit'' or analytical), or implicitly solved for by numerically zeroing the residual of an analytical optimal control condition (``implicit'' or semi-analytical).
    
    Each of these four control architectures can coexist with both \gls{acr-mdf} and \gls{acr-sand} analysis architectures, shown in \Cref{fig:control-arch}.
    The analysis architecture determines whether it is a solver or an optimization constraint which handles the dynamic residuals to converge the state, while the control architecture determines whether control variables are determined via simulation, residual solve, or simultaneous/nested optimization.
    In the ``implicit'' control diagram, the plural ``residuals'' is used rather than ``residual'' to convey that there exist residuals on the optimal control conditions to determine the control variables, in addition to those on the dynamic equations to determine the state variables.
    While the implicit-\gls{acr-sand} and simultaneous-\gls{acr-sand} block diagrams look nearly identical, the former determines the control variables via residuals of pre-determined optimality conditions expressed as constraints to the optimizer, while the latter determines them to maximize the optimization objective without a-priori knowledge of the optimal control conditions.
    
    In WecOptTool, the pseudo-spectral optimizer determines the state and control variables for a fixed design without using previously-known optimality conditions.
    Depending on whether a user adds additional design variables to the existing pseudo-spectral optimizer or in a separate loop outside, WecOptTool would be considered either simultaneous-\gls{acr-sand} or nested-\gls{acr-sand} respectively \citep{coe_initial_2020}.
    
    For MDOcean, we choose the explicit-\gls{acr-mdf} architecture.
    Like the implicit-\gls{acr-mdf} architecture, it minimizes computational cost by avoiding the need to include control variables in an optimization and lowering the number of dynamic constraints which the optimizer must consider.
    This is possible due to the semi-analytical dynamics and control module, where the optimal constrained control conditions are derived analytically.
    Explicit-\gls{acr-mdf} is chosen over implicit-\gls{acr-mdf} because it avoids the latter's numerical root-finding and is thus faster, as \Cref{sec:mod-freq-domain} will detail.
    An early version of MDOcean used the implicit-\gls{acr-mdf} architecture, which is described in \Cref{sec:appendix-qp-numerical} for completeness.
\else
    Control co-design problems are commonly classified into nested and simultaneous formulations \citep{sundarrajan_towards_2021,herber_nested_2018}.
    In the nested approach, control is optimized in an inner loop, while in the simultaneous approach, control and design variables are optimized together.
    Alternatively, the control can be obtained without numerical optimization via an implicit or explicit optimality condition. 
    An explicit optimal control condition is typically the fastest; though it is not available for all cases, its availability here (\Cref{sec:optimal-control}) makes it the natural choice.
    
    For comparison, the WecOptTool software would be considered either simultaneous-\gls{acr-sand} or nested-\gls{acr-sand} depending on whether a user adds additional design variables to the existing pseudo-spectral optimizer or in a separate loop outside, respectively \citep{coe_initial_2020}.
\fi

%%%%%%%%%%%%%%%%%%%%%%%%%%%%% XDSM and modules %%%%%%%%%%%%%%%%%%%%%%%%%%%%%%%%%%%%%%%%%%%%%%%%%%%

\subsection{Modular Decomposition and Coupling}
The simulation is composed of five coupled modules: geometry, hydrodynamics, dynamics and control, structures, and economics.
The geometry module computes derived properties such as areas, volumes, masses, centers of gravity, hydrostatics, and stability margins from bulk dimensions and structural thicknesses.
The hydrodynamics module evaluates frequency-domain hydrodynamic coefficients using a semi-analytical formulation based on the device geometry.
The dynamics and control module combines hydrodynamic coefficients, mass properties, and generator ratings to compute system response, loads, and power production using a quasi-linear frequency-domain model with approximate drag and saturation nonlinearities.
Both operational and extreme (storm) conditions are considered.
The structures module evaluates stress and factor-of-safety constraints under yield, ultimate, buckling, and fatigue criteria.
Finally, the economics module estimates capital cost (\gls{acr-pto} and structure) and levelized cost of energy using power production and material usage.

Although this work does not perform optimization, the model is developed to support system-level design optimization.
The module coupling structure is illustrated in the \gls{acr-xdsm} diagram in \Cref{fig:n2}.
Inputs and outputs are organized according to standard \gls{acr-mdo} conventions \citep{lambe_extensions_2012},
and an optimizer would update design variables $\gls{sym-vec-x}$ using objective $\gls{sym-J}$ and constraints $\gls{sym-g}$ until convergence to $\gls{sym-vec-x}^*$.

\begin{figure*}[htbp]
\centering
\includegraphics[\FigureOptions{aor:figs/from-matlab/xdsm.pdf}]{figs/from-matlab/xdsm.pdf}
\caption[Simplified \gls{acr-xdsm} diagram]{Simplified \gls{acr-xdsm} diagram. Modules are shown along the diagonal. Inputs come from the top and left sides, and outputs go out the bottom and right sides.}
\label{fig:n2}
\end{figure*}

\ifdefined\DISSERTATION
    \paragraph{Absence of Feedback Coupling}
    The lack of variables in the lower left portion of the diagram indicates that there is no feedback coupling between modules, so an external solver to enforce consistency is not necessary.
    This helps decrease the simulation computation time and is appealing for optimization.
    \paragraph{Presence of Design Variable and Feed-Forward Coupling}
    The fact that the same design variable group is an input to multiple modules indicates the presence of design variable coupling.
    Most heavily coupled are the bulk dimensions, which are input to the geometry, hydrodynamics, and structural modules.
    Thicknesses and generator ratings are each input to two modules, with thicknesses coupling geometry and structures and generator ratings coupling dynamics and economics.
    Moreover, the presence of variables in the upper right of the diagram indicates feed-forward coupling in which the output of one module directly affects subsequent modules.
    
    Analysis of coupling is important because under certain conditions, the optimization can be decomposed into separate problems to be solved in sequence or in parallel without sacrificing performance.
    For example, techniques like block coordinate descent or alternating direction method of multipliers apply if coupling is convex \citep{wright_coordinate_2015}, and monotonicity-based decomposition and constraint elimination apply if coupling is monotonic \citep{azarm_monotonicity-based_1988,papalambros_principles_2017}.
    The present coupling of \gls{acr-wec} disciplines is non-convex and non-monotonic, so concurrent optimization is required to obtain the system optimal design.
    Non-monotonicity means that perturbing a variable from one module in a certain direction does not necessarily determine the direction of the propagated change in coupling variables, objective, and constraints computed in other modules.
    Non-convexity arises from the presence of both positive and negative curvature in the coupled variable space, leading to multiple local optima.
    With an objective $\gls{sym-J}$ of \gls{acr-lcoe}, for example, while an increase in hydrodynamic damping generally increases power production (beneficial to $\gls{sym-J}$), it also increases structural loading (detrimental to $\gls{sym-g}$), and the larger bulk dimensions required to achieve that higher damping may produce higher stresses (detrimental to $\gls{sym-g}$) and increase the required structural material (detrimental to $\gls{sym-J}$).
    Thus, it would be sub-optimal to optimize the hydrodynamics module strictly for damping or power production followed by a separate optimization considering structures and economics.
    Furthermore, while even in a unified optimization it may be tempting to calculate the structural thicknesses within the structures module as the minimum required to sustain loads without failure, it is important to keep structural thicknesses as design variables because they also contribute to the hydrostatic constraints computed in the geometry module.
    If active, these constraints would make the system-level optimum material thickness larger than what is structurally necessary.
    Using a single optimizer to drive the design of all modules and consider all constraints addresses this non-monotonic coupling. 
\else
    The model exhibits strong design-variable coupling but no direct feedback coupling between modules, as expected from the prior observation that the only equation requiring iteration is within the dynamics module.
    This improves computational efficiency and is advantageous for optimization.
    Coupling occurs primarily through shared design variables, particularly bulk dimensions, which affect geometry, hydrodynamics, and structures, as well as generator ratings and structural thicknesses, which couple multiple subsystems.
    Feed-forward coupling is also present through inter-module dependencies.
    
    This coupling is both non-convex (e.g., larger bulk dimensions raise power but also raise structural loads and material requirements) and non-monotonic, motivating concurrent rather than sequential optimization.
\fi

%\clearpage
%!TeX root =../main.tex
%%%%%%%%%%%%%%%%%%%%%%%%%%%%%%%%%%%%%%%%%
%%%%%%%%%%%%%%%%%%%%%%%%%%%%%%%%%%%%%%%%%
\section{Module Details}
\label{sec:modules}
\ifdefined\DISSERTATION
    The five modules (geometry, hydrodynamics, dynamics and controls, structures, and economics) will now be described one at a time.
    For notation, unbolded unarrowed variables refer to scalars (e.g. $\gls{sym-x}$ or $X$), arrowed variables refer to vectors ($\gls{sym-vec-x}$ or $\vec{X}$), and bold variables refer to matrices ($\gls{sym-mathbf-x}$ or $\gls{sym-mathbf-X}$).
    For dynamic (time-varying) quantities, the time domain representation is denoted $\gls{sym-x}(\gls{sym-t})$, and can also be represented as the real part of a complex signal $\gls{sym-x}(\gls{sym-t})=\Re\left[\gls{sym-tilde-x}(\gls{sym-t})\right]$.
    Assuming the signal is sinusoidal, the complex signal can be expressed as $\gls{sym-tilde-x}(\gls{sym-t})=\gls{sym-hat-X}\exp(i\gls{sym-omega} \gls{sym-t})$ with complex phasor $\gls{sym-hat-X}=|\gls{sym-hat-X}|\exp(i\angle \gls{sym-hat-X})$ and angular frequency $\gls{sym-omega}$.
    Altogether, this gives $\gls{sym-x}(\gls{sym-t})=|\gls{sym-hat-X}|\cos(\gls{sym-omega} \gls{sym-t}+\angle\gls{sym-hat-X})$.
    Dots ($\dot{x}(\gls{sym-t})=\Re(\hat{\dot{X}}\exp(i\gls{sym-omega} \gls{sym-t}))=\Re(i\gls{sym-omega}\gls{sym-hat-X}\exp(i\gls{sym-omega} \gls{sym-t})))$ indicate time-derivatives.
    For static quantities, no special notation differentiates complex from real variables, as this will be indicated directly in the text or with $\gls{sym-mathbb-R}$ or $\gls{sym-mathbb-C}$ if clarification is needed.
\else
    The five modules (geometry, hydrodynamics, dynamics and controls, structures, and economics) will now be described one at a time.
    For notation, unbolded unarrowed variables refer to scalars, arrowed variables refer to vectors, and bold variables refer to matrices.
    Time-varying quantities $\gls{sym-x}(\gls{sym-t})$ are represented by complex phasor $\gls{sym-hat-X}$ such that $\gls{sym-x}(\gls{sym-t})=\Re[\gls{sym-hat-X}\exp(i\gls{sym-omega} \gls{sym-t})]$.
    \linelabel{lne:equation-cite-start}Equations drawn from prior work are cited where they appear; uncited equations are derived in this work, with full derivations in the appendices.\linelabel{lne:equation-cite-end}
    \Cref{sec:appendix-parameters} documents model parameters and hyperparameters.\linelabel{lne:appendix-parameters}
\fi
%TC:break Geometry
%%%%%%%%%%%%%%%%%%%%%%%%%%%%%%%%%%%%%%%%%
\subsection{Geometry}\label{sec:geom}
The geometry module computes submerged volume $\gls{sym-V-sub}$, structural volume $\gls{sym-V-struct}$, and associated hydrostatic properties from bulk dimensions.
\Cref{fig:dims} defines the principal geometric variables, where $\gls{sym-T-draft}$ denotes drafts, $\gls{sym-h}$ heights, and $\gls{sym-D}$ diameters.
Subscripts $f$, $s$, and $d$ refer to the float, spar, and damping plate, respectively.
Key points include the center of buoyancy $\gls{sym-B}$, center of gravity $\gls{sym-G}$, the keel $\gls{sym-K-keel}$, and metacenter $\gls{sym-M-metacenter}$ for the combined system.
Additional structural dimensions and thickness definitions are provided in \Cref{sec:structures}.

\begin{figure}[htbp]
\centering
    \includegraphics[\FigureOptions{aor:figs/from-matlab/dimensions.pdf}]{figs/from-matlab/dimensions.pdf}
\caption{Dimension labeling of system}\label{fig:dims}
\end{figure}

Static pitch stability is enforced by requiring:
\begin{equation}\label{eq:GM}
    \gls{sym-GM} = \gls{sym-KB} + \gls{sym-BM} - \gls{sym-KG} > 0,
\end{equation}
\noindent where $\gls{sym-BM} = \gls{sym-pi} \gls{sym-D-f}^4/(64 \gls{sym-V-sub})$ \citep{newman}.

Hydrostatic equilibrium is enforced by matching displaced water mass to total system mass, with any mismatch compensated using ballast with mass $\gls{sym-m-bal}$:
\begin{equation}
    \gls{sym-m-bal} = \gls{sym-rho-w} \gls{sym-V-sub} - \gls{sym-rho-M} \gls{sym-V-struct},
\end{equation}
where $\gls{sym-rho-w}$ and $\gls{sym-rho-M}$ are water and structural material densities.
To ensure feasible buoyancy, $\gls{sym-m-bal}$ must be non-negative and physically storable within the hull.
Assuming seawater ballast, this imposes a volume constraint:
\begin{equation}\label{eq:vol-constraint}
\gls{sym-V-struct} \le \gls{sym-V-sub} + \gls{sym-V-surf} - \gls{sym-V-pto},
\end{equation}
where $\gls{sym-V-surf}$ (design-dependent) and $\gls{sym-V-pto}$ (fixed) are volumes of the above-water structures and \gls{acr-pto} respectively.

%TC:break Hydrodynamic Coefficients
%%%%%%%%%%%%%%%%%%%%%%%%%%%%%%%%%%%%%%%%%
\subsection{Hydrodynamic Coefficients}\label{sec:hydro}
The first-order hydrodynamic force phasors $\gls{sym-vec-hat-F}$ for interacting floating bodies are expressed in the frequency domain as the sum of wave excitation $\gls{sym-vec-hat-F-e}$, hydrodynamic radiation $\gls{sym-vec-hat-F-rad}$, and hydrostatic restoring $\gls{sym-vec-hat-F-res}$ contributions:
\begin{equation}\label{eq:hydro-forces}
    \gls{sym-vec-hat-F} = \gls{sym-vec-hat-F-e} + \gls{sym-vec-hat-F-rad} + \gls{sym-vec-hat-F-res}
\end{equation}
where, under monochromatic forcing,
\begin{equation}\label{eq:hydro-forces-expanded}
\begin{aligned}
    \gls{sym-vec-hat-F-e} &= \gls{sym-vec-gamma}\,\gls{sym-hat-zeta},&~~
    \gls{sym-vec-hat-F-rad} &= -\gls{sym-mathbf-A-h}\gls{sym-vec-hat-ddot-xi} - \gls{sym-mathbf-B-h}\gls{sym-vec-hat-dot-xi},&~~
    \gls{sym-vec-hat-F-res} &= -\gls{sym-mathbf-K-h}\gls{sym-vec-hat-xi}.
\end{aligned}
\end{equation}
Here $\gls{sym-vec-hat-xi}$ is the body displacement phasor vector from the origin at the center of the float at the still waterline.
The hydrodynamic coefficients $\gls{sym-vec-gamma}$, $\gls{sym-mathbf-A-h}$, $\gls{sym-mathbf-B-h}$, and $\gls{sym-mathbf-K-h}$ are the wave excitation vector, added mass matrix, radiation damping matrix, and hydrostatic stiffness matrix, respectively.

The displacement vector $\gls{sym-vec-hat-xi}$ contains a state for each \gls{acr-dof}.
We define separate displacement and excitation vectors for the operational ($op$) and storm ($st$) cases.
A two-body heaving point absorber that rigidly locks the bodies in a storm following \citet{RM3} has vectors:
\begin{equation}\label{eq:define-vectors}
    \gls{sym-vec-hat-xi-op} = \begin{bmatrix}
        \gls{sym-hat-xi-f} \\ \gls{sym-hat-xi-s}
    \end{bmatrix}, ~~
    \gls{sym-vec-hat-xi-st}=\begin{bmatrix}\gls{sym-hat-xi-m}
    \end{bmatrix}, ~~
     \gls{sym-vec-gamma}_{op} = \begin{bmatrix}
        \gls{sym-gamma-f} \\ \gls{sym-gamma-s}
    \end{bmatrix},~~
     \gls{sym-vec-gamma}_{st} = \begin{bmatrix}
        \gls{sym-gamma-m}
    \end{bmatrix}.
\end{equation}
with subscripts for the float ($f$), spar ($s$), their coupling ($c$), and the merged float-spar ($m$).
Operational hydrodynamic matrices are symmetric $2\times 2$:
\begin{equation}\label{eq:op-hydro-coeffs}
\begin{aligned}
    \gls{sym-mathbf-A-h-op} = \begin{bmatrix}
        \gls{sym-A-f} & \gls{sym-A-c} \\ \gls{sym-A-c} & \gls{sym-A-s}
    \end{bmatrix},~
    \gls{sym-mathbf-B-h-op} &= \begin{bmatrix}
        \gls{sym-B-f} & \gls{sym-B-c} \\ \gls{sym-B-c} & \gls{sym-B-s}
    \end{bmatrix},
    \ifdefined\DISSERTATION
        ~
    \else
        \\
    \fi
    \gls{sym-mathbf-K-h-op} &= \begin{bmatrix}
        \gls{sym-K-f} & 0\\ 0 & \gls{sym-K-s}
    \end{bmatrix}
\end{aligned}
\end{equation}
while storm (merged-body) coefficients are scalar:
\begin{equation}\label{eq:st-hydro-coeffs}
\begin{aligned}
\gls{sym-vec-gamma}_{st}&=\gls{sym-gamma-m}=\gls{sym-Sigma}~\gls{sym-vec-gamma}_{op},&
~\gls{sym-mathbf-A-h-st}&=\gls{sym-A-m}=\gls{sym-Sigma}~\gls{sym-mathbf-A-h-op}, \\
~\gls{sym-mathbf-B-h-st}&=\gls{sym-B-m}=\gls{sym-Sigma}~\gls{sym-mathbf-B-h-op},&
~\gls{sym-mathbf-K-h-st}&=\gls{sym-K-m}=\gls{sym-Sigma}~\gls{sym-mathbf-K-h-op}.
\end{aligned}
\end{equation}
with $\gls{sym-Sigma}$ denoting the sum of all elements in the matrix or vector (``grand sum'').

The remaining task is to compute the float, spar, and coupling coefficients, with operational and storm implications discussed in \Cref{sec:design-load-cases}.
The hydrostatic stiffness terms ($\gls{sym-K-f}$, $\gls{sym-K-s}$) are frequency-independent and are computed directly from geometry (see \Cref{eq:gamma-K} in \Cref{sec:appendix-meem-details}), while the remaining coefficients require solving the frequency-domain radiation boundary value problem for the velocity potential.
Traditionally, this is performed using a \gls{acr-bem} solver, which discretizes the body surfaces and solves a large linear system.
In MDOcean, float coefficients are instead computed using the semi-analytical \gls{acr-meem}, exploiting cylindrical symmetry for improved efficiency.
Spar and coupling terms are obtained via algebraic approximations and interpolation of pre-computed \gls{acr-bem} data, and could be extended to \gls{acr-meem} in future work.
\Cref{tab:hydro-methods} summarizes the methods used for each coefficient.

%%%%%%%%%%%%%%%
\subsubsection{Float Coefficients: \gls{acr-meem}}\label{sec:hydro-meem}
\gls{acr-meem} solves the radiation and excitation problems by expanding the velocity potential in eigenfunctions within each fluid region and enforcing matching across region boundaries.
The \gls{acr-meem} radiation solution for two concentric heaving cylinders was first presented by \citet{mavrakos_hydrodynamic_2004} and detailed by \citet{chau_inertia_2012}.
In this approach, the fluid domain is partitioned into regions with separable analytical solutions, and continuity conditions are enforced at region interfaces to determine the unknown expansion coefficients.
The open-source implementation used here, summarized in \Cref{sec:appendix-meem-details}, follows the formulation by \citet{mccabe_open-source_2024}.
Forthcoming companion papers \citet{bimali_matrix_2026,best_openflash_2026} detail the numerical properties and code implementation respectively of an expanded implementation.

The primary advantage of \gls{acr-meem} is computational efficiency.
Timing comparisons (\Cref{sec:sim-runtime}) show more than an order-of-magnitude speedup over Capytaine \gls{acr-bem} at comparable convergence, significantly reducing the cost of multidisciplinary analysis.
Additional benefits include lower memory usage, elimination of meshing, and avoidance of numerical Green's function approximation.

Limitations include geometric restrictions and numerical overflow (discussed in \Cref{sec:appendix-meem-details}).
As a semi-analytical method, \gls{acr-meem} is restricted to simplified geometries; the formulation used here assumes a dual concentric cylinder (\Cref{fig:meem-geom}).
Consequently, float coefficients neglect the damping plate and approximate the float bottom as flat rather than slanted.
The damping plate enters through the spar coefficients (next subsection).
These approximations have minor impact and could be relaxed in future extensions using more general \gls{acr-meem} formulations \citep{olaya_hydrodynamic_2015,kokkinowrachos_behaviour_1986,bimali_matrix_2026}.

\begin{figure}[b!]
    \centering
     \includegraphics[\FigureOptions{aor:figs/from-matlab/MEEM_geometry.pdf}]{figs/from-matlab/MEEM_geometry.pdf}
    \caption[\gls{acr-meem} geometry]{\gls{acr-meem} geometry.
Dashed lines represent geometry that is neglected in the \gls{acr-meem} calculations}
   \label{fig:meem-geom}
\end{figure}

%%%%%%%%%%%%%%%
\subsubsection{Other coefficients: approximations}\label{sec:hydro-other}
Because the damping plate affects the spar coefficients more strongly than the float, the former are obtained by scaling existing solutions that include the damping plate rather than using \gls{acr-meem}.

\begin{table}[htbp]
    \centering
    \caption{Method of computing hydrodynamic coefficients in MDOcean}
    \label{tab:hydro-methods}
    \begin{tabular}
    {M{0.25\linewidth}M{0.63\linewidth}}
         Coefficient& Method\\ \hline
         $\gls{sym-A-f}$, $\gls{sym-B-f}$, $\gls{sym-gamma-f}$ & \gls{acr-meem} (\Cref{eq:hydro,eq:gamma-K} in \Cref{sec:appendix-meem-details})\\
         $\gls{sym-A-s}$& Approximate $\gls{sym-omega}\rightarrow\infty$ solution (\Cref{eq:As} in \Cref{sec:appendix-additional-hydro})\\
         $\gls{sym-A-c}$, $\gls{sym-B-s}$, $\gls{sym-B-c}$, $\gls{sym-gamma-s}$& Nominal \gls{acr-bem} solution, scaled with $\gls{sym-D-d}$ and $\gls{sym-T-s}$ (\Cref{eq:scale-wamit} in \Cref{sec:appendix-additional-hydro})\\
    \end{tabular}
\end{table}
%\hl{Describe how I did the WAMIT damping plate scaling for the spar excitation.}
%The MEEM implementation currently cannot take the damping plate into account, and the plate is expected to have a large impact on the spar added mass and damping. Therefore, algebraic approximations of the damping plate coefficients are used instead of the MEEM solution.
The spar added mass coefficient $\gls{sym-A-s}$ is approximated as a frequency-independent quantity based on the displaced water volume in the spherical projection of the damping plate, with the formula given in \Cref{sec:appendix-additional-hydro}.

The spar damping ($\gls{sym-B-s}$), excitation ($\gls{sym-gamma-s}$), and coupling terms ($\gls{sym-A-c}$, $\gls{sym-B-c}$) are obtained by interpolating a pre-computed \gls{acr-wamit} \gls{acr-bem} solution over the nondimensional parameter $\gls{sym-k-wavenumber}\gls{sym-D-d}$ (see \Cref{sec:appendix-additional-hydro}).
The excitation term is additionally scaled as $\exp(-\gls{sym-k-wavenumber}(\gls{sym-T-s} - \gls{sym-T-s-nom}))$ to account for depth-dependent attenuation.

%\hl{I suspect this imposes a constraint on $\gls{sym-D-d}$ needing to be similar to the nominal design, see notebook 3/20/26 referencing analysis from 4/4/25.}

Because float and spar coefficients are computed using different approximations, positive definiteness of the hydrodynamic matrices is not guaranteed (see \Cref{sec:appendix-additional-hydro}).
To ensure physicality, the added mass and damping matrices are required to remain positive definite.
Any violations are corrected by adjusting the coupling terms $\gls{sym-A-c}$ and $\gls{sym-B-c}$.

%  \subsubsection{Spar Dynamics}
% The RM3 is a two-body system where the relative motion of the float and spar produces power. However, capturing these underactuated multibody dynamics would require more complicated modeling and control \citep{faedo_principle_2022}, especially in the context of force saturation%, as well as implementation of the excitation force phase in MEEM.
% While updating these models is important future work, a fixed spar is assumed for this study. Thus the results can be considered more dynamically representative of a single-body point absorber or a two-body point absorber with bottom-mounted spar, as might be encountered in shallow water. Nonetheless, calculations still utilize the 45 m water depth of Humboldt Bay, and the results provide relevant preliminary design insights about a floating spar device, especially because the spar structures and economics are still considered. To ensure that the fixed spar assumption is reasonably accurate, the spar amplitude is required to remain small.

% Rather than checking this requirement by calculating the spar amplitude exactly, as this would require the exact modeling complications which the fixed-spar assumption intends to circumvent, the characteristics of second-order linear systems are used to develop a simple sufficient criteria for low spar motion. The peak response amplitude of a system with equation of motion $m\ddot{\gls{sym-xi}}+b\gls{sym-dot-xi}+k\gls{sym-xi}=\gls{sym-F}$ is:
% \begin{equation}
%     |X|_{pk} = \begin{cases}\frac{1}{2\zeta\sqrt{1-\zeta^2}} \frac{F}{m\omega_n^2} & 0<\zeta < 1\\
%     \frac{F}{m\omega_n^2} & \zeta \geq 1
%     \end{cases}
% \end{equation}
% where damping ratio is $\gls{sym-zeta}=\frac{\gls{sym-b}}{2\sqrt{mk}}$ and natural frequency is $\gls{sym-omega-n}=\sqrt{\frac{\gls{sym-k}}{\gls{sym-m}}}$. To keep the peak spar amplitude $|\gls{sym-X}|_{pk}$ low, both $\gls{sym-zeta}$ and $m\gls{sym-omega-n}^2/\gls{sym-F}=\gls{sym-k}/\gls{sym-F}$ should be kept high. A sufficient criteria for maintaining $|\gls{sym-X}|_{pk}$ below some nominal threshold $|\gls{sym-X}|_{pk,nom}$ is therefore to constrain $\gls{sym-zeta}=\gls{sym-zeta}_{nom}$ and $\gls{sym-k}/\gls{sym-F} > (\gls{sym-k}/\gls{sym-F})_{nom}$. While this is not a necessary condition (other designs could also maintain low spar amplitude), a sufficient condition is adequate for preliminary optimization purposes.

% Expressing these conditions in terms of the uncontrolled spar dynamic coefficients \citep{tao_low_2003}:
% \begin{equation}\label{eq:spar}
% \begin{aligned}
% b &= b_{s,rad}+b_{s,drag} \\
% k &= \frac{\pi}{4} \rho_w g  D_s^2 \\
% m &= \rho_w V_{sub,s}C_m\\
%     \omega_n^2 &= \frac{g}{C_m D_d} \frac{1}{\frac{T_s}{D_d} + \frac{h_d}{D_d}\left(\frac{D_d}{D_s}\right)^2} \leq \omega_{n,nom}^2 \\
%     \zeta &= abcd \geq \zeta_{nom}
% \end{aligned}
% \end{equation}
% where $\gls{sym-C-m}$ is the spar inertia coefficient.

% The geometry is required to maintain at least as high damping ratio and at least as low natural frequency as the nominal RM3 design. While this single-body spar analysis does not take into account the dynamic effect of the PTO force on the spar, it is sufficient to provide confidence that the spar response remains similar enough to the nominal design that the float can be analyzed assuming a fixed spar.

% \hl{Describe the frequency dependent fudge factor and add a figure showing that ratio as a function of sea state.} MDOcean scales the fixed-spar mechanical power by a period-dependent rational tuning factor
% \begin{equation}
%     P_{spar~floating} = \frac{c_1}{c_2T^2+c_3T+c_4} P_{spar~fixed}
% \end{equation}
% which adjusts only the power (not the amplitude). A fit of MDOcean and RM3 report yielded $[\gls{sym-c-1},\gls{sym-c-2},\gls{sym-c-3},\gls{sym-c}_4]=[23,1,-15,86]$, and the coefficients can be set to $[1,0,0,1]$ to disable the tuning.


%TC:break Dynamics and Control
%%%%%%%%%%%%%%%%%%%%%%%%%%%%%%%%%%%%%%%%%
\subsection{Dynamics and Control}\label{sec:dynamics}
This subsection introduces the dynamics and control modeling assumptions before detailing the frequency-domain methods and control formulations.

%%%%%%%%%%%%%% start old %%%%%%%%%%%%%%%%%%%%%%%%%%%%%%%%%%%
\ifdefined\DISSERTATION
    Methods for \gls{acr-wec} dynamics and control simulation in an optimization context were reviewed in \Cref{sec:lit} and include time-stepping (nonlinear, numerical, high computation cost), frequency domain (linear, analytical, low computation cost), pseudo-spectral (nonlinear, numerical, medium computation cost), \gls{acr-qlps} techniques (approximate nonlinear, semi-analytical, low computational cost), and Pontryagin's maximum principle (nonlinear, analytical, low computational cost).
    In wave energy, both spectral-domain numerical methods and analytical methods have generated recent research interest. \citet{bacelli_numerical_2015,coe_initial_2020,tan_computationally_2026,nie_optimal_2016,da_silva_stochastic_2023} study the former and \citet{abdulkadir_optimal_2024,lin_fast_2025} illustrate the latter.
    This section describes both the technical concept and the practical rationale for the choice of a \gls{acr-qlps} method, and contextualizes it with respect to other constrained optimal control methods in more detail than was possible in the literature review.

    MDOcean requires the ability to incorporate nonlinearities, at least approximately.
    Inherent nonlinearities such as drag significantly affect the amplitude and power production at resonance.
    Even with linearized intrinsic dynamics, maintaining optimality subject to dynamic constraints also requires a nonlinear controller.
    Dynamic constraints on amplitude, force, and peak power significantly influence the trade-off of power and cost, and amplitude limits are required to obtain a physically meaningful result.
    Besides nonlinearities, MDOcean also requires a low computation cost, since the dynamics will be evaluated many thousands of times during optimization.
    Of the methods mentioned, \gls{acr-pmp} and \gls{acr-qlps} techniques meet these requirements.
    \gls{acr-pmp} requires an analytically difficult derivation unique to the specific nonlinearities in question because wave energy dynamics have a particular form requiring so-called singular control \citep{zou_optimal_2017}.
    On the other hand, \gls{acr-qlps} techniques are simple and intuitive.
    Therefore, MDOcean pursues a semi-analytical \gls{acr-qlps}.
\else
    Standard linear frequency-domain \gls{acr-wec} analysis assumes unconstrained linear dynamics, enabling closed-form impedance-matched optimal control.
    However, real \glspl{acr-wec} are subject to (i) dynamic constraints (e.g., generator force, power, and stroke limits) and (ii) nonlinear forces, particularly viscous drag.
    Time-domain simulation handles both naturally but is computationally expensive, especially in an optimization loop.
    Pseudo-spectral methods are nonlinear and handle constraints by collocating the dynamics at specific time-steps, but require an outer numerical optimizer for both simulation and control synthesis.
\fi
MDOcean adopts a \gls{acr-qlps} approach that retains the speed of frequency-domain analysis while handling both temporal constraints and drag nonlinearities.
This formulation seeks to preserve the computational advantages and simplicity of classical linear analysis while extending its applicability to practical \gls{acr-wec} operating regimes that include saturation, drag, and dynamic constraints.
\gls{acr-qlps} combines two techniques: (a) describing functions, which quasi-linearize a nonlinear waveform by retaining only its fundamental harmonic,
and (b) constrained spectral optimal control, which solves the optimal control problem as a \gls{acr-qcqp} assuming linear dynamics and spectral-domain constraints.
\Cref{fig:venn-diagram} positions \gls{acr-qlps} relative to existing methods.
Augmenting spectral optimal control with describing functions extends its scope to approximate time-domain constraints and nonlinear dynamics.
Furthermore, because \gls{acr-qlps} assumptions under monochromatic forcing lead to a monochromatic problem, the \gls{acr-qcqp} becomes univariate for systems with a single control \gls{acr-dof}.
The optimal constrained linear controller can then be found analytically through a geometric interpretation that exploits the low dimensionality.
Its nonlinear counterpart can be reconstructed semi-analytically with inverse describing functions.
Avoiding iterative numerical optimization substantially reduces computational cost in large parametric studies and produces insight into the mathematical structure of dynamic tradeoffs to inform \gls{acr-pto} sizing.

\begin{figure}[htbp]
    \centering
    \includegraphics[\FigureOptions{aor:figs/from-matlab/control_method_venn_diagram.pdf}]{figs/from-matlab/control_method_venn_diagram.pdf}
    \caption[Venn diagram of control synthesis methods]{Venn diagram comparing different methods for control synthesis, with MDOcean's \gls{acr-qlps} approach at the intersection of frequency domain, nonlinear, and constrained optimal control.
    See text for acronym descriptions.}
    \label{fig:venn-diagram}
\end{figure}

\ifdefined\DISSERTATION
    \Cref{fig:venn-diagram} shows a Venn diagram comparing different methods for optimal control synthesis in greater detail.
    Typical frequency domain methods for control design, such as loop shaping, \gls{acr-pid} tuning, and $H_\infty$/$H_2$ control for general applications, along with complex-conjugate impedance matching for power-maximizing applications such as wave energy, cannot take into account constraints and nonlinearities.
    Meanwhile, typical methods for nonlinear control design, such as feedback linearization, sliding mode control, backstepping, and Lyapunov-based control, are performed in the time domain and do not incorporate constraints.
    Latching and declutching control are wave-energy-specific passive nonlinear bang-bang control methods that utilize the frequency-domain concept of phase-matching the fundamental harmonic of the response to the wave excitation, but they are still primarily time-domain methods and do not incorporate constraints \citep{hals_comparison_2011}.
    Methods for constrained optimal control, meanwhile, such as \gls{acr-mpc}, constrained \gls{acr-lqr}, and trajectory optimization are likewise time-domain and were originally derived for linear control problems, although nonlinear extensions exist including \gls{acr-nmpc}, constrained iterative \gls{acr-lqr} (iLQR), and differential \gls{acr-dp}.
    Other methods for nonlinear constrained optimal control in the time domain include the already-discussed \gls{acr-pmp} and the control barrier function method.
    Methods to incorporate nonlinearities in the frequency domain include harmonic balance synthesis, which retains full nonlinearity in a sinusoidal-steady state analysis so technically falls in the spectral domain rather than the frequency domain; its linearized counterpart, \gls{acr-df}-based synthesis \citep{atherton_nonlinear_1982,gelb_multiple-input_1968}; and the extension of the latter to broadband inputs, spectral linearization (also called stochastic linearization, Gaussian closure, and quasi-linear control) \citep{roberts_random_2003,ching_quasilinear_2010}.
    These methods are not optimal control methods and do not incorporate constraints, although it is possible to sub-optimally enforce constraints by applying a saturation nonlinearity to the constrained quantity.
    Finally, methods which combine frequency domain synthesis with constrained optimal control include constrained loop shaping (often enforcing performance constraints while minimizing the distance to an optimal loop shape, in this case the complex-conjugate impedance); linear spectral \acrfull{acr-oc}, which poses a trajectory optimization problem for a linear system represented with global basis functions in the spectral domain; constrained \acrfull{acr-im}, the narrow-band, frequency-domain equivalent of spectral optimal control; and convex control formulated using \glspl{acr-lmi}, typically in the context of robust control \citep{vanantwerp_tutorial_2000}.

    At the intersection of all three categories lies nonlinear spectral \gls{acr-oc}, pseudo-spectral \gls{acr-oc}, and the method used in MDOcean, which we call ``\acrlong{acr-qlps} \gls{acr-oc}.''
    The pseudospectral method uses the same global basis function representation as spectral \gls{acr-oc}, but instead of enforcing the dynamics at all points in the trajectory, it does so only at specific collocation points, which enables the representation of time-domain constraints and nonlinearities as opposed to harmonic constraints and nonlinearities \citep{elnagar_pseudospectral_1995}.
    Pseudo-spectral methods often employ Legendre or Chebyshev polynomial bases, although in wave energy Fourier bases are more typical.
    In this case, the method is equivalent to a nonlinear trajectory optimization problem that enforces both spectral-domain harmonic balance and time-domain constraints.
    The \gls{acr-qlps} method used in MDOcean borrows the pseudo-spectral concept of incorporating time-domain constraints in the spectral domain, but retains only the fundamental harmonic to collapse to the frequency domain.
    This is exactly analogous to the way that describing functions linearize the harmonic balance method.
    It can also be thought of as the combination of \glspl{acr-df} with constrained \gls{acr-im}.
\fi


%%%%%%%%%%%%%% end old %%%%%%%%%%%%%%%%%%%%%%%%%%%%%%%%%%%

\subsubsection{Describing Functions}
\label{sec:describing-functions}
\paragraph{Concept}

\linelabel{lne:df-concept-start}MDOcean's dynamics utilize the sinusoidal-input describing function method, which assumes a certain time-domain nonlinear signal shape and then calculates the fundamental amplitude of that signal for use in a typical frequency-domain simulation, assuming that the signal's higher harmonics are filtered out by the low-pass plant.
\linelabel{lne:df-concept-end}
\Cref{fig:desc-fcns} illustrates this concept, showing the nonlinear signals for saturation (\Cref{fig:sat}) and quadratic drag (\Cref{fig:drag-df}) along with their fundamental amplitude describing function approximations.
\begin{figure}[htbp]
\centering
\begin{subfigure}[t]{0.48\linewidth}
    \includegraphics[\FigureOptions{aor:figs/from-matlab/saturation_desc_fcn.pdf}]{figs/from-matlab/saturation_desc_fcn.pdf}
    \caption{Generator torque saturation}\label{fig:sat}
\end{subfigure}
\hfill
\begin{subfigure}[t]{0.48\linewidth}
    \includegraphics[\FigureOptions{aor:figs/from-matlab/drag_desc_fcn.pdf}]{figs/from-matlab/drag_desc_fcn.pdf}
    \caption{Quadratic drag}\label{fig:drag-df}
\end{subfigure}
\caption[Conceptual demonstration of describing function approximations]{Conceptual demonstration of describing function approximations. Plots show force versus time. Nonlinear signals in red; fundamental amplitudes in blue.}\label{fig:desc-fcns}
\end{figure}

Because the fundamental amplitude depends on the response, the frequency domain simulation is only quasi-linear and must either be iterated numerically or solved analytically.
MDOcean uses iteration to handle both the drag force and control force nonlinearities. % and an analytical solution for the force saturation nonlinearity.
While the describing function method used here has been known for over half a century, MDOcean represents its first application in a wave energy optimization or open-source simulation, to the authors' knowledge.
MDOcean therefore extends semi-analytical dynamics to wave energy systems, enabling efficient simulation while retaining model fidelity.

\paragraph{Assumptions}
\newcounter{asmpnumber}
Describing functions make the following assumptions:
\begin{itemize}[leftmargin=*]
    \item \assumption{nonlinearities-are-harmonics}
    \linelabel{lne:A1-start}
    \textbf{Any nonlinearities can be written as the sum of a quasi-linear transformation and harmonic generation}.
    This allows a spectral-domain model, rather than time-domain, while preserving nonlinearities.
    Using the nonlinearity classification structure of \citet{rijlaarsdam_comparative_2017,rijlaarsdam_frequency_2012},
    \asmpref{nonlinearities-are-harmonics} permits gain compression/expansion, harmonic generation, and frequency-mixing,
    including ``Volterra,'' ``special Weiner,'' or ``convergent'' system types.
    Eligible nonlinearities include drag, saturation, and nonlinear potential flow effects.
    \citet{folley_spectral-domain_2016,da_silva_stochastic_2023} discuss spectral-domain modeling for WECs.

    \item \assumption{system-low-pass}
    \textbf{The system is low-pass}.
    Most physical systems have low-pass behavior due to inertia, and WECs are no exception.
    Since the nonlinearities of \asmpref{nonlinearities-are-harmonics} considered here generate only odd harmonics,
    we define \asmpref{system-low-pass} as applicable when the third harmonic of an incident wave exceeds the system's cutoff frequency.
    \citet{mccabe_force-limited_2024} discuss further and demonstrate the assumption graphically, and \citet[Figure 3]{mccabe_wec_2025} show that RM3's hydrodynamic response is low-pass for frequencies above 1.2 rad/s.
    This gives a maximum wave period of 15.7~s,
    so \asmpref{system-low-pass} applies for the vast majority of incident energy, even in an irregular sea state.

    \item \assumption{quasi-linearized}
    \textbf{The response to higher harmonics can be neglected}, resulting in a linearized spectral-domain model.
    This is a direct consequence of \asmpref{nonlinearities-are-harmonics} and \asmpref{system-low-pass}.
    This assumption is called ``best linear approximation,'' ``quasi-linearization,'' ``statistical linearization,'' and ``Gaussian closure'' by \citet{rijlaarsdam_comparative_2017,ching_quasilinear_2010,roberts_random_2003,nie_optimal_2016} respectively.
    \citet{nie_optimal_2016,da_silva_statistical_2020,spanos_efficient_2016,tan_computationally_2026} pursue this assumption for WECs.

    \item \assumption{stationary-ergodic}
    \textbf{The input phase does not affect the analysis.}
    Our system meets this criteria since a linearized WEC is wide-sense stationary and the long-term average power is an ergodic output.
    Under this assumption, using a deterministic spectrally-equivalent multisine instead of a random-phase stochastic input for the wave elevation does not change the power and other dynamic outputs, so the former can be used for simplicity.
    Most of the spectral linearization references above implicitly assume this,
    since stochastic-steady-state rather than transient outputs are of interest.

    \item \assumption{waves-narrow-band}
    \textbf{The input wave excitation is narrow-band}.
    This implies that the deterministic multisine wave input can be replaced by a single sine wave;
    in other words, it approximates an irregular wave spectrum as an equivalent regular wave.
    \linelabel{lne:sidf-precise-start}The quasi-linearization method then collapses to the sinusoidal-input describing function method, where a nonlinearity is approximated simply as an amplitude-dependent gain and phase shift applied to a sinusoidal input.\linelabel{lne:sidf-precise-end}

    \item \assumption{signal-shape-known}
    \textbf{The signal shape of the nonlinear waveform is known.}
    Calculating the SIDF gain and phase shift requires specifying a signal shape and computing its fundamental \cite{gelb_multiple-input_1968,atherton_nonlinear_1982}.
    We assume a quadratic absolute-value nonlinearity for drag and a saturated sinusoid nonlinearity for control force.
    \linelabel{lne:A6-end}
\end{itemize}

\paragraph{Implementation}
The drag describing function formula is encoded directly into the simulation (see derivation in \Cref{sec:drag}),
while the control force describing function is implemented indirectly using knowledge of the limit cases of describing functions for multiple nonlinear control laws (see discussion in \Cref{sec:appendix-force-sat}).
\Cref{fig:assumptions} summarizes how these assumptions connect to DFs and to increasingly specialized \gls{acr-qlps} formulations.

\begin{figure*}[htbp]
    \centering
    \includegraphics[\FigureOptions{aor:figs/from-matlab/assumptions.pdf}]{figs/from-matlab/assumptions.pdf}
    \caption{Directed graph of assumptions for DFs and QLPS.
    Assumptions on system dynamics are yellow, assumptions on the optimal control objectives/constraints are green, and methods for dynamic analysis/control synthesis are blue.
    A single arrow implies necessity and all arrows imply sufficience.
    For example, A1 and A2 are each individually necessary but insufficient for A3; A1 and A2 are together sufficient for A3.}
    \label{fig:assumptions}
\end{figure*}

\subsubsection{Quasi-Linear Pseudo-Spectral Optimal Control}
\label{sec:qlps}
\paragraph{Motivation}
The describing functions of \Cref{sec:describing-functions} can evaluate a quasi-linear approximation of a nonlinear dynamic system but cannot directly synthesize an optimal controller.
For that, we introduce the \acrfull{acr-qlps} optimal control method.

\paragraph{Concept}
\gls{acr-qlps} combines the dynamic analysis concept of quasi-linearization with the pseudo-spectral concept of collocation to yield a simple yet effective method for constrained nonlinear optimal control synthesis.

\paragraph{Assumptions}
The most general form of \gls{acr-qlps} requires only a single assumption
\assumption{constraints-fundamental}:
\linelabel{lne:A7-start}
\textbf{the optimization objective and constraints can be written as a (potentially nonlinear) invertible function of the fundamental amplitude of linear dynamic system state variables}.
This follows directly from \asmpref{quasi-linearized} (linearized dynamics) and \asmpref{stationary-ergodic} (deterministic rather than stochastic analysis) when the states used to obtain the objective and constraints are the outputs of the dynamic analysis.
This assumption means that the peak value or other time-domain characteristic of a nonlinear signal can be constrained as long as it can be written in terms of the fundamental.

Here we formulate the \gls{acr-qlps} problem for a linear multiport dynamic system in which the state variables are ``effort'' and ``flow'' phasors related by linear coefficient matrices.
This form is well suited for expressing oscillatory power flow, although standard linear state-space representation yields equivalent mathematics.

Despite the linearized dynamics, the potential nonlinear relationship between the state variables and the objectives and constraints
makes the general \gls{acr-qlps} problem a nonlinear optimization problem (NLP).
This alone is not a significant computational benefit over the general pseudo-spectral method for nonlinear dynamic systems.

The computational benefits come when we introduce assumptions \assumption{quadratic-objective} and \assumption{quadratic-constraints},
which respectively state that \textbf{the objectives and constraints can be written as a quadratic function of the state variables}.
This is an easy requirement for wave energy, since quadratic metrics of a linear multiport effort-flow model include many quantities of interest: effort magnitude, flow magnitude, average power, reactive power, and linear combinations thereof.
Under this condition, the \gls{acr-nlp} becomes a \gls{acr-qcqp}.
The original \gls{acr-wec} pseudo-spectral formulation by \citet{bacelli_numerical_2015} conforms to this assumption and shows that under the additional assumption \assumption{passive-dynamics-convex} of \textbf{passive dynamics}, the \gls{acr-qcqp} is convex.

Even more computational benefits, and the novel aspect of the \gls{acr-qlps} approach, appear when two more assumptions are enforced.
\assumption{no-freq-coupling} and \assumption{no-dof-coupling} require that the system has \textbf{no dynamic coupling between different frequencies and different control DOFs} respectively.
\asmpref{no-freq-coupling} is satisfied as a result of \asmpref{waves-narrow-band} (narrow-band waves) and \asmpref{system-low-pass} (low-pass WEC dynamics).
\asmpref{no-dof-coupling} is satisfied because the \gls{acr-rm3} contains only a single controlled \gls{acr-dof}.
\asmpref{no-dof-coupling} permits coupling between the body \glspl{acr-dof} as long as there is a single controlled \gls{acr-dof} as is present here,
or it permits multiple controlled \glspl{acr-dof} only if they are uncoupled when viewed from each control load port at which power is to be optimized
(i.e. an array of non-interacting multi-\gls{acr-dof} \glspl{acr-wec} for which off-diagonal terms of the intrinsic impedance matrix are negligible).
\linelabel{lne:A12-end}

Under \asmpref{no-freq-coupling} and \asmpref{no-dof-coupling}, the \gls{acr-qcqp} for all frequencies and \glspl{acr-dof} can be separated into a set of single-order \glspl{acr-qcqp},
one for each frequency and \gls{acr-dof},
which can be solved semi-analytically using the geometric properties of circles.
In fact, this order reduction enables a semi-analytical solution even in the presence of certain nonlinear constraints or nonconvexities,
in violation of \asmpref{quadratic-constraints} and \asmpref{passive-dynamics-convex}.

\paragraph{Implementation}
\Cref{fig:mod-freq-domain-synthesis,fig:mod-freq-domain-evaluation,fig:mod-freq-domain-combined} distinguish three application workflows for the \gls{acr-qlps} method:
control synthesis (the design of a potentially nonlinear controller to obey constraints and maximize performance of the quasi-linearized system),
evaluation (simulation of an arbitrary controller's power performance, where the controller, plant, or both are nonlinear), and both.
MDOcean uses \gls{acr-qlps} for both control synthesis and evaluation.
\begin{figure*}[htbp]
\centering
\caption[Flowchart depicting the \gls{acr-qlps} method (M1)]{Flowchart depicting the \gls{acr-qlps} method (M1). The \acrlong{acr-df}s for response/constraint, controller, and plant are shown in yellow, pink, and blue respectively.}
\begin{subfigure}[t]{.9\linewidth}
    \includegraphics[\FigureOptions{aor:figs/from-matlab/mod_freq_domain_ctrl_synthesis.pdf}]{figs/from-matlab/mod_freq_domain_ctrl_synthesis.pdf}
    \caption{Control synthesis - applies to (M1.2) only.}
    \label{fig:mod-freq-domain-synthesis}
\end{subfigure}
\vfill
\begin{subfigure}[t]{.9\linewidth}
    \includegraphics[\FigureOptions{aor:figs/from-matlab/mod_freq_domain_ctrl_evaluation.pdf}]{figs/from-matlab/mod_freq_domain_ctrl_evaluation.pdf}
    \caption{Control evaluation - applies to both (M1.1) and (M1.2).}
    \label{fig:mod-freq-domain-evaluation}
\end{subfigure}
\vfill
\begin{subfigure}[t]{.9\linewidth}
    \includegraphics[\FigureOptions{aor:figs/from-matlab/mod_freq_domain_ctrl_combined.pdf}]{figs/from-matlab/mod_freq_domain_ctrl_combined.pdf}
    \caption{Combined synthesis and evaluation.}
    \label{fig:mod-freq-domain-combined}
\end{subfigure}
\end{figure*}

\Cref{sec:optimal-control} derives the overall optimal control problem, and \Cref{sec:appendix-qp-solution} derives the semi-analytical solution for the single-order QCQP.

\subsubsection{Approximate \glsxtrshort{acr-qlps}: Modified Frequency-Domain Methods}
\label{sec:mod-freq-domain}
\paragraph{Motivation}
Viewing DFs (\Cref{sec:describing-functions}) and \gls{acr-qlps} (\Cref{sec:qlps}) as part of the broader family of \acrfull{acr-mfd} methods provides greater context and facilitates comparison to other work.
In prior wave energy studies, related approaches have been used for evaluation, with simplified treatments of constraints (see F+ and PI+ entries in \Cref{tab:lit}).
We refer to these as the \gls{acr-mfd} family of methods, of which \gls{acr-qlps} is the most accurate and optimal.
In contrast, the present work applies \gls{acr-qlps} and other \gls{acr-mfd} methods to control synthesis, enabling nonlinear controller design within a frequency-domain formulation.
This differs from linearization approaches such as stochastic linearization, which are generally limited to linear controller design.

\gls{acr-mfd} methods besides \gls{acr-qlps} can be considerably easier to implement and may be of interest to practitioners seeking simpler alternatives to assess the effect of dynamic constraints,
or to provide a path forward when the assumptions of \gls{acr-qlps} are not fully satisfied but the computational cost of time-domain optimal control is prohibitive.

\paragraph{Concept}
When handling dynamic constraints with a general \gls{acr-mfd} method, first the system is (quasi-) linearized and the unconstrained optimal controller (and the response signals corresponding to that controller) are determined in the frequency domain.
Then the unconstrained response is checked for constraint violations.
This check can be performed in the frequency domain if the time-domain constraint can be adequately linearized (via describing functions) as a harmonic constraint, as shown in the top of \Cref{fig:mod-freq-domain-synthesis}, or the response can be reconstructed in the time domain via the inverse describing function if the constraint is not easily representable via harmonics.
If the response for a given sea state does not violate any constraints, then the unconstrained controller and corresponding power is used for that sea state as the frequency-domain linear optimal controller.
Otherwise, a variety of approaches are possible to address the violated constraint, summarized in \Cref{tab:constraint-approaches} and described below.
Methods M1--M2 can be used for both control synthesis and evaluation, while methods M4--M5 can only be used for evaluation since they do not attempt to find a controller that satisfies the constraint.
Method M3 can, in general, only be used for evaluation.
However, in the special case where the only time-domain constraint is on the control force, the method can be used for control synthesis.

\begin{table}[htbp]
\centering
\caption[Constraint satisfaction approaches for \glsxtrlong{acr-mfd} methods]{Constraint satisfaction approaches for \glsxtrlong{acr-mfd} methods. Filled stars indicate better performance.}
\label{tab:constraint-approaches}
\begin{tabular}{M{.03\linewidth} >{\raggedright\arraybackslash}m{0.45\linewidth}l l}
&\textbf{Method} & \textbf{Accuracy} & \textbf{Optimality} \\ \toprule
M1 & Nonlinear optimal (M1.1) or near-optimal (M1.2) controller, quasi-linearized via describing functions& \Stars{3} &\Stars{3} \\ \hline
M2 & Linear optimal controller& \Stars{3} & \Stars{2} \\ \hline
M3 & Saturate or zero the times with constraint violations, neglecting the effect on other signals& \Stars{2} & \Stars{1}\\ \hline
M4 & Saturate or zero entire sea state& \Stars{1} & \Stars{0.5} \\ \hline
M5 &Mark design as infeasible& N/A& \Stars{0} \\
\end{tabular}
\end{table}

% --- M1 ---
\paragraph{M1: \gls{acr-qlps} with nonlinear control}
The \gls{acr-mfd} approach with the highest accuracy and optimality (M1) is to synthesize a new nonlinear optimal (or near-optimal) time-domain controller that obeys the constraint, taking into account the linearized effect of that nonlinear controller on the frequency domain response.
This synthesis can occur in a number of ways.
If a time-domain optimal control law (e.g. those that \citet{abdulkadir_optimal_2024,lin_fast_2025} develop via \gls{acr-pmp}) is available for the constraint under consideration (M1.1), it can be used as the nonlinear optimal controller directly, essentially bypassing the use of \gls{acr-qlps} in the control synthesis step (\Cref{fig:mod-freq-domain-synthesis}) and using it only for evaluation (\Cref{fig:mod-freq-domain-evaluation}).
\ifdefined\DISSERTATION
    This differs from the standard \gls{acr-pmp} method in that the evaluation of the response and power corresponding to the time-domain \gls{acr-pmp} controller is determined in the frequency domain, rather than via numerical time-integration (as in the method of \citet{zou_optimal_2017}) or analytical time-integration (as in the method of \citet{lin_fast_2025}).
    This is advantageous if the computational cost or mathematical labor of the numerical or analytical \gls{acr-pmp} evaluation methods, respectively, are of concern.
\fi
For situations without an explicit \gls{acr-pmp} control law (M1.2), a nonlinear near-optimal controller can be constructed by first deriving the optimal linear controller via constrained optimization
(M2) and then adjusting it to be nonlinear using insights from signal saturation and filtering.
M1.2 should be considered the primary formulation of \gls{acr-qlps}, and the other \gls{acr-mfd} methods (M2)--(M5) in \Cref{tab:constraint-approaches} act as approximations of \gls{acr-qlps} with varying levels of accuracy and optimality.
M1.1 and M1.2 both linearize the nonlinear controller via the describing function method.
In addition to the control nonlinearity, a describing function is also utilized to linearize the plant drag nonlinearity.

% --- M2 ---
\paragraph{M2: \gls{acr-qlps} with linear control}
If a nonlinear controller or the corresponding describing function to linearize it are not available, the most optimal linear controller that obeys the constraint can be used instead (M2).
This approach accurately evaluates the maximum power that a linear controller could produce, although this power will be lower (less optimal) than that of the optimal nonlinear controller.

\ifdefined\DISSERTATION
    The constrained optimization to identify the controller in (M1.2) and (M2) can be performed numerically or analytically by analyzing an impedance-mismatched linear system in the frequency domain and will be discussed in detail in \Cref{sec:optimal-control}.
    The current MDOcean software uses the numerical approach, which requires more computation time but is simpler to implement and equally accurate.
    \Cref{sec:optimal-control} derives the analytical approach, while \citet{mccabe_force-limited_2024} and \Cref{sec:appendix-qp-solution} provide details and visualizations.
\else
    The constrained optimization used to identify the controller in (M1.2) and (M2) can be performed numerically or analytically by analyzing an impedance-mismatched linear system in the frequency domain;
    the analytical derivation is provided in \Cref{sec:optimal-control}, with additional details in \Cref{sec:appendix-qp-solution}.
\fi

% --- M3 ---
\paragraph{M3: Saturate or Zero Violations}
If neither a numerical nor an analytical implementation is available to identify the optimal constrained linear controller, then the violating portion of the time-domain signals can be saturated or zeroed for the sake of power calculation (M3).
This approach neglects the corresponding effect of the saturation/zeroing on other quantities in the system, which lowers the accuracy, and the controller does not have the opportunity to even approximately consider the constraint, resulting in a power production potentially much lower than the constrained optimal.
\ifdefined\DISSERTATION
    This simple approach does not require any derivations to implement.
    It has been used in several previous \gls{acr-wec} optimization studies for various constraints \citep{garcia-teruel_design_2022,garcia-teruel_reliability-based_2021,cotten_multi-objective_2022,mccabe_constrained_2013} and is evidently the most common \gls{acr-mfd} method in the marine energy field.
\else
    This simple approach has been used in several previous \gls{acr-wec} optimization studies for various constraints in the studies \citep{garcia-teruel_design_2022,garcia-teruel_reliability-based_2021,cotten_multi-objective_2022,mccabe_constrained_2013}
    and is the most common \gls{acr-mfd} method in the marine energy field.
\fi

% ---- M4 -- M5 ---
\paragraph{M4--M5: Penalize Violations}
As a last resort, the energy in that sea state can be zeroed (M4),
essentially assuming that the device enters survival mode in those cases to avoid constraint violation, or the entire design can be marked infeasible with a constraint in the optimization (M5), which \citet{mccabe_constrained_2013} pursues for a radiation limit constraint.
These approaches do not attempt to find dynamic solutions which satisfy the limits, and essentially deal with violations by penalizing the optimization through either the objective or constraints.

%Approach (M5) is also pursued here for amplitude limits in the operational design load case, in the interest of consistency and implementation expedience.

\paragraph{Method Choice for Each Nonlinearity}\label{sec:mfd-method-choice}
\begin{table}[htbp]
    \centering
    \caption{Dynamic nonlinearities and approach for each}
    \label{tab:nonlinearities}
\begin{tabular}{L{0.25\linewidth}L{0.7\linewidth}}
         \textbf{Dynamic nonlinearity/limit}& \parbox{\linewidth}{\centering\textbf{Method}}\\ \hline
         Drag& Describing function (M1.2) for $\sin(\gls{sym-omega} \gls{sym-t})|\sin(\gls{sym-omega} \gls{sym-t})|$ with iteration to find effect on response and optimal controller\\
         Force limit& Describing function (M1.2) for $\text{sat}(\sin(\gls{sym-omega} \gls{sym-t}))$ with iteration to find effect on response and optimal controller\\
        Operational amplitude limit& Optimal linear controller (M2)\\
         Power limit& Average value of $\text{sat}(\sin(\gls{sym-omega} \gls{sym-t})+\gls{sym-a})$, neglecting any effect on response and optimal controller (M3)\\
        Storm amplitude limit&Optimizer considers limit violations as infeasible designs (M5)\\
    \end{tabular}
\end{table}
\Cref{tab:nonlinearities} summarizes the dynamic nonlinearities and limits and the method MDOcean currently uses to handle each.

For the control force limit, the corresponding describing function (saturated-sine) is simple to implement, and prior \gls{acr-rm3} optimizations \citep{mccabe_multidisciplinary_2022,gaebele_tpl_2025,mcgilton_optimal_2024} show the constraint to be frequently active and impact results strongly, justifying the use of M1.2 (full \gls{acr-qlps}).
For amplitude limits in the operational design load case, M1.2 produces the same result as M2 due to the low-pass assumption that makes the constrained quantity a pure sinusoid, so we use the simpler M2 without loss of optimality.

For a power limit, under cases where the power limit is the only constraint, applying M2 is straightforward, and applying M1.2 is possible but more cumbersome:
assuming that the optimal power signal is a saturated sinusoid implies that the corresponding force signal is the quotient of the truncated power sinusoid and the amplitude sinusoid, which requires deriving an additional describing function.
When a power limit is applied simultaneously with a force limit, both methods break down.
M2 would be incompatible with the previous choice of M1.2 for the co-occurring force limit, and M1.2 creates ambiguity about the nonlinear signal shapes,
since it is not possible for the force and power signals to both be saturated sinusoids while assuming sinusoidal amplitude.
To avoid this additional complexity, the power limit uses M3, although future work should consider the use of M1.2 or M2 to improve accuracy and optimality.

For slamming amplitude limits in the storm design load case, approach M5 is actually the only option. 
This is because the device is already in survival mode (preventing method M4), is not producing power (preventing methods M3 and M4), and is not applying a control input (preventing methods M1 and M2).

%%%%%%%%%%%%%%%
\FloatBarrier
\subsubsection{Linearized Dynamics}\label{sec:eom}
\ifdefined\DISSERTATION
    Before applying the \gls{acr-qlps} method described above, we first present the equations of motion for a standard linear \gls{acr-wec} model, starting in the time domain and quickly moving to the frequency domain.
    Specifically, we combine the bi-conjugate network model of \citet{coe_co-design_2025} with the underactuated multi-\gls{acr-dof} model of \citet{faedo_principle_2022} to arrive at a formulation suitable for the impedance matching of a multibody \gls{acr-wec} subject to wave forcing, powertrain kinematics and dynamics, and drag.
\else
    Before applying the \gls{acr-qlps} method described above,
    we first present the dynamics for a standard linear \gls{acr-wec},
    combining the bi-conjugate network model of \citet{coe_co-design_2025} with the underactuated multi-\gls{acr-dof} model of \citet{faedo_principle_2022} to express the equations of motion for a multibody \gls{acr-wec} subject to wave forcing, powertrain kinematics and dynamics, and drag.
    \linelabel{lne:single-freq}Note that \Cref{sec:eom,sec:optimal-control,sec:drag} consider a single monochromatic sea state; \Cref{sec:energy} performs annual aggregation.
\fi

\paragraph{Multiport Structure}
The dynamics portion of \Cref{fig:design-scope} represents this system as a multiport circuit, with effort 
variables $\gls{sym-vec-e}$ and flow variables $\gls{sym-vec-q}$ at each port.
Ports include the wave input, body mechanical, \gls{acr-pto} mechanical, generator mechanical, and generator electrical.
Between each set of ports are the intrinsic body dynamics, \gls{acr-pto} kinematics, drivetrain dynamics, and generator dynamics respectively.
The overall relationship between the \gls{acr-pto} mechanical port and generator electrical port is collectively referred to as the \gls{acr-pto} dynamics.

\Cref{sec:appendix-multiport} derives three matrix forms to describe various aspects of the port relationships:
impedance ($\gls{sym-mathbf-Z}$) groups all ports together and is used for the intrinsic body dynamics;
cascade ($\gls{sym-mathbf-a}$/\gls{acr-abcd}) transmits efforts and flows sequentially between port groups and is used for the \gls{acr-pto} dynamics;
hybrid ($\gls{sym-mathbf-h}$) mixes the two and is used for the \gls{acr-pto} kinematics.
Full conventions and conversion formulas are given in \Cref{sec:appendix-multiport-conventions}.

\paragraph{Floating Body Intrinsic Dynamics}
The float and spar are each modeled as a floating rigid body, and the two bodies are kinematically coupled to transmit force between them.
Forces include wave excitation, hydrodynamic radiation, hydrostatic restoring, drag, and power takeoff.
The excitation, radiation, and restoring forces are linear; the drag and \gls{acr-pto} forces are generally nonlinear but represented linearly in the multiport model assuming they have been quasi-linearized.
\Cref{sec:appendix-body-dynamics} provides the equations and derivations for the intrinsic body dynamics, which encompasses the body inertia and all external forces except the \gls{acr-pto} forces.

%%%%%%%%%%%%%%%
\paragraph{\glsxtrshort{acr-pto} Model}
\label{sec:pto-model-overview}
The \gls{acr-pto} model is a linear multiport effort-flow model that is valid for multiple \gls{acr-dof} devices across various energy domains.
It consists of \gls{acr-pto} kinematics
(a hybrid matrix relating body \glspl{acr-dof} to the \gls{acr-pto} port, accommodating underactuation)
and \gls{acr-pto} dynamics
(cascade matrices for a mechanical drivetrain and permanent magnet generator,
relating the \gls{acr-pto} port to the generator electrical port,
with generator voltage $\gls{sym-hat-V}$ and current $\gls{sym-hat-I}$ as the effort and flow variables).
The Thévenin equivalent circuit method is used to find the effective intrinsic impedance $\gls{sym-Z-s-th}$ and open-circuit source voltage $\gls{sym-hat-V-s-th}$ viewed from the generator electrical port, as required for impedance matching.
Details and circuit diagrams are provided in \Cref{sec:appendix-pto-kinematics,sec:appendix-combine-intrinsic-kinematics,sec:appendix-pto-dynamics}.

%%%%%%%%%%%%%%%
\paragraph{Power Production}
\label{sec:power}
The instantaneous time-domain power at any port is simply the product of the port effort and flow variables, $\gls{sym-p}(\gls{sym-t}) = \gls{sym-e-t}~\gls{sym-q-t}$.
\ifdefined\DISSERTATION
    Substituting $\{\gls{sym-e},\gls{sym-q}\}=\{\gls{sym-F-p},\gls{sym-dot-xi}\}$ yields the floating body mechanical power,
    $\{\gls{sym-e},\gls{sym-q}\}=\{\gls{sym-F-PTO},\gls{sym-dot-X-PTO}\}$ yields the mechanical power at the \gls{acr-pto} port,
    $\{\gls{sym-e},\gls{sym-q}\}=\{\gls{sym-tau},\gls{sym-Omega}\}$ yields the generator mechanical power,
    and $\{\gls{sym-e},\gls{sym-q}\}=\{\gls{sym-V},\gls{sym-I}\}$ yields the generator electrical power.
\else
    The same expression is evaluated at the floating body, \gls{acr-pto}, generator mechanical, and generator electrical ports by substituting the corresponding effort/flow pair $\{\gls{sym-e},\gls{sym-q}\}$.
\fi

In the frequency domain assuming sinusoidal waveforms, the active power $\gls{sym-P}$, reactive power $\gls{sym-Q}$, complex power $\gls{sym-S}$, and apparent power $|\gls{sym-S}|$ from electrical engineering \citep{saadat_power_1999} are
\begin{equation}\label{eq:power-PQS}
\begin{aligned}
    \gls{sym-P} &= \tfrac{1}{2}\Re (\gls{sym-hat-e}\gls{sym-hat-q}^*) = \tfrac{1}{2}\Re (\gls{sym-Z-term})|\gls{sym-hat-q}|^2,\\
    \gls{sym-Q} &= \tfrac{1}{2}\Im (\gls{sym-hat-e}\gls{sym-hat-q}^*) = \tfrac{1}{2}\Im (\gls{sym-Z-term})|\gls{sym-hat-q}|^2, \\
    \gls{sym-S} &= \gls{sym-P} + i\gls{sym-Q} = \tfrac{1}{2}\gls{sym-hat-e}\gls{sym-hat-q}^* = \tfrac{1}{2}\gls{sym-Z-term}|\gls{sym-hat-q}|^2, \\
    |\gls{sym-S}| &= \tfrac{1}{2}|\gls{sym-Z-term}||\gls{sym-hat-q}|^2,
\end{aligned}
\end{equation}
where the port is terminated with complex impedance $\gls{sym-Z-term}$ and $(\,)^*$ denotes the conjugate transpose.

The average and peak powers over a wave period, derived in \Cref{sec:appendix-power-time}, are
\begin{equation}\label{eq:power-avg-peak}
    \gls{sym-p-avg} = \gls{sym-P}, \qquad p_{pk} = \gls{sym-P} + |\gls{sym-S}|.
\end{equation}
The corresponding minimum power is $\gls{sym-p}_{\text{min}} = \gls{sym-P} - |\gls{sym-S}|$, negative when $|\gls{sym-S}|>\gls{sym-P}$ (i.e., reactive control).

To capture non-dynamic power losses such as those in an electrical drive, a constant efficiency $\gls{sym-eta}$ describes the gain between average generator electrical power and average absorbed electrical power.
\ifdefined\DISSERTATION
    If the generator cascade matrix is not used (i.e.\ $\gls{sym-mathbf-a-tau-Omega-leftarrow-VI} = \gls{sym-mathbf-I}$), this efficiency can also approximate losses in the generator itself.
    Note that this is not strictly equivalent to assuming the generator has constant efficiency, since controllers with reactive terms ($\gls{sym-K-l}\neq 0$) spend some time expending energy, which reverses the direction of the generator efficiency.
    The average efficiency is assumed to also apply to the peak powers, though this is not guaranteed in general for the same reason.
\fi
Validation in \Cref{sec:validation-benchmarking} sets $\gls{sym-eta}=80\%$ to match the \gls{acr-rm3} design assumption of \citet{RM3}.

The full expansion of $\gls{sym-p-avg-elec}$ and $\gls{sym-p-pk-elec}$ in terms of the optimal-control variables $\gls{sym-B-l}$, $\gls{sym-K-l}$, and the electrical-port current phasor $\gls{sym-hat-I}$ is given in \Cref{sec:appendix-power-elec}.

% Additionally, the drivetrain mechanical cascade matrix $[\gls{sym-a}]_{\gls{sym-m}}$ is set to the identity matrix, which represents a static effective gear radius $\gls{sym-R-gear}=1$:
% \begin{equation}
% [a]_m = \begin{bmatrix}
% 1/\gls{sym-R-gear} &~ 0 \\
% 0 &~ \gls{sym-R-gear}
% \end{bmatrix} = \begin{bmatrix}
% 1 &~ 0 \\
% 0 &~ 1
% \end{bmatrix}
% \end{equation}
% With this $[\gls{sym-a}]_{\gls{sym-m}}$, the control damping and stiffness $\gls{sym-B-u}$ and $K_u$ equal the PTO damping and stiffness $\gls{sym-B-p}$ and $\gls{sym-K-p}$.
% This will simplify the impedance matching procedure described below, requiring matching at only one interface.

%The time-domain mechanical power is $\gls{sym-P}_{\text{mech}}(\gls{sym-t}) = \gls{sym-tau}(\gls{sym-t}) \gls{sym-Omega}(\gls{sym-t}) = (\gls{sym-B-u} \gls{sym-Omega}(\gls{sym-t}) + \gls{sym-K-u} \gls{sym-theta}(\gls{sym-t}) ) \gls{sym-Omega}(\gls{sym-t}) $, where $\gls{sym-theta}(\gls{sym-t})$ is the generator rotational position. Assuming sinusoidal waveforms, the time-average and peak mechanical power in a given sea state are then

%%%%%%%%%%%%%%%
\subsubsection{Optimal Control}\label{sec:optimal-control}
This subsection presents the unconstrained and constrained optimal-control solutions used in the model,
starting with the linearized dynamics of \Cref{sec:eom} and building up in complexity to arrive back at the \gls{acr-qlps} formulation laid out in \Cref{sec:qlps}.

\paragraph{Linearized, Unconstrained}
In the absence of constraints, maximizing power at a given port requires impedance matching at that port due to the maximum power transfer theorem.
We apply matching at the generator electrical port, so the electrical load impedance $\gls{sym-Z-l}$ must be the complex conjugate of the Th\'{e}venin equivalent source impedance $\gls{sym-Z-s-th}$ seen from that port.
This is called reactive control, with subscript $\text{reac}$:
\begin{equation}\label{eq:matched-load}
    \gls{sym-Z-l-reac} = \gls{sym-Z-s-th}^*
\end{equation}

However, some \glspl{acr-pto} lack reactive capability (the ability to function as a motor and add energy to the system at some times), such as rectifying hydraulic systems or winches where cable tension must be maintained.
Avoiding reactive power ($\gls{sym-Q}=0$) at the electrical port restricts the system to so-called damping control, with subscript $\text{damp}$, and requires $\gls{sym-K-l-damp}=0$.
The power-maximizing \gls{acr-pto} damping $\gls{sym-B-l-damp}$ is then:
\begin{equation}\label{eq:damping-control}
    \gls{sym-B-l-damp} = \sqrt{ (\gls{sym-B-l-reac})^2 + (\gls{sym-K-l-reac}/\gls{sym-omega})^2}
\end{equation}
\ifdefined\DISSERTATION
    When a port other than the electrical port must avoid reactive power, the constrained solution detailed later in this section can be used to enforce $\gls{sym-Q}=0$ at that port.

    Both of these controllers frequently result in excessively large forces and thus structural and powertrain requirements, especially in energetic sea states or in reactive control when $\gls{sym-K-l}$ is far from zero.
    $\gls{sym-K-l-reac}$ is often a large negative number because cost minimization favors a low mass, and the negative \gls{acr-pto} stiffness reduces the natural frequency to better match the ocean wave frequency with a smaller device mass.
    Control laws \eqref{eq:matched-load} and \eqref{eq:damping-control} provide no mechanism to reduce this torque, so below we formulate a control law with the ability to obey constraints.
\else
    Both controllers frequently produce excessively large forces in energetic sea states or under reactive control, and provide no mechanism to obey limits.
    We therefore formulate a control law that handles constraints below.
\fi

\paragraph{Linearized, Constrained}
Now we incorporate constraints on the linearized system, yielding the method we term ``constrained impedance matching.''
We first compile expressions for all metrics $\gls{sym-underline-M}$ which must be constrained.
\ifdefined\DISSERTATION
    For the time being, we keep these expressions generic so they apply to any port.
    To facilitate the calculations in transmission form, we write the metrics as the quadratic product of the port variable vector with a suitable matrix $\gls{sym-mathbf-Q-M}$:
\else
    \linelabel{lne:quadratic-metrics-start}
    Metrics such as active power, reactive power, port variable magnitude, and linear combinations thereof can be formulated as a quadratic product of the complex port variable vector with a matrix $\gls{sym-mathbf-Q-M}$:
    \linelabel{lne:quadratic-metrics-end}
\fi
\begin{equation}\label{eq:metrics}
    \gls{sym-underline-M} = \frac{1}{4} \begin{bmatrix}
    \gls{sym-hat-e} \\ \gls{sym-hat-q}
    \end{bmatrix}^*
    \gls{sym-mathbf-Q-M}
    \begin{bmatrix}
    \gls{sym-hat-e} \\ \gls{sym-hat-q}
    \end{bmatrix}
\end{equation}
The matrices for various metrics are given in \Cref{tab:power-metrics}.
\begin{table}[bhp]
\centering
\caption[Matrices for calculating port metrics]{Matrices for calculating various metrics at a port via \Cref{eq:metrics}}
\label{tab:power-metrics}
\begin{tabular}{L{0.1\linewidth}|ccccc}
    Metric $\gls{sym-underline-M}$ & $\gls{sym-p}_{\text{avg}}=\gls{sym-P}$ & $\gls{sym-Q}$ & $\gls{sym-S}$ & $|\gls{sym-hat-e}|^2$ & $|\gls{sym-hat-q}|^2$ \\ \hline
    Matrix $\gls{sym-mathbf-Q-M}$ &
    $\begin{bmatrix}
        0 & 1 \\
        1 & 0
    \end{bmatrix}$
    &
    $\frac{1}{i}
    \begin{bmatrix}
        0 & -1 \\
        1 & 0
    \end{bmatrix}$
    &
    $\begin{bmatrix}
        0 & 0 \\
        2 & 0
    \end{bmatrix}$
    &
    $\begin{bmatrix}
        1 & 0 \\
        0 & 0
    \end{bmatrix}$
    &
    $\begin{bmatrix}
        0 & 0 \\
        0 & 1
    \end{bmatrix}$
\end{tabular}
\end{table}

Note that the apparent power $|\gls{sym-S}|$, and therefore the peak power from \Cref{eq:power-avg-peak}, cannot be written in the quadratic form of \Cref{eq:metrics} and is instead fourth-order:
\begin{equation}
    |\gls{sym-S}|^2 = \gls{sym-S} \gls{sym-S}^* 
    = \frac{1}{16}
    \begin{bmatrix} \gls{sym-hat-e} \\ \gls{sym-hat-q} \end{bmatrix}^*
    \gls{sym-mathbf-Q-S}
    \begin{bmatrix} \gls{sym-hat-e} \\ \gls{sym-hat-q} \end{bmatrix}
    \begin{bmatrix} \gls{sym-hat-e} \\ \gls{sym-hat-q} \end{bmatrix}^*
    \gls{sym-mathbf-Q-S}^{\top}
    \begin{bmatrix} \gls{sym-hat-e} \\ \gls{sym-hat-q} \end{bmatrix}
\end{equation}
Peak power constraints use a more complicated approach discussed in \Cref{sec:appendix-mag-S-constraints} and are left out of the formulation below for simplicity.

To implement the constraints, we create an intentional linear impedance mismatch, sacrificing power to satisfy the limits (method M2 in \Cref{tab:constraint-approaches}).
For select constraints where a describing function is available (currently the torque limit), we additionally use nonlinear control to satisfy the limit with less of a power sacrifice (method M1.2 in \Cref{tab:constraint-approaches}).

For the linear impedance mismatch, the average electrical power from \Cref{eq:power-avg-peak} can be maximized subject to quadratic constraints on the effort and flow variables.
This is a \gls{acr-qcqp}:
\begin{equation}\label{eq:opt-problem}
\begin{aligned}
& \max_{\gls{sym-vec-x} \in \gls{sym-mathbb-C}^2} & \gls{sym-p-avg-VI} &= \frac{1}{4}\gls{sym-vec-x}^{\,*} \gls{sym-mathbf-Q-P} \gls{sym-vec-x} \\[3pt]
& \text{\gls{sym-s}.\gls{sym-t}.}
& \gls{sym-vec-x}^{\,*}\gls{sym-mathbf-Q-mu}\gls{sym-vec-x}
    + 2\Re\!\left\{\gls{sym-vec-a-mu}^{\,*}\gls{sym-vec-x}\right\}
    &\le \gls{sym-b-mu},
    \quad \forall \gls{sym-mu} \in [1,\cdots,\gls{sym-N}] \\[3pt]
& & \gls{sym-vec-c}^{\,*}\gls{sym-vec-x} &= d \\[3pt]
& \text{where}
& \gls{sym-vec-x} &= \begin{bmatrix} \gls{sym-hat-V} & \gls{sym-hat-I} \end{bmatrix}^{\top} \\[3pt]
& & \gls{sym-mathbf-Q-mu} &\in \gls{sym-mathbb-C}^{2\times 2}, \quad \gls{sym-mathbf-Q-mu} = \gls{sym-mathbf-Q-mu}^{\,*} \\[3pt]
& & \gls{sym-vec-a-mu} &\in \gls{sym-mathbb-C}^2 \\[3pt]
& & \gls{sym-b-mu} &\in \gls{sym-mathbb-R} \\[3pt]
& & \gls{sym-vec-c}^{\,*} &= \begin{bmatrix} 1 & \gls{sym-Z-s-th} \end{bmatrix} \\[3pt]
& & d &= \gls{sym-hat-V-s-th}
\end{aligned}
\end{equation}

The linear equality constraint $\gls{sym-vec-c}^{\,*}\gls{sym-vec-x} = d$ enforces the \gls{acr-wec} and \gls{acr-pto} dynamics through a voltage-current relationship at the generator electrical port, while the quadratic inequality constraints encode limits on linear or quadratic combinations of effort and flow at any port.
\ifdefined\DISSERTATION
    The specific quantities constrained in this case are shown in \Cref{tab:qp-constraints}.
    More precisely, this simulation implements constraints on the maximum allowed generator torque and mechanical power, float and spar body amplitudes, and \gls{acr-pto} amplitude, as well as a positivity constraint on the electrical power.
    In addition to the amplitude limits shown in the first three rows of \Cref{tab:qp-constraints}, a differently-structured quadratic amplitude constraint will be presented in \Cref{sec:design-load-cases} (see \Cref{eq:slam-constraint-quadratic}) representing a coupled amplitude-phase criterion for the bodies to not exit the water or be fully submerged.
\else
    Constraints implemented in MDOcean include generator torque, mechanical power, float and spar amplitudes, \gls{acr-pto} amplitude, and electrical power positivity.
    Quadratic coefficients $\{\gls{sym-mathbf-Q-mu}, \gls{sym-vec-a-mu}, \gls{sym-b-mu}\}$ for each are given in \Cref{tab:qp-constraints} (\Cref{sec:appendix-qp-constraints-table}).
    A differently-structured quadratic amplitude constraint for float and spar slamming/submersion is presented in \Cref{sec:design-load-cases}.
\fi

The solution to \Cref{eq:opt-problem} can be written as:
\begin{equation}\label{eq:constrained-qp-solution}
\begin{aligned}
    \gls{sym-vec-x-opt} &=
    \frac{\gls{sym-hat-V-s-th}}{2\Re (\gls{sym-Z-s-th})}
        \begin{bmatrix} Z_{s,th}^* (1+\gls{sym-Gamma-opt}) \\ 1-\gls{sym-Gamma-opt} \end{bmatrix} \\
    \gls{sym-p-avg-VI-opt} &= \frac{|\gls{sym-hat-V-s-th}|^2}{8\Re (\gls{sym-Z-s-th})}~ (1 - |\gls{sym-Gamma-opt}|^2)
\end{aligned}
\end{equation}
where $\gls{sym-Gamma-opt}$ is a complex scalar with magnitude typically between 0 and 1 called the optimal reflection coefficient, with formula given in \Cref{eq:gamma-opt-quadratic,eq:candidate-points,eq:circle-intersection} of \Cref{sec:appendix-qp-solution}.
The reflection coefficient represents the degree to which the constraints introduce an impedance mismatch.
$\gls{sym-Gamma-opt}=0$ corresponds to the impedance-matched solution (no active inequality constraints).

\ifdefined\DISSERTATION
    As \Cref{sec:mod-freq-domain} describes, the solution for $\gls{sym-Gamma-opt}$ can be found numerically or analytically, and we further distinguish that the numerical approach can be accomplished with either direct or indirect optimal control.
    Direct methods solve the optimization problem itself, while indirect methods first derive the optimality conditions and then solve those conditions.
    MDOcean currently implements both the analytical approach, which corresponds to the explicit architecture in \Cref{fig:control-arch}, and the numerical indirect approach, which corresponds to the implicit architecture in \Cref{fig:control-arch}.
    The analytical approach is described in \Cref{sec:appendix-qp-solution}, and the numerical indirect approach in \Cref{sec:appendix-qp-numerical}.

    The numerical direct approach can correspond to either the nested architecture or the simultaneous architecture of \Cref{fig:control-arch}.
    While the nested architecture solves the quadratic program \Cref{eq:opt-problem}, the simultaneous architecture integrates the \gls{acr-qcqp} into the nonlinear design optimization, which does not take advantage of the efficiency available with the \gls{acr-qcqp} formulation.
    It is unknown whether the numerical direct approach to solving the \gls{acr-qcqp} would be more efficient than the numerical indirect approach, but both are expected to be significantly less efficient than the analytical approach, so the analytical approach is used for all results and benchmarks presented here.
\else
    Both analytical and numerical-indirect solutions for $\gls{sym-Gamma-opt}$ are implemented in MDOcean.
    This study focuses on the superior analytical approach described in \Cref{sec:appendix-qp-solution}, which has a geometric interpretation as a minimum-distance problem on the complex plane of $\gls{sym-Gamma-opt}$ and is used for all results presented here.
    \citet[Appendix A.3.3]{mccabe_dissertation_2026} describes the numerical-indirect approach.
\fi

\paragraph{Constrained, QLPS with Linear Control (M2)}
The previous section on constrained impedance matching describes how to achieve optimal power transfer for a linear dynamic system subject to quadratic constraints.
Our system is quasi-linear, so the linear constrained impedance matching approach is applied iteratively with updated quasi-linear dynamic coefficients and updated constraint limits to approximate the optimal solution.
\Cref{sec:appendix-dynamic-limit} covers the relationship between quasi-linear constraint limits and time-domain constraint limits for each constraint type.
Then, we implement an iterative loop around the dynamics and control module to update the quasi-linear coefficients (in the case of linear control, the only nonlinear coefficient is drag) and constraint limits until convergence is achieved.

\paragraph{Constrained, QLPS with Nonlinear Control (M1.2)}
\ifdefined\DISSERTATION
    For the nonlinear controller, the maximum generator torque constraint is considered separately from the others.
    This constraint changes the optimal torque signal from a pure sinusoid to a saturated waveform, approaching a square wave for sufficiently low torque limit.
    As is, this nonlinear signal could not be handled with a frequency-domain transfer function and would require time-domain simulation.
    Conveniently, through the describing function method, the saturated signal can be represented with Fourier analysis as a superposition of harmonics.
    Furthermore, because the plant is a second order oscillator and filters out high frequency inputs, harmonics other than the fundamental contain very little energy and can be safely neglected.
    \Cref{fig:desc-fcns}(a) shows this approximation graphically for a saturated-sine force signal.
\else
    To extend to a nonlinear controller, the nonlinear coefficients now include both drag and the \gls{acr-pto} coefficients, and there is an additional relationship needed between the quasi-linear torque limit and the time-domain torque limit.
\fi

Details of the calculation are provided in \Cref{sec:appendix-force-sat}.
The result is updated effective values for control gains $\gls{sym-B-l}$ and $\gls{sym-K-l}$ that represent the nonlinear saturated controller with a torque that obeys the constraint and is used in place of \Cref{eq:constrained-qp-solution}.

%%%%%%%%%%%%%%%
\subsubsection{Drag}\label{sec:drag}
\ifdefined\DISSERTATION
    The hydrodynamics model of \Cref{sec:hydro} does not capture all fluid forces on the body due to its linear, irrotational, and inviscid assumptions.
    More specifically, viscous shear stresses (skin friction),
    higher-order nonlinear potential flow effects, and
    steady and unsteady viscous separation phenomena are all neglected.
    In this section, we model the form drag due to viscous flow separation, including steady wake losses, turbulence, and cycle-averaged vortex shedding.
    \citet{quartier_influence_2021} finds that these effects are the dominant drag force for heaving \glspl{acr-wec} and are important to avoid unrealistically high resonant peaks in the response amplitude which would lead to overestimates of power production.
\else
    The linear, inviscid hydrodynamics model of \Cref{sec:hydro} omits viscous flow separation, which is the dominant drag mechanism for heaving \glspl{acr-wec} \citep{quartier_influence_2021}.
    Without modeling drag, the predicted resonant amplitude is unrealistically high and power is overestimated.
\fi

% Maha question: should we just jump to the result instead of deriving p_d approximation? (5/18/26)
Drag force $\gls{sym-vec-F-d}$ is proportional to the product of the relative velocity of the \gls{acr-wec} and incident wave, $\gls{sym-vec-v-rel}(\gls{sym-y},\gls{sym-t})=\gls{sym-vec-dot-xi}(\gls{sym-t})-\gls{sym-vec-v-wave}(\gls{sym-y},\gls{sym-t})$, with its absolute value.
The resulting quadratic form $|\gls{sym-vec-F-d}|\sim|\gls{sym-vec-v-rel}|^2$ implies a nonlinear input-output response that would normally require time-domain simulation.
To improve computational efficiency and allow solution in the frequency domain, the nonlinearity is quasi-linearized with a describing function, following \citet{quartier_influence_2021}.
The pressure along the bottom surface of each body due to drag $\gls{sym-vec-p-d}$ is modeled with empirical drag coefficient vector $\gls{sym-vec-C-d}$:
\begin{equation}\label{eq:drag-pressure}
    \gls{sym-vec-p-d} = \frac{1}{2}\gls{sym-rho-w} \gls{sym-vec-C-d} ~\gls{sym-vec-v-rel} (\gls{sym-y},\gls{sym-t}) |\gls{sym-vec-v-rel} (\gls{sym-y},\gls{sym-t})|
\end{equation}
\Cref{eq:drag-pressure-approx} expands the velocity term then applies the describing function approximation:
\begin{equation}\label{eq:drag-pressure-approx}
\begin{aligned}
    \left(\gls{sym-vec-v-rel} (\gls{sym-y},\gls{sym-t}) ~\cdot
    \ifdefined\DISSERTATION
    \else
        \right.\\
        \left.  
    \fi
    \left|\gls{sym-vec-v-rel} (\gls{sym-y},\gls{sym-t}) \right|\right)
    &= |\gls{sym-vec-hat-v-rel} (\gls{sym-y})|^2 \cos (\gls{sym-omega} \gls{sym-t}+\angle \gls{sym-vec-hat-v-rel} (\gls{sym-y}))
    \left| \cos (\gls{sym-omega} \gls{sym-t}+\angle \gls{sym-vec-hat-v-rel} (\gls{sym-y}))\right|\\
    &\approx~|\gls{sym-vec-hat-v-rel} (\gls{sym-y})|^2\frac{8}{3\gls{sym-pi}} \cos (\gls{sym-omega} \gls{sym-t}+\angle \gls{sym-vec-hat-v-rel} (\gls{sym-y}))
\end{aligned}
\end{equation}
The relative velocity uses the standard linear-wave velocity phasor evaluated at body draft; the explicit expression is in \Cref{sec:appendix-drag} (\Cref{eq:wave-velocity-phasor}).

The time-domain nonlinear waveform and describing function approximation for the drag pressure/force are compared in \Cref{fig:desc-fcns}(b).
The amplitude of the sinusoidal approximation equals $\frac{8}{3\gls{sym-pi}}\approx0.85$ times the peak of the nonlinear waveform.

Integrating the drag pressure over the wetted surface via strip theory (\Cref{sec:appendix-drag}) yields a quasi-linearized drag force consisting of a damping term in phase with the body motion and an excitation term potentially out of phase:
\begin{equation}\label{eq:drag-damping-excitation}
    \gls{sym-vec-hat-F-d} = -\gls{sym-mathbf-B-d} \gls{sym-vec-hat-dot-xi}
    + \gls{sym-vec-gamma-d} \gls{sym-hat-zeta}
\end{equation}
The excitation term arises from the use of the relative velocity rather than direct \gls{acr-wec} velocity, which is often overlooked in \gls{acr-wec} models and has important implications for the phase of the drag force.
The damping-excitation grouping of the drag terms, as well as the dependence of the relative velocity on the direction of wave propagation $\gls{sym-y}$, improves on the approach of \citet{quartier_influence_2021} and is discussed more in \Cref{sec:appendix-drag}.
Explicit expressions for $\gls{sym-mathbf-B-d}$ and $\gls{sym-vec-gamma-d}$ in terms of nondimensional strip-theory integrals over body geometry are given in \Cref{sec:appendix-drag} (\Cref{eq:drag-coeffs}).

These numerical integrals are precomputed and stored as a lookup table to support fast optimization.
The resulting magnitude and phase of \gls{acr-wec} motion are obtained by numerically converging the state-dependent coefficients $\gls{sym-mathbf-B-d}$ and $\gls{sym-vec-gamma-d}$.
\Cref{sec:appendix-drag} discusses details and control implications.

%\subsubsection{Power Saturation}
%Equation~\eqref{eq:mech-power} shows that the power peak to maximum ratio is at minimum two for a damping controller,
%\hl{describe}

%%%%%%%%%%%%%%%
\subsubsection{Energy Production}
\label{sec:energy}
The analysis thus far either is sea-state agnostic (\Cref{sec:geom}) or applies separately to each individual sea state (\Cref{sec:hydro,sec:dynamics}).
We now combine the results across sea states to find the long-term energy statistics in a given wave environment.

\ifdefined\DISSERTATION
    While matrices in \Cref{sec:hydro,sec:dynamics} represent coupled relationships between vector output variables (i.e.\ $\gls{sym-vec-hat-xi}$) for systems with multiple bodies or \glspl{acr-dof}, in the current section each output element is treated as a scalar, either representing a single element of an output vector (i.e.\ $\gls{sym-hat-xi-f}$) or the sum across all \glspl{acr-dof} (i.e.\ power in systems with multiple actuated \glspl{acr-dof}).
    These scalars are assembled into ``matrices'' (i.e.\ $\gls{sym-boldsymbol-hat-xi-H-T-f}$) indexed by sea state rather than by \gls{acr-dof}, indicated by the $\gls{sym-H},\gls{sym-T-period}$ superscript.
    Such ``matrices'' are incompatible with matrix multiplication and any matrix operations in this section should be interpreted as elementwise operations across sea states.
    To distinguish from the average and maximum signal values over a wave period denoted with avg and pk subscripts in \Cref{sec:power}, the long-term average across sea states is denoted with an overbar $\overline{(\cdot)}$ and long-term maximum with subscript $\text{max}$.
    This intentionally matches the $\text{max}$ subscript used in \Cref{tab:qp-constraints} for upper limits within each sea state because MDOcean currently applies the same limit value for a given constraint across all sea states, though this need not be the case in general.
    An exception is the slamming and submergence constraints, which will be discussed separately in \Cref{sec:design-load-cases}.
\else
    Quantities indexed by sea state are written as ``power matrices'' $(\cdot)^{\gls{sym-H},\gls{sym-T-period}}$ (e.g.\ \gls{sym-mathbf-P-H-T-elec}),
    with operations interpreted elementwise across sea states rather than as matrix multiplication.
    Long-term averages across sea states are denoted with an overbar $\overline{(\cdot)}$ and long-term maxima with subscript $\text{max}$,
    distinguishing them from the within-period averages and peaks (avg, pk) introduced in \Cref{sec:power}.
    A separate quadratic amplitude limit for slamming and submergence, which varies across sea states, is treated in \Cref{sec:design-load-cases}.
\fi

To find the long-term average power production \gls{sym-overline-P-elec} in a location with a given distribution of sea states, the power matrix \gls{sym-mathbf-P-H-T-elec} for each sea state is weighted by that sea state's probability using a \gls{acr-jpd} matrix and then summed.
This method is illustrated in \Cref{fig:JPD-multiply}.
The \gls{acr-jpd} used in this analysis represents wave conditions in Humboldt Bay, \gls{acr-ca} and is taken from the study by \citet{janzou_system_2020}.

\begin{figure}[b!]
\centering
\includegraphics[\FigureOptions{aor:figs/manual/jpd_mult.pdf}]{figs/manual/jpd_mult.pdf}
\caption{Power matrix multiplication}
\label{fig:JPD-multiply}
\end{figure}

The annual energy production $\gls{acr-aep}$ in kWh/year for a farm of multiple devices is the product of the number of devices, the average power per device, an efficiency $\gls{sym-eta-array}$ representing transmission losses and array downtime, and an appropriate unit conversion:
\begin{equation}\label{eq:AEP}
    \gls{acr-aep} = \gls{sym-N-WEC} ~ \gls{sym-overline-P-elec} ~\gls{sym-eta-array} \frac{8766}{1000}.
\end{equation}

%Finally, the standard deviation $\gls{sym-sigma}$ of power across all sea states $\gls{sym-P}(\gls{sym-omega})$ is computed, and normalized by the average power $\gls{sym-mu}$ to get the coefficient of variation (see equation \eqref{cv}).

%%%%%%%%%%%%%%%
\subsubsection{Dynamic Limits and \glspl{acr-dlc}}
\label{sec:design-load-cases}
This study considers two design load cases: (1) cyclic operational loading in sea states where power is produced, and (2) storm loading in sea states where the device enters survival mode, where the \gls{acr-pto} locks and stops producing power.
See \Cref{sec:appendix-dynamic-limit} for further discussion and derivation.

\paragraph{Relating Peak Forces and Amplitudes to Fundamental}
In the operational load case,
the peak \gls{acr-pto} force is found from the fundamental assuming that the \gls{acr-pto} force signal shape is a saturated sine wave and the \gls{acr-pto} force constructively interferes with shaft torque (a conservative assumption).
The peak fluid force and peak displacement are assumed to equal their respective fundamental amplitudes.
This is a reasonable approximation for the displacements due to the low-pass nature of the system, and could be an under- or over-estimate for the fluid force.

In the storm case, the peak is assumed to equal the fundamental amplitude for the \gls{acr-pto} force, fluid force, and displacements.
The assumption applies here for the \gls{acr-pto} force since the storm regime uses a brake not a generator, so there is no saturation nonlinearity.
\ifdefined\DISSERTATION
    Note that the analysis of this section continues to use standard linear wave theory with equivalent regular waves,
    even for the storm condition.
    This is a questionable assumption,
    but it is necessary because optimization-suitable models for nonlinear wave forces do not currently exist.
    \Cref{sec:discussion} will discuss possible future alternatives to this modeling choice.
\else
    Both load cases use standard linear wave theory with equivalent regular waves,
    a questionable approximation whose limitations and possible future alternatives are discussed in \Cref{sec:unmodeled-effects}.
\fi

\paragraph{Permissible Force and Amplitude Limits}
Permissible amplitudes must meet the following criteria:
\begin{itemize}
\item Geometric clearance between the float and spar tops and bottoms to prevent overtravel
\item Geometric clearance of the tubular structure connecting the float and spar
\item Linear wave-body interaction to maintain compatibility with the linear potential flow theory approach of \Cref{sec:hydro}
\item Prevent slamming, where the bottom surface rises out of the water and impacts on reentry
\item Prevent submersion, where the top of the \gls{acr-wec} becomes fully submerged
\end{itemize}

\ifdefined\DISSERTATION
    The values $\gls{sym-h-fs-clear}$, $\gls{sym-h-fs-up}$, $\gls{sym-h-fs-down}$,
    $\gls{sym-xi-f-lin}$, and $\gls{sym-xi-s-lin}$
    are all upper limits and are independent of sea state,
    while $\gls{sym-xi-slam}$ depends on the sea state
    and can separately act as both an upper and lower limit on the oscillation magnitude,
    as shown in \Cref{sec:appendix-slam}.
\else
    The first two limits are static and independent of sea state;
    the third limit is static but depends on the sea state;
    the last two limits are dynamic (depend on the body motion as well as the sea state).
    The last two limits, together called the free surface constraints, ensure that the \gls{acr-wec} surface does not transiently pierce the wave free surface since doing so would cause extremely high loads.
    Each free surface constraint can separately act as both an upper and lower limit on the oscillation magnitude.
    These limits depend on the phase of the response, not just the magnitude, so they do not use the quadratic coefficients of \Cref{tab:qp-constraints} but can nonetheless be expressed in a quadratic form compatible with \Cref{eq:opt-problem}.
\fi

\ifdefined\DISSERTATION
    In early design studies that do not directly compute the hydrodynamic coefficients and associated power production,
    the amplitude limits proposed here could be plugged into Budal's upper bound,
    a common analytical expression that combines an upper bound on the excitation force with the maximum permissible body velocity
    to yield an upper bound on power production \citep{zou_practical_2023}.
\fi

Permissible forces are determined based on structural criteria to be explained in \Cref{sec:structures}.

% Todo: This \creftab:sea-state} is \hl{possibly out of date now:}
% \begin{table}[htbp]
% \centering

% \begin{tabular}{| l  |L{0.35\linewidth}|L{0.35\linewidth}|}
% \hline
% \textbf{Value} & \textbf{Sea state to use} & \textbf{Saturation to meet it?} \\
% \hline
% $\gls{sym-F}_{\text{heave},\text{storm}}$& Storm sea state & Not relevant (controller is off during storm) \\
% \hline
% $\gls{sym-F}_{\text{heave},op}$& Operating sea states with nonzero JPD & Ideally, but requires more work \\
% \hline
% $\gls{sym-F}_{pto}$& Operating sea states with nonzero JPD & Analytical \\
% \hline
% $\gls{sym-X}_{\text{linear}}$& Operating sea states with nonzero JPD & Ideally, but requires more work \\
% \hline
% $\gls{sym-X}_{\text{slam}}$& Storm sea state & Not relevant (controller is off during storm) \\ \hline
%  $\gls{sym-X}_{fs}$& Operating sea states with nonzero JPD&Ideally, but requires more work\\\hline
% $\gls{sym-P}_{\text{max}}$& Operating sea states with nonzero JPD & Ideally, but requires more work \\
% \hline

% \end{tabular}
% \caption{Sea states}
% \label{tab:sea-state}
% \end{table}

%%%%%%%%%%%%%%%
\subsubsection{Irregular Waves}\label{sec:irregular-waves}
MDOcean uses a \gls{acr-jpd} matrix to capture the long-term variation in wave conditions and accurately estimate annual average power and fatigue load.
Additionally, it represents each irregular wave condition within the \gls{acr-jpd} as an equivalent sinusoidal regular wave.
Regular wave height $\gls{sym-H}$ and period $\gls{sym-T-period}$ are selected to provide equal energy as the irregular sea state:
\begin{equation}\label{eq:regular-wave-equiv}
    \gls{sym-H} = \frac{\gls{sym-H-s}}{\sqrt{2}},
    \qquad
    \gls{sym-T-period} = \gls{sym-T-e} \approx 0.857 ~\gls{sym-T-p},
\end{equation}
where the irregular sea state has significant wave height $\gls{sym-H-s}$, energy period $\gls{sym-T-e}$, and (for a Pierson-Moskowitz spectrum) peak period $\gls{sym-T-p}$.
This enables the analytical solution of the optimal control problem.
% moved to appendix
Further discussion is in \Cref{sec:appendix-narrowband}.

\ifdefined\DISSERTATION
    Nonetheless, irregular time-domain transient peaks and the variation of hydrodynamic coefficients over the spectral width of each individual sea state are not considered.
    Therefore, MDOcean likely underestimates the sensitivity to peak force, power, and amplitude constraints.
    \Cref{sec:discussion} will discuss possible future extensions to incorporate irregular wave effects into the analysis.
\else
    Transient peaks and spectral coefficient variation within sea states are not captured, so MDOcean likely underestimates sensitivity to peak force, power, and amplitude constraints.
    Stochastic linearization is a candidate for future extension; see \Cref{sec:unmodeled-effects}.
\fi

%TC:break Structures
%%%%%%%%%%%%%%%%%%%%%%%%%%%%%%%%%%%%%%%%%
\FloatBarrier
\subsection{Structures}\label{sec:structures}
The MDOcean structures module uses forces from the dynamics module as well as area and volume outputs from the geometry module to calculate the factors of safety for various structural failure modes.
The factors of safety represent the multiplicative increase in force before stress would surpass peak limits for storm loads and the endurance limit under operational loads.

\ifdefined\DISSERTATION
    Again aligning with the modeling philosophy of investing the up-front development time to create a simulation that runs extremely quickly with moderate accuracy, MDOcean integrates a combination of explicit algebraic and tabulated semi-analytical structural models that are well-established in the structures community yet rarely utilized for \gls{acr-wec} design.
    The structures module therefore runs significantly faster than traditional \gls{acr-fem} simulations, though it lacks the ability to easily generalize to other geometries, boundary conditions, and load cases as a \gls{acr-fem} model could.

    More specifically, the peak limits analyzed include yield, global buckling, local buckling, and ultimate where appropriate, while the endurance limit assesses high-cycle fatigue.
    The latter is potentially conservative, and future work could implement Miner's law to better represent infrequent operational loads (low-cycle fatigue) \citep{ove_arup__partners_ltd_structural_2016}.
\else
    Peak-limit checks include yield, global buckling, local buckling, and ultimate where appropriate;
    the endurance limit assesses high-cycle fatigue.
\fi

The factor of safety is directly calculated for the most heavily loaded subcomponent in each of the three components: the bottom plate in the float, the cylindrical shell of the spar, and the annular disc of the damping plate.
Thicknesses of other subcomponents are either held constant or scaled with the dimensions of subcomponents that are directly assessed, as outlined in \Cref{tab:struct-calc-vs-scale}.
\ifdefined\DISSERTATION
    In the future, extending the structures module to directly calculate the factor of safety for each subcomponent would allow for more realistic structural design.
\fi

\begin{table}[t!]
    \centering
    \caption{Structural analysis methodology by subcomponent}
    \label{tab:struct-calc-vs-scale}
\begin{tabular}{M{0.15\linewidth} M{0.4\linewidth} M{0.3\linewidth}}
         \textbf{Component}&  \textbf{Structural Subcomponent}&  \textbf{Method of Analysis}\\ \hline
         \multirow[c]{5}{*}{\centering Float}&  \begin{itemize}\item Bottom trapezoidal plate\end{itemize}&  Directly calculated\\
         &  \begin{itemize}\item Top trapezoidal plate\end{itemize}&  Scaled from bottom plate\\
         &  \begin{itemize}\item Radial rectangular plate\end{itemize}&  Scaled from bottom plate\\
        & \begin{itemize}\item Circumferential rectangular plates\end{itemize}& Scaled from bottom plate\\
        & \begin{itemize}\item \gls{acr-pto} connection tubes\end{itemize}& Held constant\\ \hline
        Spar& Cylindrical shell& Directly calculated\\ \hline
        \multirow[c]{2}{*}{\centering Damping plate}& \begin{itemize}\item Annular plate\end{itemize}& Directly calculated\\
        & \begin{itemize}\item Spar connection tubes\end{itemize}& Held constant\\
    \end{tabular}
    \end{table}

\Cref{fig:FBDs} shows a free body diagram of the applied load on each of the three structural units and the locations where they are modeled as structurally fixed.

\begin{figure}[t!]
\centering
\includegraphics[\FigureOptions{aor:figs/from-matlab/FBD.pdf}]{figs/from-matlab/FBD.pdf}
\caption{Applied loads and fixed points of each structure}\label{fig:FBDs}
\end{figure}

Each of the twelve float sections is modeled as a stiffened shell subject to distributed bottom loading and top fixity at the tube support weld joint.
\ifdefined\DISSERTATION
    To simplify the representation of load transfer between plate elements, the bottom %and top
    float plate is analyzed individually, assuming fixed edges and solving for the reaction loads at these edges.
    This means that the %top and bottom plate are each
    plate is conservatively sized assuming the other plate elements are perfectly rigid. %The fixed-edge reaction loads from these top and bottom solutions are then applied to the central shell element modeled as a cross section subject to bending.
\else
    The bottom plate is analyzed individually with fixed-edge boundary conditions, conservatively sizing it as if the other plate elements were perfectly rigid.
\fi

The cylindrical shell of the spar is modeled as a short column under compression and hydrostatic hoop stress.
The column's short length requires an intermediate solution between pure compression and Euler buckling, a consideration that the original \gls{acr-rm3} design did not account for \citep{previsic_reference_2011}.

The damping plate is a thin stiffened annular plate subject to a vertical distributed force on its bottom surface, reacted by the welded connection to the column along the column's circumference and by four welded tubular supports.
\ifdefined\DISSERTATION
    Two analytical solutions for a thin annular (cylindrical with central hole) plate fixed at its inner radius are superimposed: a distributed pressure and four point loads at the edge, with appropriate modifications to account for the stiffeners.
\else
    Two analytical solutions for a fixed-inner-radius annular plate (distributed pressure and four point loads at the edge) are superimposed, with modifications for the stiffeners.
\fi

Further details of all structural calculations performed in the simulation are provided in \Cref{sec:appendix-structures}.
For illustrative purposes only, assuming a fixed maximum stress and considering only dominant terms for geometries similar to the nominal design, the scaling of the required structural thicknesses with bulk dimensions simplifies to the following:
\begin{equation}\label{eq:struct-scaling}
\begin{aligned}
    \gls{sym-t-f-b} h_{\text{stiff},f}^2 &\sim \frac{F_{\text{heave}} \gls{sym-D-f}}{\left(\frac{\gls{sym-D-f}}{\gls{sym-D-s}}-1\right)^2} \\
    \gls{sym-t-s-r}^4 &\sim F_{\text{heave}} \gls{sym-h-s}^2 \\
    \gls{sym-t-d} h_{\text{stiff},d}^2 &\sim \frac{F_{\text{heave}}\gls{sym-D-d}^2}{\gls{sym-D-s}} \\
\end{aligned}
\end{equation}

\ifdefined\DISSERTATION
    Since structural material scales linearly with both the thickness $\gls{sym-thick}$ and the stiffener height $\gls{sym-h-stiff}$, the squaring of stiffener heights in \Cref{eq:struct-scaling} means that increasing the stiffener height is generally more structurally efficient than increasing the thickness.
    We therefore expect that an optimization would potentially drive the thicknesses to their minimum values, and thus it is also required to model the minimum value of $\gls{sym-thick}/\gls{sym-h-stiff}$ for which the stiffened plate assumptions hold.
    These limits are also derived in \Cref{sec:appendix-structures}.
\else
    Because structural material volume scales linearly in both $\gls{sym-thick}$ and $\gls{sym-h-stiff}$, the squared dependence on stiffener height in \Cref{eq:struct-scaling} implies that taller stiffeners are more material-efficient than thicker plates.
    Optimization therefore tends to drive thicknesses toward the minimum value of $\gls{sym-thick}/\gls{sym-h-stiff}$ for which the stiffened-plate assumptions hold;
    these limits are derived in \Cref{sec:appendix-structures}.
\fi

%TC:break Economics
%%%%%%%%%%%%%%%%%%%%%%%%%%%%%%%%%%%%%%%%%
\subsection{Economics}\label{sec:econ}
The common economic metric $LCOE$ is calculated as the ratio of annualized expenditures to the annual energy production $\gls{acr-aep}$:
\begin{equation}\label{eq:lcoe}
    LCOE = \frac{CAPEX \cdot \acrshort{acr-fcr} + OPEX}{\gls{acr-aep}}
\end{equation}
where the numerator consists of the up-front capital expenditure $CAPEX$, annualized via a fixed charge rate $\gls{acr-fcr}$, plus an annual operational expenditure $OPEX$.
\ifdefined\DISSERTATION
    The fixed charge rate incorporates the effect of device lifetime as well as a variety of financial values affecting the time-value of money such as tax rate, interest rate, depreciation, and how much of the project is financed with debt (loans) compared to equity (investors).
\fi

\ifdefined\DISSERTATION
    The MDOcean economics model aims not to predict the precise cost breakdown of a given design but to accurately reflect the scaling of cost with design variables to facilitate design optimization and tradeoff analysis.
    With this in mind, the model lumps costs into categories that are either constant or scale with design, and the constant costs and scale factors are tuned to match more detailed cost breakdowns.
    The cost data chosen for tuning is the \gls{acr-rm3} \gls{acr-cbs}, first published by \citet{neary_reference_2014} and reproduced in \Cref{tab:CBS} of \Cref{sec:appendix-econ}.
\else
    The MDOcean economics model prioritizes capturing the scaling of cost with design variables over predicting absolute cost.
    Constant and design-dependent cost categories are calibrated against the \gls{acr-rm3} \gls{acr-cbs} \citep{neary_reference_2014} (\Cref{tab:CBS} in \Cref{sec:appendix-econ}).
\fi

Two \gls{acr-cbs} categories scale with design:
structural cost (category 1.4) scales linearly with the volume of structural material $\gls{sym-V-struct}$,
and \gls{acr-pto} cost (category 1.5) scales with the peak electrical power $\gls{sym-p-pk-elec}$ and the force limit $\gls{sym-F-max}$.
All other categories are held constant per device but decrease with the number of devices $\gls{sym-N-WEC}$
via a power law calibrated against the \gls{acr-cbs} estimates for
$\gls{sym-N-WEC} = (1, 10, 50, 100)$ \citep{neary_reference_2014}.

\ifdefined\DISSERTATION
    Some cost components assumed constant in the model would scale with design in reality, and modifying the model to reflect this is an area for future work.
    For example, installation and decommissioning costs (categories 1.7 and 1.8) could scale with the device maximum dimension since a larger device requires a larger ship, but this is not included due to the difficulty of determining dimension thresholds.
    Likewise, the mooring and foundation costs (category 1.3) could scale with storm surge force, and the replacement parts and consumables costs (categories 2.5 and 2.6) could scale with the \gls{acr-pto} power and force as the \gls{acr-pto} capital costs do.
\else
    Several categories that would scale with design in reality (installation by device size, mooring by storm surge force, replacement parts and consumables by \gls{acr-pto} ratings) are held constant here and flagged in \Cref{sec:unmodeled-effects} as future work.
\fi

The chosen scalings are a power law for the number of devices to reflect economies of scale, direct proportionality between structural cost and material volume, and constant plus linear for the \gls{acr-pto} cost since the main components of the \gls{acr-pto} (generator, drivetrain, and structural support) are generally sized linearly with power and torque (\Cref{sec:appendix-econ}).
    %the cheaper of the induction and permanent magnet machines is chosen with a torque threshold:
    % \begin{equation}
    %     C_{gen} = \begin{cases}
    %     80.86 \tau + 1032, & 0~\leq\tau <47~\text{Nm, ~~~~induction generator} \\
    %     9.151 \tau + 4386, & 47\leq\tau \leq 4000~\text{Nm, permanent magnet generator}
    %     \end{cases}
    % \end{equation}
This yields:
\begin{equation}\label{eq:unit-cost}
\begin{aligned}
    \gls{sym-C-unit}   &= \gls{sym-C-pto} + \gls{sym-C-struct} + \gls{sym-C-fixed} \\
    \gls{sym-C-pto}    &= \gls{sym-C-pto-constant} + \gls{sym-p-P} \gls{sym-P-max} + \gls{sym-p-F} \gls{sym-F-max} \\
    \gls{sym-C-struct} &= \gls{sym-p-s} \gls{sym-V-struct}
\end{aligned}
\end{equation}
where $\gls{sym-p-P}$, $\gls{sym-p-F}$, and $\gls{sym-p-s}$ are the prices of power, force, and structural material in units of \linelabel{lne:2012-prices}2012\$/kW, 2012\$/N, and 2012\$/m$^3$ respectively.
See \Cref{sec:appendix-econ-power-law} for more details.

The per-\gls{acr-wec} unit costs and prices $\gls{sym-C-fixed}$, $\gls{sym-C-pto-constant}$, \gls{sym-p-P}, \gls{sym-p-F}, and \gls{sym-p-s} all decrease with the number of devices $\gls{sym-N-WEC}$ via a power law (\Cref{eq:cost-power-law} in \Cref{sec:appendix-econ}), with curve-fit parameters given in \Cref{tab:econ-model-values}.

The total \gls{acr-capex} and \gls{acr-opex} are found by multiplying the unit costs by the number of devices $\gls{sym-N-WEC}$.
For consistency with prior references \citep{neary_reference_2014,RM3}, all costs are intentionally kept in units of 2012 \gls{acr-usd} without adjusting for inflation.

%\clearpage
%%%%%%%%%%%%%%%%%%%%%%%%%%%%%%%%%%%%%%%%%
%%%%%%%%%%%%%%%%%%%%%%%%%%%%%%%%%%%%%%%%%
\section{Validation and Benchmarking}
\label{sec:validation-benchmarking}
This section demonstrates MDOcean's accuracy and speed through verification (checking that the code reliably implements the specified assumptions), calibration (tuning the model to replicate specified outputs), validation (checking that the specified assumptions are realistic), and benchmarking (measuring model runtime).
The majority of checks are implemented as unit tests in a continuous integration workflow reproducibility pipeline following software engineering best practices.

\subsection{Verification, Calibration, and Validation: Hydrodynamics Module}
\label{sec:hydro-checks}
To verify and validate MDOcean's hydrodynamics module, we compare against analytically expected behavior, previously published MEEM results for various geometries, and existing BEM solutions for RM3 \cite{RM3}.
Hydrodynamic calibration is not necessary.

\subsubsection{Hydrodynamics Verification}
MDOcean's results satisfy analytically expected behaviors for a set of 6,250 geometries representing a full-factorial design sweep.
We confirm that:
damping goes to zero at high and low frequency;
added mass is logarithmically singular at low frequency (16.8\% of geometries have at least one frequency where a non-finite value is capped) and constant at high frequency;
added mass and damping matrices are positive definite
(6.9\% of geometries require truncating the coupling terms as described in \Cref{sec:appendix-additional-hydro});
potential and radial velocity match at the boundary of adjacent fluid regions
(see \cite[Fig.~7]{mccabe_open-source_2024});
and
hydrodynamic coefficients have converged with respect to chosen series truncation order
(see \cite[Fig.~6]{mccabe_open-source_2024}).
Additionally, potential and hydrodynamic coefficients match previously published MEEM results for a reference geometry
(see \cite[Fig.~4]{mccabe_open-source_2024}).

\subsubsection{Hydrodynamics Validation}
Hydrodynamic validation for the float, spar/coupling, and system radiation impedance shows that MDOcean agrees well with \gls{acr-bem} and damping-plate-inclusive \gls{acr-meem}.
\ifdefined\DISSERTATION
    The \gls{acr-meem} hydrodynamic coefficients for the nominal \gls{acr-rm3} geometry match \gls{acr-wamit} \gls{acr-bem} results closely, as previously shown in Fig.~5 of \citep{mccabe_open-source_2024}, except at frequencies below 0.3~rad/s.
    This is expected because the \gls{acr-bem} solution assumes infinite depth, which requires a finite added mass at low frequency, whereas \gls{acr-meem} uses finite depth, where the added mass must grow logarithmically toward infinity at zero frequency \citep{mciver_added_1991}.
    For the site under consideration, the lowest-frequency sea state containing any energy is 0.4~rad/s, so the discrepancy is not a concern.
    After the publication of reference~\citet{mccabe_open-source_2024}, excitation phase was added to the \gls{acr-meem} implementation.
    It also matches \gls{acr-wamit} results well.
\else
    Validating the assumption that the slanted surfaces of the float can be approximated as flat,
    the \gls{acr-meem} hydrodynamic coefficients for the nominal \gls{acr-rm3} geometry match \gls{acr-wamit} \gls{acr-bem} results closely for $\gls{sym-omega} \geq 0.3$~rad/s (Fig.~5 of \citet{mccabe_open-source_2024}).
    The divergence at lower frequencies arises from differing finite-depth versus infinite-depth assumptions \citep{mciver_added_1991} and is not a concern for the site considered (minimum sea-state frequency 0.4~rad/s).
\fi
Meanwhile, the spar and coupling coefficients for the RM3 geometry match the \gls{acr-wamit} \gls{acr-bem} results perfectly due to the interpolation scheme.
Therefore, spar and coupling coefficient validation is performed with the data used to determine dimensional scalings \citep{olaya_hydrodynamic_2015}.
%\Cref{fig:olaya-validation} shows the results.
At the system level, it is the magnitude and phase of the system radiation impedance that determines how hydrodynamic errors propagate to power and amplitude.
The third subfigure of \Cref{fig:error-histogram} shows the corresponding power and amplitude errors after running the coefficients through the dynamics and controls module, turning drag off to eliminate the effect of any dynamics linearization errors.
Prior studies have successfully validated \gls{acr-meem} \cite{} and \gls{acr-bem} \cite{} coefficients (and their assumptions of linear wave theory) against experimental response data.

%%%%%%%%%%%%%%%%%%%%%%%%%%%%%%%%%%%%%%%

\subsection{Verification, Calibration, and Validation: Dynamics and Control Module}
To verify, calibrate, and validate MDOcean's dynamics and controls module, we compare against analytical solutions, the software WEC-Sim \citep{ruehl_wec-simwec-sim_2024}, and to the reference model report power matrix \cite{RM3}.

\subsubsection{Dynamics Verification}
Numerical checks confirm agreement with expected behavior:
the power, motion, and forces for any \gls{acr-wec} with an unconstrained controller match the analytical impedance-matched solution;
the power output for a \gls{acr-wec} with no drag, lossless \gls{acr-pto}, and unconstrained control does not exceed the radiation limit;
the power output for a \gls{acr-wec} with no drag, lossless \gls{acr-pto}, and amplitude-constrained control does not exceed the Budal upper bound;
and
the response for any \gls{acr-wec} with constrained control does not violate constraint limits.

The first subfigure in \Cref{fig:error-histogram} shows that the power and motion amplitude for \gls{acr-rm3} with no drag, lossless \gls{acr-pto}, and a reactive scheme
match those predicted by WEC-Sim with identical hydrodynamic inputs,
verifying correct implementation of the dynamics for a linear system.
Promisingly, the mode of the histogram is less than one percent error,
although a minority of sea states exhibit substantially larger errors.
These large errors vanish when drag is turned on or when control is switched to a damping scheme.
We theorize that this represents an error in the WEC-Sim results, rather than in MDOcean, arising from transient normal modes that have not had enough time to decay in the highly-undamped system.
Future work should confirm this hypothesis by extending the time-domain simulation period past the current 200~s duration.
Because the same percent error can have disproportionate impact in different sea states,
vertical lines indicate the most relevant summary statistic for each metric:
maximum float amplitude and annual average power (\resultsAOR[wecsimAvgPowerErrorBestCase] error). %(\resultsAOR[wecsimMaxAmpErrorBestCase] error).
Damping control results are similar (not shown).

\subsubsection{Dynamics Calibration}
Two dynamics outputs require calibration to ensure agreement with reference models and experimental data.

Prior to calibration, the structural force in storm waves differs significantly from \citet{RM3}.
\ifdefined\DISSERTATION
    This is because load cases in \citet{RM3} are derived experimentally from wave tank tests rather than a linear model.
    For this study, a tuned scale factor on force of \resultsAOR[forceScaleFactor] is used to account for the discrepancy.
    Improving the force model to align with the wave tank data is an area of future work, with suggestions discussed in \Cref{sec:unmodeled-effects}.
\else
    The latter computes loads by scaling wave-tank test results rather than a linear model.
    MDOcean applies a tuned scale factor of \resultsAOR[forceScaleFactor] to obtain agreement.
    \Cref{sec:unmodeled-effects} discusses future work to improve this.
\fi

\ifdefined\DISSERTATION
    As with force, a tuned scale factor of \resultsAOR[powerScaleFactor] is also applied to the power when the geometry is set to the nominal \gls{acr-rm3} geometry to obtain an equal average power to the reference model report (86~kW).
    Unlike the force, the power discrepancy is not attributed to a modeling error because the separate power validation against WEC-Sim tracks extremely well as described in \Cref{sec:dynamic-validation}.
    Rather, the discrepancy is attributed to the fact that the reference model report uses a different dynamic model that is not described in detail, and does not document the hydrodynamic coefficients, water depth, drag coefficient, and other necessary input parameters that would be required to obtain a close match.
\else
    A power scale factor of \resultsAOR[powerScaleFactor] is similarly applied to match the reference report's 86~kW average power.
    The discrepancy is attributed to the reference report's use of an undocumented dynamic model rather than to a modeling error in MDOcean
    since the WEC-Sim validation of \Cref{sec:dynamic-validation} tracks closely.
\fi
Ultimately, the scaled quantities match within 1\% of the reference model report.

\subsubsection{Dynamics Validation}
Validation tests the assumptions introduced in \Cref{sec:dynamics}, focusing on those that are non-standard in \gls{acr-wec} modeling.
Despite a poor match of the drag force magnitude, the overall system dynamics and fundamental dynamic assumptions agree well with the WEC-Sim results.

\paragraph{\glspl{acr-df}: Excellent Agreement}
To validate \asmpref{system-low-pass}, we confirm that the float displacement waveform is very nearly sinusoidal, with a \gls{acr-thd} of less than 1\% in all sea states.
We also check \asmpref{signal-shape-known}, finding excellent agreement. %an RMS error of \resultsAOR[dragShapeRmsErrorPercent] averaged across sea states. see slides 152-153 3/22/26 in MDOcean debugging slides
\Cref{fig:wecsim-thd} and \Cref{fig:wecsim-drag-waveform} in \Cref{sec:appendix-dynamic-validation} provide detailed results.

\paragraph{System Response: Good Agreement}
The absolute error in average power compared to the WEC-Sim power is less than \resultsAOR[wecsimAvgPowerErrorBestCase] in the best case and \resultsAOR[wecsimAvgPowerErrorWorstCase] in the worst case,
with an error breakdown for all simulation scenarios and sea states shown in \Cref{fig:error-histogram}.
\begin{figure*}[htbp]
    \centering
    \includegraphics[\FigureOptions{aor:figs/from-matlab/wecsim_wcsm_multi__histogram.pdf}]{figs/from-matlab/wecsim_wcsm_multi__histogram.pdf}
    \caption[WEC-Sim validation error breakdown]{Error breakdown based on WEC-Sim Validation Runs}
    \label{fig:error-histogram}
\end{figure*}
% Though this seems large, the extremely low errors when enforcing the same hydrodynamic coefficients as WEC-Sim and with zero drag
% (\resultsAOR[wecsimAvgPowerErrorBestCase] and 1.9\% in power and amplitude respectively) indicate that this is not an error in the dynamic model itself,
% but in the way that the model amplifies errors in drag and hydrodynamic coefficient inputs.
\Cref{sec:appendix-dynamic-validation} provides errors of additional outputs as a function of sea state.

\paragraph{Drag Force: Poor Agreement}
Validation of various effort and flow magnitudes all show good agreement with WEC-Sim, with the exception of drag force.
The systematic trend in phase error is particularly concerning, suggesting a sign error in the drag force derivation or implementation.
MDOcean users should avoid drawing conclusions specifically about the effect of drag until the discrepancy is understood and eliminated,
although thankfully overall power or amplitude predictions remain reliable.
Other gaps in validation include the lack of direct testing of \asmpref{waves-narrow-band} (narrowband wave assumption) and comparison of the \gls{acr-qcqp} controller to a time-domain numerical optimal controller.
Even given this possible error, the authors believe that for early-stage design and tradeoff analysis, MDOcean is preferable over other dynamic models,
%simple damping-only drag models
since it captures nonlinearities and coupling effects and does not compromise power and amplitude predictions.

%\Gls{acr-thd} measures how a signal's higher harmonics compare to its fundamental;
%pure sinusoids have a \gls{acr-thd} of 0 and square waves have a \gls{acr-thd} of approximately 48.3\%.
% \Cref{sec:appendix-dynamic-validation} also validates the describing-function approximation itself,
% showing total harmonic distortion below 1\% in the worst sea state
% and excellent agreement between the assumed and actual drag force waveforms at all four corners of the \gls{acr-jpd}.

%  --- Fig 37: drag numerical amplitude and phase are not very correct (but doesn't matter too much)

% - the power for RM3 with drag, lossless PTO, and damping control match that predicted by the reference model report

% - validation = how accurate are my assumptions?
%  --- subfigures 2,3,4 of figure 12

%  Assumption: use meem for hydro coeffs
%  Assumption: simultaneous interaction of (use describing function for drag) and (use meem for hydro coeffs)
 
% Assumption: linearization
% -- linearization of drag: validated above
% -- linearization of PTO saturation: not directly validated, but the drag linearization validates the concept that it stands on (low pass filter)
% Assumption: regular waves (narrowband)
% -- not yet validated, but refer to 92\% study
% Assumption: peak fluid force uses the fundamental
% -- not yet validated
% Assumption: peak amplitude uses the fundamental
% -- not directly validated, but the low drag THD validates the concept that it stands on (low pass filter)
% Assumption: you only care about sinusoidal steady state, no transients
% -- not directly validated, but the drag linearization validates the concept that it stands on (low pass filter bc transients are high frequency)


% Table of assumptions, whether I validate it, and if I don't validate it if that's because it doesn't need validating or because it's not validated yet.

%%%%%%%%%%%%%%%%%%%%%%%%%%%%%%%%%%%%%%%%%
\subsection{Dynamic Validation Using WEC-Sim}\label{sec:dynamic-validation}
The popular time-domain hydrodynamic simulation software WEC-Sim \citep{ruehl_wec-simwec-sim_2024} is used to validate the dynamics module.
The WEC-Sim \gls{acr-rm3} example is run with regular waves and with the device constrained to oscillate only in heave.
\ifdefined\DISSERTATION
    Notably, the \gls{acr-rm3} geometry provided in WEC-Sim differs slightly from the published dimensions in the \gls{acr-rm3} report, so for validation the dimensions input to MDOcean are adjusted to match WEC-Sim.
\fi

\ifdefined\DISSERTATION
    WEC-Sim runs utilize hydrodynamic coefficients obtained with the \gls{acr-wamit} \gls{acr-bem} for dynamics, and control coefficients calculated with MDOcean for consistency.
    MDOcean is run with \gls{acr-meem} as usual, and separately also run with the \gls{acr-wamit} coefficients to distinguish differences caused by disparate hydrodynamic coefficients from those caused by the underlying dynamics.
\else
    Two MDOcean configurations are compared against WEC-Sim:
    one using \gls{acr-meem} hydrodynamic coefficients (the default) and one using \gls{acr-wamit} \gls{acr-bem} coefficients matching WEC-Sim's,
    to separate dynamics-model error from hydrodynamic-coefficient error.
\fi


These errors are deemed acceptable for the purposes of this study,
since the goal is to demonstrate the value of simultaneously analyzing multiple disciplines
and the ability to quickly evaluate a large number of preliminary design options, rather than to conduct late-stage refinement of a known design.
%%%%%%%%%%%%%%%%%%%%%%%%%%%%%%%%%%%%%%%%%

\subsection{Verification, Calibration, and Validation: Other Modules}
\label{sec:static-validation}

\subsubsection{Verification}

The overall model is validated by comparing simulated structural forces, power, mass, and cost results to the nominal values in \citet{RM3}, as shown in \Cref{tab:validation}.
\ifdefined\DISSERTATION
    Volume and center of buoyancy results are also compared to the WEC-Sim \gls{acr-rm3} tutorial geometry, which differs slightly from the nominal \gls{acr-rm3} geometry.
    Structural calculations are additionally compared to \citet{previsic_reference_2011}, an unpublished report by an author of the reference model report providing more detail on the structural calculations.
\fi
The mass, power, and cost track well.

%\hl{Describe the radiation and amplitude limit and show they aren't violated} \citep{zou_practical_2023}.

% \begin{table}[htbp]
% \centering
% \begin{tabular}{l|l|l|l|l}
%                                                      & \multicolumn{2}{l|}{WEC-Sim RM3 Design}                                      & \multicolumn{2}{l}{DOE Report RM3 Design \citep{RM3}}                               \\
%                                                      Variable& MDOcean                    & Actual                     & MDOcean                    & Actual                            \\ \hline
% Float mass                     &                            &                            & 202.4 Mg                   & 208 Mg                            \\
% Vertical column mass           &                            &                            & \hl{fail}                  &                                   \\
% Reaction plate mass            &                            &                            & 230.9 Mg                   & 245 Mg                            \\
% Total mass                     &                            &                            & 614.3 Mg                   & 680 Mg                            \\
% Float volume                   & 726.8 m\textsuperscript{3} & 725.8 m\textsuperscript{3} & 701.9m\textsuperscript{3}  & 701.97m\textsuperscript{3}        \\
% Spar volume                    & 887.8 m\textsuperscript{3} & 886.7 m\textsuperscript{3} & 1,007 m\textsuperscript{3} & 1,008 m\textsuperscript{3}        \\
% CAPEX                          &                            &                            & \hl{fail}                  &                                   \\
% OPEX                           &                            &                            & [1.0 e6, \hl{...}]         & [1.2, 3.3, 6.6, 9.4] $\cdot 10^6$ \\
% LCOE                           & \hl{fail}                  &                            & \hl{fail}                  &                                   \\
% Average power                  &                            &                            &                            &                                   \\
% Heave force                    &                            &                            & 8.497 MN                   & 8.5 MN                            \\
% Spar factor of safety          &                            &                            & 11.894                     & 11.1\textsuperscript{*}           \\
% Power coeff. of variation& \hl{fail}                  &                            & 76.7\%                     & 71.1\%                            \\
% Float center of buoyancy       & 1.293 m                    & 1.293 m                    &                            &                                   \\
% Float center of gravity        & 0.283 m                    & 0.283 m                    &                            &                                  
% \end{tabular}
% \caption{Validation}
% \label{tab:validation}
% \end{table}

    \begin{table*}[htbp]
        \centering
        \ifdefined\JOURNALTWOCOLUMN
            \input{tables/from-matlab/validation_smaller.tex}
        \else
            \input{tables/from-matlab/validation_larger.tex}
        \fi
        \caption{Verification}
        \label{tab:validation}
    \end{table*}

%\hl{Explain sources of error and rough uncertainty and the implications of what we can trust}

The scaling behavior of economic outputs against the number of \glspl{acr-wec} is validated by sweeping $\gls{sym-N-WEC}$ from 1 to 100 and comparing the simulated \gls{acr-capex}, \gls{acr-opex}, and \gls{acr-lcoe} against the values in the reference model report \citep{RM3}.
The average percent errors are included in \Cref{tab:validation}, with additional detail in \Cref{fig:econ-nwec-validate} of \Cref{sec:appendix-econ-validation}.

\subsubsection{Calibration}
\ifdefined\DISSERTATION
    A final tuned scale factor of \resultsAOR[massScaleFactor] is applied to the structural mass to account for the fact that not all structural components are modeled, including spar stiffeners, damping plate circumferential stiffeners, and miscellaneous features like flanges, enclosures, and access mounts.
\else
    A structural-mass scale factor (of \resultsAOR[massScaleFactor]) accounts for unmodeled subcomponents (spar stiffeners, damping plate circumferential stiffeners, flanges, enclosures, access mounts).
\fi

% Assumption: surge force is not contributing to the structural stress
% -- not yet validated

%%%%%%%%%%%%%%%%%%%%%%%%%%%%%%%%%%%%%%%%%
\subsection{Runtime Benchmarking}
\label{sec:sim-runtime}
Benchmarking the runtime of the MDOcean simulation is important to verify it achieves the desired speed to facilitate rapid design optimization.
An initial speed requirement order of magnitude of 100~ms for all modules was set to enable a 100-iteration finite difference optimization with 12 design variables to complete in approximately two minutes.
Ultimately, each simulation run takes around \resultsAOR[simRuntime], missing the initial goal but still a significant improvement over alternatives.
\ifdefined\DISSERTATION
    See \Cref{diss:sec:conclusion-future-work} for a discussion of future work to improve runtime.
\fi
The timings in this section are performed on an Ubuntu 20.04 \gls{acr-lts} server with a 14-core Intel Core i9-10940X \gls{acr-cpu} (3.3~GHz base clock) and 256~GB of DDR4 \gls{acr-ram} at 3200~MHz, running MATLAB R2024b.

\Cref{fig:runtime-modules} visualizes the breakdown of runtime between modules.
The \gls{acr-meem} hydrodynamics module takes the majority (\resultsAOR[pctRuntimeMEEM]) of the time and is broken down in \Cref{fig:runtime-hydro}.
\ifdefined\DISSERTATION
    The biggest portion is dedicated to evaluating Bessel functions in the semi-analytical solution, another large portion to unpacking variables from the cell data structure, and a smaller portion to solving the imaginary modes of the dispersion relation and the linear matrix equation.
\fi
The simulation is an order of magnitude faster than the Capytaine \gls{acr-bem} solver for similar convergence levels.
\ifdefined\DISSERTATION
    See \Cref{ch:hydrodynamic} for a more detailed analysis of the comparative accuracy and runtime of \gls{acr-meem} and Capytaine.
\fi

\begin{figure}[b!]
\centering
\includegraphics[\FigureOptions{aor:figs/from-matlab/sim_runtime.pdf}]{figs/from-matlab/sim_runtime.pdf}
\caption[Simulation runtime breakdown between modules]{Bar chart showing simulation runtime breakdown between modules}\label{fig:runtime-modules}
\end{figure}

\begin{figure}[t!]
\centering
\includegraphics[\FigureOptions{aor:figs/from-matlab/hydro_runtime_logscale.pdf}]{figs/from-matlab/hydro_runtime_logscale.pdf}
\caption[Speed improvement of MDOcean's hydro module over Capytaine]{Bar chart demonstrating the speed improvement of MDOcean's hydro module over baseline solver Capytaine}\label{fig:runtime-hydro}
\end{figure}

The dynamics and controls module takes the next longest (\resultsAOR[pctRuntimeDynamics], enlarged in \Cref{fig:runtime-dynamics}), with contributions from force saturation, spar analysis, drag linearization, and evaluation of the motion transfer function.
This represents a three-order-of-magnitude improvement over the equivalent regular-wave WEC-Sim simulation run in parallel.
\ifdefined\DISSERTATION
    Simplifying the dynamics to a single \gls{acr-dof} achieves another order of magnitude speedup, although the optimization and benchmarking results presented here utilize the 2-\gls{acr-dof} model.
\fi

\begin{figure}[t!]
\centering
\includegraphics[\FigureOptions{aor:figs/from-matlab/dynamics_runtime.pdf}]{figs/from-matlab/dynamics_runtime.pdf}
\caption[Speed improvement of MDOcean's dynamics module over WEC-Sim]{Bar chart demonstrating the speed improvement of MDOcean's dynamics module over baseline solver WEC-Sim}\label{fig:runtime-dynamics}
\end{figure}

The structures, geometry, and economics modules are not computationally expensive and together compose the remaining \resultsAOR[pctRuntimeOther]~of the total runtime.

% say that it's 210 sea states, that wecsim is parallelized across X cores, that MEEM is using N=M=K=10

%%%%%%%%%%%%%%%%%%%%%%%%%%%%%%%%%%%%%%%%%
\subsection{Discussion}
\ifdefined\DISSERTATION
    \Cref{sec:validation-benchmarking} demonstrates that MDOcean achieves accuracies generally within single-digit-percent \gls{acr-jpd}-weighted annual average power under realistic conditions, despite worst-case per-sea-state errors that can be larger due to drag and inter-body phase sensitivity in the 2-\gls{acr-dof} model %within 10\% of established models and benchmarks both at the individual subsystem level and at the system level.
    Just three scale factors are used to tune results across disciplines and achieve an interdisciplinary model fully consistent with all relevant outputs in the reference model report, providing a valuable benchmark to compare technology improvements while retaining the flexibility to adjust scale factors as needed to represent other scenarios.
    The model is quite accurate for inputs near the nominal \gls{acr-rm3} design and therefore its ability to capture multidisciplinary interactions in a single software package while running 1-3 orders of magnitude faster than established software is considered a major success.
    Still, the model is not perfect, and although only minor tuning was performed, the model's development centers around the Reference Model 3 and there may be unidentified discrepancies for geometries far from the reference design.
    We emphasize that the model is not intended to replace detailed models for final design and analysis, but rather to enable rapid design space exploration and optimization in the early stages of design.
    Since the values without scale factors are of the correct order of magnitude and the trends are reasonable, the use of scale factors to resolve discrepancies in the force, power, and mass models is deemed acceptable for the purposes of this study.
\else
    MDOcean reproduces the dynamics of an established time-domain solver within a few percent under matched modeling assumptions and within single-digit-percent \gls{acr-jpd}-weighted annual average power under realistic conditions, despite worst-case per-sea-state errors that can be larger due to drag and inter-body phase sensitivity in the 2-\gls{acr-dof} model.
    Three scale factors are used to tune force, power, and mass outputs across disciplines for consistency with the reference model report.
    The \resultsAOR[simRuntime] runtime is 1-3 orders of magnitude faster than established software, enabling multidisciplinary optimization workflows that would otherwise be prohibitive.
    The model is intended to support rapid design-space exploration in early design, not to replace detailed \gls{acr-fem} or \gls{acr-bem} simulations for final analysis.
    Accuracy may degrade for geometries far from the reference design; users should validate against detailed models before drawing conclusions about previously-unexplored regions of the design space.
\fi
%\clearpage
%%%%%%%%%%%%%%%%%%%%%%%%%%%%%%%%%%%%%%%%%
%%%%%%%%%%%%%%%%%%%%%%%%%%%%%%%%%%%%%%%%%
\section{Insights and Discussion}
\label{sec:discussion}

%%%%%%%%%%%%%%%%%%%%%%%%%%%%%%%%%%%%%%%%%
\subsection{Variable Sweeps}
This subsection summarizes insights from variable-sweep analyses before presenting specific trend breakdowns.

%%%%%%%%%%%%%%%%%%%
\subsubsection{Effect of Bulk Dimensions on Hydrodynamic Efficiency}
\ifdefined\DISSERTATION
    A brute-force sweep of nondimensionalized bulk geometry variables and the nondimensional wavenumber $\gls{sym-m-0}\gls{sym-h}$ allows for the identification of design trends.
    MDOcean performs a six-dimensional sweep with \resultsAOR[meemSweepSize] combinations in only \resultsAOR[meemSweepTime], an average of 16~ms per combination.
    Using only the hydrodynamics and dynamics modules for the present, we examine nondimensional outputs including radiation efficiency, capture width ratio, and nondimensionalized surface area.
    These serve respectively as rough proxies for power, power per unit cost, and cost, allowing for the visualization of tradeoffs between power and cost across the design space.
    Radiation efficiency was defined in \Cref{eq:power-matrix-decomposition} as the ratio of the capture width to the maximum radiation-limited capture width, and represents how well the device's geometry allows it to absorb power at a given frequency.
    Capture width ratio is the ratio of capture width to a characteristic length scale, often taken as the device diameter but taken here as the square root of the wetted surface area to better serve as a proxy for cost.
    (The cube root of submerged volume was also considered as a characteristic length scale; results were qualitatively similar to the square root of surface area and are not shown here.
    If wavelength were used as the characteristic length scale, the resulting capture width ratio would be a constant multiple of the radiation efficiency).
    For the final output, to nondimensionalize the wetted surface area, we divide it by the wavelength squared.
\else
    A brute-force sweep of nondimensionalized bulk geometry variables and the nondimensional wavenumber $\gls{sym-m-0}\gls{sym-h}$ identifies design trends.
    MDOcean performs the six-dimensional sweep with \resultsAOR[meemSweepSize] combinations in \resultsAOR[meemSweepTime].
    We examine three nondimensional outputs: radiation efficiency (proxy for power), capture width ratio (proxy for power per unit cost, normalized by $\sqrt{\gls{sym-A-wet}}$), and nondimensional wetted surface area (proxy for cost).
\fi

\ifdefined\DISSERTATION
    Visualizing this many-dimensional dataset (six inputs, three outputs) can pose a challenge.
    \Cref{fig:meem-sweep-m0h} will provide a view of all six input dimensions simultaneously, focusing centrally on the nondimensional wavenumber $\gls{sym-m-0}\gls{sym-h}$, and examine two outputs of radiation efficiency and capture width ratio.
    Meanwhile, \Cref{fig:meem-sweep-pareto} will show the tradeoff of radiation efficiency with nondimensionalized surface area and denote five of the six input dimensions, excluding only $\gls{sym-m-0}\gls{sym-h}$.
    In both figures, marker size and shape represent $\gls{sym-a-2}/\gls{sym-h}$ and $\gls{sym-a-3}/\gls{sym-a-1}$, which are swept over a range of 0-3 and 1-4 respectively.
    Ratios $\gls{sym-a-1}/\gls{sym-a-2}$, $\gls{sym-d-1}/\gls{sym-h}$, and $\gls{sym-d-2}/\gls{sym-d-1}$ must fall between zero and one and are represented as the red, blue, and green decimal values of the \gls{acr-rgb} color space, respectively.
\else
    Visualizing this many-dimensional dataset (six inputs, three outputs) can pose a challenge.
    \Cref{fig:meem-sweep-m0h} will provide a view of all six input dimensions simultaneously, focusing centrally on the nondimensional wavenumber $\gls{sym-m-0}\gls{sym-h}$, and examine two outputs of radiation efficiency and capture width ratio.
    Meanwhile, \cref{fig:meem-sweep-pareto} will show the tradeoff of radiation efficiency with nondimensionalized surface area and denote five of the six input dimensions, excluding only $\gls{sym-m-0}\gls{sym-h}$.
    In both figures, marker size and shape represent $\gls{sym-a-2}/\gls{sym-h}$ and $\gls{sym-a-3}/\gls{sym-a-1}$, which are swept over a range of 0-3 and 1-4 respectively.
    Ratios $\gls{sym-a-1}/\gls{sym-a-2}$, $\gls{sym-d-1}/\gls{sym-h}$, and $\gls{sym-d-2}/\gls{sym-d-1}$ must fall between zero and one and are represented as the red, blue, and green decimal values of the \gls{acr-rgb} color space, respectively.
\fi

\paragraph{Best Design as a function of Wave Environment}
\begin{figure*}[htbp]
    \centering
    \begin{subfigure}[t]{0.48\linewidth}
        \includegraphics[\FigureOptions{aor:figs/from-matlab/sweep_geoms_line.pdf}]{figs/from-matlab/sweep_geoms_line.pdf}
        \caption{Radiation efficiency}
        \label{fig:meem-sweep-eff}
    \end{subfigure}
    \hfill
    \begin{subfigure}[t]{0.48\linewidth}
        \includegraphics[\FigureOptions{aor:figs/from-matlab/sweep_geoms_line_cwr_sa.pdf}]{figs/from-matlab/sweep_geoms_line_cwr_sa.pdf}
        \caption{Capture width ratio}
        \label{fig:meem-sweep-cwr}
    \end{subfigure}
    \caption[Effect of wave environment and design variables on efficiency and capture width]{Effect of wave environment and hydrodynamic design variables on (a) radiation efficiency and (b) capture width ratio}
    \label{fig:meem-sweep-m0h}
\end{figure*}

\ifdefined\DISSERTATION
    Starting with \Cref{fig:meem-sweep-eff}, we observe that the most efficient design depends on the frequency-depth regime $\gls{sym-m-0}\gls{sym-h}$ of the wave environment, shown on the x-axis.
    At low $\gls{sym-m-0}\gls{sym-h}$ (shallow water, slow waves), the designs with the highest radiation efficiency are those with $\gls{sym-a-1}/\gls{sym-a-2} \approx 1$, $\gls{sym-a-2}/\gls{sym-h} \ll 1$, and $\gls{sym-a-3}/\gls{sym-a-1}\approx 1$ with any value of $\gls{sym-d-1}/\gls{sym-h}$ (the blue/green/purple curves on the left).
    This represents a pure cylinder (float, spar, and damping plate all equal in diameter) with a small radius and any draft.
    
    The radiation limit is marked with a horizontal dashed line, and efficiencies above the limit are obtained around $\gls{sym-m-0}\gls{sym-h}\approx 1, \gls{sym-a-1}/\gls{sym-a-2}\approx 1, \gls{sym-a-2}/\gls{sym-h}\approx 3, \gls{sym-a-3}/\gls{sym-a-1}\approx 1$ (the pink/orange/yellow curves in the center of the figure).
    This indicates an issue with the simulation and data over the limit is disregarded.
    At this moderate $\gls{sym-m-0}\gls{sym-h}$ regime, the most efficient designs that stay within the radiation limit are green/blue, generally indicating a low $\gls{sym-a-1}/\gls{sym-a-2}$ (small spar).
    At high $\gls{sym-m-0}\gls{sym-h}$ (deep water, fast waves), the most efficient designs (pink/red) have
\else
    The most efficient design depends on the frequency-depth regime $\gls{sym-m-0}\gls{sym-h}$.
    The radiation limit is marked with a horizontal dashed line; efficiencies above the limit indicate simulation errors and are disregarded.
    Low $\gls{sym-m-0}\gls{sym-h}$ (shallow water, slow waves) favors a pure cylinder geometry with small radius.
    Moderate $\gls{sym-m-0}\gls{sym-h}$ favors low $\gls{sym-a-1}/\gls{sym-a-2}$ (a small spar relative to the float).
    High $\gls{sym-m-0}\gls{sym-h}$ favors qualitatively different geometries.
\fi

\paragraph{Design Cost-Power Tradeoffs}
\Cref{fig:meem-sweep-pareto} shows a Pareto front for the tradeoff of radiation efficiency (power proxy) and normalized surface area (cost proxy) across the six geometric design variables.
\ifdefined\DISSERTATION
    The normalized surface area does not depend on the seventh input variable $\gls{sym-m-0}\gls{sym-h}$,
    and the largest radiation efficiency for that geometry across all $\gls{sym-m-0}\gls{sym-h}$ is plotted,
    essentially showing the best possible radiation efficiency for each geometry across all wave environments.
    Black squares denote the non-dominated Pareto-optimal points:
    the geometries that achieve the highest radiation efficiency for a given normalized surface area.
    %Designs with \hl{some quality} tend to perform well, with \hl{some quality} being optimal for the high-efficiency, high-cost end of the Pareto front and \hl{some quality} being optimal for the low-efficiency, low-cost end of the Pareto front.
\else
    The plot shows only the largest radiation efficiency across all $\gls{sym-m-0}\gls{sym-h}$ for each geometry.
    Black squares denote Pareto-optimal points.
\fi

\begin{figure}[htbp]
    \centering
    \includegraphics[\FigureOptions{aor:figs/from-matlab/sweep_geoms_pareto_nondim.pdf}]{figs/from-matlab/sweep_geoms_pareto_nondim.pdf}
    \caption[Radiation efficiency vs. power per unit surface area tradeoff]{Effect of hydrodynamic design variables on radiation efficiency and power per unit surface area}
    \label{fig:meem-sweep-pareto}
\end{figure}

%%%%%%%%%%%%%%%%%%%
\subsubsection{Damping Plate Size}
\Cref{fig:damping-plate-maxs} shows that damping plate stress can be reduced substantially by increasing the plate inner radius.
This is because the same force has a lower lever arm to the column and therefore creates less bending moment.
\begin{figure}[htbp]
    \centering
    \includegraphics[\FigureOptions{aor:figs/from-matlab/damping_plate_aspect_ratio.pdf}]{figs/from-matlab/damping_plate_aspect_ratio.pdf}
    \caption[Effect of damping plate aspect ratio on maximum stress and deflection]{Normalized effect of damping plate aspect ratio on maximum stress and deflection.
The dashed lines indicate the nominal design point at $\gls{sym-b-inner-damping-plate}/\gls{sym-a-outer-damping-plate} = 0.2$.}
    \label{fig:damping-plate-maxs}
\end{figure}

%%%%%%%%%%%%%%%%%%%
\subsubsection{Effect of \gls{acr-pto} Force and Power Limits}
\ifdefined\DISSERTATION
    Previous work has shown that capping the \gls{acr-pto} force can substantially reduce \gls{acr-pto} size and cost with minimal decrease in power.
    \citet{devin_high-dimensional_2024} finds that impedance matching at resonance has a large effect on power and \gls{acr-pto} sizing has little effect.
    This suggests that if cost were included, it may be preferable to have an undersized \gls{acr-pto} that frequently saturates its control force.
    This agrees with the conceptual argument of \citep{coe_maybe_2021} and the numerical findings of \citep{mcgilton_optimal_2024}.
    The latter shows that reducing the maximum \gls{acr-pto} force by 70\% leads only to a 3\% decrease in annual energy production, a trend consistent across two device archetypes (\gls{acr-rm3} and RM5), four locations, and several dimensional scales.
    The recent \gls{acr-rm3} optimization \citep{gaebele_tpl_2025} finds that the device should be scaled down by a factor of 0.55 to a new float diameter of 11~m in order to maximize the power per unit surface area, and that the continuous torque constraint is often active, illustrating the importance of including this constraint in the model.
    \citet{mccabe_force-limited_2024} shows analytically that for a given sea state, power decreases quadratically with decreasing force limit for the worst-case scenario of zero intrinsic reactance, and that highly reactive devices are even less sensitive to force limits.
    Intuitively, the sensitivity of annual average power to force limit should be even lower when considering a range of realistic sea states, because large sea states that require higher forces are comparatively rare.
    Using MDOcean, we investigate the effect of \gls{acr-pto} force limit on average power and structural load in \figureautorefname~\ref{fig:force-limit}.
\else
    Prior work consistently shows that capping the \gls{acr-pto} force can substantially reduce \gls{acr-pto} size and cost with minimal decrease in power \citep{devin_high-dimensional_2024,coe_maybe_2021,mcgilton_optimal_2024,gaebele_tpl_2025,mccabe_force-limited_2024}.
\fi
Using MDOcean, we investigate the effect of \gls{acr-pto} force limit on average power and structural load in \Cref{fig:force-limit}.

\begin{figure}[htbp]
\centering
\includegraphics[\FigureOptions{aor:figs/from-matlab/power_force_sensitivity.pdf}]{figs/from-matlab/power_force_sensitivity.pdf}
\caption[Effect of force limit on power and structural load]{Effect of Force Limit on Annual Average Power and Peak Structural Load}
\label{fig:force-limit}
\end{figure}

Decreasing the force limit has almost no effect on power until the force limit is around 50\% of the nominal value, after which power falls off steeply.
\ifdefined\DISSERTATION
    From theory (see \Cref{sec:appendix-constraint-sensitivity}), we expect a smooth nonlinear relationship between power and force limit when considering a single sea state, and a piecewise continuous nonlinear relationship when considering multiple sea states as in \Cref{fig:force-limit}.
\fi
On the other hand, the structural load for the operational design load case scales nearly linearly with the force limit.
\ifdefined\DISSERTATION
    For structural load, we expect a linear relationship when considering a single sea state and a piecewise linear relationship when considering multiple sea states.
    The fact that the linear relationship holds when considering multiple sea states implies that both the worst-case sea state and its quasi-linearized drag coefficients hardly change as the force limit is swept.
    If this trend holds true for other designs, it could suggest a more efficient analysis strategy where the worst-case sea state is identified at the nominal design and then maintained (or only infrequently updated) throughout a design optimization, rather than performing a full multi-sea-state structural analysis for each design being considered.
\fi

We next move to a simultaneous sweep of both the force limit and the power limit, which is more relevant for \gls{acr-pto} design optimization since generator selection dictates both limits.
\ifdefined\DISSERTATION
    Here we hold structural thicknesses constant, so any cost savings only reflect reduced direct cost of the \gls{acr-pto} and not the cost savings from the reduced structural material that would be possible with lower loads.
\fi
\Cref{fig:force-power-limit-sweep} shows the results for both power and \gls{acr-lcoe}, with hatched shading to indicate regions that are infeasible.
\ifdefined\DISSERTATION
    The unsaturated solution is in the top right corner, past which power is constant because the unconstrained reactive controller satisfies the limit in all sea states.
\fi

\begin{figure*}[htbp]
\centering
\includegraphics[\FigureOptions{aor:figs/manual/pto_sweep_3.pdf}]{figs/manual/pto_sweep_3.pdf}
\caption[Effect of force and power limits on power and \gls{acr-lcoe}]{Effect of Generator Force Limit and Power Limit on Annual Average Power and \gls{acr-lcoe}. Cost values are in 2012 USD.}
\label{fig:force-power-limit-sweep}
\end{figure*}
Notably, the minimum \gls{acr-lcoe} point occurs at a force limit far below the unsaturated solution because the cost savings from a smaller \gls{acr-pto} outweigh the loss in power.
\ifdefined\DISSERTATION
    If a full design optimization were performed to adjust the structural thicknesses to take advantage of the reduced structural loads, we expect an even lower force limit to be optimal.
\fi
This illustrates the importance of including the force-cost tradeoff in an optimization model rather than assuming an unconstrained solution.
\ifdefined\DISSERTATION
    We also observe that power and \gls{acr-lcoe} both appear to be convex with respect to the force and power limits, and that the feasible region is convex as well.
    The theoretical analysis in \Cref{sec:appendix-constraint-sensitivity} derives mathematical conditions for when this is the case.
    Convexity suggests that local optimization methods should be effective for finding the optimal \gls{acr-pto} design.
\else
    Power and \gls{acr-lcoe} both appear to be convex with respect to the force and power limits across the feasible region, a structural property that future optimizations can exploit.
\fi

%%%%%%%%%%%%%%%%%%%
\subsubsection{Design Space Exploration}
\ifdefined\DISSERTATION
    The \gls{acr-oat} design sweep, shown in \Cref{fig:experiments}, immediately draws attention to the float diameter $\gls{sym-D-f}$, shown with a red dashed line, as the variable with the strongest effect on the objectives.
    A slight reduction in the diameter from its nominal value causes the cost to decrease close to linearly but the \gls{acr-lcoe} to skyrocket due to a substantial reduction in power generation.
    Meanwhile, slightly increasing the diameter causes the power production to grow significantly but increases loads, leading to fatigue failure of the damping plate.
    This is shown as blank on the diagram to indicate infeasibility.
    A larger $\sim30\%$ diameter increase restores feasibility of the damping plate and achieves a roughly 15\textcent/kWh reduction in \gls{acr-lcoe}, representing the best design of the \gls{acr-oat} sweep.
    The results also point to the possibility of an even lower \gls{acr-lcoe} design with float diameter perhaps $\sim20\%$ larger than nominal, if the damping plate can be made feasible by increasing $\gls{sym-t-d}$ or $h_{1,\text{stiff},d}$ without a major cost penalty.
\else
    The one-at-a-time (\gls{acr-oat}) design sweep (\Cref{fig:experiments}) reveals that float diameter $\gls{sym-D-f}$ has the strongest effect on objectives.
    The feasible design space is highly fragmented: an ``island'' of low-\gls{acr-lcoe} designs at $1.3 < \gls{sym-D-f}/\gls{sym-D-f-nom} < 2.1$ is separated from the nominal design by the damping plate fatigue constraint.
\fi

\begin{figure}[htbp]
\centering
\includegraphics[\FigureOptions{aor:figs/from-matlab/experiments_ratios.pdf}]{figs/from-matlab/experiments_ratios.pdf}
\caption{Design of experiments}\label{fig:experiments}
\end{figure}

\ifdefined\DISSERTATION
    Besides $\gls{sym-D-f}$, the other design variable displaying a non-monotonic relationship with $LCOE$ is the force limit $F_{\text{max}}$, with an optimal at around half the nominal force limit.
    Together, these insights inform the selection of the starting design for optimization, $\gls{sym-vec-x-0}$.
    The \gls{acr-oat} sweep also does not reveal any multimodality in the objectives, a promising sign that only a limited number of multi-starts may be necessary, although objective multimodality could still exist in other parts of the design space.
    Finally, the \gls{acr-oat} design sweep shows the highly constrained nature of the problem, with large regions of the design space infeasible.
    In particular, the feasible region $1.3<\gls{sym-D-f}/\gls{sym-D-f-nom}<2.1$ is an ``island'' in the feasible design space, separated from the other feasible points at $\gls{sym-D-f}/\gls{sym-D-f-nom}<1 $ by the damping plate fatigue constraint.
    Thus even without objective multimodality, multimodality exists in the Lagrangian (objective plus weighted active constraints), indicating the need for a multi-start optimization procedure as planned.
\else
    Force limit $F_{\text{max}}$ also shows non-monotonic \gls{acr-lcoe} dependence, with an optimum near half the nominal limit.
    Although no objective multimodality is observed within the swept region, the fragmented feasible space introduces multimodality in the Lagrangian, indicating the need for a multi-start optimization procedure \citep{mccabe_leveraging_2026}.
\fi

%%%%%%%%%%%%%%%%%%%%%%%%%%%%%%%%%%%%%%%%%
\subsection{Multidisciplinary Insights}
\label{sec:discussion-multidisciplinary}
This section leverages the analytical multidisciplinary nature of the model to draw intuitive insights on limit cases and tradeoffs, and observe nondimensional relationships and scaling laws that would not be readily apparent in a purely numerical or single-discipline model.

\paragraph{Power Matrix Drivers}
Similar to \Cref{fig:JPD-multiply} showing the multiplication of the power matrix by the \gls{acr-jpd} to obtain average power,
\Cref{fig:power-matrix-decomposition} further decomposes the power matrix into component matrices
showing the impact of shape, drag, maximum force, and \gls{acr-pto} losses, with the following breakdown:
\begin{equation}\label{eq:power-matrix-decomposition}
\gls{sym-mathbf-P-H-T-elec} = \gls{sym-mathbf-J-H-T-wave} ~\gls{sym-mathbf-CW-H-T-max} ~\gls{sym-boldsymbol-eta-H-T-rad} ~\gls{sym-boldsymbol-eta-H-T-drag} ~\gls{sym-boldsymbol-eta-H-T-limits} ~\gls{sym-boldsymbol-eta-H-T-elec}
\end{equation}

Here $\mathbf{J}^{H,T}_{\text{wave}}$ is the wave power density (design-independent) and $\mathbf{CW}^{\gls{sym-H},\gls{sym-T-period}}_{\text{max}}$ is the radiation-limit capture width (depends only on the mode of motion).
The four efficiency matrices isolate the design dependencies: $\gls{sym-boldsymbol-eta-H-T-rad}$ for hull-shape radiation efficiency, $\mathbf{\eta}^{H,T}_{\text{drag}}$ for drag, $\mathbf{\eta}^{H,T}_{\text{limits}}$ for dynamic constraints, and $\mathbf{\eta}^{H,T}_{\text{elec}}$ for \gls{acr-pto} losses and dynamics.
Explicit definitions are given in \Cref{sec:appendix-power-decomposition}.

\ifdefined\DISSERTATION
    Two features of this decomposition deserve emphasis.
    First, $\mathbf{\eta}^{H,T}_{\text{drag}}$ may exceed 1 in sea states where the drag-free system is so undamped that the unconstrained optimal controller would induce instability and a sub-optimal controller is used instead --- this reflects a denominator artifact, not an actual energy gain from drag.
    Second, the multiplicative structure depends on the order chosen: alternative factorizations (e.g., defining drag effect relative to the saturated rather than the ideal system) would represent independent sensitivities and could not be multiplied to recover the overall power matrix.
\fi


This multiplicative decomposition allows isolation of the effect of each design variable and parameter on power, intuitively identifying which design variables are most important for improving power.
\Cref{tab:power-matrix-dependence} maps the dependence explicitly.
{
\renewcommand{\arraystretch}{1.35}
\begin{table}[h]
\centering
\caption{Dependence of efficiencies on inputs}\label{tab:power-matrix-dependence}
\begin{tabular}{lcccccc}
\hline
 & $x_{\text{shape}}$ & $\gls{sym-C-d-f}$ & $\gls{sym-C-d-s}$ & $\gls{sym-F-max}$ & $\gls{sym-P-max}$ & $\gls{sym-eta}$ \\
\hline
$\gls{sym-boldsymbol-eta-H-T-rad}$   & $\checkmark$ &  &  &  &  &  \\
$\gls{sym-boldsymbol-eta-H-T-drag}$  & $\checkmark$ & $\checkmark$ & $\checkmark$ &  &  &  \\
$\gls{sym-boldsymbol-eta-H-T-limits}$ & $\checkmark$ & $\checkmark$ & $\checkmark$ & $\checkmark$ &  &  \\
$\gls{sym-boldsymbol-eta-H-T-elec}$  & $\checkmark$ & $\checkmark$ & $\checkmark$ & $\checkmark$ & $\checkmark$ & $\checkmark$ \\
\hline
\end{tabular}
\end{table}
}

\ifdefined\DISSERTATION
    Importantly, the later matrices in the product depend on the design variables and parameters that affect the earlier matrices.
    This coupling means that the effect of the drag coefficients, for example, on power in a device with \gls{acr-pto} limits and dynamics must be understood not only through $\mathbf{\eta}^{H,T}_{\text{drag}}$ but also through their effect on $\mathbf{\eta}^{H,T}_{\text{limits}}$ and $\mathbf{\eta}^{H,T}_{\text{elec}}$.
    While it would be possible to formulate variants such as $\mathbf{\eta}^{H,T}_{\text{drag},var} = \mathbf{P}_{IV}/\mathbf{P}_{IV, no \text{drag}}$ that capture the full effect of drag on the saturated system, such matrices would represent independent sensitivities and could not be multiplied together to obtain the overall power matrix, which is a key advantage of the current formulation.
\else
    The later matrices depend on the design variables affecting earlier matrices, so the total impact of (e.g.) drag must be assessed through its effects on $\mathbf{\eta}^{H,T}_{\text{drag}}$, $\mathbf{\eta}^{H,T}_{\text{limits}}$, and $\mathbf{\eta}^{H,T}_{\text{elec}}$ jointly.
\fi

\Cref{fig:power-matrix-decomposition} shows the power matrix decomposition for the nominal \gls{acr-rm3} design.
\ifdefined\DISSERTATION
    Aligning with \citet{zou_practical_2023}, low-period sea states achieve mechanical powers closer to the radiation limit (higher radiation efficiency).
    High-period sea states have lower radiation efficiencies due to the amplitude limit yet still contribute significantly due to their higher energy content.
    Despite the most probable wave period being 6-8 seconds, 10 second waves contribute the most to energy production for this reason, though even longer waves contribute little due to their rarity.
\else
    Low-period sea states achieve higher radiation efficiency, while high-period sea states have lower radiation efficiency but contribute significantly to total production due to their higher energy content; despite the most probable period being 6--8~s, 10~s waves contribute the most to annual energy production \citep{zou_practical_2023}.
\fi

\begin{figure*}[htbp]
\centering
\includegraphics[\FigureOptions{aor:figs/from-matlab/nominal_power_matrix.pdf}]{figs/from-matlab/nominal_power_matrix.pdf}
\caption{Power matrix decomposition}
\label{fig:power-matrix-decomposition}
\end{figure*}
The effect of drag is most significant in the 11-12 second range, corresponding to the spar's natural frequency.
\ifdefined\DISSERTATION
    Interestingly, this is also where the force limit has the strongest effect, indicating that the large amplitudes at the natural frequency drive up the force more than the larger stiffnesses required for reactive control at frequencies far from the natural frequency.
\else
    Interestingly, this is also where the force limit has the strongest effect, indicating that the large amplitudes at the natural frequency drive up the force more than the larger stiffnesses required for reactive control at frequencies far from the natural frequency.
\fi

\paragraph{Constraint Sensitivities}
The structure of the constrained optimal control problem allows for quantification of the sensitivity of optimal power to each constraint limit, formalizing the intuition behind $\gls{sym-boldsymbol-eta-H-T-limits}$.
From \Cref{eq:constrained-qp-solution,eq:power-elec}, the average electrical power (the \gls{acr-lcoe} denominator) is a sea-state-weighted sum of the impedance-matched power in each sea state, each scaled down by the factor $(1-|\Gamma_{\beta,\text{opt}}|^2)$ that captures the effect of active dynamic constraints:
\begin{equation}\label{eq:power-elec-sum}
\begin{aligned}
    \gls{sym-overline-P-elec} &= \gls{sym-eta} \sum_{\gls{sym-beta}=1}^{\gls{sym-N-sea}} \gls{sym-w-beta} \cdot
    \frac{|\gls{sym-hat-V-s-th-beta}|^2}{8\Re(\gls{sym-Z-s-th-beta})}~ (1 - |\gls{sym-Gamma-beta-opt}|^2)
    \\ \text{\gls{sym-s}.\gls{sym-t}.}
\quad
    |\gls{sym-hat-F-p-beta}| &\leq \gls{sym-F-max},\quad
    |\gls{sym-p-pk-V-I-beta}| \leq \gls{sym-P-max},\\
    |\gls{sym-hat-xi-f-beta}| &\leq \gls{sym-xi-f-max},\quad
    |\gls{sym-hat-xi-s-beta}| \leq \gls{sym-xi-s-max}\\
\end{aligned}
\end{equation}

Power is maximized when $|\gls{sym-Gamma-beta-opt}|=0$ in every sea state, which corresponds to unconstrained impedance-matched control.
When constraints are active, the optimal $|\gls{sym-Gamma-beta-opt}|>0$ represents the impedance mismatch required to satisfy them, and power decreases quadratically with $|\gls{sym-Gamma-beta-opt}|$.

\Cref{sec:appendix-constraint-sensitivity} uses this relationship to derive a closed-form scaling law for the optimal annual average power $\gls{sym-overline-P-elec}$ as a function of the constraint coefficients $\gls{sym-b-mu}$ (which are affine or square functions of the constraint limits $\gls{sym-tau-max}$, $\gls{sym-P-tau-Omega-max}$, $\gls{sym-xi-f-max}$, $\gls{sym-xi-s-max}$, and $\gls{sym-X-PTO-max}$ defined in \Cref{tab:qp-constraints}).

The result is:
\begin{equation}\label{eq:power-double-sum-main}
\begin{aligned}
    \gls{sym-overline-P-elec} =~& \gls{sym-overline-P-elec}_{,0}
    - \gls{sym-eta} \sum_{\gls{sym-mu}\gls{sym-nu}\gls{sym-beta}}
    \left[ 
        \sqrt{\textrm{aff}_{\gls{sym-mu}\gls{sym-nu}\gls{sym-beta}}(\gls{sym-b-mu})}
        ~+ \right. \\
        & \left. \textrm{quad}_{\gls{sym-mu}\gls{sym-nu}\gls{sym-beta}}(\gls{sym-b-mu}, \gls{sym-b-nu})
    \right.+
     \\
    &\left.
        \left(\textrm{aff}_{\gls{sym-mu}\gls{sym-nu}\gls{sym-beta}}(\gls{sym-b-mu}) + \textrm{aff}_{\gls{sym-mu}\gls{sym-nu}\gls{sym-beta}}(\gls{sym-b-nu})\right)
        \sqrt{\textrm{quad}_{\gls{sym-mu}\gls{sym-nu}\gls{sym-beta}}(\gls{sym-b-mu}, \gls{sym-b-nu})}~
    \right]
\end{aligned}\end{equation}
where $\gls{sym-overline-P-elec}_{,0}$ is the unconstrained power and $\textrm{aff}_{\gls{sym-mu}\gls{sym-nu}\gls{sym-beta}}$, $\textrm{quad}_{\gls{sym-mu}\gls{sym-nu}\gls{sym-beta}}$ are affine and quadratic functions of $\gls{sym-b-mu}$ and $\gls{sym-b-nu}$, respectively.
$\gls{sym-mu}$ and $\gls{sym-nu}$ are constraint indices, expanded in \Cref{sec:appendix-constraint-sensitivity}.

Under conditions derived in \Cref{sec:appendix-constraint-sensitivity}, this scaling law is convex in the constraint coefficients $\gls{sym-b-mu}$.
While satisfaction of these conditions is not guaranteed, convexity is observed to hold across the \gls{acr-pto} sweep of \Cref{fig:force-power-limit-sweep}.

The scaling law also displays coupling in the constraint variables: the effect of one constraint limit depends on the values of the other constraint limits.
This illustrates the importance of a multidisciplinary optimization considering all interactions, rather than a sequential optimization that optimizes one subsystem at a time holding the others fixed.

\paragraph{\gls{acr-lcoe} Drivers}
To better see the dependence of \gls{acr-lcoe} on design for a fixed $\gls{sym-N-WEC}$ and $\gls{acr-fcr}$, one can divide out the constant factors from \Cref{eq:lcoe}, leaving normalized price constants $\gls{sym-p-prime-star}$:
\begin{equation}\label{eq:LCOE-scale}
  \gls{acr-lcoe} = \frac{
    \gls{sym-p-prime-F} \gls{sym-F-max} + \gls{sym-p-prime-P}\gls{sym-P-max} + \gls{sym-p-prime-s} \gls{sym-V-struct} + \gls{sym-p-prime-0}
    }{\gls{sym-overline-P-elec}}
\end{equation}
The numerator is affine in the design inputs $\gls{sym-F-max}$ and $\gls{sym-P-max}$, while $\gls{sym-V-struct}$ is the product of structural thicknesses with their corresponding areas, which depend quadratically on the bulk dimensions.
The fact that decreasing a constraint limit causes the numerator of the \gls{acr-lcoe} to monotonically increase and the denominator to monotonically decrease implies that at system optimality, each constraint will be active in at least one sea state.

\ifdefined\DISSERTATION
    The \gls{acr-lcoe} minimization with design variables of $\gls{sym-F-max}$, $\gls{sym-P-max}$, and thicknesses is a fractional quadratic program, a structure common in resource allocation problems.
    When the denominator is convex with respect to the \gls{acr-pto} limits, this means that \gls{acr-lcoe} is ``quasi-convex'' in the \gls{acr-pto} limits for fixed geometry \citep{agrawal_disciplined_2020}.
    This quasi-convexity can be leveraged by using a subgradient method or by solving a small sequence of convex problems, which simplifies the optimization process significantly compared to a general nonlinear solver.
    This motivates a three-level nested optimization strategy:
    at the innermost level, the optimal control problem for fixed geometry and \gls{acr-pto} limits is solved as an analytical \gls{acr-qp};
    at the mid-level, the optimal \gls{acr-pto} limits are found by exploiting the quasi-convex fractional structure when possible;
    and at the outer level, the bulk geometry is optimized with a general nonlinear program.
\else
    The fractional-quadratic structure of \Cref{eq:LCOE-scale} admits an efficient nested optimization that exploits convexity properties; see the companion paper \citep{mccabe_leveraging_2026} for the optimization strategy.
\fi

%%%%%%%%%%%%%%%%%%%%%%%%%%%%%%%%%%%%%%%%%
\subsection{Limitations and Future Work}\label{sec:unmodeled-effects}
\ifdefined\DISSERTATION
    Broadly, future work must both address limitations of the model and apply it to answer design questions.
    As MDOcean is open source, the authors encourage the wave energy community to contribute to both areas, and are themselves actively pursuing the latter with a multi-objective optimization of the \gls{acr-rm3} \gls{acr-wec} \citep{mccabe_leveraging_2026}.
    \Cref{tab:future-work} summarizes potential future improvements to the model, distinguishing between model improvements that would enhance the accuracy or realism of studies that can be conducted with the present model and those that would unlock the ability to answer design questions that the current model cannot.
    This section describes the relevance and possible implementation paths for each development area.
\else
    \Cref{tab:future-work} summarizes principal limitations and future work, distinguishing model improvements to enhance the accuracy of currently-achievable studies from extensions to unlock new design questions.
    MDOcean is open-source \citep{mccabe_mdocean_2024}; community contributions are welcome, and the authors are pursuing the multi-objective \gls{acr-rm3} optimization in the companion paper \citep{mccabe_leveraging_2026}.
\fi

\newcommand{\modelTrustBuilders}{
    \begin{enumerate}[leftmargin=*]
        \item Surge force and mooring cost
        \item Nonlinear storm wave forces
        \item Irregular waves
        \item Structure: thinner plates, better validation
        \item Additional \gls{acr-pto} nonlinearities
    \end{enumerate}
}

\newcommand{\modelStudyEnablers}{
    \begin{enumerate}[leftmargin=*]
        \item Spectral power, load, amplitude
        \item Different \gls{acr-wec} archetypes
        \item Lifetime and sea state contours
        \item \gls{acr-meem} extensions for damping plate and 3+ cylinders
        \item Generator magnetics and cost
        \item Other structural materials
    \end{enumerate}
}

\begin{table}[htbp]
\caption{Future model improvements}
    \label{tab:future-work}
    \expanded{
        \noexpand\tabular{\TableColumnSpecs{aor:future-work}}
    }
Enhance Trust in Achievable Studies  & Unlock New Studies  \\ \hline
\modelTrustBuilders                & \modelStudyEnablers \\
    \endtabular
    \end{table}

\ifdefined\DISSERTATION
    \paragraph{Neglect of Surge Force and Mooring Cost}
    \Cref{sec:design-load-cases} establishes that although the dynamics model estimates the surge force, the structures model uses only the heave force to calculate the factors of safety.
    The neglect of surge force means that geometries with significant lateral areas that would experience high surge force appear to perform better than they should.
    This is not only because of the impact on the structures module, but also because the surge force drives mooring cost, an effect which is not modeled in this study.
    The model predicts that the nominal \gls{acr-rm3} design experiences a storm surge force of \resultsAOR[surgeForceFloatNominal] on the float and \resultsAOR[surgeForceSparNominal] on the spar, compared to the \resultsAOR[heaveForceNominal] nominal heave force.
    With mooring and foundation comprising 12\% of the capital cost (see \Cref{tab:CBS}), this effect may be substantial.
    Incorporating surge force into the structures module and adding a mooring cost module are recommended for future work.
    If these effects are implemented and confirmed to be important, then it may also be desired to obtain a more accurate estimate of the surge force across frequencies by utilizing the surge hydrodynamic coefficients instead of the long-wavelength approximation in \Cref{eq:surge-force}. 
    
    \paragraph{Nonlinear Storm Loads}
    As \Cref{sec:irregular-waves} describes, MDOcean uses standard linear wave theory with equivalent regular waves, even for the storm condition.
    However, storm waves are nonlinear, and it is generally recommended to utilize higher order potential flow, \gls{acr-cfd} simulations, or wave tank tests to obtain storm loads \citep{coe_survey_2018}.
    Additionally, even when storm wave analysis uses linear or quasi-linear hydrodynamics, other design procedures typically utilize a time-domain ``design wave elevation'' signal to capture transient peaks.
    The regular wave analysis used in MDOcean cannot capture nonsinusoidal peaks.
    Therefore, determination of design loads from a given extreme sea state represents a major area of uncertainty in this model and a challenge for future development.
    Open-source tools for \gls{acr-wec} extreme response, such as MHKiT\footnote{\url{https://mhkit-software.github.io/MHKiT/}} and the \gls{acr-dlc} Generator\footnote{\url{https://dlc.primre.org/DLCGenerator}}, apply extreme statistics to facilitate the creation of a design wave elevation but do not calculate loads from this elevation.
    Finding a modeling methodology that captures hydrodynamic nonlinearities and transient peaks without optimization-prohibitive computational expense remains an open problem.
    Second-order \gls{acr-meem} can potentially address the former, but nonlinearities require solving the radiation and diffraction problems separately, and corner discontinuities present new analytical challenges \citep{cong_novel_2020,mavrakos_second-order_2009}.
    Another option is a slender-body approximation for second order loads that was recently implemented in the frequency-domain quasi-linear hydrodynamics solver \gls{acr-raft} \citep{carmo_slender-body_2025}.
    Meanwhile, to address the transient dynamics of an irregular design wave elevation, it requires more investigation to determine whether the semi-analytical dynamics model used here can be adequately extended.
    Alternatives include integrating second-order hydrodynamics into pseudo-spectral methods \citep{coe_initial_2020}, or an effort to speed up time-domain models where second-order hydrodynamics have recently been integrated \citep{issoglio_second-order_2026}.
    
    \paragraph{Irregular Waves}
    Besides storm loads, the regular wave assumption also affects the operational power and load calculations.
    In the future, the interaction of irregular waves with dynamic constraints could be roughly approximated in MDOcean with minimal implementation effort by computing the probability density function for each signal and saturating any amplitude above the constraint threshold.
    To more fully capture nonlinearities in irregular waves, the describing function mentioned earlier (which quantifies the fundamental amplitude of the response to a deterministic sinusoidal input) could be replaced with its probabilistic counterpart, stochastic linearization (which quantifies the expected value of the fundamental over the spectral input).
    This technique has been explored in several \gls{acr-wec} papers \citep{da_silva_statistical_2020,da_silva_stochastic_2023,kluger_synergistic_2017,folley_spectral-domain_2016,spanos_efficient_2016}, including one that performs multi-objective design optimization \citep{neshat_enhancing_2024}.
    
    \paragraph{Probabilistic Variation in Load, Amplitude, and Power}
    Once implemented, an irregular wave formulation would allow for the quantification of the short-term variability in power, load, and response amplitude within a sea state.
    Specifically, the Wiener-Khinchin theorem describes the spectrum of the output of a linear system in terms of the input spectrum and the transfer function.
    This opens the door for investigations in other disciplines, including the development of a spectral fatigue model, insight into broadband impedance matching and dynamic constraints, and analysis of grid integration.
    A broadband matching study could consider the complementary $\gls{sym-Z-i}(\gls{sym-omega})$ shapes of different \gls{acr-pto} types, \gls{acr-wec} shapes, and wave spectra.
    Use of probabilistic amplitude constraints rather than the regular (or equivalent wave-by-wave) analysis assumed in MDOcean would allow a designer to encode a desired level of risk tolerance for extreme events, which may be more appropriate than the current approach of using the maximum nonzero joint probability density for the operational amplitude constraints. 
    
    Additionally, while MDOcean accounts for the effect of the power limit on \gls{acr-pto} cost and models long-term power variation across sea states, its failure to model short-term power variation within a sea state means that it cannot reliably be used for sizing energy storage or grid connection.
    In particular, the temporal alignment of power production with grid demand can affect the value of the produced electricity, and therefore the economic viability.
    The authors are actively pursuing this extension, with an initial modeling process described in the study \citep{mccabe_wec_2025} and ongoing refinements in future work.
    Going further, future studies could investigate whether the addition of local or grid-connected electrical energy storage or the modulation of hydrodynamic design to smooth the dynamic response is more cost-effective at reducing power variability.
    
    \paragraph{Multi-cylinder \gls{acr-meem} with damping plate}
    Recall from \Cref{sec:hydro-meem} that the present study utilizes a two-cylinder \gls{acr-meem} formulation that neglects the slanted float bottom and damping plate.
    The multi-cylinder \gls{acr-meem} has recently been implemented \citep{best_openflash_2026,bimali_matrix_2026}, but has not yet been integrated into MDOcean.
    Doing so would not only improve the accuracy of the hydrodynamic coefficients, allowing calculation of the spar coefficients without relying on the approximations of \Cref{sec:hydro-other}, but also expand the reachable design space to include qualitatively different frequency-dependence of the hydrodynamic coefficients.
    Implementing the damping plate in the \gls{acr-meem} model would further expand the reachable design space to include spar designs that deviate from the nominal aspect ratios imposed by the interpolation in \Cref{sec:hydro-other}.
    %would let me actually optimize the spar dimensions, which is important since I'm up against that constraint of min damping plate diameter 
    
    \paragraph{Additional \gls{acr-pto} nonlinearities}
    The current model includes describing functions for a force saturation nonlinearity and a drag nonlinearity.
    Other relevant nonlinearities that could be added include drivetrain static friction, hydraulic check valve rectification, and power saturation.
    It would also be possible to utilize the higher harmonics of the existing describing functions if the fundamental approximation is found to be insufficient and a fully nonlinear model is desired.
    This would require a major overhaul to the dynamics model to transition from the quasi-linear frequency domain analysis to a spectral or pseudo-spectral analysis, and might best be implemented in an existing tool such as WecOptTool.
    
    \paragraph{Generator magnetics and cost model}
    Building off the \gls{acr-pto} \gls{acr-ccd} study that is easily realizable with the current model, a generator model would unlock \gls{acr-ccd}  with the generator itself, for example determining the conditions under which an expensive high-torque generator outperforms a cheap low-torque generator.
    The torque limit would essentially impose a constraint on the $F_{\text{max}}$ design var, the core losses introduce nonlinear damping, and impose a relation between generator torque limit and max/min generator inertia.
    This analysis could leverage a simplified generator model as explored in the \gls{acr-rm3} \gls{acr-ccd} study \citep{anderson_re-imagining_2024}, or a full generator model as in the offshore wind \gls{acr-mdo} study \citep{barter_beyond_2023}.
    
    \paragraph{Structural model extensions}
    The structures module could be improved for geometries with more prominent stiffeners and thinner skins by adding local buckling failure criteria and effective breadth considerations.
    Confidence in the semi-analytical structures model would be enhanced by performing a systematic validation study using \gls{acr-fea}, rather than relying on minimally-documented \gls{acr-fea} results in the reference model report.
    
    Extending the structures model to other types of materials would allow the comparison of novel \gls{acr-wec} materials with traditional steel designs.
    The existing relations between force and stress should still hold, but the relation between stress and factor of safety require modification for materials with different limit states.
    This is of interest because steel accounts for a significant portion of device cost, and the assumption of steel construction may be holding the \gls{acr-wec} industry back from economic viability.
    For example, \citet{roberts_bringing_2021} finds inflatable polyurethane coated nylon fabric to be more optimal than steel, reinforced concrete, fiberglass, and rubber based on basic density and cost considerations.
    However, unlike ductile metals, composites require anisotropic analysis, concrete requires brittle fracture analysis such as Mohr's circle, and inflatables require tensioned membrane modeling.
    
    \paragraph{Lifetime and storm contour modeling}
    MDOcean's conservative use of the endurance limit to model fatigue implies an infinite structural lifetime.
    A less conservative fatigue analysis that quantifies the fatigue lifetime, such as Miner's law, could in turn couple with the economics module and wave statistics to understand the tradeoff between capital cost and lifetime.
    Explicitly modeling the storm sea state contour as a function of location and device lifetime is possible by integrating with \gls{acr-wdrt} or a similar tool.
    This is of interest in design scenarios where the \gls{acr-wec} lifetime is a design variable.
    This may prove useful to explore the tradeoffs of single-use (set-and-forget) \glspl{acr-wec} for far-offshore sensing applications where maintenance is particularly difficult or costly.
    
    \paragraph{Different \gls{acr-wec} archetypes}
    Extending the model beyond 2-body heaving point absorbers to other \gls{acr-wec} archetypes such as oscillating water columns, overtopping devices, or surge/pitch devices would allow systematic quantification of the comparative benefits and drawbacks of each \gls{acr-wec} archetype in a consistent way.
    This capability is especially valuable in the early concept design phase to inform architectural decisions and could facilitate design convergence for the field as a whole, which is considered an industry bottleneck \citep{caio_tackling_2019}.
    Generalizing MDOcean's dynamics/controls and economics modules to other archetypes would require minor adaptations including alternate degrees of freedom and cost fits, but not structural modification of the core models.
    On the other hand, the hydrodynamics and structures models would require significant modification since separate semi-analytical continuum models would need to be implemented for each new geometry.
    Before initiating such an undertaking, it would be important to confirm that numerical alternatives remain too inefficient for optimization.
    Since development of MDOcean began in 2021, new numerical tools have emerged that could potentially reduce the computational cost of mesh-based geometry-agnostic computations.
    For example, a differentiable \gls{acr-bem} solver in Julia \texttt{MarineHydro.jl} was released in 2025 and is in its infancy \citep{khanal_fully_2025}, while a lightweight \gls{acr-fea} tool with plate/shell elements \texttt{pynite} was released in 2021 and has seen widespread adoption.
    Whether semi-analytical or numerical, extension to other archetypes would require significant development effort but unlock new comparative insights.
\else
    The most consequential current limitations are:
\begin{itemize}[leftmargin=*]
    \item \textbf{Neglect of surge force in the structures module and absence of a mooring cost model.} Surge forces on the nominal \gls{acr-rm3} are \resultsAOR[surgeForceFloatNominal] (float) and \resultsAOR[surgeForceSparNominal] (spar), and mooring/foundation accounts for 12\% of \gls{acr-capex} (\Cref{tab:CBS}); incorporating these would likely affect optimal designs.
    \item \textbf{Regular-wave assumption in storm load cases.} Storm waves are nonlinear, and the regular-wave equivalent cannot capture transient peaks; second-order \gls{acr-meem} \citep{cong_novel_2020,mavrakos_second-order_2009} or the slender-body approximation in \gls{acr-raft} \citep{carmo_slender-body_2025} are candidate extensions.
    \item \textbf{Regular-wave assumption in operational loads.} Stochastic linearization \citep{da_silva_statistical_2020,da_silva_stochastic_2023,kluger_synergistic_2017,folley_spectral-domain_2016,spanos_efficient_2016,neshat_enhancing_2024} could replace the describing function to handle spectral inputs and enable spectral fatigue, grid-integration, and storage-sizing analyses \citep{mccabe_wec_2025}.
    \item \textbf{Single \gls{acr-wec} archetype.} Extending \gls{acr-meem} to handle the damping plate or 3+ concentric cylinders \citep{best_openflash_2026,bimali_matrix_2026} would broaden applicability, as would generalizing the framework to non-point-absorber archetypes.
\end{itemize}
    
    Additional model refinements that would build trust without unlocking new study types include local buckling and effective-breadth criteria in the structures module, \gls{acr-fea}-based validation of the structural model, additional \gls{acr-pto} nonlinearities (drivetrain friction, check-valve rectification, power saturation), and Miner's-law fatigue analysis coupled with storm-sea-state contour modeling for finite-lifetime designs.
    Per-item rationale and implementation paths are in \citet{mccabe_dissertation_2026}.
\fi

\ifdefined\DISSERTATION
    \paragraph{Runtime Improvements}
    The runtime of the simulation could also be improved, either by further simplifying the model
    or by implementing code optimizations to avoid unnecessary computations.
    Early MDOcean versions (v0.2.0 to v0.3.1) ran in under \resultsAOR[simRuntimeOld],
    about one-%
    \makeatletter%
    \phantom{\rlap{% make upcoming print invisible and zero width (hack because you can't retrieve value without also printing it)
    \resultsAOR[multiplierOldRuntime] % print value, which locally sets resultsAOR@out
    \global\let\SavedMultiplier\resultsAOR@out % save resultsAOR@out globally so it can be accessed outside the phantom environment
    }}%
    \makeatother%
    \ordinalstringnum{\SavedMultiplier}
    of the time of the current version.
    These versions feature largely the same disciplinary scope but use fixed-point iteration rather than a solver to converge the drag loop.
    This results in far fewer evaluations of the dynamics but an inability to converge the drag loop for some sea states under reactive control.
    Other instances where returning to a simpler model that could be justified in certain instances and would improve runtime include
    turning off the \gls{acr-qlps} constrained optimal control in favor of a suboptimal control scheme
    and reverting to a simpler drag formulation that neglects the strip-theory integration of the drag force.
    
    To increase speed while preserving the current features, recommended areas for code optimization include:
    \begin{enumerate}
        \item Share calculations between the multibody dynamics evaluation function and the \gls{acr-qcqp} function
        \item Refactor to enhance vectorization
        \item Strategically balance complex-valued and real-valued operations in the dynamics module: a given transfer function may be implemented in a small number of expensive complex operations or a large number of inexpensive real operations, and the preferred balance has not yet been systematically explored. Deferring the introduction of complex numbers until later in the calculation may also be beneficial.
        \item Decrease the number of iterations required to converge the drag loop (e.g., improve the initial guess, tune solver parameters, or solve the nonlinear equations analyticaly to avoid iteration entirely).
    \end{enumerate}
\else
    The runtime of the simulation could also be improved by revisiting earlier model versions that achieved
    even faster (<40~ms) runtimes but exhibited difficulty converging,
    and by code optimizations such as sharing calculations between the dynamics and controls,
    enhancing vectorization, balancing complex and real operations,
    and solving the drag equation analytically to obviate the need for iteration entirely.
\fi

\section{Conclusion}
% Summary:
This paper develops MDOcean, an open-source semi-analytical modeling framework for two-body axisymmetric point-absorber wave energy converters that integrates hydrodynamics, dynamics, control, structures, and economics in a single fast simulation.

% Model
\ifdefined\DISSERTATION
    The geometry module calculates volumes, areas, and the pitch stability margin with simple algebraic relations, and the economics module calculates design-dependent cost from structural material volume and powertrain force and power ratings.
    The hydrodynamics module implements the \gls{acr-meem} for a dual concentric cylinder geometry to obtain the float radiation coefficients, and interpolates existing \gls{acr-bem} solutions to obtain the spar radiation coefficients due to the lack of a damping plate in the current \gls{acr-meem} model.
    This achieves an order of magnitude speedup compared to the popular \gls{acr-bem} solver Capytaine.
    The dynamics module employs describing functions to extend the typical frequency-domain formulation to capture nonlinearities and dynamic constraints, including viscous drag and generator force saturation, and includes a novel analytical \gls{acr-qcqp} for the constrained linear controller.
    This achieves a three-order-of-magnitude computational speedup over the common time-domain solver WEC-Sim.
    The structures module uses an equivalent-thickness model of stiffened plates and a library of canonical solutions to structural boundary-value problems to assess the factor of safety of the float bottom plate, spar shell, and damping plate in cyclic fatigue and storm load cases.
    This represents, to the authors' knowledge, the only public \gls{acr-wec} structural simulation regardless of modeling technique.
    Finally, the economics module calculates the design-dependent cost based on structural material volume and powertrain force and power rating, performing a straightforward \gls{acr-lcoe} calculation.
\else
    The geometry module handles hydrostatics.
    The hydrodynamics module uses \gls{acr-meem} for the float, achieving an order-of-magnitude speedup over \gls{acr-bem}.
    The dynamics module employs describing functions and a \gls{acr-qcqp} for constrained linear control,
    which we demonstrate can be solved analytically to achieve a three-order-of-magnitude speedup over time-domain solvers.
    This module represents a relevant methodological contribution to the broader \gls{acr-wec} modeling community,
    introducing and contextualizing \gls{acr-qlps} and \gls{acr-mfd} within constrained optimal control approaches.
    The structures module uses semi-analytical stiffened-plate models with realistic failure criteria for fatigue and storm load cases---to the authors' knowledge, the only public \gls{acr-wec} structural simulation.
    The economics module includes a calibrated cost model.
\fi

Validation against WEC-Sim and the \gls{acr-rm3} reference design demonstrates that MDOcean reproduces the dynamics of an established time-domain solver to within a few percent under matched modeling assumptions and within single-digit-percent \gls{acr-jpd}-weighted annual average power under realistic conditions.
Three tuned scale factors account for assumption differences between MDOcean and the reference model, reducing error to under 1\%.
Total simulation runtime is \resultsAOR[simRuntime], which is 1-3 orders of magnitude faster than established baselines.

These results suggest that semi-analytical multidisciplinary modeling provides a useful intermediate between simplified conceptual studies and computationally intensive high-fidelity simulation,
improving physical insight during early-stage \gls{acr-wec} development
and unlocking future \linelabel{lne:conclusion-mdo-start}multidiscplinary optimization workflows that would otherwise be computationally prohibitive (such as the RM3 optimization performed in the companion paper \citep{mccabe_leveraging_2026}).\linelabel{lne:conclusion-mdo-end}

% Discussion
\ifdefined\DISSERTATION
    A key feature of this study that other \gls{acr-wec} models lack is the inclusion of more realistic effects such as shape-informed structural stress, force saturation, amplitude limits, and a nonlinear drag model.
    Nonetheless, the model has limitations: the structural loads and stresses have non-trivial uncertainty due to the regular-wave assumption in storm load cases; the impact of irregular waves on operational power and constraints is approximated rather than directly modeled; surge force is not propagated to the structures module and no mooring cost model is included; and only a single \gls{acr-wec} archetype (two-body axisymmetric point absorber) is considered.
    Future modeling work should pursue these enhancements alongside extensions to additional \gls{acr-wec} archetypes, generator magnetics, and probabilistic / spectral analyses for fatigue, grid integration, and storage sizing.
\else
    Principal limitations include the linear potential flow assumption in storm load cases,
    approximation of irregular waves as equivalent regular waves,
    and restriction to a single \gls{acr-wec} archetype.
    \Cref{sec:unmodeled-effects} gives detailed future work directions.
\fi

% Impact of findings:
MDOcean is \linelabel{lne:open-source-conclusion}open-source and is intended to support continued community development of fast multidisciplinary \gls{acr-wec} modeling tools that integrate hydrodynamics, controls, structures, and economics in a unified framework.

\section*{Acknowledgements}

The authors thank Kapil Khanal, Yinghui Bimali, En Lo, and John Fernandez for assistance with hydrodynamics;
Fabien Royer for guidance on structures;
Ryan Coe, Jacob Mays, Patrick Reed, Nate DeGeode, and Alaa Ahmed for technical feedback on a draft manuscript;
Nola McCabe for proofreading support;
Pete Bachant for support on the reproducible computation pipeline;
and Jeff Scruggs for suggestion of resources on quasi-linearization.

Rebecca McCabe acknowledges funding from the National Science Foundation Graduate Research Fellowship Program, Cornell Provost's Diversity Fellowship, and Cornell Mechanical and Aerospace Engineering McMullen Fellowship.
Madison Dietrich acknowledges funding from the Cornell Duffield Engineering Fund for Undergraduate Research on Solutions to Climate Change and the Bill Nye '77 Award in Undergraduate Research.

This material is based on work supported by the National Science Foundation Graduate Research Fellowship under Grant No.~DGE--2139899.
Any opinions, findings, conclusions, or recommendations expressed in this material are those of the authors and do not necessarily reflect the views of the National Science Foundation.

%The graphical abstract uses icons from \url{www.flaticon.com} in accordance with the Flaticon license.

\section*{Data availability statement}

The MATLAB code for all simulation, analysis, and visualization to fully reproduce this work is available open-source via the MDOcean project at
\url{https://github.com/symbiotic-engineering/MDOcean} \citep{mccabe_mdocean_2024}.
The version used in this paper is v1.5.3.
Questions and contributions via GitHub issues and pull requests are welcomed.
%A Python version using the \texttt{OpenMDAO} package is forthcoming.
All computational environments, results, and \LaTeX~artifacts can be reproduced via a data version control pipeline built with the \texttt{Calkit} package
and accessed at \url{https://calkit.io/symbiotic-engineering/mdocean}.
Sphinx API documentation is available at \url{https://symbiotic-engineering.github.io/MDOcean/}.
See the documentation for full instructions for reproducing results and using the MDOcean code.

\section*{Declaration of interests}
Rebecca McCabe reports administrative support was provided by New York State Education Department Adult Career and Continuing Education Services Vocational Rehabilitation Services.
Rebecca McCabe reports travel was provided by \gls{acr-umerc}.
Maha Haji reports a relationship with \gls{acr-umerc} that includes board membership. 
Madison Dietrich declares no known competing financial interests or personal relationships that could have appeared to influence the work reported in this paper.

\section*{Declaration of generative AI and AI-assisted technologies in the manuscript preparation process}
During the preparation of this work, the authors used the following generative AI tools:
Anthropic Claude (accessed May 2026);
OpenAI ChatGPT (GPT-5.5, GPT-5.3, GPT-5.3-mini, and GPT-5-mini: accessed June 2025--June 2026);
GitHub Copilot Cloud Agent (GPT-5.3 Codex: accessed March 2026--June 2026);
GitHub Copilot Chat Extension for Visual Studio Code (Claude Sonnet 4, Claude Sonnet 4.5, GPT-4.1, GPT-5.4 mini, and GPT-5 mini: accessed February 2024--June 2026); and
Google Search AI Overview (Gemini 1.0, Gemini 1.5, Gemini 2.0, Gemini 2.5, Gemini 3, Gemini 3.5: accessed May 2024--June 2026).
AI tools were dominantly used for the generation and debugging of code and computational workflows not core to the technical contributions of the work.
This includes generation of changelogs, pull request reviews, and continuous integration pipeline features;
execution of merges, dependency version bumps, and compatibility updates;
refactoring of code for improved organization and readability;
syntax reference and usage examples for MATLAB, Python, Bash, Mermaid, \LaTeX, and their associated packages;
resolution of errors and warnings encountered in the above languages during human-led code development and manuscript preparation;
and the execution of specific formatting refinements to figures without altering their core content.

The following exceptions describe instances where AI tools were used to assist with technical development, mathematical derivations, or figure generation related to the core technical contributions of the work:
\begin{itemize}[leftmargin=*]
\item Based on author-provided descriptions of the desired content and formatting, ChatGPT generated the initial source code for the taxonomy diagram (Mermaid), control synthesis flowchart (Mermaid), control Venn diagram (TikZ), and \gls{acr-meem} geometry diagram (TikZ).
\item GitHub Copilot Cloud Agent checked that the notation and equations were consistent between the manuscript and the codebase. The authors resolved all identified inconsistencies.
\item Based on author-provided descriptions of the desired content and formatting, GitHub Copilot Cloud Agent generated the initial source code for the geometry sweep Pareto front, WEC-Sim total harmonic distortion figure, the \gls{acr-pto} sweep contour plot, and the overlay of \gls{acr-wamit} results on the existing spar excitation \gls{acr-meem} result plot.
\item After the authors developed and debugged the optimal control code for generic quadratic constraints and implemented the force constraint, GitHub Copilot Cloud Agent integrated the amplitude and power constraints following the same pattern, and implemented some changes to reduce runtime while preserving identical outputs.
\item ChatGPT was used to conduct targeted literature review to identify relevant optimization algorithms for systems with structured coupling, optimal control techniques to populate the Venn diagram, and references describing convex optimization sensitivity analysis.
\item With the intent of balancing technical accuracy and approachability of the terminology for broad audiences, the authors prompted Github Copilot Cloud Agent to generate the two sentences that conceptually introduce convexity in \Cref{sec:appendix-qp-convexity}.
% thesis version only: chatgpt was used for in the same way for the sentences describing types of nonlinear pressures in the drag section.
\item ChatGPT was used to assist with the algebraic manipulation, simplification, and/or initial typesetting of the following equations, which were then verified by hand and/or with symbolic math software:
the drag integral limit formulas in \Cref{sec:appendix-drag};
the circle/ellipse intersection formulas in \Cref{sec:appendix-qp-solution};
the sensitivity formulas for the \gls{acr-pto} limit in \Cref{sec:appendix-constraint-sensitivity};
the curvature of the function in \Cref{sec:appendix-qp-convexity};
the trigonometric simplifications in the slamming angle formula in \Cref{sec:appendix-slam}.
\item ChatGPT was used to assist with the mathematical reasoning and interpretation of the following technical concepts, which were then verified via code implementation: the passivity, stability, and cardinality of underactuated multiport systems in \Cref{sec:appendix-pto-dynamics}, and curve fit structures that maintain desired limiting properties for the spar excitation nondimensionalization in \Cref{sec:appendix-additional-hydro}.
\item Claude and ChatGPT were used editorially to improve clarity, grammar, style, and readability of the text.
\end{itemize}
After using these tools, the authors reviewed and edited the content as needed and take full responsibility for the content of the published article.
Author-generated continuous-integration tests and validation cases were used to verify the correctness of simulation code, both AI- and human-generated.

% GitHub copilot chat models:
%   From ~/AppData/Roaming/Code/User/workspaceStorage/*/chatSessions/*.json:
%     copilot/auto
%     copilot/claude-sonnet-4
%     copilot/gpt-4.1
%     copilot/gpt-5-mini
%     github.copilot-chat/claude-sonnet-4
%   From github useage history:
%     claude-sonnet-4.5
%     gpt-5.4 mini
%     gpt-5 mini

\appendix

\ifdefined\JOURNALTWOCOLUMN
    \onecolumn%
\fi

%\appendix

\section{\acrlong{acr-meem}}
\label{sec:appendix-meem-details}
As introduced in \Cref{sec:hydro-meem}, the hydrodynamic coefficients are computed semi-analytically using \gls{acr-meem}.
This section explains the \gls{acr-meem} formulation and solution methodology for the radiation of a dual truncated concentric cylinder geometry, originally presented by \citet{mavrakos_hydrodynamic_2004,chau_inertia_2010,chau_inertia_2012} as an extension of the single-cylinder \gls{acr-meem} radiation solution in the study by \citet{yeung_added_1981}.
Further numeric and realization details of the authors' implementation may be found in references \citet{bimali_matrix_2026,best_openflash_2026,mccabe_open-source_2024}.
The computation involves splitting the fluid domain into regions, approximating an infinite series by truncation, and solving a matrix equation to enforce the continuity of potential and velocity across regions.

\subsection{Linear Hydrodynamics and Eigenfunctions}
The dynamics of a floating body in water waves are well-described by linear potential flow theory, a simplification of the Navier-Stokes equation.
This theory states that the fluid velocity field is the gradient of some complex potential $\gls{sym-phi}$, $\gls{sym-vec-v}=\nabla\gls{sym-phi}$, and $\gls{sym-phi}$ satisfies the Laplace equation, $\nabla^2\gls{sym-phi}=0$.
Adding the free surface condition, far-field or incident waves, and body surface and sea-bed conditions detailed in \citet{chatjigeorgiou_analytical_2018} yields a boundary value problem.
When boundary conditions correspond to the heave radiation problem (body moving vertically, no incident waves), solving for the potential $\gls{sym-phi}(\gls{sym-r},\gls{sym-theta},\gls{sym-z})$ determines the heave added mass and damping $\gls{sym-A-h}$ and $\gls{sym-B-h}$, hereafter ``hydro coefficients.”

For appropriate geometries, the partial differential equation is separable and $\gls{sym-phi}$ can be expressed as the product of radial, vertical, and circumferential basis functions called eigenfunctions.
In cylindrically symmetric problems, the radial eigenfunctions are a family of transcendental functions called Bessel functions.
The fluid is then divided into cylindrical regions.
Arbitrarily many fluid regions can exist, so the method applies to any axisymmetric geometry, including multiple concentric bodies that oscillate independently.
Here, two concentric cylinders and thus three fluid regions are demonstrated.
Extension to many regions is discussed in section 3.4.
\Cref{fig:meem-regions} illustrates the regions and dimensions: two interior regions \textit{i1} and \textit{i2}, and an exterior region \textit{e} extending to infinity.

\begin{figure}[htb]
\centering
\includegraphics[\FigureOptions{aor:figs/from-matlab/meem_dims.pdf}]{figs/from-matlab/meem_dims.pdf}
\caption{Fluid regions and dimensions used in the dual concentric cylinder \gls{acr-meem}}
\label{fig:meem-regions}
\end{figure}

 The potential in each region is split into a homogeneous part for the unforced solution and a particular part due to body motion: $\gls{sym-phi}=\gls{sym-phi-h}+\gls{sym-phi-p}$.
Boundary conditions dictate $\gls{sym-phi-p}$ and the eigenfunctions for $\gls{sym-phi-h}$ in each region, which the textbook \citet{chatjigeorgiou_analytical_2018} describes in detail.
\Cref{tab:MEEM-eigenfunctions} shows the equations originally presented in the studies by \citet{chau_inertia_2010,chau_inertia_2012} for the potential and eigenfunctions in each region, which include infinitely many unknown eigencoefficients $\gls{sym-C-1n-i1}$, $\gls{sym-C-1m-i2}$, $\gls{sym-C-2m-i2}$ and $\gls{sym-B-k-e}$.
By construction, this potential obeys all boundary conditions except for zero radial velocity on radial body surfaces.
The unknown coefficients must be computed to enforce this final condition as well as continuity across regions, which will be the subject of \Cref{sec:appendix-meem-matching}.

\newcommand{\dispersion}{
    \begin{cases}
        \gls{sym-m-0}\tanh (\gls{sym-m-0}\gls{sym-h})
            =\gls{sym-omega}^{2}/\gls{sym-g},
            & \gls{sym-k-index}=0,\\
        \gls{sym-m-k}\tan (\gls{sym-m-k}\gls{sym-h})
            =-\gls{sym-omega}^{2}/\gls{sym-g},
            & \gls{sym-k-index}\geq 1
    \end{cases}
}

\begin{table}[htbp]
    \centering
    \expanded{
        \noexpand\tabular{\TableColumnSpecs{aor:meem-eigenfunctions}}
    }
    \hline 
    Region
        &
        $i1$
        &
        $i2$
        &
        $e$
        \\
        \hline 
    Homog. potential $\gls{sym-phi-h}(\gls{sym-r},\gls{sym-z})$
        &
        $\displaystyle\sum_{\gls{sym-n}}
            \gls{sym-C-1n-i1}
            \gls{sym-R-1n-i1} (\gls{sym-r})
            \gls{sym-Z-n-i1} (\gls{sym-z})$
        &
        $\displaystyle\sum_{\gls{sym-m}}\left(\gls{sym-C-1m-i2}\,
            \gls{sym-R-1m-i2} (\gls{sym-r})
            +
            \gls{sym-C-2m-i2}\,
            \gls{sym-R-2m-i2} (\gls{sym-r})
        \right)
        \times~
        \gls{sym-Z-m-i2} (\gls{sym-z})$
        &
        $\displaystyle\sum_{\gls{sym-k-index}}
            \gls{sym-B-k-e}
            \gls{sym-Lambda-k-e} (\gls{sym-r})
            \gls{sym-Z-k-e} (\gls{sym-z})$ \\
    Partic. potential $\gls{sym-phi-p}(\gls{sym-r},\gls{sym-z})$
        &
        $\displaystyle
        \frac{1}{2 (\gls{sym-h}-\gls{sym-d-1})}
        \left[
            (\gls{sym-z}+\gls{sym-h})^2
            -
            \frac{\gls{sym-r}^2}{2}
        \right]$
        &
        $\displaystyle
        \frac{1}{2 (\gls{sym-h}-\gls{sym-d-2})}
        \left[
            (\gls{sym-z}+\gls{sym-h})^2
            -
            \frac{\gls{sym-r}^2}{2}
        \right]$
        &
        $0$
        \\
        \hline 
    Radial eigen-function $\gls{sym-R-r}$
        &  $\gls{sym-R-1n-i1} (\gls{sym-r})
            =
            \begin{cases}
                \frac12
                    & \gls{sym-n}=0\\[1em]
                \displaystyle
                \frac{
                    \gls{sym-mathrm-I-0}
                    (\gls{sym-lambda-n-i1}\gls{sym-r})
                }{
                    \gls{sym-mathrm-I-0}
                    (\gls{sym-lambda-n-i1}\gls{sym-a-2})
                }
                    & \gls{sym-n}\ge1
            \end{cases}$
        &  \shortstack{
            $
            \gls{sym-R-1m-i2} (\gls{sym-r})
            =
            \begin{cases}
                \frac12
                    & \gls{sym-m}=0\\[1em]
                \displaystyle
                \frac{
                    \gls{sym-mathrm-I-0}
                    (\gls{sym-lambda-m-i2}\gls{sym-r})
                }{
                    \gls{sym-mathrm-I-0}
                    (\gls{sym-lambda-m-i2}\gls{sym-a-2})
                }
                    & \gls{sym-m}\ge1
            \end{cases}
            $
            \\
            $
            \gls{sym-R-2m-i2} (\gls{sym-r})
            =
            \begin{cases}
                \frac12
                \ln\!\left(\frac{\gls{sym-r}}
                        {\gls{sym-a-2}}
                \right)
                    & \gls{sym-m}=0\\
                \displaystyle
                \frac{
                    \gls{sym-mathrm-K-0}
                    (\gls{sym-lambda-m-i2}\gls{sym-r})
                }{
                    \gls{sym-mathrm-K-0}
                    (\gls{sym-lambda-m-i2}\gls{sym-a-2})
                }
                    & \gls{sym-m}\ge1
            \end{cases}
            $
            }
        &
        $\gls{sym-Lambda-k-e} (\gls{sym-r})
            =
            \begin{cases}
                \displaystyle
                \frac{
                    \gls{sym-mathrm-H-0-1}
                    (\gls{sym-m-0}\gls{sym-r})
                }{
                    \gls{sym-mathrm-H-0-1}
                    (\gls{sym-m-0}\gls{sym-a-2})
                }
                    & \gls{sym-k-index}=0\\[1em]
                \displaystyle
                \frac{
                    \gls{sym-mathrm-K-0}
                    (\gls{sym-m-k}\gls{sym-r})
                }{
                    \gls{sym-mathrm-K-0}
                    (\gls{sym-m-k}\gls{sym-a-2})
                }
                    & \gls{sym-k-index}\ge1
            \end{cases}$
            \\
            \hline 
    Vertical eigen-function $\gls{sym-Z}(\gls{sym-z})$
        &
        $\gls{sym-Z-n-i1} (\gls{sym-z})
        =
        \sqrt{2}\times
        \begin{cases}
            1/\sqrt{2}
                & \gls{sym-n}=0\\[1em]
            \cos\!\left(\gls{sym-lambda-n-i1}
                (\gls{sym-z}+\gls{sym-h})
            \right)
                & \gls{sym-n}\ge1
        \end{cases}$
        &
        $\gls{sym-Z-m-i2} (\gls{sym-z})
        =
        \sqrt{2}\times
        \begin{cases}
            1/\sqrt{2}
                & \gls{sym-m}=0\\[1em]
            \cos\!\left(\gls{sym-lambda-m-i2}
                (\gls{sym-z}+\gls{sym-h})
            \right)
                & \gls{sym-m}\ge1
        \end{cases}$
        &
        $\gls{sym-Z-k-e} (\gls{sym-z})
        =
        \gls{sym-N-k}^{-\frac12}
        \times
        \begin{cases}
            \cosh\!\left(\gls{sym-m-0}
                (\gls{sym-z}+\gls{sym-h})
            \right)
                & \gls{sym-k-index}=0\\[1em]
            \cos\!\left(\gls{sym-m-k}
                (\gls{sym-z}+\gls{sym-h})
            \right)
                & \gls{sym-k-index}\ge1
        \end{cases}$
        \\
        \hline
    Eigen-value& 
        $\displaystyle
        \gls{sym-lambda-n-i1}
        =
        \frac{\gls{sym-n}\gls{sym-pi}}
            {\gls{sym-h}-\gls{sym-d-1}},
        \quad
        \gls{sym-n}\ge1$
        &
        $\displaystyle
        \gls{sym-lambda-m-i2}
        =
        \frac{\gls{sym-m}\gls{sym-pi}}
            {\gls{sym-h}-\gls{sym-d-2}},
        \quad
        \gls{sym-m}\ge1$
        &
        \ifdefined\DISSERTATION
            $\gls{sym-m-k}$ (see \Cref{aor:eq:dispersion})
        \else
            $\displaystyle
            \dispersion
            $
        \fi
    \\
    \hline
    \endtabular
    \caption{Equations for potential (homogeneous and particular) and eigenfunctions (radial and vertical) for each region.}
    \label{tab:MEEM-eigenfunctions}
\end{table}
\ifdefined\DISSERTATION
    The dispersion relation is a nonlinear equation that implicitly gives the exterior-region eigenvalues $\gls{sym-m-k}$:
    \begin{equation}
        \label{aor:eq:dispersion}
    \dispersion
    \end{equation}
\fi
In \Cref{tab:MEEM-eigenfunctions},  $\gls{sym-mathrm-I-0}$,  $\gls{sym-mathrm-K-0}$, and $\gls{sym-mathrm-H-0-1}$ are different Bessel functions of order zero; 1M and 2M mean body 1 and 2 (spar and float) are moving respectively, while 1S and 2S mean each is stationary; and the $\gls{sym-N-k}$ expression is defined as:

\begin{equation}
    \gls{sym-N-k} = \frac{1}{2}\left(1+\frac{\gls{sym-f-k}}{2\gls{sym-m-k}\gls{sym-h}}  \right)~
    \textrm{ where }
    \gls{sym-f-k} =
    \begin{cases}
        \sinh(2\gls{sym-m-0}\gls{sym-h}), & \gls{sym-k-index}=0 \\ \sin(2\gls{sym-m-k}\gls{sym-h}), & \gls{sym-k-index}\geq1
    \end{cases}
\end{equation}

 %%%%%%%%%%%%%%%%%%%%%%%%%%%%%%%%%%%%%%%%%%%%%%%%%%%%%%%%%%%%%%% begin equation helper functions %%%%%%%%%%%%%%%%%%%%%%%%%%%%%%%%%%%%%%%%%%%%%%%%%%%%%%%%%%%
%% integral definitions
\newcommand{\RintOneDefn}{
    \shortstack{
        $\displaystyle
            \gls{sym-boldsymbol-mathcal-R}_{1j} =
            \int\limits_{\gls{sym-a-in}}^{\gls{sym-a-out}} \gls{sym-vec-R}_{1j} (\gls{sym-r})\gls{sym-r} dr$
        \\
            for $\gls{sym-j}= (\gls{sym-n},\gls{sym-m})
        $
    }
}

\newcommand{\RintTwoDefn}{
    $\displaystyle \gls{sym-boldsymbol-mathcal-R}_{2\gls{sym-m}} = \int\limits_{\gls{sym-a-in}}^{\gls{sym-a-out}} \gls{sym-vec-R}_{2\gls{sym-m}}(\gls{sym-r})rdr$
}

\newcommand{\ZmnDefn}{
    \shortstack{
        $\displaystyle \gls{sym-boldsymbol-mathcal-Z}_{nm} = \gls{sym-boldsymbol-mathcal-Z}_{mn}^\top =$
        \\
        $\displaystyle \int_{-\gls{sym-h}}^{-\gls{sym-d-1}} \gls{sym-vec-Z-n-i1}^{\top}\gls{sym-vec-Z-m-i2} dz$
    }
}
\newcommand{\ZmkDefn}{
    \shortstack{
        $\displaystyle \gls{sym-boldsymbol-mathcal-Z}_{mk} = \gls{sym-boldsymbol-mathcal-Z}_{km}^\top = $
        \\
        $\displaystyle \int_{-\gls{sym-h}}^{-\gls{sym-d-2}} \gls{sym-vec-Z-m-i2}^{\top} \gls{sym-vec-Z-k-e} dz $
    }
}

%% R definitions
\newcommand{\RintOneJzero}{
    $\displaystyle\frac{\gls{sym-a-out}^2-\gls{sym-a-in}^2}{4}$
}

\newcommand{\RintOneJOne}{
    $\displaystyle
        \frac
            {\gls{sym-a-out} \gls{sym-mathrm-I-1} (\gls{sym-vec-lambda}_{\gls{sym-j}} \gls{sym-a-out}) - \gls{sym-a-in}\gls{sym-mathrm-I-1} (\gls{sym-vec-lambda}_{\gls{sym-j}} \gls{sym-a-in})}
            {\gls{sym-vec-lambda}_{\gls{sym-j}} \gls{sym-mathrm-I-0} (\gls{sym-vec-lambda}_{\gls{sym-j}} \gls{sym-a-2})}
    $
}

\newcommand{\RintTwoJzero}{
    $\displaystyle
        \frac
        {
            \gls{sym-f-2m} \gls{sym-a-in}^2-\gls{sym-a-out}^2
        }
        {8}
    $
    where $
        \gls{sym-f-2m} = 1 + 2\ln\frac{\gls{sym-a-out}}{\gls{sym-a-in}}
    $
}

\newcommand{\RintTwoJOne}{
    $\displaystyle
    \frac
    {
        \gls{sym-a-in}\,\gls{sym-mathrm-K-1} (\gls{sym-vec-lambda-m} \gls{sym-a-in})
        -\gls{sym-a-out}\,\gls{sym-mathrm-K-1} (\gls{sym-vec-lambda-m} \gls{sym-a-out})
    }
    {
        \gls{sym-vec-lambda-m} \,\gls{sym-mathrm-K-0} (\gls{sym-vec-lambda-m} \gls{sym-a-2})
    }
    $
}

% Znm formulas
\newcommand{\ZnZeroMZero}{
    $\gls{sym-h}-\gls{sym-d-1}$
}

\newcommand{\ZnZeroMOne}{
    $\displaystyle
     \frac
                {\sqrt{2}\sin\left(
				\gls{sym-vec-lambda-m} (\gls{sym-h} - \gls{sym-d-1})
			\right)}
                { \gls{sym-vec-lambda-m}}
    $
}

\newcommand{\ZnOneMZero}{
    $\displaystyle 0
     % this more complicated formula is only needed for multi-region
     %\frac
      %          {\sqrt{2}\sin (\gls{sym-vec-lambda-n}^{\top} (\gls{sym-h} - \gls{sym-d-1}))}
        %        { \gls{sym-vec-lambda-n}^{\top}}
    $
}

\newcommand{\ZnOneMOne}{
$\displaystyle
\begin{cases}
\displaystyle\frac{
            2 (-1)^{\gls{sym-vec-n}}
            }
            {
                    1 - \left(\frac{\gls{sym-vec-lambda-n}^{\top}}{  \gls{sym-vec-lambda-m}}
                        \right)^2
            }
\displaystyle\frac{
            \sin\left(
                \gls{sym-vec-lambda-m} (\gls{sym-h}-\gls{sym-d-1})
         	\right)
            }
            {
            \gls{sym-vec-lambda-m}
            }
,&  \gls{sym-lambda-n-i1} \neq \gls{sym-lambda-m-i2}\\
\gls{sym-h} -\gls{sym-d-1}, & \gls{sym-lambda-n-i1} = \gls{sym-lambda-m-i2}
\end{cases}
$
% this more complicated formula is only needed for multi-region
% \shortstack{
%     $\displaystyle
%     \frac
%         {\sin\left(% 		(\gls{sym-vec-lambda-n}^{\top} + \gls{sym-vec-lambda-m}) (\gls{sym-h}-\gls{sym-d-1})
% 	\right)}
%         {\gls{sym-vec-lambda-n}^{\top} + \gls{sym-vec-lambda-m}}
%     + \gls{sym-Z}_{nm}^-$
%     \\
%     where $\gls{sym-Z}_{nm}^-=
%     \begin{cases}
%         \frac
%             {\sin\left(% 		(\gls{sym-vec-lambda-n}^{\top} - \gls{sym-vec-lambda-m}) (\gls{sym-h}-\gls{sym-d-1})
% 	\right)}
%             {(\gls{sym-vec-lambda-n}^{\top} - \gls{sym-vec-lambda-m})},
%         & \gls{sym-lambda-n-i1} \neq \gls{sym-lambda-m-i2}
%         \\
%         \gls{sym-h}-\gls{sym-d-1},
%         & \gls{sym-lambda-n-i1} = \gls{sym-lambda-m-i2}
%     \end{cases}
%   $
%  }
}

% Zmk formulas
\newcommand{\ZmZeroKZero}{
    $\displaystyle \frac{\sinh(\gls{sym-m-0}(\gls{sym-h}-\gls{sym-d-2}))}{\gls{sym-m-0} \sqrt{\gls{sym-N-0}}}$
}

\newcommand{\ZmOneKZero}{
    $\displaystyle \frac{\sqrt{2}~\gls{sym-m-0} (-1)^{\vec{\gls{sym-m}}} \sinh(\gls{sym-m-0}(\gls{sym-h}-\gls{sym-d-2}))}{\sqrt{\gls{sym-N-0}} \left(\gls{sym-m-0}^2+{(\gls{sym-vec-lambda-m}^{\top})}^2\right)}$
}

\newcommand{\ZkOneMZero}{
    $\displaystyle\frac{
                        \sin (\gls{sym-vec-m-k} (\gls{sym-h}-\gls{sym-d-2}))
                        }
                        {
                        \gls{sym-vec-m-k}
                        \sqrt{\vec{\gls{sym-N-k}}}
                        }$
}

\newcommand{\ZkOneMOne}{
    $\displaystyle
    \begin{cases}
            \displaystyle\frac{1}{\sqrt{2\vec{\gls{sym-N-k}}}}
            \left(
                \displaystyle\frac
                    {
                        \sin\left((\gls{sym-h}-\gls{sym-d-2})
                            (\gls{sym-vec-m-k} + \gls{sym-vec-lambda-m}^{\top})
                        \right)
                    }
                    {
                        \gls{sym-vec-m-k}
                        + \gls{sym-vec-lambda-m}^{\top}
                    }
            + \right.
                 \\
                \left.~~~~~~~~~
                \displaystyle\frac
                    {\sin\left((\gls{sym-h}-\gls{sym-d-2}) (\gls{sym-vec-m-k}- \gls{sym-vec-lambda-m}^{\top})\right)}
                    {\gls{sym-vec-m-k}- \gls{sym-vec-lambda-m}^{\top}}
            \right),
        & |\gls{sym-m-k}| \neq \gls{sym-lambda-m}
        \\
        \displaystyle\frac{\gls{sym-h}-\gls{sym-d-2}}{2},
        & |\gls{sym-m-k}| = \gls{sym-lambda-m}
    \end{cases}$
}

 %%%%%%%%%%%%%%%%%%%%%%%%%%%%%%%%%%%%%%%%%%%%%%%%%%%%%%%%%%%%%%% end equation helper functions %%%%%%%%%%%%%%%%%%%%%%%%%%%%%%%%%%%%%%%%%%%%%%%%%%%%%%%%%%%

\Cref{tab:meem-integrals} lists several integrals of the radial and vertical eigenfunctions, $\gls{sym-boldsymbol-mathcal-R}$ and $\boldsymbol{\mathcal{\gls{sym-Z}}}$ respectively, that will be needed in the calculations to follow.

\begin{table}[htbp]
    \centering
    \caption{Eigenfunction integrals}
    \label{tab:meem-integrals}
    \begin{tabular}{|M{0.16\linewidth}|M{0.05\linewidth}|M{0.20\linewidth}|M{0.48\linewidth}|} \hline 
        \multicolumn{2}{|M{0.21\linewidth}|}{Integral}           & $\gls{sym-m},\gls{sym-j}=0$                       & $\gls{sym-m},\gls{sym-j}\geq1$            \\ \hline 
        \multicolumn{2}{|M{0.21\linewidth}|}{\RintOneDefn}       & \RintOneJzero                 & \RintOneJOne          \\ \hline 
        \multicolumn{2}{|M{0.21\linewidth}|}{\RintTwoDefn}       & \RintTwoJzero                 & \RintTwoJOne          \\ \hline 
        \multirow{6}{*}{\ZmnDefn}   & \multirow{2}{*}{$\gls{sym-n}=0$}     & \multirow{2}{*}{\ZnZeroMZero} & \ZnZeroMOne           \\ \cline{2-4} 
                                    & \multirow{3}{*}{$\gls{sym-n}\geq1$}  & \multirow{3}{*}{\ZnOneMZero}  & \ZnOneMOne            \\ \hline
        \multirow{9}{*}{\ZmkDefn}   & \multirow{2}{*}{$\gls{sym-k-generic-index}=0$}     & \multirow{2}{*}{\ZmZeroKZero} & \ZmOneKZero           \\ \cline{2-4}
                                    & \multirow{4}{*}{$\gls{sym-k-generic-index}\geq1$}  & \multirow{4}{*}{\ZkOneMZero}  & \ZkOneMOne            \\ \hline
    \end{tabular}
\end{table}


\subsection{Matching Across Fluid Boundaries}
\label{sec:appendix-meem-matching}

The eigencoefficients must be selected to enforce the radial velocity body boundary condition and the matching of the potentials and radial velocities at the edges of each region, earning this technique the name \gls{acr-meem}.
The radiation problem was first solved this way for a floating cylinder in 1980 \citep{yeung_added_1981}.

First, the infinite sums in $\gls{sym-phi-h}$ must be truncated.
Assuming truncation to $\gls{sym-N}$ terms in \textit{i1}, $\gls{sym-M}$ terms in \textit{i2}, and $\gls{sym-K}$ terms in \textit{e}, the total number of eigencoefficients to solve for is $\gls{sym-N}+2\gls{sym-M}+\gls{sym-K}$.
For a 3-region problem, there are 2 boundaries.
Thus there are four matching equations: (1) potential at $\gls{sym-a-1}$, (2) potential at $\gls{sym-a-2}$, (3) velocity at $\gls{sym-a-1}$, and (4) velocity at $\gls{sym-a-2}$.
As-is, this is not enough equations ($4 < \gls{sym-N}+2\gls{sym-M}+\gls{sym-K}$).
We must leverage eigenfunction orthogonality to get enough equations.
The first equation will turn into $\gls{sym-N}$ equations; the second and third each give $\gls{sym-M}$; the fourth $\gls{sym-K}$.
The transformation uses the following property of orthogonality.
Consider a generic function $\gls{sym-Y}(\gls{sym-x})$ expressed as a series with coefficients $\gls{sym-alpha-generic}$ and basis functions $e(\gls{sym-x})$: $\gls{sym-Y}(\gls{sym-x})=\sum_i \gls{sym-alpha-generic}_i \gls{sym-e-i}(\gls{sym-x})$.
If $\gls{sym-e-j}(\gls{sym-x})$ is orthogonal to $\gls{sym-e-i}(\gls{sym-x})$ from $\gls{sym-x} = \gls{sym-a}$ to $\gls{sym-b}$, then:

\begin{equation}
\begin{aligned}
       & \int\limits_{\gls{sym-a}}^{\gls{sym-b}} \gls{sym-Y} (\gls{sym-x})\gls{sym-e-j} (\gls{sym-x})dx = (\gls{sym-b}-\gls{sym-a}) <\gls{sym-Y},\gls{sym-e-j}>= (\gls{sym-b}-\gls{sym-a})<\sum_i \gls{sym-alpha-generic}_i \gls{sym-e-i}, \gls{sym-e-j}> \\ &= (\gls{sym-b}-\gls{sym-a}) \sum_i \gls{sym-alpha-generic}_i <\gls{sym-e-i},\gls{sym-e-j}> = (\gls{sym-b}-\gls{sym-a}) \sum_i \gls{sym-alpha-generic}_i \gls{sym-delta-ij} = (\gls{sym-b}-\gls{sym-a}) \gls{sym-alpha-generic}_{\gls{sym-j}}
\end{aligned}
\end{equation}

where $<\cdot,\cdot>$ is the inner product and $\gls{sym-delta-ij}$ is Kronecker's delta.
In the current hydrodynamics problem, the basis functions are the vertical eigenfunctions $\gls{sym-Z-n-i1}$, $\gls{sym-Z-m-i2}$, and $\gls{sym-Z-k-e}$.
Orthogonality of each eigenfunction can be verified with the inner product.
In the first region, for example, $<\gls{sym-Z-n-1-i-1},Z_{n_2}^{i_1}>=\delta_{n_1 n_2}$.
Note that eigenfunctions of different domains are not orthogonal, and their inner products will be expressed as coupling integrals in \Cref{tab:meem-integrals}.

For each of the four matching equations, the property of orthogonality applies only after multiplying by the appropriate eigenfunction and integrating over appropriate bounds.
For the potential matching equations, multiply both sides by the eigenfunction of the region with smaller fluid height (so $\gls{sym-Z-n-i1}$ at $\gls{sym-a-1}$ and $\gls{sym-Z-m-i2}$ at $\gls{sym-a-2}$).
Then integrate over that fluid height ($\gls{sym-z}=-\gls{sym-h}$ to $-\gls{sym-d-1}$ at $\gls{sym-a-1}$, and $-\gls{sym-h}$ to $-\gls{sym-d-2}$ at $\gls{sym-a-2}$).
For velocity matching, multiply instead by the eigenfunction corresponding to the larger region, while still integrating over the smaller region.
In velocity matching, an extra step is required to incorporate the boundary condition of zero radial velocity along the radial surface of the body.
Since it is zero-valued, the integral of this velocity may be added to one side of the equation (the one corresponding to the velocity of the larger region) to change the integration bounds only on that side.
This manifests in the bounds of the coupling integrals to be presented in \Cref{tab:meem-integrals}.
Other combinations of eigenfunction multiplication or integration besides those described above are not useful since they result in integrating a quantity on a region where it is undefined, or a form unsuitable for the application of the orthogonality property.

\subsection{Block Matrix Structure}

Once orthogonality is applied, the matching equations create a linear system $\gls{sym-mathbf-A}\gls{sym-vec-x}=\gls{sym-vec-b}$ where $\gls{sym-mathbf-A}$ is a complex sparse $(\gls{sym-N}+2\gls{sym-M}+\gls{sym-K})$x$(\gls{sym-N}+2\gls{sym-M}+\gls{sym-K})$ square matrix corresponding to the homogeneous case, $\gls{sym-vec-x}=[\gls{sym-vec-C-1n-i1}, \gls{sym-vec-C-1m-i2}, \gls{sym-vec-C-2m-i2}, \gls{sym-vec-B-k-e}]$ is the complex eigencoefficient vector, and $\gls{sym-vec-b}$ is the real boundary condition vector corresponding to the particular case.
We elaborate on the block structure of the A-matrix and b-vector, an implementation detail that prior discussion of \gls{acr-meem} overlooks.
The A-matrix and b-vector block structures are shown in \Cref{tab:MEEM-A-matrix,tab:MEEM-b-vector} respectively.
They are written in compact notation using row vectors of basis functions, so $\gls{sym-vec-R-1n-i1}=[R_{10}^{i1}, R_{11}^{i1},\dots, R_{1(N-1)}^{i1}]$ and so on.
Each basis function is evaluated at the radius described to the left of its row in the table. $0_{ij}$ and $1_{ij}$ are the $i$ x $\gls{sym-j}$ matrices of zeros and ones respectively; diag($\cdot$) constructs a diagonal matrix from a vector; and $\odot$ is the Hadamard (element-wise) product.

\begin{table}[htbp]
    \centering
    \caption{\gls{acr-meem} A-matrix}
    \label{tab:MEEM-A-matrix}
    \begin{tabular}{|M{0.12\linewidth}|c||c|c|c|c|}\hline
 & & $\gls{sym-vec-C-1n-i1}$& $\gls{sym-vec-C-1m-i2}$& $\gls{sym-vec-C-2m-i2}$&$\gls{sym-vec-B-k-e}$\\\hline
          &size&  N&  M&  M& K
          \\ \hline \hline
            \shortstack{
                            $\gls{sym-phi-i1}=\gls{sym-phi-i2}$
                        \\
                            at $\gls{sym-r}=\gls{sym-a-1}$
                        }
            &N
            &  $(\gls{sym-h}-\gls{sym-d-1})~\gls{sym-mathrm-diag}(\gls{sym-vec-R-1n-i1})$
            &  $-\gls{sym-boldsymbol-mathcal-Z}_{nm}\odot 1_{\gls{sym-N}1}\gls{sym-vec-R-1m-i2}$
            &  $-\gls{sym-boldsymbol-mathcal-Z}_{nm}\odot 1_{\gls{sym-N}1}\gls{sym-vec-R-2m-i2}$& $0_{NK}$
          \\ \hline
            \shortstack{
                            $\gls{sym-phi-i2}=\gls{sym-phi-e}$
                        \\
                            at $\gls{sym-r}=\gls{sym-a-2}$
                        }
            &M
            &  $0_{MN}$&  $(\gls{sym-h}-\gls{sym-d-2})~\gls{sym-mathrm-diag}(\gls{sym-vec-R-1m-i2})$
            &  $(\gls{sym-h}-\gls{sym-d-2})~\gls{sym-mathrm-diag}(\gls{sym-vec-R-2m-i2})$
            & $-\gls{sym-boldsymbol-mathcal-Z}_{mk}\odot 1_{\gls{sym-M}1}\vec{\Lambda}_{\gls{sym-k-index}}$
          \\ \hline
            \shortstack{
                            $\frac{\partial}{\partial \gls{sym-r}}
                            \gls{sym-phi-i1}=\frac{\partial}{\partial \gls{sym-r}}
                            \gls{sym-phi-i2}$
                        \\
                            at $\gls{sym-r}=\gls{sym-a-1}$
                        }
            &M
            &  $- \gls{sym-boldsymbol-mathcal-Z}_{mn} \odot 1_{\gls{sym-M}1} \frac{\partial}{\partial \gls{sym-r}}\gls{sym-vec-R-1n-i1}$
            &  $(\gls{sym-h}-\gls{sym-d-2})~\gls{sym-mathrm-diag}(\frac{\partial}{\partial \gls{sym-r}}\gls{sym-vec-R-1m-i2})$
            &  $(\gls{sym-h}-\gls{sym-d-2})~\gls{sym-mathrm-diag}(\frac{\partial}{\partial \gls{sym-r}}\gls{sym-vec-R-2m-i2})$
            & $0_{MK}$
          \\ \hline
            \shortstack{
                            $\frac{\partial}{\partial \gls{sym-r}}
                            \gls{sym-phi-i2}=\frac{\partial}{\partial \gls{sym-r}}
                            \gls{sym-phi-e}$
                        \\
                            at $\gls{sym-r}=\gls{sym-a-2}$
                        }
                &K
            &  $0_{KN}$
            &  $-\gls{sym-boldsymbol-mathcal-Z}_{km} \odot 1_{\gls{sym-K}1}\frac{\partial}{\partial \gls{sym-r}}\gls{sym-vec-R-1m-i2}$
            &  $-\gls{sym-boldsymbol-mathcal-Z}_{km} \odot 1_{\gls{sym-K}1}\frac{\partial}{\partial \gls{sym-r}}\gls{sym-vec-R-2m-i2}$
            & $\gls{sym-h}~\gls{sym-mathrm-diag}
            (\frac{\partial}{\partial \gls{sym-r}} \gls{sym-vec-Lambda-k-e})$
        \\ \hline
    \end{tabular}
\end{table}

\begin{table}[htbp]
    \centering
    \caption{\gls{acr-meem} b-vector}
    \label{tab:MEEM-b-vector}
    \begin{tabular}{|c|c|} \hline
         N& $\displaystyle
\int_{-\gls{sym-h}}^{-\gls{sym-d-1}} (\gls{sym-phi-p}^{i2} - \gls{sym-phi-p}^{i1})\gls{sym-vec-Z-n-i1}^{~\top} dz
$ \\ \hline
         M& $\displaystyle-
\int_{-\gls{sym-h}}^{-\gls{sym-d-2}} \gls{sym-phi-p}^{i2} \gls{sym-vec-Z-m-i2}^{~\top} dz
$\\ \hline
         M& $\displaystyle
\int_{-\gls{sym-h}}^{-\gls{sym-d-1}} \frac{\partial}{\partial \gls{sym-r}}\gls{sym-phi-p}^{i1}\gls{sym-vec-Z-m-i2}^{i2~\top} dz - \int_{-\gls{sym-h}}^{-\gls{sym-d-2}}\frac{\partial}{\partial \gls{sym-r}}\gls{sym-phi-p}^{i2}\gls{sym-vec-Z-m-i2}^{~\top} dz
$\\ \hline
         K& $\displaystyle
\int_{-\gls{sym-h}}^{-\gls{sym-d-2}} \frac{\partial}{\partial \gls{sym-r}}\gls{sym-phi-p}^{i2} \vec{\gls{sym-Z}}_{\gls{sym-k-index}}^{e~\top} dz
$\\ \hline
    \end{tabular}
\end{table}

The dense blocks contain coupling integrals $\gls{sym-boldsymbol-mathcal-Z}$ of the vertical eigenfunctions.

Of the sixteen blocks that make up the matrix, six are diagonal, four are zero, and six are dense, resulting in the sparsity pattern shown in \Cref{fig:sparsity}.
An even sparser matrix could be obtained with the alternate eigenfunction scaling for the second region described in \citet{chau_inertia_2012}.
\begin{figure}[htbp]
    \centering
    \includegraphics[\FigureOptions{aor:figs/from-matlab/meem_sparsity.pdf}]{figs/from-matlab/meem_sparsity.pdf}
    \caption{A-matrix sparsity pattern, shown for $\gls{sym-N}=\gls{sym-M}=\gls{sym-K}=4$}
    \label{fig:sparsity}
\end{figure}


\subsection{Calculation of Outputs}

Once the eigencoefficients have been calculated by solving the linear system for $\gls{sym-vec-x}$, the hydrodynamic radiation coefficients $\gls{sym-A-h}$ and $\gls{sym-B-h}$ are found by integrating the potentials and velocities over the body surface:
\begin{alignat}{5}\label{eq:meem-c-derivation}
    \gls{sym-A-h} + \frac{i\gls{sym-B-h}}{\gls{sym-omega}}&=\gls{sym-rho-w} \gls{sym-h}^3  \int_{\gls{sym-theta}=0}^{2\gls{sym-pi}}
    & \int_{\gls{sym-r}=\gls{sym-a-in}}^{\gls{sym-a-out}}
    &\gls{sym-phi} (\gls{sym-r},\gls{sym-z}) \frac{\partial \gls{sym-phi} (\gls{sym-r},\gls{sym-z})}{\partial \gls{sym-z}}
    \gls{sym-r} d\gls{sym-r} d\gls{sym-theta}  \nonumber \\
    %
    %
    &= 2\pi \rho h^3
    & \int_{r=a_{in}}^{a_{out}}&
   \left(\phi_p (r,z) +
         \sum_{j_1} C_{j_1} R_{j_1} (r)Z_{j_1} (z)
   \right) \nonumber \\
   %
  & & &  \left(\frac{\partial \phi_p (r,z)}{\partial z}
               + \sum_{j_2} C_{j_2} R_{j_2} (r)
               \frac{\partial Z_{j_2} (z)} {\partial z}
        \right)r dr \nonumber \\
    %
    %
    &= 2\pi \rho h^3 & \left[
    % first term
    \int_{r=a_{in}}^{a_{out}} \right. &
    \left. \phi_p (r,z)
        \frac{\partial \phi_p (r,z)}{\partial z} rdr
    \right.
    \\
    % second term
    & &  + \sum_{j_1} &C_{j_1} Z_{j_1} (z)
    \int_{r=a_{in}}^{a_{out}}R_{j_1} (r)
    \frac{\partial \phi_p (r,z)}{\partial z}rdr \nonumber
    \\
    % third term
    &&   + \sum_{j_2}&
    \frac{\partial Z_{j_2} (z)}{\partial z} C_{j_2}
    \int_{r=a_{in}}^{a_{out}}
    \phi_p (r,z)  R_{j_2} (r)rdr \nonumber
    \\
    % fourth term
    & & +  \sum_{j_1}&
    \left.
        C_{j_1}Z_{j_1} (z) \sum_{j_2} C_{j_2}
        \frac{\partial Z_{j_2} (z)}{\partial z}
        \int_{r=a_{in}}^{a_{out}}
        R_{j_1} (r)R_{j_2} (r)rdr
    \right] \nonumber
\end{alignat}
where the generic indices $\gls{sym-j}_1$ and $\gls{sym-j}_2$ are used to represent $\gls{sym-n}$ or $\gls{sym-m}$ depending on the region.
For regions with multiple radial eigenfunctions or surfaces covering multiple fluid regions, summation over each eigenfunction and region respectively is implied.
After substituting the eigenfunctions from \Cref{tab:MEEM-eigenfunctions} into \Cref{eq:meem-c-derivation}, all $\gls{sym-z}$-dependent quantities are evaluated at the body draft on the integration surface, $\gls{sym-z}=-d$.
This means $\frac{\partial \gls{sym-phi-p}}{\partial \gls{sym-z}}=1$, and $\gls{sym-Z-j}(\gls{sym-z})$ simplifies to
\begin{equation}
    \gls{sym-Z-j} (\gls{sym-z}=-d) =
    \begin{cases}
        1,& \gls{sym-j}=0
        \\
        \sqrt2(-1)^{\gls{sym-j}},& \gls{sym-j} \geq 1
    \end{cases}
\end{equation}
The third and fourth integrals of \Cref{eq:meem-c-derivation} vanish, and the second integral is the radial eigenfunction integral $\gls{sym-boldsymbol-mathcal-R}$ expressed previously in \Cref{tab:meem-integrals}.

The first and second term are independent of and scale linearly with the eigencoefficients $\gls{sym-vec-x}$, respectively.
The radiation coefficients are thus computed as
\begin{equation}\label{eq:hydro}
    \gls{sym-A-h} + \frac{i\gls{sym-B-h}}{\gls{sym-omega}}=2\gls{sym-pi} \gls{sym-rho-w} \gls{sym-h}^3 (\gls{sym-c-0} + \gls{sym-vec-c}\cdot\gls{sym-vec-x}) =2\gls{sym-pi} \gls{sym-rho-w} \gls{sym-h}^3 (\gls{sym-c-0} + \gls{sym-vec-c} \cdot \gls{sym-mathbf-A}^{-1} \gls{sym-vec-b})
\end{equation}
with output constant $\gls{sym-c-0}$ and output row vector $\gls{sym-vec-c}$ defined in \Cref{eq:c0,tab:MEEM-c-vector} respectively:
\begin{equation}\label{eq:c0}
    \gls{sym-c-0} =  \frac{\left({\gls{sym-a-out}}-{\gls{sym-a-in}}\right)\,\left(-{\gls{sym-a-in}}-{\gls{sym-a-out}}+4 (\gls{sym-h}-d)^2\right)}{16\,\left(\gls{sym-h} - d\right)}
\end{equation}
\begin{table}[h]
    \centering
    \caption{Output vector $\gls{sym-vec-c}$}
    \label{tab:MEEM-c-vector}
\begin{tabular}{|l|c|c|c|c|} \hline
          Body&N&  M&  M& K\\ \hline
          Float&$0_{1\gls{sym-N}}$&  $\gls{sym-vec-Z-m-i2}(\gls{sym-z}=-d) \odot\gls{sym-boldsymbol-mathcal-R}_{1\gls{sym-m}}$&  $\gls{sym-vec-Z-m-i2}(\gls{sym-z}=-d) \odot\gls{sym-boldsymbol-mathcal-R}_{2\gls{sym-m}}$& $0_{1\gls{sym-K}}$\\ \hline
          Spar& $\gls{sym-vec-Z-n-i1}(\gls{sym-z}=-d) \odot\gls{sym-boldsymbol-mathcal-R}_{1\gls{sym-n}}$& $0_{1\gls{sym-M}}$& $0_{1\gls{sym-M}}$&$0_{1\gls{sym-K}}$\\\hline
    \end{tabular}
    \end{table}
The appropriate dimensions are
\begin{equation}
    \gls{sym-a-out} = \begin{cases}
        \gls{sym-a-2}, & \text{float} \\
        \gls{sym-a-1}, & \text{spar}
    \end{cases}
    \qquad
    \gls{sym-a-in} = \begin{cases}
        \gls{sym-a-1}, & \text{float} \\
        0, & \text{spar}
    \end{cases}
    \qquad
   d = \begin{cases}
        -\gls{sym-d-2}, & \text{float} \\
        -\gls{sym-d-1}, & \text{spar}
    \end{cases}
\end{equation}
to obtain the hydrodynamic coefficients of the float or the spar respectively.

Using the A-matrix and b- and c-vectors, the \gls{acr-meem} solution directly calculates radiation coefficients $\gls{sym-B-h}$ and $\gls{sym-A-h}$.
The hydrostatic stiffness $\gls{sym-K-h}$ and the excitation coefficient $\gls{sym-gamma}$ are calculated as follows:
\begin{equation}\label{eq:gamma-K}
    |\gls{sym-gamma}|  = \sqrt{\frac{ 4 \gls{sym-rho-w} \gls{sym-g} \gls{sym-V-g}  \gls{sym-B-h}} {\gls{sym-m-0}}}, \quad % excitation
   \angle \gls{sym-gamma} = -\frac{\gls{sym-pi}}{2} + \angle\frac{ {\gls{sym-B-k-e}}_{=0}}{\gls{sym-mathrm-H-0-1} (\gls{sym-m-0} \gls{sym-a-2})},\quad
    \gls{sym-K-h}       = \gls{sym-rho-w} \gls{sym-g} \underbrace{\frac{\gls{sym-pi}}{4} (\gls{sym-D-f}^2 - \gls{sym-D-s}^2)}_{\gls{sym-A-w-f}},\quad  % hydrostatic stiffness
    \gls{sym-K-s} = \gls{sym-rho-w} \gls{sym-g} \underbrace{\frac{\gls{sym-pi}}{4} \gls{sym-D-s}^2}_{\gls{sym-A-w-s}}
\end{equation}
where $\gls{sym-g}$ is the acceleration due to gravity, $\gls{sym-rho-w}$ is the density of water, $\gls{sym-m-0}$ is the wavenumber, $\gls{sym-V-g}$ is the finite depth group velocity, $\gls{sym-mathrm-H-0-1}$ is the zeroth-order Hankel function of the first kind, and $\gls{sym-A-w}$ is the waterplane area \citep{newman}.
This method of calculating excitation from damping is the well-known Haskind relation.
Note that while the excitation magnitude $|\gls{sym-gamma}|$ depends on the radiation damping $\gls{sym-B-h}$ which in turn depends on all the inner region eigencoefficients ($\gls{sym-vec-C-1m-i2}$ and $\gls{sym-vec-C-2m-i2}$ for float excitation and $\gls{sym-vec-C-1n-i1}$ for the spar excitation), the excitation phase $\angle\gls{sym-gamma}$ depends only on the first exterior eigencoefficient, ${\gls{sym-B-k-e}}_{=0}$.

\subsection{Verification and Validation}

Hydro coefficient results are verified by comparing to a benchmark shallow-water concentric-cylinder \gls{acr-meem} solution in \citet{chau_inertia_2012}.
Excellent agreement is observed, shown in \Cref{fig:meem-yeung-validation}.
The results are also experimentally validated in the studies \citet{chau_inertia_2012,son_performance_2016}.
Previously in \Cref{sec:dynamic-validation}, the N=M=K=11 case was compared to \gls{acr-wamit} results for \gls{acr-rm3} in deep water.

\begin{figure}[htbp]
    \centering
    \includegraphics[\FigureOptions{aor:figs/from-matlab/meem_validation.pdf}]{figs/from-matlab/meem_validation.pdf}
    \caption{Nondimensional added mass and damping coefficient validation against \citep{chau_inertia_2012}}
    \label{fig:meem-yeung-validation}
\end{figure}

\subsection{Convergence}

As $\gls{sym-N},\gls{sym-M},\gls{sym-K}\rightarrow\infty$, matching quality improves, and hydro coefficients converge toward their true values.
Previous \gls{acr-meem} papers use $\gls{sym-N}=\gls{sym-M}=\gls{sym-K}=50$ to obtain 4-digit matching accuracy without elaborating on convergence properties \citep{chau_inertia_2012}.
We observe that potential matching converges faster than velocity matching.
\Cref{fig:meem-matching} shows the matching behavior for $\gls{sym-N}=\gls{sym-M}=\gls{sym-K}=11$, where potential matches well but velocity still has noticeable mismatch.
\begin{figure}[htbp]
    \centering
    \includegraphics[\FigureOptions{aor:figs/from-matlab/meem_matching.pdf}]{figs/from-matlab/meem_matching.pdf}
    \caption{Matching for $\gls{sym-N}=\gls{sym-M}=\gls{sym-K}=11$ for benchmark geometry}
    \label{fig:meem-matching}
\end{figure}

Hydro coefficient convergence depends on the geometry: the benchmark shallow-water geometry of \citep{chau_inertia_2012} converges to within 0.25\% with only $\gls{sym-N}=\gls{sym-M}=\gls{sym-K}=4$, but \gls{acr-rm3} requires $\gls{sym-N}=\gls{sym-M}=\gls{sym-K}>10$, shown in \Cref{fig:meem-convergence}.
There, damping converges well at low frequencies but requires more harmonics at higher frequencies, while added mass has similar convergence across frequencies.
\begin{figure}[htbp]
    \centering
    \includegraphics[\FigureOptions{aor:figs/from-matlab/meem_convergence_vs_omega.pdf}]{figs/from-matlab/meem_convergence_vs_omega.pdf}
    \caption{Convergence for $\gls{sym-N}=\gls{sym-M}=\gls{sym-K}=(5,10,20,30)$ for \gls{acr-rm3}}
    \label{fig:meem-convergence}
\end{figure}


\subsection{Numerical Notes}\label{sec:appendix-meem-numerics}
Note that radial eigenfunctions $\gls{sym-R-1n-i1}(\gls{sym-r})$ and $\gls{sym-R-1m-i2}(\gls{sym-r})$ contain the modified Bessel function of the first kind $\gls{sym-mathrm-I-0}(\gls{sym-chi})$, which diverges for $\gls{sym-chi}\rightarrow \infty$, resulting in numeric overflow.
To prevent overflow, each fluid region must exceed a minimum height $\gls{sym-Delta-z-min}$ that is proportional to the region diameter $\gls{sym-D}$ for both the spar ($\gls{sym-Delta-z-s}, \gls{sym-D-s}$) and the float ($\gls{sym-Delta-z-f}, \gls{sym-D-f}$), with variables shown in \Cref{fig:meem-geom}.
The constant of proportionality between $\gls{sym-Delta-z-min}$ and $\gls{sym-D}$ is a function of the number of harmonics $\gls{sym-N-harmonics}$ used in that region ($\gls{sym-N}$ for spar, $\gls{sym-M}$ for float), as well as the maximum argument to the \texttt{besseli} function in MATLAB before overflow, $\gls{sym-chi-max}$:
\begin{equation}\label{eq:delta-z-min}
    \frac{\gls{sym-Delta-z-min}}{\gls{sym-D}} = \frac{\gls{sym-pi}\gls{sym-N-harmonics}}{2\gls{sym-chi-max}} \approx \frac{\gls{sym-N-harmonics}}{446}.
\end{equation}
By trial, $\gls{sym-chi-max}$ is found to be $\approx 700.5$, close to the theoretical value of \texttt{log (realmax)} $\approx 709.8$ for exponential scaling.
Since $\gls{sym-N-harmonics}=10$ gives adequate convergence for most geometries, this condition is trivially satisfied for nearly all floating bodies of practical relevance, but must still be added as a constraint to prevent the optimizer from exploiting the numerical divergence.
In future work, if high-accuracy solutions are desired for large bodies close to the sea floor, exponentially scaled Bessel functions could be used instead, with symbolic cancellation of the exponential scaler from the eigenfunction numerator and denominator.

Likewise, the vertical eigenfunction $\gls{sym-Z-k-e}$ for $\gls{sym-k-index}=0$ contains the $\cosh$ and $\sinh$ functions, which diverge for large values of $\gls{sym-m-0}\gls{sym-h}$ (high frequencies or deep water).
Since the largest relevant value of $\gls{sym-m-0}\gls{sym-h}$ depends on the site rather than on the \gls{acr-wec} design, it is not possible to add geometric constraints to prevent overflow as it was above.
Therefore, the limit is derived analytically. %\hl{This formula hasn't been numerically checked yet}
\begin{equation}
    \lim_{\gls{sym-m-0}\gls{sym-h}\rightarrow\infty} \gls{sym-Z-0-e} (\gls{sym-z})= \frac{\cosh^2 (\gls{sym-m-0}\gls{sym-h})}{\sqrt{2\gls{sym-m-0}\gls{sym-h}}}\exp\left(1+\frac{\gls{sym-z}}{\gls{sym-h}}\right)
\end{equation}
Plugging this into the first element of the bottom block of the b-vector (position $\gls{sym-N}+2\gls{sym-M}+1$) results in
\begin{equation}
    \lim_{\gls{sym-m-0}\gls{sym-h}\rightarrow\infty}\gls{sym-b-N-2M-1}=\frac{-\gls{sym-a-2}}{\gls{sym-h}-\gls{sym-d-2}} \sqrt{\frac{\gls{sym-h}}{2\gls{sym-m-0}}}\exp(-\gls{sym-d-2}\gls{sym-m-0})
\end{equation}
\begin{figure}[htbp]
    \centering
    \includegraphics[\FigureOptions{aor:figs/from-matlab/asymptotic_b_vector.pdf}]{figs/from-matlab/asymptotic_b_vector.pdf}
    \caption{Asymptotic b-vector for large $\gls{sym-m-0} \gls{sym-h}$}
    \label{fig:meem-b-limit}
\end{figure}

\Cref{fig:meem-b-limit} shows the asymptotic approximation as well as the standard expression for the relevant b-vector entry as a function of nondimensional ratio $\gls{sym-m-0}\gls{sym-h}/\textrm{acosh(realmax})$.
With the standard expression, values for inputs above $\frac{1}{2}$ become \texttt{NaN}, which poses problems in deep water.
With the asymptotic approximation, we are able to avoid \texttt{NaN}s for the zone $\frac{1}{2}<\gls{sym-m-0}\gls{sym-h}/\textrm{acosh(realmax})<\frac{\gls{sym-h}}{\gls{sym-d-2}}$.
Above this threshold, \texttt{NaN}s occur even with the approximation, and the value in the b-vector is set to zero.
This introduces negligible error because in this range, the magnitude of the true value is less than $10^{-300}$ for the nominal \gls{acr-rm3} geometry.
The asymptotic approximation is accurate for values of $\gls{sym-m-0}\gls{sym-h}/\textrm{acosh(realmax)}$ as low as $10^{-2}$,
though it is only applied for values above $\frac{1}{2}$.


The corresponding limit for the vertical coupling integral is:
\begin{equation}
\begin{aligned}
    \lim_{\gls{sym-m-0}\gls{sym-h}\rightarrow\infty}
    \gls{sym-boldsymbol-mathcal-Z}_{\gls{sym-m},\gls{sym-k-generic-index}=0}
    &
    %= \int_{-h}^{-d_2} \vec{Z}_m^{i2 ~T}\vec{Z}_0^{\textit{e}} dz
    = \gls{sym-h}~\frac{
                            \cosh^2 (\gls{sym-m-0}\gls{sym-h})
                        }
                        {
                            \sqrt{2\gls{sym-m-0}\gls{sym-h}}
                        }
            \cdot \frac{
                            -1+(-1)^{\gls{sym-m}}
                            \exp(
                                1-\frac{\gls{sym-d-2}} {\gls{sym-h}}
                                )
                        }
                        {
                            \gls{sym-f-m}
                        }
    \\
    \text{where}~~\gls{sym-f-m}&=
    \begin{cases}
        1,                                   & \gls{sym-m}=0 \\
        \gls{sym-h}^2\gls{sym-lambda-m}^2+1, & \gls{sym-m} \geq 1
    \end{cases}
\end{aligned}
\end{equation}
%\hl{check this formula}

A final numerical subtlety worth discussing is finite precision effects in calculating $\gls{sym-m-k}$.
Bounds of $180^\textrm{o}\cdot[\gls{sym-k-index}-\frac{1}{2}, \gls{sym-k-index}]$
are placed on $\gls{sym-m-k}\gls{sym-h}$ in a root-finding algorithm to ensure the $\gls{sym-k-index}$th root is identified.
Without these bounds, sometimes the \texttt{fsolve} solver would wrongly converge to values where the residual approaches
$+\infty$ from one side and $-\infty$ from the other
(asymptotes of $\tan()$) rather than true zeroes.
Degrees are used instead of radians so asymptotes occur at rational values.

\subsection{Runtime and Computational Cost Scaling}

The runtime of the \gls{acr-meem} method is the time required to find the eigencoefficients, then obtain the hydrodynamic coefficients from eigencoefficients.
First, a nonlinear root-finding algorithm runs $\gls{sym-K}-1$ times to generate the $\gls{sym-m-k}$ inputs used in the A-matrix and b-vector.
Then $3\gls{sym-N}+8\gls{sym-M}+2\gls{sym-K}-11$ Bessel functions must be evaluated for the radial terms of the A-matrix.
 (Note that this would reduce to $2\gls{sym-N}+8\gls{sym-M}+2\gls{sym-K}-10$ evaluations if the definition of the $\gls{sym-R-1n-i1}(\gls{sym-r})$ eigenfunction were modified to use a scale factor based on $\gls{sym-a-1}$ rather than $\gls{sym-a-2}$, although this enhancement is not pursued here to maintain consistency with previous work).
% if I change scaling of i2 from a2 to (a1+a2)/2, it increases by 2* (M-1). So total would be 2N+10M+2K-12. see notebook p51.
The cost of evaluating vertical coupling integrals (\Cref{tab:meem-integrals}) is negligible since they are trigonometric.
Linear solves scale almost cubically with matrix size, so this step scales with $(\gls{sym-N}+2\gls{sym-M}+\gls{sym-K})^3$.
The radial integrals for the c-vector do not require evaluating Bessel functions with any arguments that were not already evaluated for the A-matrix.
For $\gls{sym-N}=\gls{sym-M}=\gls{sym-K}=10$, the simulation of a single frequency averages \resultsAOR[MEEMRuntime] on the hardware described in \Cref{sec:sim-runtime}.
\Cref{fig:runtime-hydro} shows the time breakdown.
Most of the time is spent evaluating Bessel functions, so future code optimization should focus on speeding up Bessel evaluations, such as with lookup tables.
\citet{chau_inertia_2012} proposes using the sparsity pattern to reduce matrix size from N+2M+K to 2M, but this seems low impact since the linear solve only takes a few percent of compute time and the reduced matrix would require an equal number of Bessel evaluations.
On the other hand, matrix size in a boundary element method solver is much larger (meshes may have 1000s of panels) and the linear solve can drive computation cost.
On the same machine, Capytaine boundary element method for the same geometry takes an average of 323 ms for a 710 panel mesh (1\% convergence). %\hl{the 323 ms capy number came from my laptop not from the lab computer}
Thus, \gls{acr-meem} achieves a 10x time reduction over Capytaine.

%%%%%%%%%%%%%%%%%%%%%%%%%%%%%%%%%%%%%%%%%
%%%%%%%%%%%%%%%%%%%%%%%%%%%%%%%%%%%%%%%%%
\section{Additional Hydrodynamics Details}
\label{sec:appendix-additional-hydro}
\subsection{Spar Coefficients}
As \Cref{sec:hydro-other} introduces, because of the importance of the damping plate, the spar hydrodynamic coefficients do utilize other methods besides \gls{acr-meem}.
We approximate the spar added mass coefficient $\gls{sym-A-s}$ as a frequency-independent quantity using a formula that \citet{philip_damping_2012} proposes based on the displaced water volume in the spherical projection of the damping plate:
\begin{equation}\label{eq:As}
 \gls{sym-A-s} =\frac{1}{3}\gls{sym-rho-w} \gls{sym-D-d}^3\left(1 - \frac{3}{4}\gls{sym-sigma-ratio}^2\sqrt{1-\gls{sym-sigma-ratio}^2} - \frac{1}{4}(1-\sqrt{1-\gls{sym-sigma-ratio}^2})^2(2+\sqrt{1-\gls{sym-sigma-ratio}^2})\right),
\end{equation}
where $\gls{sym-sigma-ratio} = \gls{sym-D-s}/\gls{sym-D-d}$.
This expression smoothly transitions from a pure plate limit ($\frac{1}{3}\gls{sym-rho-w} \gls{sym-D-d}^3$ at $\gls{sym-sigma-ratio}=0$) to a pure column limit ($\frac{1}{6}\gls{sym-rho-w} \gls{sym-D-d}^3$ at $\gls{sym-sigma-ratio}=1$).
\ifdefined\DISSERTATION
    The trend is shown in \figureautorefname~\ref{fig:spar_added_mass}.
    \begin{figure}
        \centering
        \includegraphics[width=.6\linewidth]{figs/from-matlab/spar_added_mass.pdf}
        \caption[Nondimensional spar added mass vs. diameter ratio]{Nondimensional spar added mass as a function of column to damping plate diameter ratio}
        \label{fig:spar_added_mass}
    \end{figure}
\fi

The other spar coefficients are calculated by fitting and interpolating data from \gls{acr-wamit} for the nominal design, then adjusting for known dimensional scalings based on differences from the nominal design.
To derive the dimensional scalings, we use the numerical results that \citet{olaya_hydrodynamic_2015} present for a variant of \gls{acr-meem} that incorporates the damping plate.
By plotting the data for a variety of geometries on a single plot and scaling the axes by dimensionless geometric ratios until the points begin to overlap as a single trend,
we empirically arrive at the following form for the excitation coefficient magnitude:
\begin{equation}\label{eq:scale-wamit}
|\gls{sym-gamma-s}| = \gls{sym-rho-w} \gls{sym-g} \gls{sym-pi} \gls{sym-D-s}^2/4 |\gls{sym-mathrm-H-0-1}(\gls{sym-k-wavenumber}\gls{sym-R-x})| \exp(-\gls{sym-k-wavenumber} \gls{sym-e-1}) \gls{sym-k-wavenumber}\gls{sym-h} \gls{sym-alpha} \max(1,\gls{sym-alpha})
\sqrt{\frac{\gls{sym-Upsilon}\left(\frac{\gls{sym-k-wavenumber}\gls{sym-h} \gls{sym-R-x} \gls{sym-alpha}^3}{\gls{sym-beta-olaya} \gls{sym-R-p}}\right)}
           {\gls{sym-beta-olaya}\left(1+\frac{\gls{sym-beta-olaya}}{\gls{sym-alpha}}\right)}}
\end{equation}
where $\gls{sym-Upsilon}(\cdot)$ indicates a nonlinear relationship to be interpolated from data, and $\gls{sym-alpha}$, $\gls{sym-beta-olaya}$,  $\gls{sym-e-1}$, $\gls{sym-R-p}$, and $\gls{sym-R-x}$ are dimensions and nondimensional ratios defined in \citet{olaya_hydrodynamic_2015}.
\Cref{fig:olaya-hydro-data} shows this nonlinear relationship, which has collapsed remarkably well for the variety of $\gls{sym-alpha}$ and $\gls{sym-beta-olaya}$ values.
Case numbers in the legend refer to the numerical parameters used by \citet[Section~V.]{olaya_hydrodynamic_2015}
\begin{figure}[htbp]
\centering
    \includegraphics[\FigureOptions{aor:figs/from-matlab/case4_v1_auto_case2_overlay_loglog_khRxRpa2b_fv1.pdf}]{figs/from-matlab/case4_v1_auto_case2_overlay_loglog_khRxRpa2b_fv1.pdf}
\caption[Nondimensionalized hydrodynamic data used to interpolate spar coefficients]{Nondimensionalized hydrodynamic data used to interpolate spar coefficients, showing the rebound effect.}
\label{fig:olaya-hydro-data}
\end{figure}
The dip at mid frequencies shows the excitation and damping coefficients experiencing ``rebound'' effects where rather than slowly asymptoting to zero like those of a single cylinder, they go quickly to zero, increase to a peak, then slowly asymptote down to zero again.

\subsection{Positive-definiteness}
Computing the float coefficients in a different way than the spar and coupling coefficients introduces the possibility of violating physical requirements of the solution.
In particular, the damping matrix must be positive definite because it represents a dissipative force where outgoing waves remove energy from the system, and the added mass matrix must be positive definite because it represents the kinetic energy of the fluid which must be non-negative.
\begin{equation}\label{eq:open-loop-stabilty}
    \det(\gls{sym-mathbf-A-h-op})=\gls{sym-A-f} \gls{sym-A-s} - \gls{sym-A-c}^2>0 \,, \quad
    \det(\gls{sym-mathbf-B-h-op})=\gls{sym-B-f} \gls{sym-B-s} - \gls{sym-B-c}^2>0
\end{equation}
This positive-definiteness may be violated for certain combinations of coefficients and can cause the system dynamics to become unstable.
Even when other forces such as body inertia, drag, friction, and control stabilize the system, violation of the conditions of \Cref{eq:open-loop-stabilty} can artificially inflate power production.
These conditions must be checked for each sea state and the coefficients adjusted if needed.
MDOcean chooses to decrease the magnitude of the coupling coefficients $|\gls{sym-A-c}|$ and $|\gls{sym-B-c}|$ if a violation occurs because the coupling coefficients are hypothesized to correlate less strongly than the float and spar coefficients with the float and spar dimensions respectively.
This intends to lower the likelihood of creating saturation behavior (where a dimension can be perturbed without influencing power production) that could interfere with a gradient-based optimization routine.

%%%%%%%%%%%%%%%%%%%%%%%%%%%%%%%%%%%%%%%%%
%%%%%%%%%%%%%%%%%%%%%%%%%%%%%%%%%%%%%%%%%
\section{Dynamics and Control Module Details}
\label{sec:appendix-dynamics}
This section provides supplementary derivations and implementation details for the dynamics and control module.
See
\Cref{sec:appendix-force-sat} for describing functions,
\Cref{sec:appendix-multiport} for the multiport model,
\Cref{sec:appendix-optimal-control} for optimal control,
\Cref{sec:appendix-drag} for the drag model,
\Cref{sec:appendix-dynamic-limit} for dynamic limits,
\Cref{sec:appendix-slam} for slamming and submersion criteria,
\Cref{sec:appendix-narrowband} for the narrow-band assumption,
\Cref{sec:appendix-power-decomposition} for power matrix decomposition, and
\Cref{sec:appendix-dynamic-validation} for dynamic validation.

\subsection{Describing Function Details}
\label{sec:appendix-force-sat}

\Cref{sec:dynamics} introduces the use of describing functions to quantify the fundamental amplitude of non-sinusoidal force waveforms, which can achieve higher powers than the equivalent sinusoidal waveform for the same force limit.
\citet{hall_open-source_2022} documents the \gls{acr-raft} open-source Python package for offshore wind turbines, which uses the drag describing function with iteration but does not include describing functions for saturation.
\ifdefined\DISSERTATION
    The indirect implementation of describing functions for the force saturation nonlinearity is possible because the describing function \gls{acr-qlps} method is used for both controller synthesis and evaluation, so the time-domain nonlinear optimal controller that is output from the right of \Cref{fig:mod-freq-domain-synthesis} is fed into the top left of \Cref{fig:mod-freq-domain-evaluation},
    resulting in a chaining of the inverse control describing function with the forward control describing function, which cancel out.
    Unless intentionally made different to study robustness, the linearized plant model is identical for the synthesis and evaluation steps.
    In this case, two workflows essentially become one, with no further computation required after the optimal control step.
    Reconstruction of the nonlinear controller is not necessary for the evaluation of the power corresponding to that controller.
    The power can be calculated directly from the frequency-domain response and the describing function of the nonlinear controller.
    The time-domain representation of the nonlinear controller is only necessary for hardware implementation, or if the synthesized controller is to be evaluated with a different method such as time-stepping.
\else
    Because the same describing function is applied in both synthesis and evaluation, the inverse and forward describing functions cancel;
    the time-domain nonlinear controller need not be reconstructed for power evaluation, only for hardware implementation.
\fi
The simulation permits the nonsinusoidal waveform to have a maximum fundamental a factor of $\frac{4}{\gls{sym-pi}}$ above the force limit.
This coincides with the fundamental to peak ratio of a square wave, which is the limiting case both for the bang-singular-bang controller (a piecewise-discontinuous sine-square wave combination that \citet{hendrikx_optimal_2017} finds as the analytically optimal solution), and for the saturated sine controller (a piecewise-continuous sine-square wave combination that \citet{coe_initial_2020} finds as numerically optimal).
In fact, the optimization need not specify the exact nonlinear waveform, since under the describing function filtering hypothesis, only the fundamental matters for simulation purposes.
For hardware implementation purposes, the nonlinear controller can be obtained for the saturated-sine controller by the process outlined in the study by \citet{mccabe_force-limited_2024}.

The appropriate complex scale factor from the optimal unsaturated controller (damping or reactive respectively) is found as described in \Cref{sec:appendix-qp-solution} and will be real for the damping control case and complex in the reactive control case.

%Mathematically, a saturation factor $0< \gls{sym-f}_{sat}\leq 1$ is computed as the ratio between the fundamental amplitude of the saturated torque and the original unsaturated torque signal,
%\begin{equation}
%    f_{sat} = \min\left(\alpha\frac{\tau_{max}}{\tau}, 1\right),
%\end{equation}
%where $\gls{sym-tau}_{\text{max}}$ is the maximum permissible generator torque and $\gls{sym-alpha}$ is the ratio between the fundamental amplitude of the saturated waveform and the maximum torque $\gls{sym-tau}_{\text{max}}$, using the describing function derived in \citep{mccabe_force-limited_2024},
%\begin{equation}
%    \alpha = \frac{2}{\pi} \left(\frac{F_p}{\tau_{max}}\sin^{-1}\frac{\tau_{max}}{F_p} +  \sqrt{1 - \left(\frac{\tau_{max}}{F_p}\right)^2} \right).
%\end{equation}

%When $\frac{\gls{sym-F-p}}{\gls{sym-tau}_{\text{max}}}$ = 1 (no saturation), then $\gls{sym-alpha}= 1$;
%when $\frac{\gls{sym-F-p}}{\gls{sym-tau}_{\text{max}}} \rightarrow \infty$ (full saturation), then $\gls{sym-alpha}= \frac{4}{\gls{sym-pi}}$.
%This makes sense because $\frac{4}{\gls{sym-pi}}$ is the harmonic amplitude of a unit square wave,
%and in full saturation the saturated signal approaches a square wave.
%$\gls{sym-alpha}$ is plotted in \Cref{fig:alpha-fit}.

%The factor $\gls{sym-f-sat}$ describes the ratio of the saturated and unsaturated torque signal.
%Previous work \citet{mccabe_multidisciplinary_2022} implies that this is also the ratio of the saturated and unsaturated control gains, $\gls{sym-m-sat}$, but this is not the case because the \gls{acr-wec} response also changes with saturation.

% %Substituting $\gls{sym-f-sat}$ back into the equation of motion reveals that the control gain saturation multiplier $\gls{sym-m-sat}$ is the solution to a quadratic equation $\gls{sym-a-q} \gls{sym-m-sat}^2 + \gls{sym-b-q} \gls{sym-m-sat} + \gls{sym-c-q} = 0$. If the uncontrolled dynamics have a natural frequency $\gls{sym-omega}_{\gls{sym-n},\gls{sym-u}}$ and damping ratio $\gls{sym-zeta-u}$, and the controlled unsaturated dynamics have natural frequency $\gls{sym-omega-n}$ and damping ratio $\gls{sym-zeta}$, then the quadratic coefficients can be expressed as:

% \begin{equation}\label{eq:abc-matrix}
%     \begin{bmatrix}
%         a_q \\ b_q \\ c_q
%     \end{bmatrix}
%     =
%     \begingroup % keep the row spacing change local
%     \setlength\arraycolsep{2pt}
%     \begin{bmatrix} -4 & \frac{4}{f_{sat}^2} & \frac{1}{f_{sat}^2} & -1 & 0 & 0\\
%     8 & 0 & 0 & 2 & 8 & 2\\
%     -4 & -4 & -1 & -1 & -8 & -2 \end{bmatrix}
%     \endgroup
%     \begin{bmatrix}
%  (\theta_1)^2  \\
%   (\theta_2)^2 \\
%    (\theta_3)^2 \\
%    (\theta_4)^2  \\
%    \theta_1 \theta_2 \\
%     \theta_3 \theta_4
%     \end{bmatrix}
% \end{equation}
% where the basis vector $\vec{\gls{sym-theta}}$ is a dimensionless function of the natural frequencies and damping ratios:
% \begin{equation}\label{eq:theta-basis}
%         \begin{bmatrix}
%   \theta_1  \\
%   \theta_2 \\
%    \theta_3 \\
%    \theta_4  \\
%     \end{bmatrix} =     \begin{bmatrix}
% ~ \overline{\omega} (\zeta_u - \overline{\omega}_n\zeta)  \\
%   \overline{\omega} ~\overline{\omega}_n \zeta \\
%   \overline{\omega}^2 - \overline{\omega}_n^2\\
%    \overline{\omega}_n^2 - 1  \\
%     \end{bmatrix}
% \end{equation}
% and the overlined frequencies indicate nondimensionalization by the uncontrolled natural frequency: $\overline{\gls{sym-omega}} = \frac{\gls{sym-omega}}{\gls{sym-omega}_{\gls{sym-n},\gls{sym-u}}}, \overline{\gls{sym-omega}}_{\gls{sym-n}} = \frac{\gls{sym-omega-n}}{\gls{sym-omega}_{\gls{sym-n},\gls{sym-u}}}$. The quadratic equation is then solved for $\gls{sym-m-sat}$, and of the two roots, the one that is real-valued and between 0 and 1 is used.

% The no saturation case ($\gls{sym-f-sat} = 1$) for \eqref{abc-matrix} yields $\gls{sym-m-sat}=1$. This confirms the expected behavior that when no force reduction is required, the powertrain coefficients can remain unchanged.

% Equations \eqref{eq:abc-matrix} and \eqref{eq:theta-basis} are an equivalent reformulation of \citet{mccabe_force-limited_2024}, which uses expressions for the complex reflection coefficient $\gls{sym-Gamma}$ common in electrical engineering, rather than the damping ratio and natural frequency more common in mechanical engineering. The equivalence is:
% \begin{equation}
%     \Gamma = \frac{z-1}{z+1}, \quad z = \frac{2\theta_1 + i~ \theta_4}{-2 (\theta_1+\theta_2) + i~(\theta_3+\theta_4)}
% \end{equation}
% % I am inputting zeta and omega_n corresponding to mechanical dynamcis when calculating m_sat, but I really should be inputting them corresponding to electrical dynamics. Maybe describe that?

% The gain saturation multiplier $\gls{sym-m-sat}$ is then used to modify the powertrain coefficients,
% \begin{align}
%     B_{p,sat} &= m_{sat} B_p, & K_{p,sat} &= m_{sat} K_p.
% \end{align}
% and these new control gains are then plugged into equation \eqref{eq:eom-freq-domain} to find the saturated displacement amplitude $\gls{sym-X}_{sat}$.


\subsection{Multiport Model Details}
\label{sec:appendix-multiport}
To implement the multiport model of \Cref{sec:dynamics}, we must specify its
conventions, body dynamics, \gls{acr-pto} kinematics, \gls{acr-pto} dynamics, power relations, and ultimate solution procedure.
This section provides a general structure in effort-flow form applicable across multiple energy domains for any number of \glspl{acr-dof}, as well as specific equations for a common \gls{acr-pto} structure consisting of a mechanical drivetrain and a permanent magnet generator.

%%%%%%%%%%%%%%%
\subsubsection{Multiport Circuit Conventions}
\label{sec:appendix-multiport-conventions}

The intrinsic, powertrain, and Th\'evenin-equivalent representations used throughout \Cref{sec:dynamics} are all instances of multiport circuit models,
with effort variables $\gls{sym-vec-e}$ and flow variables $\gls{sym-vec-q}$ at each port.
Following \citet{reveyrand_multiport_2018}, three matrix forms describe the relationships between port variables:
impedance ($\gls{sym-mathbf-Z}$), cascade ($\gls{sym-mathbf-a}$, also called \gls{acr-abcd}), and hybrid ($\gls{sym-mathbf-h}$):
\begin{equation}\label{eq:multiport-matrix-definitions}
    \vec{e} = \gls{sym-mathbf-Z}\,\gls{sym-vec-q},
\qquad
    \begin{bmatrix}
        \vec{e}_2 \\ \vec{q}_2
    \end{bmatrix}
    = [\gls{sym-mathbf-a}]_{2\leftarrow 1}
    \begin{bmatrix}
        \vec{e}_1 \\ \vec{q}_1
    \end{bmatrix},
\qquad
    \begin{bmatrix}
        \vec{e}_1 \\ \vec{q}_2
    \end{bmatrix}
    = [\gls{sym-mathbf-h}]_{1,2}
    \begin{bmatrix}
        \gls{sym-vec-q}_1 \\ \vec{e}_2
    \end{bmatrix}.
\end{equation}
These equations apply to both time-domain and phasor representations of each port.
The three forms differ in how the ports are grouped:
\begin{itemize}
    \item \textbf{Impedance} ($\gls{sym-mathbf-Z}$)
        groups all ports together, describing the effort required to produce a unit of flow.
        The intrinsic dynamics matrix $\gls{sym-mathbf-Z-i}$ and the \gls{acr-pto} impedance
        $\gls{sym-mathbf-Z-p}$ in \Cref{eq:eom-freq-domain,eq:intrinsic-impedance} both use this form.
    \item \textbf{Cascade} ($\gls{sym-mathbf-a}$)
        divides the ports into two groups and describes how effort and flow transmit from one group to the other.
        The arrow subscript indicates direction:
        $[\gls{sym-mathbf-a}]_{2\leftarrow 1}$ maps port group 1 to port group 2.
        Reversing direction inverts the matrix,
        \begin{equation}\label{eq:cascade-inverse}
            [\gls{sym-mathbf-a}]_{1\leftarrow 2} = [\gls{sym-mathbf-a}]_{2\leftarrow 1}^{-1},
        \end{equation}
        and sequential multiplication represents serial composition of two-port elements.
        This form is used for the \gls{acr-pto} dynamics (drivetrain followed by generator).
    \item \textbf{Hybrid} ($\gls{sym-mathbf-h}$)
        mixes the two conventions, expressing one effort and one flow at the two port groups as functions of the complementary variables.
        This form is used for the \gls{acr-pto} kinematics, where it naturally accommodates underactuation.
\end{itemize}

Conversion between the three forms is accomplished using the formulas proposed by \citet{reveyrand_multiport_2018}.
Some port relationships are physically well-defined in one convention but degenerate in another (for example, an ideal current source has no impedance representation but a well-defined hybrid representation) so the convention is chosen to match the physics of each subsystem.

% Verify against \citet{coe_co-design_2025} before submission and edit
% if the convention differs.
The sign convention adopted here treats both flows as positive when entering the multiport,
matching the suggestion of \citet{coe_co-design_2025}.

\subsubsection{Body Dynamics}
\label{sec:appendix-body-dynamics}
\Cref{fig:multiport-circuit-intrinsic} presents the multiport circuit representation of the floating body intrinsic dynamics.
Forces include wave excitation $\gls{sym-vec-F-e}$, hydrodynamic radiation $\gls{sym-vec-F-rad}$, hydrostatic restoring $\gls{sym-vec-F-res}$, drag $\gls{sym-vec-F-d}$, and power takeoff $\gls{sym-vec-F-p}$.

\begin{figure}[htbp]
\centering
    \includegraphics[\FigureOptions{aor:figs/from-matlab/circuit_intrinsic.pdf}]{figs/from-matlab/circuit_intrinsic.pdf}
\caption[Multiport circuit for two hydrodynamic \glspl{acr-dof}]{Multiport circuit for two hydrodynamic \glspl{acr-dof} with intrinsic and powertrain impedances $\gls{sym-mathbf-Z-i}$ and $\gls{sym-mathbf-Z-p}$}
\label{fig:multiport-circuit-intrinsic}
\end{figure}

In \Cref{fig:multiport-circuit-intrinsic}, excitation forces are represented as dependent voltage sources (diamonds) driven by a single wave elevation source (circle), and the intrinsic and powertrain dynamics are each represented as a two-port impedance matrix.
The dependent sources and intrinsic impedance matrix are inherently linked through the Haskind relation, which relates excitation coefficient magnitudes $|\gls{sym-gamma}|$ to diagonal elements of the radiation damping $\gls{sym-B-h}$.
This bounds the maximum absorbed power for any geometry.

\ifdefined\DISSERTATION
    The MDOcean dynamic model assumes regular waves for reasons to be detailed in \Cref{sec:irregular-waves}, so the system experiences monochromatic wave forcing.
    Combined with the linearity of the model, this implies that all signals are sinusoids.
    Therefore, the time-domain radiation force is a simple product rather than a convolution integral.
    With these considerations, the equation of motion can be expressed in the time domain as:
\else
    Under the regular-wave assumption justified in \Cref{sec:irregular-waves}, all signals are sinusoidal
    and the radiation force reduces to a simple product rather than a convolution integral, yielding the time-domain equation of motion:
\fi
% todo: this is time domain so gamma is not allowed to be complex
\begin{equation}\label{eq:eom}
    \gls{sym-mathbf-M}\gls{sym-ddot-vec-xi} (\gls{sym-t})
    = \gls{sym-vec-F-e} + \gls{sym-vec-F-rad} + \gls{sym-vec-F-res} + \gls{sym-vec-F-d} + \gls{sym-vec-F-p} (\gls{sym-t})
\end{equation}
where, in the time domain under monochromatic forcing,
\begin{equation}\label{eq:eom-forces}
\begin{aligned}
    \gls{sym-vec-F-e}   &= \Re\!\left(\gls{sym-vec-gamma}\,\gls{sym-tilde-zeta} (\gls{sym-t})\right),
    & \gls{sym-vec-F-rad} &= -\gls{sym-mathbf-A-h}\gls{sym-ddot-vec-xi} (\gls{sym-t}) - \gls{sym-mathbf-B-h}\gls{sym-dot-vec-xi} (\gls{sym-t}), \\
    \gls{sym-vec-F-res} &= -\gls{sym-mathbf-K-h}\gls{sym-vec-xi} (\gls{sym-t}),
    & \gls{sym-vec-F-d}   &= -\gls{sym-mathbf-B-d}\gls{sym-dot-vec-xi} (\gls{sym-t}) + \Re\!\left(\gls{sym-vec-gamma-d}\,\gls{sym-tilde-zeta} (\gls{sym-t})\right), \\
    \gls{sym-vec-F-p} (\gls{sym-t})   &= -\gls{sym-mathbf-B-p}\gls{sym-dot-vec-xi} (\gls{sym-t}) - \gls{sym-mathbf-K-p}\gls{sym-vec-xi} (\gls{sym-t}).
\end{aligned}
\end{equation}
for both operational and storm conditions.
$\gls{sym-vec-xi}(\gls{sym-t})$ (body position) and $\gls{sym-zeta}$ (wave elevation) were defined in \Cref{sec:hydro}.
The mass matrix for the operational and storm cases is given by:
\begin{equation}
    \gls{sym-M-op}=
    \begin{bmatrix}
        \gls{sym-m-f} & 0 \\ 0 & \gls{sym-m-s}
    \end{bmatrix},
    \quad
    \gls{sym-M-storm}=
    \begin{bmatrix}
        \gls{sym-m-f} + \gls{sym-m-s}
    \end{bmatrix}
\end{equation}

\ifdefined\DISSERTATION
    For each force component, a linear or quasi-linear model is adopted, with coefficients $\gls{sym-vec-gamma}$, $\gls{sym-mathbf-A-h}$, $\gls{sym-mathbf-B-h}$, $\gls{sym-mathbf-K-h}$ from the hydrodynamics module and $\gls{sym-mathbf-B-d}$, $\gls{sym-vec-gamma-d}$, $\gls{sym-mathbf-B-p}$, and $\gls{sym-mathbf-K-p}$ to be derived in the following sections.
    The excitation coefficients $\gls{sym-vec-gamma}$ and $\gls{sym-vec-gamma-d}$ are complex, and all other coefficients are real.
\else
    Hydrodynamic coefficients $\gls{sym-vec-gamma}$, $\gls{sym-mathbf-A-h}$, $\gls{sym-mathbf-B-h}$, $\gls{sym-mathbf-K-h}$ come from \Cref{sec:hydro};
    drag coefficients $\gls{sym-mathbf-B-d}$, $\gls{sym-vec-gamma-d}$ are derived in \Cref{sec:drag};
    \gls{acr-pto} coefficients $\gls{sym-mathbf-B-p}$, $\gls{sym-mathbf-K-p}$ are derived in \Cref{sec:appendix-pto-kinematics,sec:appendix-combine-intrinsic-kinematics,sec:appendix-pto-dynamics,sec:optimal-control}.
    The excitation coefficients $\gls{sym-vec-gamma}$ and $\gls{sym-vec-gamma-d}$ are complex;
    all others are real.
\fi

In the frequency domain, the complex transfer function from real regular wave height $\gls{sym-H}$ to the vector of complex body velocity phasors $\gls{sym-vec-hat-dot-xi}$ is then:
\begin{equation}\label{eq:eom-freq-domain}
    \frac{\gls{sym-vec-hat-dot-xi}}{\gls{sym-H}} = \frac{1}{2}\left[\gls{sym-mathbf-Z-i}+\gls{sym-mathbf-Z-p}\right]^{-1} (\gls{sym-vec-gamma}+\gls{sym-vec-gamma-d})
\end{equation}
where the inertial, hydrodynamic radiation, hydrostatic restoring, and drag damping terms have been combined into a single intrinsic impedance matrix $\gls{sym-mathbf-Z}_{i}$ for compactness:
\begin{equation}\label{eq:intrinsic-impedance}
\gls{sym-mathbf-Z-i} = -i\gls{sym-omega} (\gls{sym-mathbf-M}+\gls{sym-mathbf-A-h})+ (\gls{sym-mathbf-B-h}+\gls{sym-mathbf-B-d})+\frac{1}{i\gls{sym-omega}} (\gls{sym-mathbf-K-h})
\end{equation}
The \gls{acr-pto} kinematics and dynamics are yet to be specified, currently represented by generic impedance matrix $\gls{sym-mathbf-Z-p} = \gls{sym-mathbf-B-p}+\frac{\gls{sym-mathbf-K-p}}{i\gls{sym-omega}}$.

%%%%%%%%%%%%%%%
\subsubsection{\gls{acr-pto} Kinematics and Underactuation}
\label{sec:appendix-pto-kinematics}
In this section and the next, we turn our attention to the \gls{acr-pto} kinematics and dynamics, to specify the generic impedance matrix $\gls{sym-mathbf-Z-p}$ in \Cref{eq:eom-freq-domain}.

A \gls{acr-pto} model must capture the dynamics of any elements that transmit, transduce, or condition power between the rigid-body mechanical energy and the end energy product, in this case electricity.
Relevant dynamics include the inertia, stiffness, static friction, mechanical advantage, losses, and electromagnetic coupling of flywheels, shafts, springs, gears, hydraulic circuits, generators, and power electronics.
\citet{penalba_review_2016} review \gls{acr-wec} \gls{acr-pto} models and \citet{coe_co-design_2025} synthesize these into a multiport impedance matching network framework to facilitate \gls{acr-ccd}.
\citet{coe_co-design_2025} present a \gls{acr-pto} model in cascade matrix form for a single \gls{acr-dof} system.
The \gls{acr-pto} kinematics describe the relationship between \gls{acr-wec} body displacements $\gls{sym-vec-hat-xi}$ to \gls{acr-pto} displacements $\gls{sym-vec-hat-X-PTO}$,
determined by linkages or similar mechanisms and not including subsequent \gls{acr-pto} elements such as a mechanical drivetrain or generator.
In general form, the \gls{acr-pto} kinematics can be expressed as:
\begin{equation}\label{eq:pto-kinematics-motion}
\gls{sym-vec-hat-X-PTO} = \gls{sym-mathbf-T-kin}~ \gls{sym-vec-hat-xi}
\end{equation}
with kinematic transformation matrix $\gls{sym-mathbf-T-kin}$ that may depend on the body displacements $\gls{sym-vec-hat-xi}$ if the linkages are nonlinear,
but hereafter is assumed constant for simplicity.
This assumption means that \Cref{eq:pto-kinematics-motion} can be differentiated trivially to yield the same matrix relationship for velocities.
Assuming power conservation across the linkages,
the kinematics describe the relationship between not only motions but also forces:
\begin{equation}\label{eq:pto-kinematics-force}
\gls{sym-vec-hat-F-p} = [\gls{sym-mathbf-T-kin}]^\top  \gls{sym-vec-hat-F-PTO}
\end{equation}
To avoid confusion, subscript $\gls{acr-pto}$ is used for forces in the \gls{acr-pto} \glspl{acr-dof} (post-kinematics),
while subscript $\gls{sym-p}$ denotes forces in the hydrodynamic body \glspl{acr-dof} (pre-kinematics).
\Cref{eq:pto-kinematics-motion,eq:pto-kinematics-force} imply rigid lossless massless linkages,
and any linkage deflection, dissipation, and inertia would be included in the \gls{acr-pto} dynamics rather than kinematics.

Multibody \glspl{acr-wec} are frequently underactuated,
meaning that the \gls{acr-pto} has fewer controlled \glspl{acr-dof} than the hydrodynamic bodies,
and $\gls{sym-mathbf-T-kin}$ has more columns than rows.
Underactuated systems receive considerable attention in robotics \citep{tedrake_underactuated_2024}
but have only recently been explored in the context of wave energy conversion \citep{faedo_principle_2022}.
Cascade matrices require an identical number of input and output ports
and thus cannot be used to describe the unbalanced ports that underactuated kinematics create
\citep{reveyrand_multiport_2018}.
Non-square impedance matrices may be used to capture unbalanced port relations in general,
but the impedance is undefined for the specific relations of
\Cref{eq:pto-kinematics-motion,eq:pto-kinematics-force}.
Instead, the kinematics can be packaged into a square hybrid matrix $\gls{sym-mathbf-h-kin}$
which accurately represents the underactuated system:
\begin{equation}\label{eq:h-matrix-kinematics}
    \begin{bmatrix}
        \gls{sym-vec-hat-F-p} \\
        \gls{sym-vec-hat-dot-X-PTO}
    \end{bmatrix}
    = \underbrace{
        \begin{bmatrix}
            \mathbf{0}             & \gls{sym-mathbf-T-kin}^{\top} \\
            \gls{sym-mathbf-T-kin} & \mathbf{0}
        \end{bmatrix}}
        _{\gls{sym-mathbf-h-kin}}
    \begin{bmatrix}
        \gls{sym-vec-hat-dot-xi} \\ \gls{sym-vec-hat-F-PTO}
    \end{bmatrix}
\end{equation}

Here, we model two hydrodynamic \glspl{acr-dof} (heave of the float and spar respectively) and a single \gls{acr-pto} \gls{acr-dof} corresponding to the relative heave motion between the float and spar.
The kinematics matrix is therefore:
\begin{equation}
    \label{eq:T-kin-specific}
    \gls{sym-mathbf-T-kin} = \begin{bmatrix}1 & -1\end{bmatrix}
\end{equation}
Thus, even absent any mechanism with ``kinematics'' in the traditional sense,
the kinematics matrix is filled with zeroes, ones, and negative ones to represent which \glspl{acr-dof} are controlled.

In the circuit representation with two hydrodynamic and one \gls{acr-pto} \gls{acr-dof},
the hybrid matrix \Cref{eq:h-matrix-kinematics} can be represented as a three-port element,
as shown in \Cref{fig:multiport-circuit-kinematics}.
\begin{figure}[htbp]
    \centering
    \includegraphics[\FigureOptions{aor:figs/from-matlab/circuit_kinematics.pdf}]{figs/from-matlab/circuit_kinematics.pdf}
    \caption[Multiport circuit with powertrain kinematics and Th\'{e}venin equivalent]{Multiport circuit with powertrain kinematics represented as a hybrid matrix, followed by the Th\'{e}venin equivalent of the intrinsic dynamics and kinematics}
    \label{fig:multiport-circuit-kinematics}
\end{figure}

%%%%%%%%%%%%%%%
\subsubsection{Combining Intrinsic Dynamics with \gls{acr-pto} Kinematics}
\label{sec:appendix-combine-intrinsic-kinematics}
Taking the Th\'{e}venin equivalent of the intrinsic dynamics and kinematics at the \gls{acr-pto} port yields a scalar equivalent circuit for each \gls{acr-pto} \gls{acr-dof},
with force sources $\gls{sym-vec-hat-F-th}$ and mechanical impedances $\gls{sym-vec-Z-i-th}$,
shown in the bottom portion of \Cref{fig:multiport-circuit-kinematics} and computed with the following expressions:
\begin{equation}
\begin{aligned}
    \gls{sym-vec-hat-F-th} &= (\gls{sym-mathbf-T-kin}\gls{sym-mathbf-T-kin}^{\top})^{-1}
                               \gls{sym-mathbf-T-kin} \gls{sym-mathbf-Z-i}^{-1} \gls{sym-vec-gamma} \gls{sym-hat-zeta}
    = \frac{1}{2} [1,~-1] \gls{sym-mathbf-Z-i}^{-1} \gls{sym-vec-gamma} \gls{sym-hat-zeta} \\
    \gls{sym-vec-Z-i-th} &= (-\gls{sym-mathbf-T-kin}\gls{sym-mathbf-Z-i}^{-1} \gls{sym-mathbf-T-kin}^{\top})^{-1}
                          = -\frac{\det(\gls{sym-mathbf-Z-i})}{\gls{sym-Sigma} \gls{sym-mathbf-Z-i}}
                          = -\frac{\gls{sym-Z-i-f}\gls{sym-Z-i-s}-\gls{sym-Z-i-c}^2}
                                  {\gls{sym-Z-i-f}+\gls{sym-Z-i-s}+2\gls{sym-Z-i-c}}
\end{aligned}
\end{equation}
where the vector expressions for the general case are shown first,
followed by the scalar result for the $\gls{sym-mathbf-T-kin}$ in \Cref{eq:T-kin-specific} with 1 controlled \gls{acr-dof}.

Although, as previously mentioned, a true cascade matrix must have the same number of input and output ports, it is possible to formulate a cascade-like matrix that can transform from the \gls{acr-pto} port to the body ports.
However, it cannot transform in the other direction due to underactuation.

Even in the favorable direction, this matrix cannot transform to the body ports directly.
Due to the wave excitation, body velocities can be nonzero even when the \gls{acr-pto} force and velocity are both zero.
Hence, there is no way to express $\gls{sym-hat-dot-xi}$ as a weighted sum of $\gls{sym-hat-F-PTO}$ and $\gls{sym-dot-hat-X-PTO}$.
$\gls{sym-hat-dot-xi}$ could be expressed as a function of the \gls{acr-pto} variables by substituting
$\gls{sym-mathbf-Z-p}=\gls{sym-mathbf-T-kin}^{\top} (\gls{sym-vec-hat-F-PTO}/\gls{sym-vec-hat-dot-X-PTO}) \gls{sym-mathbf-T-kin}$
into \Cref{eq:eom-freq-domain}, but this yields a nonlinear expression of limited utility.
Instead, we can use cascade-like matrices to solve for a modified variable $\gls{sym-vec-dot-xi-Delta}$ defined as the change in the body velocities caused by the \gls{acr-pto},
in other words the difference between the actual body velocities $\gls{sym-vec-dot-xi}$
and their values if the \gls{acr-pto} were disconnected with $\gls{sym-F-PTO}=0$, $\gls{sym-vec-dot-xi-0}$:
\begin{equation}\label{eq:modified-body-velocity}
    \gls{sym-vec-dot-xi-Delta} = \gls{sym-vec-dot-xi} - \gls{sym-vec-dot-xi-0}
    \quad \textrm{ where } \quad
    \gls{sym-vec-dot-xi-0} = \gls{sym-mathbf-Z-i}^{-1}~\gls{sym-vec-gamma}~\frac{\gls{sym-H}}{2}
\end{equation}
The cascade-like matrix from the \gls{acr-pto} port $\{\gls{sym-F-PTO}, \gls{sym-dot-X-PTO}\}$
to the modified body port $\{\gls{sym-F-p}, \gls{sym-dot-xi-Delta}\}$ is then:
\begin{equation}\label{eq:body-transmission-matrix}
    [\gls{sym-mathbf-a}]_{\gls{sym-F-p}\gls{sym-xi-Delta}\leftarrow \gls{sym-F-PTO}\gls{sym-dot-X}} =
    \begin{bmatrix}
        \gls{sym-mathbf-I}^{2\times2} \\
        -\gls{sym-mathbf-Z-i}^{-1}
    \end{bmatrix}
    \gls{sym-mathbf-T-kin}^{\top}
    \begin{bmatrix}
        1 & 0
    \end{bmatrix}
\end{equation}
noting that this matrix is non-square ($\in \gls{sym-mathbb-C}^{4\times2}$) since there are more hydrodynamic than \gls{acr-pto} \glspl{acr-dof},
and that the matrix label drops the $PTO$ subscripts for compactness.
\Cref{fig:multiport-circuit-kinematics} represents this in circuit form,
noting that the excitation source must be duplicated to create $\gls{sym-vec-hat-dot-xi-0}$.

%%%%%%%%%%%%%%%
\subsubsection{\gls{acr-pto} Dynamics}
\label{sec:appendix-pto-dynamics}
In addition to the \gls{acr-pto} port already discussed, we define two additional ports of relevance:
the generator mechanical port (torque $\gls{sym-tau}$ and rotation speed $\gls{sym-Omega}$)
and the generator electrical port (voltage $\gls{sym-V}$ and current $\gls{sym-I}$).
Assuming a drivetrain consisting of a linear to rotational gear ratio with equivalent radius $\gls{sym-R-gear}$ followed by a driveshaft with impedance $\gls{sym-Z-shaft}$,
the \gls{acr-pto} mechanical dynamics are expressed with cascade matrix
$[\gls{sym-mathbf-a-F-dot-X-leftarrow-tau-Omega}]$:
\begin{equation}\label{eq:mech-transmission-matrix}
\begin{bmatrix}
\gls{sym-F-PTO} \\ \gls{sym-dot-X-PTO}
\end{bmatrix}
=
\underbrace{\begin{bmatrix}
\gls{sym-Z-shaft}/\gls{sym-R-gear} & 1/\gls{sym-R-gear} \\
\gls{sym-R-gear} & 0
\end{bmatrix}}
_{[\gls{sym-mathbf-a-F-dot-X-leftarrow-tau-Omega}]}
\begin{bmatrix}
 \gls{sym-tau}\\ \gls{sym-Omega}
\end{bmatrix}
\end{equation}

Meanwhile, a synchronous permanent magnet generator modeled as an ideal gyrator with torque constant $\gls{sym-k-tau}$
plus a winding impedance $\gls{sym-Z-w}$ has cascade matrix $[\gls{sym-mathbf-a-tau-Omega-leftarrow-VI}]$:
\begin{equation}\label{eq:elec-transmission-matrix}
    \begin{bmatrix}
        \gls{sym-tau} \\ \gls{sym-Omega}
    \end{bmatrix}
    =
    \underbrace{\begin{bmatrix}
    0 & \gls{sym-k-tau} \\
    1/\gls{sym-k-tau} & \gls{sym-Z-w}/\gls{sym-k-tau}
    \end{bmatrix}}
    _{[\gls{sym-mathbf-a-tau-Omega-leftarrow-VI}]}
    \begin{bmatrix}
        \gls{sym-V} \\ \gls{sym-I}
    \end{bmatrix}
\end{equation}
with Park-transformed quadrature electrical voltage and current $\gls{sym-V}$ and $\gls{sym-I}$.

The system operating point is determined by a controller that applies some electrical load to the generator.
The controller can be formulated as an impedance at either the generator mechanical or electrical port.
The former is more convenient for hardware implementation, because modern drives typically expect torque commands,
while the latter is more convenient for electrical power calculations and optimal control.
%Here we formulate both as linear impedances with damping and stiffness terms, where
The mechanical control impedance $\gls{sym-Z-u}$ and electrical load impedance $\gls{sym-Z-l}$ are defined as follows:
\begin{equation}\label{eq:controller-impedances}
    \gls{sym-Z-l} = \frac{\gls{sym-V}}{\gls{sym-I}} = \gls{sym-B-l} + \frac{\gls{sym-K-l}}{i\gls{sym-omega}}
    ,  \qquad
    \gls{sym-Z-u} = \frac{\gls{sym-tau}}{\gls{sym-Omega}}% = B_u + \frac{K_u}{s}
\end{equation}
Combining \Cref{eq:controller-impedances,eq:elec-transmission-matrix} gives the following expression for $\gls{sym-Z-u}$ as a function of $\gls{sym-Z-l}$:
\begin{equation}\label{eq:control-impedance}
    \gls{sym-Z-u} =
    \frac{
    \begin{bmatrix}
    1 & 0
    \end{bmatrix}
    [\gls{sym-mathbf-a-tau-Omega-leftarrow-VI}]
    \begin{bmatrix}
    \gls{sym-Z-l} \\ 1
    \end{bmatrix}
}{
    \begin{bmatrix}
    0 & 1
    \end{bmatrix}
    [\gls{sym-mathbf-a-tau-Omega-leftarrow-VI}]
    \begin{bmatrix}
    \gls{sym-Z-l} \\ 1
    \end{bmatrix}
}
\end{equation}
This equation allows the $\gls{sym-Z-l}$ identified with optimal control to be implemented as a corresponding $\gls{sym-Z-u}$ in hardware.
%Incorporating \Cref{eq:mech-transmission-matrix} yields the Th\'{e}venin equivalent scalar PTO impedance at the hydrodynamic body port:
% \begin{equation}\label{eq:pto-impedance}
% Z_{PTO,th} =
% \frac{
%     \begin{bmatrix}
%     1 & 0
%     \end{bmatrix}
%     [\mathbf{a}]_{F\dot{X}\leftarrow \tau \Omega}
%     [\mathbf{a}]_{\tau \Omega \leftarrow V I}
%     \begin{bmatrix}
%     Z_l \\ 1
%     \end{bmatrix}
% }{
%     \begin{bmatrix}
%     0 & 1
%     \end{bmatrix}
%     [\mathbf{a}]_{F\dot{X}\leftarrow \tau \Omega} [\mathbf{a}]_{\tau \Omega \leftarrow V I}
%     \begin{bmatrix}
%     Z_l \\ 1
%     \end{bmatrix}
% }
% \end{equation}
% which is useful for calculating the mechanical response of each body.

The \gls{acr-pto} dynamics can be represented as a multiport circuit, as shown in \Cref{fig:multiport-circuit-pto}.
\begin{figure}[htbp]
\centering
    \includegraphics[\FigureOptions{aor:figs/from-matlab/circuit_pto.pdf}]{figs/from-matlab/circuit_pto.pdf}
\caption[\gls{acr-pto} dynamics represented as a multiport circuit]{\gls{acr-pto} dynamics represented as a multiport circuit and electrical Th\'{e}venin equivalent}
\label{fig:multiport-circuit-pto}
\end{figure}

The intrinsic dynamics, \gls{acr-pto} kinematics, and \gls{acr-pto} dynamics can be combined into a Th\'{e}venin equivalent source voltage and impedance $\gls{sym-hat-V-s-th}$ and $\gls{sym-Z-s-th}$ as follows:
\begin{equation}\label{eq:thevenin-electrical}
    \gls{sym-Z-s-th} = \frac{
        \begin{bmatrix}1 & 0\end{bmatrix}
        \gls{sym-mathbf-a-VI-leftarrow-F-dot-X}
        \begin{bmatrix}-\gls{sym-Z-i-th} \\ 1\end{bmatrix}
    }
    {
        \begin{bmatrix}0 & 1\end{bmatrix}
        \gls{sym-mathbf-a-VI-leftarrow-F-dot-X}
        \begin{bmatrix}-\gls{sym-Z-i-th} \\ 1\end{bmatrix}
    } = \frac{
                \gls{sym-B-a} - \gls{sym-A-a} \gls{sym-Z-i-th}
            }
            {
                \gls{sym-D-a}-\gls{sym-C-a} \gls{sym-Z-i-th}
            }
    \qquad
    \gls{sym-hat-V-s-th} =
        \gls{sym-hat-F-th}
        (\gls{sym-A-a} - \gls{sym-C-a} \gls{sym-Z-s-th})
\end{equation}
where $\gls{sym-A-a},\gls{sym-B-a},\gls{sym-C-a},\gls{sym-D-a}$
are the elements of
$\gls{sym-mathbf-a-VI-leftarrow-F-dot-X}$.
This is shown in the bottom portion of \Cref{fig:multiport-circuit-pto}.
The load voltage and current can then be found as:
\begin{equation}\label{eq:voltage-current}
    \gls{sym-hat-V} = \gls{sym-hat-V-s-th}\frac{
                                                    \gls{sym-Z-l}
                                                }
                                                {
                                                    \gls{sym-Z-s-th} + \gls{sym-Z-l}
                                                },
            \qquad \gls{sym-hat-I} = \frac{
                                                \gls{sym-hat-V-s-th}
                                            }
                                            {
                                                \gls{sym-Z-s-th}+\gls{sym-Z-l}
                                            }
\end{equation}

%%%%%%%%%%%%%%%

\subsubsection{Time-Domain Expansion of Instantaneous Power}
\label{sec:appendix-power-time}

Equation~\eqref{eq:power-PQS} can be used to rewrite the instantaneous power $\gls{sym-p}(\gls{sym-t})$ at any port as \citep{saadat_power_1999}
\begin{equation}\label{eq:power-time-PQS}
    \gls{sym-p}(\gls{sym-t}) =
    \gls{sym-P} + |\gls{sym-S}|
    \cos\!\left(
                2(\gls{sym-omega} \gls{sym-t} + \angle\hat{e})
                - \tan^{-1}\!\left(
                                    \frac{\gls{sym-Q}} {\gls{sym-P}}
                            \right)
        \right),
\end{equation}
a \gls{acr-dc} offset of $\gls{sym-P}$ plus an oscillation at $2\gls{sym-omega}$ with amplitude $|\gls{sym-S}|$.
The average and peak powers in \Cref{eq:power-avg-peak} follow directly:
$\gls{sym-p-avg} = \gls{sym-P}$ since the oscillation averages to zero,
and $\gls{sym-p-pk} = \gls{sym-P} + |\gls{sym-S}|$ at the cosine maximum.
The minimum $\gls{sym-p-min} = \gls{sym-P} - |\gls{sym-S}|$ is negative when $|\gls{sym-S}|>\gls{sym-P}$,
corresponding to instants when energy flows from the load back into the source under reactive control.

%%%%%%%%%%%%%%%
\subsubsection{Electrical-Port Power in Optimal-Control Variables}
\label{sec:appendix-power-elec}

Expressing the average and peak powers at the electrical port in terms of the constrained-control parameters introduced in \Cref{sec:optimal-control}
(control damping $\gls{sym-B-l}$, control stiffness $\gls{sym-K-l}$, and current phasor $\gls{sym-hat-I}$),
and applying the constant electrical efficiency $\gls{sym-eta}$,
\begin{equation}\label{eq:power-elec}
    \begin{bmatrix}
        \gls{sym-p-avg-elec} \\
        \gls{sym-p-pk-elec}
    \end{bmatrix}
    =
    \gls{sym-eta}
    \begin{bmatrix}
        \gls{sym-p-avg-VI} \\
        \gls{sym-p-pk-VI}
    \end{bmatrix}
    =
    \gls{sym-eta} \tfrac{1}{2} \gls{sym-B-l} |\gls{sym-hat-I}|^2
    \begin{bmatrix}
        1 \\
        1 + \sqrt{1 + \left(
                        \tfrac{\gls{sym-K-l}} {\gls{sym-B-l}\gls{sym-omega}}
                    \right)^2}
    \end{bmatrix}.
\end{equation}
The peak-power form is used as the apparent-power constraint in the \gls{acr-qcqp} of \Cref{sec:optimal-control}.

%%%%%%%%%%%%%%%
% Additionally, the drivetrain mechanical cascade matrix $[\gls{sym-a}]_{\gls{sym-m}}$ is set to the identity matrix, which represents a static effective gear radius $\gls{sym-R-gear}=1$:
% \begin{equation}
% [a]_m = \begin{bmatrix}
% 1/\gls{sym-R-gear} &~ 0 \\
% 0 &~ \gls{sym-R-gear}
% \end{bmatrix} = \begin{bmatrix}
% 1 &~ 0 \\
% 0 &~ 1
% \end{bmatrix}
% \end{equation}
% With this $[\gls{sym-a}]_{\gls{sym-m}}$, the control damping and stiffness $\gls{sym-B-u}$ and $\gls{sym-K-u}$ equal the PTO damping and stiffness $\gls{sym-B-p}$ and $\gls{sym-K-p}$.
% This will simplify the impedance matching procedure described below, requiring matching at only one interface.

%The time-domain mechanical power is $\gls{sym-P}_{\text{mech}}(\gls{sym-t}) = \gls{sym-tau}(\gls{sym-t}) \gls{sym-Omega}(\gls{sym-t}) = (\gls{sym-B-u} \gls{sym-Omega}(\gls{sym-t}) + \gls{sym-K-u} \gls{sym-theta}(\gls{sym-t}) ) \gls{sym-Omega}(\gls{sym-t}) $, where $\gls{sym-theta}(\gls{sym-t})$ is the generator rotational position. Assuming sinusoidal waveforms, the time-average and peak mechanical power in a given sea state are then
%%%%%%%%%%%%%%%

%%%%%%%%%%%%%%%%%%%%%%%%%%%%%%%%%%%%%%%%%
\subsubsection{Linear Solution Procedure}
\label{sec:appendix-dynamics-solution}
\Cref{sec:eom,sec:pto-model-overview} introduce linear models for the rigid body dynamics, \gls{acr-pto} kinematics, and \gls{acr-pto} dynamics
using impedance, hybrid, and cascade matrices respectively.
Here we unify the formulation and present a clear system-level solution procedure for all port variables:
\begin{enumerate}
    \item Use \Cref{eq:thevenin-electrical} to assemble the Th\'{e}venin source voltage and impedance
        $\gls{sym-hat-V-s-th}$ and $\gls{sym-Z-s-th}$
        as seen from the port on which power should be maximized
        (in this case, the electrical generator port
        $\{\gls{sym-V},\gls{sym-I}\}$).
    \item Set the load impedance $\gls{sym-Z-l}$.
        Any load impedance can be simulated,
        but typically it is selected as a function of $\gls{sym-Z-s-th}$ to maximize the power at this port,
        possibly subject to constraints on various ports.
        \Cref{sec:optimal-control} frames the load impedance selection as a constrained optimization problem.
    \item From  $\gls{sym-hat-V-s-th}$, $\gls{sym-Z-s-th}$, and $\gls{sym-Z-l}$,
        calculate the $\{\gls{sym-V},\gls{sym-I}\}$ port variables from \Cref{eq:voltage-current}.
    \item Use cascade matrices \Cref{eq:mech-transmission-matrix,eq:elec-transmission-matrix,eq:body-transmission-matrix}
        to propagate the variables back toward the rigid body,
        yielding the $\{\gls{sym-tau},\gls{sym-Omega}\}$, $\{\gls{sym-F-PTO},\gls{sym-dot-X-PTO}\}$,
        and $\{\gls{sym-F-p}, \gls{sym-dot-xi-Delta}\}$ port variables.
    \item Transform from $\gls{sym-dot-xi-Delta}$ to $\gls{sym-dot-xi}$ via \Cref{eq:modified-body-velocity}.
\end{enumerate}
Once all the port variables have been calculated,
power at each port can be calculated according to \Cref{sec:power}.
This serves as the solution to a linear subproblem,
which is then re-solved using updated quasi-linear coefficients until drag and force saturation nonlinearities have converged.


\FloatBarrier
\subsection{Optimal Control Details}
\label{sec:appendix-optimal-control}
%%%%%%%%%%%%%%%%%%%%%%%%%%%%%%%%%%%%%%%%%
\ifdefined\DISSERTATION
    \subsubsection{Indirect Optimal Control Solution Procedure}
    \label{sec:appendix-qp-numerical}
    The indirect optimal control approach applies optimality conditions and then numerically solves the resulting equations.
    The optimality conditions depend on the active set of constraints.
    In the most straightforward case where simply one constraint exists, the constraint is active if and only if the unconstrained solution violates that constraint.
    For example, if only a maximum force constraint exists, the constrained solution has a force $\gls{sym-hat-F-p-constr}$ that is the minimum of the unconstrained force $\gls{sym-hat-F-p-unconstrained}$ and the maximum allowed force fundamental $\gls{sym-F-max}$:
    \begin{equation}
    \gls{sym-hat-F-p-constr} = \min\left(\gls{sym-hat-F-p-unconstrained}, \gls{sym-F-max}\right)
    \end{equation}
    The solver guesses a value of $\gls{sym-Gamma}$, calculates the resulting force $\gls{sym-hat-F-p-guess}$, and then uses $\gls{sym-hat-F-p-guess}-\gls{sym-hat-F-p-constr}$ as an error signal to adjust its next guess of $\gls{sym-Gamma}$, iterating until the error converges to zero.
    This solver-based optimal control performs numerical root-finding, not numerical optimization, and can happen in an inner loop of the drag-linearization solver iteration (where $\gls{sym-hat-F-p-unconstrained}$ stays constant each control iteration) or simultaneously with the drag-linearization iteration (where $\gls{sym-hat-F-p-unconstrained}$ and the system linearization are updated each control iteration).
    MDOcean implements the latter, which simplifies the setup by requiring only one solver.
    Comparing the efficiency of the two approaches is left for future work.

    When multiple constraints exist, the situation is more complex, since it is possible for a given constraint to be violated by the unconstrained solution but not active in the constrained solution, or for multiple constraints to be active simultaneously.
    In the special case where the unconstrained solution violates no more than one constraint, that constraint is active and the solver can proceed as described above.
    When the unconstrained solution violates multiple constraints, the solver must try all subsets of 1-2 of those constraint(s) as the active set, find which of these active sets produces the largest power while satisfying all constraints, and then set the error as the deviation of the guessed solution from the constrained solution for that active set.
    Since the active set cannot be calculated a-priori as in the simpler case, this approach ends up looking more like a numerical optimization than a numerical root-finding, although the optimization is over a finite set of active set possibilities instead of a continuous control domain as it would be in a direct optimal control method.
    The main advantage of the indirect approach is its ability to handle arbitrary nonlinear\footnote{Nonlinear meaning greater than second order in an optimization sense, not in a dynamic sense, because the dynamics are already quasi-linearized here} constraints, although for the quadratic constraints considered here, the analytical approach described in \Cref{sec:appendix-qp-solution} is more efficient.
\fi
%%%%%%%%%%%%%%%%%%%%%%%%%%%%%%%%%%%%%%%%%
\subsubsection{Analytical \gls{acr-qp} for Constrained Power Maximization}
\label{sec:appendix-qp-solution}
This section describes the analytical solution process for the constrained optimal control problem \Cref{eq:opt-problem} introduced in \Cref{sec:optimal-control}.

%%%%%%%%%%%%%%%
\paragraph{Quadratic Constraint Coefficients}
\label{sec:appendix-qp-constraints-table}

The quadratic coefficients $\{\gls{sym-mathbf-Q-mu}, \gls{sym-vec-a-mu}, \gls{sym-b-mu}\}$ in \Cref{eq:opt-problem} for each constraint implemented in MDOcean are given in \Cref{tab:qp-constraints}.

\begin{table}[b!]
\centering
\caption{Quadratic constraint coefficients}
\label{tab:qp-constraints}
\begin{tabular}{lccc}
    Constraint & $\gls{sym-mathbf-Q-mu}$ & $\gls{sym-vec-a-mu}$ & $\gls{sym-b-mu}$ \\ \hline
    $|\gls{sym-hat-xi-f}| \leq \gls{sym-xi-f-max}$ &
        $\frac{1}{\gls{sym-omega}^2}
        [\gls{sym-mathbf-a-F-p-dot-xi-f-leftarrow-VI}]^*
        \gls{sym-mathbf-Q-hat-q-2}
        [\gls{sym-mathbf-a-F-p-dot-xi-f-leftarrow-VI}]$ &
        $[0,~ \gls{sym-hat-xi-f-0}]^\top$ & $\gls{sym-xi-f-max} - |\gls{sym-hat-xi-f-0}|$ \\
    $|\gls{sym-hat-xi-s}| \leq \gls{sym-xi-s-max}$ &
        $\frac{1}{\gls{sym-omega}^2}
        [\gls{sym-mathbf-a-F-p-dot-xi-s-leftarrow-VI}]^*
        \gls{sym-mathbf-Q-hat-q-2}
        [\gls{sym-mathbf-a-F-p-dot-xi-s-leftarrow-VI}]$ &
        $[0,~ \gls{sym-hat-xi-s-0}]^\top$ & $\gls{sym-xi-s-max} - |\gls{sym-hat-xi-s-0}|$ \\
    $|\gls{sym-hat-X-PTO}| \leq \gls{sym-X-PTO-max}$ &
        $\frac{1}{\gls{sym-omega}^2}
        [\gls{sym-mathbf-a-F-dot-X-leftarrow-VI}]^*
        \gls{sym-mathbf-Q-hat-q-2}
        [\gls{sym-mathbf-a-F-dot-X-leftarrow-VI}]$ &
        $\gls{sym-vec-zero}$ & $\gls{sym-X-PTO-max}^2$ \\
    $|\gls{sym-hat-tau}| \leq \gls{sym-tau-max}$ &
        $[\gls{sym-mathbf-a-tau-Omega-leftarrow-VI}]^*
        \gls{sym-mathbf-Q-hat-e-2}
        [\gls{sym-mathbf-a-tau-Omega-leftarrow-VI}]$ &
        $\gls{sym-vec-zero}$ & $\gls{sym-tau-max}^2$ \\
    $\gls{sym-P-tau-Omega} \leq \gls{sym-P-tau-Omega-max}$ &
        $[\gls{sym-mathbf-a-tau-Omega-leftarrow-VI}]^*
        \gls{sym-mathbf-Q-P}
        [\gls{sym-mathbf-a-tau-Omega-leftarrow-VI}]$ &
        $\gls{sym-vec-zero}$ & $\gls{sym-P-tau-Omega-max}$ \\
    $\gls{sym-P-VI} \geq 0$ &
        $-\gls{sym-mathbf-Q-P}$ &
        $\gls{sym-vec-zero}$ & $0$ \\
\end{tabular}
\end{table}

%%%%%%%%%%%%%%%
\paragraph{Reduced Dimension Formulation}
The linear equality constraint $\gls{sym-vec-c}^{\,*}\gls{sym-vec-x} = d$ reduces the optimization problem dimension from $\gls{sym-vec-x} \in \gls{sym-mathbb-C}^2$ to $\gls{sym-mathbb-C}^1$.
Specifically, the equality constraint requires that the solution be a particular solution plus a scalar multiple of $\gls{sym-vec-n}$, the null space vector of $\gls{sym-vec-c}^{\,*}$:
\begin{equation}\label{eq:x-p-plus-nullspace}
    \gls{sym-vec-x} = \underbrace{
        \frac{\gls{sym-hat-V-s-th}}{2\Re(\gls{sym-Z-s-th})}
        \begin{bmatrix} \ Z_{s,th}^* \\ 1 \end{bmatrix}
    }_{
        \textrm{particular solution } \gls{sym-vec-x-p}}
    -~ \gls{sym-Gamma} \cdot \frac{\gls{sym-hat-V-s-th}}{2\Re(\gls{sym-Z-s-th})}
    \underbrace{
        \begin{bmatrix} -\ Z_{s,th}^* \\ 1 \end{bmatrix}
    }_{
        \textrm{null space } \gls{sym-vec-n}}
\end{equation}
where $\gls{sym-Gamma}$ is a nondimensional complex scalar representing the remaining \gls{acr-dof} in the optimization after enforcing the equality constraint.
For the particular solution, we have chosen to use the optimal solution in the absence of inequality constraints, which is identical to the solution implied by the impedance-matching condition \Cref{eq:matched-load} and can be derived from the quadratic program as follows:
\begin{equation}\label{eq:x-p}
    \gls{sym-vec-x-p} = \frac{d \gls{sym-mathbf-Q-P}\gls{sym-vec-c}}{\gls{sym-vec-c}^{\,*}\gls{sym-mathbf-Q-P}\gls{sym-vec-c}}
\end{equation}
This selection for $\gls{sym-vec-x-p}$ means that the \gls{acr-dof} $\gls{sym-Gamma}$ represents the deviation of the solution from the impedance-matched solution.
Using $\gls{sym-hat-I-p}$ for brevity to denote the current in the particular solution, \Cref{eq:x-p-plus-nullspace} can be rewritten as:
\begin{equation}\label{eq:x-p-plus-nullspace-simplified}
\gls{sym-vec-x} = \gls{sym-vec-x-p} - \gls{sym-Gamma} \gls{sym-hat-I-p} \gls{sym-vec-n}
    %&= \begin{bmatrix}
    %    1+\Gamma &  0 \\
    %    0 & 1-\Gamma
    %\end{bmatrix}
    %\vec{x}_p
\end{equation}

Inequality constraints
$\gls{sym-vec-x}^{\,*}\gls{sym-mathbf-Q-mu}\gls{sym-vec-x}
+ 2\Re\!\left\{\gls{sym-vec-a-mu}^{\,*}\gls{sym-vec-x}\right\}
\le \gls{sym-b-mu}$
further reduce the feasible design space, and the presence of multiple inequality constraints introduces the possibility of \Cref{eq:opt-problem} having no solutions due to infeasibility.
While constrained quadratic programs like \Cref{eq:opt-problem} generally require numerical solutions, the low dimensionality allows an analytical solution with geometric intuition, even if the problem is non-convex (see \Cref{sec:appendix-qp-convexity} for a discussion on convexity).
First, we rewrite the objective and inequality constraints in terms of $\gls{sym-Gamma}$ by plugging \Cref{eq:x-p-plus-nullspace-simplified} into \Cref{eq:opt-problem}:
\begin{subequations}
\begin{align}
\gls{sym-p-avg-VI}
&= \frac{|\gls{sym-hat-V-s-th}|^2}{8\Re (\gls{sym-Z-s-th})}~ (1 - |\gls{sym-Gamma}|^2)
\label{eq:obj-gamma}
\\[6pt]
\gls{sym-A-q-mu}|\gls{sym-Gamma}|^2 +&
2\Re\!\left\{
    \gls{sym-B-q-mu}\gls{sym-Gamma}
\right\}+\gls{sym-C-q-mu}
\le
0
\quad \forall \gls{sym-mu}
\label{eq:ineq-gamma}
\end{align}
\end{subequations}
where quadratic coefficients $\gls{sym-A-q-mu}$, $\gls{sym-B-q-mu}$, and $\gls{sym-C-q-mu}$ from \Cref{eq:ineq-gamma} are defined as:
\begin{equation}
\begin{aligned}
    \gls{sym-A-q-mu} &= |\gls{sym-hat-I-p}|^{2},
    \gls{sym-vec-n}^{\,*}\gls{sym-mathbf-Q-mu}\gls{sym-vec-n}& &\in \gls{sym-mathbb-R} \\
    \gls{sym-B-q-mu} &= -\gls{sym-hat-I-p},
    \left(\gls{sym-vec-x-p}^{\,*}\gls{sym-mathbf-Q-mu} + \gls{sym-vec-a-mu}^{\,*}\right)\gls{sym-vec-n}& &\in \gls{sym-mathbb-C} \\
    \gls{sym-C-q-mu} &=
        \gls{sym-vec-x-p}^{\,*}\gls{sym-mathbf-Q-mu}\gls{sym-vec-x-p}
        +
        2\Re\!\left\{\gls{sym-vec-a-mu}^{\,*}\gls{sym-vec-x-p}\right\}
        -\gls{sym-b-mu}&
    &\in \gls{sym-mathbb-R}
\end{aligned}
\end{equation}
Since \Cref{eq:obj-gamma} shows that power decreases with increasing $|\gls{sym-Gamma}|$, the power maximization problem \Cref{eq:opt-problem} can be recast as minimizing the norm of $\gls{sym-Gamma}$:
\begin{subequations}\label{eq:opt-problem-gamma}
\begin{alignat}{2}
& \min_{\gls{sym-Gamma} \in \gls{sym-mathbb-C}} &  |\gls{sym-Gamma}|^2& \\[3pt]
& \text{s.t.} &
\gls{sym-S-mu}|\gls{sym-Gamma} - \gls{sym-Gamma-c-mu}|^2 &\le S_{\mu} r_{\mu}^2,
\quad \forall \mu \label{eq:disc-constraint}\\
& \text{where} &
\gls{sym-Gamma-c-mu}  &\in \gls{sym-mathbb-C} \nonumber\\
& & r_{\mu}^2 &\in \gls{sym-mathbb-R} \nonumber \\
& & \gls{sym-S-mu} &\in \left\{-1,1\right\} \nonumber
\end{alignat}
\end{subequations}
where the constraint \Cref{eq:disc-constraint} is derived from \Cref{eq:ineq-gamma} by completing the square.
More precisely, when $\gls{sym-A-q-mu} \neq 0$, the $\gls{sym-mu}$th inequality constraint defines a circle in the complex plane of $\gls{sym-Gamma}$ centered at $\gls{sym-Gamma-c-mu}$ with radius $\gls{sym-r-mu}$.
The constraint requires that $\gls{sym-Gamma}$ lie within this circle when $\gls{sym-A-q-mu} > 0$ ($\gls{sym-S}_{\gls{sym-mu}}=1$) or outside this circle when $\gls{sym-A-q-mu} < 0$ ($\gls{sym-S}_{\gls{sym-mu}}=-1$).
Parameters $\gls{sym-Gamma-c-mu}$, $\gls{sym-r-mu}$, and $\gls{sym-S-mu}$ are defined as:
\begin{subequations}\label{eq:gamma-disc-params}
\begin{align}
\gls{sym-Gamma-c-mu}
&= \frac{-\gls{sym-B-q-mu}^{\,*}}{\gls{sym-A-q-mu}} =
\frac{
    \gls{sym-hat-I-p}^{\,*}\,
    \gls{sym-vec-n}^{\,*}
    \left(\gls{sym-mathbf-Q-mu}\gls{sym-vec-x-p} + \gls{sym-vec-a-mu}\right)
}{
    |\gls{sym-hat-I-p}|^{2}\,
    \gls{sym-vec-n}^{\,*}\gls{sym-mathbf-Q-mu}\gls{sym-vec-n}
}
\label{eq:gamma-disc-center}
\\[6pt]
\gls{sym-r-mu}^2
&=
\frac{|\gls{sym-B-q-mu}|^2-\gls{sym-A-q-mu}\gls{sym-C-q-mu}}{\gls{sym-A-q-mu}^2} =
\frac{
\gls{sym-b-mu}
-
\gls{sym-vec-x-p}^{\,*}\gls{sym-mathbf-Q-mu}\gls{sym-vec-x-p}
-
2\Re\!\left\{\gls{sym-vec-a-mu}^{\,*}\gls{sym-vec-x-p}\right\}
}{
|\gls{sym-hat-I-p}|^{2}\,
\gls{sym-vec-n}^{\,*}\gls{sym-mathbf-Q-mu}\gls{sym-vec-n}
}
+
|\gls{sym-Gamma-c-mu}|^2
\label{eq:gamma-disc-radius}
\\
\gls{sym-S-mu} &= \textrm{sign} (\gls{sym-A-q-mu}) = \textrm{sign}\left(\gls{sym-vec-n}^{\,*}\gls{sym-mathbf-Q-mu}\gls{sym-vec-n}\right)
\label{eq:gamma-disc-sign}
\end{align}
\end{subequations}

%%%%%%%%%%%%%%%
\paragraph{Solution}
\Cref{eq:opt-problem-gamma} can be interpreted graphically as finding the point closest to the origin that is interior to all constraint circles for which $\gls{sym-S-mu}=1$ and exterior to all circles for which $\gls{sym-S-mu}=-1$.
The unconstrained optimal solution of \Cref{eq:opt-problem-gamma} is the origin ($\gls{sym-Gamma}=0$), and the constrained optimal solution will also be zero in the case where the origin is feasible (no active constraints), which occurs when $ \gls{sym-S-mu} \gls{sym-C-q-mu} \leq 0 \quad \forall \gls{sym-mu}$.
If the origin is \textrm{infeasible}, the constrained optimal solution will lie on the boundary of one or more of the circles, specifically either at an intersection of two or more circles or at any circle's point closest to the origin.
We analytically compute all such candidate points and select the one with the lowest $|\gls{sym-Gamma}|$ that satisfies all constraints.
For $\gls{sym-N}$ constraints, there are up to $\gls{sym-N}^2$ candidate points ($\gls{sym-N}^2-\gls{sym-N}$ at intersections and $\gls{sym-N}$ at closest points).
This leads to the following expression for the constrained optimal solution (valid when $\gls{sym-A-q-mu} \neq 0 ~\forall \gls{sym-mu}$):
\begin{equation}\label{eq:gamma-opt-quadratic}
\gls{sym-Gamma-opt} =
\begin{cases}
 0 & \textrm{if } \gls{sym-S-mu} \gls{sym-C-q-mu} \leq 0 \quad \forall \gls{sym-mu} \\[3pt]
 \arg\min\limits_{\gls{sym-Gamma} \in \gls{sym-mathcal-F} \bigcap \gls{sym-mathcal-C}} |\gls{sym-Gamma}|
 & \textrm{otherwise} \\[3pt]
 \textrm{infeasible} & \textrm{if } \gls{sym-mathcal-F} \bigcap \gls{sym-mathcal-C} = \emptyset
\end{cases}
\end{equation}
where $\gls{sym-mathcal-F} = \bigcap_{\gls{sym-mu}} \left\{\gls{sym-Gamma}: \gls{sym-S}_{\gls{sym-mu}}|\gls{sym-Gamma} - \gls{sym-Gamma-c-mu}| \le \gls{sym-S}_{\gls{sym-mu}} \gls{sym-r-mu} \right\}$ is the set of all feasible points and $\gls{sym-mathcal-C}$ is the set of candidate points. $\gls{sym-mathcal-C}$ is the union of $\gls{sym-mathcal-C-int}$, the candidates at circle intersections, and $\gls{sym-mathcal-C-cl}$, the candidates at the closest point on each circle to the origin:
\begin{equation}\label{eq:candidate-points}
\begin{aligned}
\gls{sym-mathcal-C-int} &= \left\{
    \gls{sym-Gamma-mu-nu}, \quad \forall \gls{sym-mu} \neq \gls{sym-nu}
\right\}.
\\
\gls{sym-mathcal-C-cl} &= \left\{
    (|\gls{sym-Gamma-c-mu}| - \gls{sym-r-mu})\exp (i\angle\,\gls{sym-Gamma-c-mu}), \quad \forall \gls{sym-mu}
\right\}.
\end{aligned}
\end{equation}
The intersection points of the $\gls{sym-mu}$th and $\gls{sym-nu}$th circles are given by:
\begin{equation}\label{eq:circle-intersection}
    \gls{sym-Gamma-mu-nu} =
    \gls{sym-Gamma-c-mu} + \frac{\gls{sym-Gamma-c-nu} - \gls{sym-Gamma-c-mu}}{\gls{sym-d-mu-nu}}
    \left(        \gls{sym-a-mu-nu}
        \pm
        1i \sqrt{\gls{sym-r-mu}^2-\gls{sym-a-mu-nu}^2}
    \right)
\end{equation}
where $\gls{sym-d-mu-nu} = |\gls{sym-Gamma-c-nu} - \gls{sym-Gamma-c-mu}|$ is the distance between the circle centers,
$\gls{sym-a-mu-nu} = \frac{1}{2}(\gls{sym-r-mu}^2-\gls{sym-r-nu}^2+\gls{sym-d-mu-nu}^2)/\gls{sym-d-mu-nu}$ is the distance from
the center of the $\gls{sym-mu}$th circle to the midpoint of the two intersection points,
and $1i$ is the imaginary unit.
The closest-point candidates $\gls{sym-mathcal-C-cl}$ match the solution that \citet[\equationautorefname~4]{mccabe_force-limited_2024} derive for problems with only one constraint.

Importantly, we need not compute the entire feasible space $\gls{sym-mathcal-F}$, only compute $\gls{sym-mathcal-C}$ and then check whether each candidate point is feasible.
\Cref{fig:qp-solution-geometry} illustrates the various cases for feasibility and optimality for a toy problem with $\gls{sym-r-mu}=4$, $\gls{sym-S-mu}=1$, and various $\gls{sym-Gamma-c-mu}$.
\begin{figure}[htbp]
    \centering
    \includegraphics[\FigureOptions{aor:figs/from-matlab/qp_circles.pdf}]{figs/from-matlab/qp_circles.pdf}
    \caption[Visualization of the optimization problem in the complex plane of $\gls{sym-Gamma}$]{Visualization of the optimization problem \Cref{eq:opt-problem-gamma} in the complex plane of $\gls{sym-Gamma}$.
Toy problems with various constraint circle centers $\gls{sym-Gamma-c-mu}$, all with radius $\gls{sym-r-mu}=4$ and $\gls{sym-S-mu}=1$, demonstrate the various solution cases in \Cref{eq:gamma-opt-quadratic}.}
    \label{fig:qp-solution-geometry}
\end{figure}

If $\gls{sym-A-q-mu} = 0$ for any $\gls{sym-mu}$, the $\gls{sym-mu}$th inequality constraint defines a line in the complex plane of $\gls{sym-Gamma}$ rather than a circle.
This can occur when $\gls{sym-mathbf-Q-mu} = 0$ (the constraint on $\gls{sym-vec-x}$ is linear rather than quadratic, which is not the case for any of the constraints in \Cref{tab:qp-constraints}) or when $\gls{sym-vec-n}^{\, *}\gls{sym-mathbf-Q-mu}\gls{sym-vec-n} = 0$ ($\gls{sym-mathbf-Q-mu}$ is rank-deficient and $\gls{sym-vec-n}$ is in its nullspace).
The same types of critical points (now consisting of line-line, circle-line, and circle-circle intersections, as well as the closest point on each circle and line) apply, but for implementation simplicity, we replace $\gls{sym-A-q-mu}$ with a small positive $\gls{sym-epsilon}$ in this case and use the quadratic implementation.

%%%%%%%%%%%%%%%
\paragraph{Quartic Apparent Power Constraints}
\label{sec:appendix-mag-S-constraints}
While constraints on effort and flow magnitudes and real and reactive power are quadratic in $\gls{sym-vec-x}$, constraints on the magnitude of apparent power $|\gls{sym-S}|<|\gls{sym-S}_{\text{max}}|$, and therefore the maximum and minimum instantaneous power, are fourth-order in $\gls{sym-vec-x}$ and thus $\gls{sym-Gamma}$.
This results in a ``Cassini oval'' curve \citep{weisstein_cassini_2026} on the complex plane of $\gls{sym-Gamma}$, which consists of two closed curves when $|\gls{sym-S}_{\text{max}}|<|\gls{sym-S-0}|$ and one curve when $|\gls{sym-S}_{\text{max}}|\geq|\gls{sym-S-0}|$, where $|\gls{sym-S-0}|$ is the apparent power of the unconstrained optimum.
Though the optimization is no longer a quadratically constrained quadratic program, the same geometric solution procedure applies, adding circle-oval intersections $\gls{sym-mathcal-C-int-oval-circ}$, oval-oval intersections, and oval closest points $\gls{sym-mathcal-C-cl-oval}$ to the set of candidate points $\gls{sym-mathcal-C}$.
The new candidate points are:
\begin{equation}
\begin{aligned}
\gls{sym-mathcal-C-int-oval-circ} &= \left\{
    \gls{sym-Gamma}:
    |\gls{sym-Gamma} - \gls{sym-Gamma-c-mu}| = \gls{sym-r-mu}, ~(1-|\gls{sym-Gamma}|^2)^2 + 4\Im (\gls{sym-Gamma}) = \frac{|\gls{sym-S-max-k}|}{|\gls{sym-S-0}|}, \quad \forall \gls{sym-mu},\gls{sym-k-generic-index}
\right\}.
\\
\gls{sym-mathcal-C-cl-oval} &= \left\{
    \gls{sym-Gamma}:
    \gls{sym-Gamma} = \begin{cases}
        \pm \sqrt{1-\frac{|\gls{sym-S-max-k}|}{|\gls{sym-S-0}|}}, \quad \forall \gls{sym-k-generic-index}~ |\gls{sym-S-max-k}|<|\gls{sym-S-0}| \\
        \pm \sqrt{1+\frac{|\gls{sym-S-max-k}|}{|\gls{sym-S-0}|}}, \quad \forall \gls{sym-k-generic-index}~ |\gls{sym-S-max-k}|>|\gls{sym-S-0}|
    \end{cases}
\right\}.
\end{aligned}
\end{equation}
where $\gls{sym-mu}$ indexes circle constraints and $\gls{sym-k-generic-index}$ indexes oval constraints.
While the earlier circle-circle intersections of \Cref{eq:candidate-points} result in a quadratic equation that is solved analytically for the coordinates of $\gls{sym-Gamma}$, oval-circle intersections lead to quartic equations that are solved numerically with standard polynomial root-finding routines.
Typically there is only one apparent power constraint, so oval-oval intersections are not derived, but if present would likewise yield a fourth-order polynomial in $\gls{sym-Gamma}$.
Therefore, the presence of apparent power constraints makes the optimal control method semi-analytical rather than analytical.

%%%%%%%%%%%%%%%
\paragraph{Extension to Higher Dimensions and Convexity}
\label{sec:appendix-qp-convexity}
Future work may wish to extend the constrained power maximization to the full nonlinear dynamics, or to linear dynamics with irregular wave forcing as in the study by \citet{bacelli_numerical_2015}.
In the terminology of \Cref{fig:venn-diagram}, these would be known as nonlinear and linear spectral optimal control respectively.
Both of these cases solve the full harmonic-balance equation, producing a higher-dimensional analog of \Cref{eq:opt-problem-gamma}.
The problem dimension equals twice the number of frequencies retained in the spectral expansion, preventing the use of the geometric analytical solution used in the monochromatic case.
Instead, a numerical solver must be used, and the computational scaling of that solver depends critically on whether the problem is \emph{convex}.

A convex optimization problem is one in which both the feasible set and the objective function are convex---informally, the feasible region has no holes or concavities, and the objective is bowl-shaped with no local valleys.
The key practical implication is that any locally optimal solution is also globally optimal, so gradient-based solvers converge to the true optimum without exhaustive search, and solution time scales polynomially with the number of decision variables rather than exponentially.

After reducing to the $\gls{sym-Gamma}$ parameterization, the power objective $|\gls{sym-Gamma}|^2$ is a sum of squares and therefore convex.
Even though the original power matrix $\gls{sym-mathbf-Q-P}$ (see \Cref{tab:power-metrics}) has eigenvalues $\{-1,1\}$ and is indefinite, \citet{bacelli_numerical_2015} show that substituting the linear dynamics equality constraint into the objective---which is exactly what the reduction to $\gls{sym-Gamma}$ in \Cref{eq:obj-gamma} accomplishes---renders the objective convex for both regular and irregular wave forcing.
Nonlinear dynamics can reintroduce non-convexity and should be checked on a case-by-case basis.

The inequality constraints are more problematic.
A constraint is convex if the set of points satisfying it forms a convex region.
In the $\gls{sym-Gamma}$ plane, a convex constraint corresponds to a circle whose feasible region is the \emph{interior} ($\gls{sym-S-mu}=1$ in \Cref{eq:disc-constraint}), while a non-convex constraint corresponds to a circle whose feasible region is the \emph{exterior} ($\gls{sym-S-mu}=-1$).
Constraints on the effort and flow magnitude (e.g., force limits $|\gls{sym-hat-F-PTO} |\leq \gls{sym-F-max}$) can be convex or non-convex, depending on the complex angle of the Th\'{e}venin equivalent of the relevant source impedance and on the ratio of the constraint limit to its unconstrained value.
\citet{mccabe_force-limited_2024} visualize the feasible regions for families of such constraints on the complex $\gls{sym-Gamma}$ plane.
They show that the constraints are always convex for purely real impedances; that for reactive plants, nonconvexities occur when the value of the constraint limit exceeds some threshold; and that this threshold lowers as the reactivity increases.
The quadratic matrices $\gls{sym-mathbf-Q-mu}$ in \Cref{tab:qp-constraints} confirm this: $\gls{sym-mathbf-Q-hat-q-2}$ and $\gls{sym-mathbf-Q-hat-e-2}$ are only positive \emph{semi}definite, indicating constraints that define half-spaces rather than bounded ellipsoids.
For the 2D problem studied here, non-convexity is handled by the exhaustive enumeration of candidate points in \Cref{eq:gamma-opt-quadratic}, which is efficient in two dimensions but does not scale to higher dimensions.
Future work should investigate semidefinite relaxations of these non-convex constraints using the S-procedure \citep{vanantwerp_tutorial_2000}, which replaces each non-convex constraint with a conservative but convex outer approximation.
It is expected that the optimal design will have low reactivity because reactivity increases signal peak-to-average ratios, implying that a semi-definite relaxation would likely improve in accuracy as it approaches the optimal.

%%%%%%%%%%%%%%%
\paragraph{Comparison to Other \gls{acr-wec} Constrained Control Formulations}
Notably, $\gls{sym-Gamma}$ is equal to the reflection coefficient at the generator electrical port, which is commonly used in microwave circuit design to quantify impedance mismatch.
\citet{mccabe_force-limited_2024,coe_co-design_2025} further discuss the reflection coefficient for \glspl{acr-wec}.

Similar graphical interpretations of constraints for wave energy converters are presented by \citet{mccabe_force-limited_2024,merigaud_geometrical_2023,bacelli_geometric_2013}, with all three visualizing quadratic constraints as circles on the 2D plane.
\citet{mccabe_force-limited_2024} focus on force/current limits but discuss amplitude, phase voltage, and apparent power constraints as well.
\citet{merigaud_geometrical_2023} consider exclusively amplitude limits while \citet{bacelli_geometric_2013} cover amplitude and force limits.
\citet{mccabe_force-limited_2024} use Smith charts to show constraints directly in $\gls{sym-Gamma}$ space as we do here, while \citet{merigaud_geometrical_2023} work in the complex plane of the hydrodynamic transmission coefficient and \citet{bacelli_geometric_2013} in that of the \gls{acr-pto} force.
\citet{bacelli_geometric_2013} also extend the formulation to capture higher harmonics (nonsinusoidal forces), where the exact constraint cannot be easily represented geometrically but sufficient conditions for constraint violation and satisfaction in the higher dimensional space can be represented as hyperspheres and hyperellipsoids.

%%%%%%%%%%%%%%%%%%%%%%%%%%%%%%%%%%%%%%%%%

\subsubsection{Constraint Sensitivity Analysis via Parametric Programming}
\label{sec:appendix-constraint-sensitivity}
To support the discussion in \Cref{sec:discussion}, this section derives the scaling of optimal average power with the constraint limits and demonstrates its conditions for convexity.

The constraint limits are
$\gls{sym-tau}_{\text{max}}$, $\gls{sym-P}_{\gls{sym-tau}\gls{sym-Omega},\text{max}}$, $\gls{sym-xi-f-max}$, $\gls{sym-xi-s-max}$, and $\gls{sym-X-PTO-max}$, which we collectively denote as $\gls{sym-L-mu}$.
The dependence of power on the first two limits (\gls{acr-pto} sizing variables) is of more immediate relevance than the last three (amplitude limits) because the former can be independently manipulated through \gls{acr-pto} design, while the latter can only be manipulated by changing the bulk dimensions.
The sensitivity of power to bulk dimensions is substantially more complex than the sensitivity to constraint limits for reasons discussed in \Cref{sec:appendix-dimension-sensitivity}, so we focus on the \gls{acr-pto} sizing sensitivities here.

In a given sea state, the global sensitivity of optimal power to a constraint limit can be quantified using parametric programming theory, where a local sensitivity is found parametrically and then integrated over the parameter space to find the global parameter relationship.
The local sensitivity of power to a constraint parameter
$\frac{\partial \gls{sym-p}_{\text{avg},VI,\text{opt}}}{\partial \gls{sym-L-mu}}$
equals the product of the Lagrange multiplier for that constraint, $\gls{sym-lambda-mu}$,
and the partial derivative of the constraint function with respect to the parameter,
$\frac{\partial \gls{sym-g-mu}}{\partial \gls{sym-L-mu}}$ \cite[Section~5.9.3]{boyd_convex_2004}.
Since the constraint function is defined as
$\gls{sym-g-mu} (\gls{sym-vec-x}) = \gls{sym-vec-x}^{\,*}\gls{sym-mathbf-Q-mu}\gls{sym-vec-x}
                 + 2\Re\!\left\{\gls{sym-vec-a-mu}^{\,*}\gls{sym-vec-x}\right\} - \gls{sym-b-mu}\leq 0$,
and \Cref{tab:qp-constraints} shows that each constraint limit $\gls{sym-L-mu}$ enters $\gls{sym-g-mu}$ only through $\gls{sym-b-mu}$, we have
$\frac{\partial \gls{sym-g-mu}}{\partial \gls{sym-L-mu}} = -\frac{\partial \gls{sym-b-mu}}{\partial \gls{sym-L-mu}}$.
This derivative can be trivially computed for all constraints in \Cref{tab:qp-constraints}.
The power sensitivity simplifies to:
\begin{equation}
    \frac{\partial \gls{sym-p-avg-VI-opt}}{\partial \gls{sym-L-mu}} = -\gls{sym-lambda-mu} \frac{\partial \gls{sym-b-mu}}{\partial \gls{sym-L-mu}}
\end{equation}
If the optimization problem were solved with a numerical solver,
the Lagrange multipliers would be directly available as part of the solution.
In our case, we have a semi-analytical solution for the optimal reflection coefficient at each sea state in terms of the constraint centers and radii at that sea state,
$\gls{sym-Gamma-opt}(\gls{sym-r-mu}, \gls{sym-Gamma-c-mu})$,
so we can compute the Lagrange multipliers as follows:
\begin{equation}\label{eq:lagrange}
\begin{aligned}
   \gls{sym-lambda-mu}
    &\equiv
    \frac{\partial \gls{sym-p-avg-VI-opt}}{\partial \gls{sym-g-mu}}
    \\
    &= - \frac{\partial \gls{sym-p-avg-VI-opt}}{\partial \gls{sym-b-mu}}
    \\
    &= -
    \frac{\partial \gls{sym-p-avg-VI-opt}}{\partial |\gls{sym-Gamma-opt}|}
    \frac{\partial |\gls{sym-Gamma-opt}|}{\partial \gls{sym-b-mu}}
    \\
    &= -
    \frac{\partial \gls{sym-p-avg-VI-opt}}{\partial |\gls{sym-Gamma-opt}|}
    \left(        \frac{\partial |\gls{sym-Gamma-opt}|}{\partial \gls{sym-r-mu}}
        \frac{\partial \gls{sym-r-mu}}{\partial \gls{sym-b-mu}}
        +
        \frac{\partial |\gls{sym-Gamma-opt}|}{\partial \gls{sym-Gamma-c-mu}}
        \frac{\partial \gls{sym-Gamma-c-mu}}{\partial \gls{sym-b-mu}}
    \right)
\end{aligned}
\end{equation}
The computation of each of the five partials in the final expression of \Cref{eq:lagrange} is discussed next.

%%%%%%%%%%%%%%%
\paragraph{Effect of Constraint Limits on Circle Radii}
Recall from \Cref{eq:gamma-disc-radius} that each inequality constraint maps to a circle in the complex $\gls{sym-Gamma}$ plane
with center $\gls{sym-Gamma-c-mu}$ and radius $\gls{sym-r-mu}$ given by:
\begin{equation}\label{eq:radius-full}
    \gls{sym-r-mu}^2
    = \frac{\gls{sym-b-mu} - \gls{sym-vec-x-p}^{\,*}\gls{sym-mathbf-Q-mu}\gls{sym-vec-x-p}
            - 2\Re\!\left\{\gls{sym-vec-a-mu}^{\,*}\gls{sym-vec-x-p}\right\}}
           {|\gls{sym-hat-I-p}|^{2}\,\gls{sym-vec-n}^{\,*}\gls{sym-mathbf-Q-mu}\gls{sym-vec-n}}
    + |\gls{sym-Gamma-c-mu}|^2
\end{equation}
The constraint limit $\gls{sym-L-mu}$ enters $\gls{sym-r-mu}$ only through $\gls{sym-b-mu}$
 (the right-hand side of the quadratic inequality in \Cref{tab:qp-constraints}),
while $\gls{sym-vec-x-p}$, $\gls{sym-vec-n}$, $\gls{sym-hat-I-p}$, $\gls{sym-mathbf-Q-mu}$, $\gls{sym-vec-a-mu}$, and $\gls{sym-Gamma-c-mu}$ depend only on the Th\'{e}venin impedance $\gls{sym-Z-s-th}$ and excitation $\gls{sym-hat-V-s-th}$ for the sea state,
which are independent of the constraint limits.
Defining $\gls{sym-d-mu} = |\gls{sym-hat-I-p}|^{2}\,\gls{sym-vec-n}^{\,*}\gls{sym-mathbf-Q-mu}\gls{sym-vec-n}$
and $\gls{sym-c-mu} = \gls{sym-d-mu}|\gls{sym-Gamma-c-mu}|^2 - \gls{sym-vec-x-p}^{\,*}\gls{sym-mathbf-Q-mu}\gls{sym-vec-x-p} - 2\Re\!\left\{\gls{sym-vec-a-mu}^{\,*}\gls{sym-vec-x-p}\right\}$,
both independent of the constraint limits, \Cref{eq:radius-full} and its derivative simplify to:
\begin{equation}\label{eq:radius-simplified}
    \gls{sym-r-mu}^2 = \frac{\gls{sym-b-mu} + \gls{sym-c-mu}}{\gls{sym-d-mu}}
, \quad
    \frac{\partial \gls{sym-r-mu}}{\partial \gls{sym-b-mu}} = \frac{1}{2\sqrt{\gls{sym-d-mu} (\gls{sym-b-mu} + \gls{sym-c-mu})}}
\end{equation}
Since $\gls{sym-r-mu}\in\gls{sym-mathbb-R}$ (circle radii cannot have an imaginary part), then $\gls{sym-b-mu} + \gls{sym-c-mu}$ and $\gls{sym-d-mu}$ must have the same sign, so $\frac{\partial \gls{sym-r-mu}}{\partial \gls{sym-b-mu}}$ is always positive.
In other words, increasing $\gls{sym-b-mu}$ always enlarges the circle radius, regardless of whether the constraint represents the interior or exterior of the circle.
Meanwhile, \Cref{eq:gamma-disc-center} shows that the center $\gls{sym-Gamma-c-mu}$ does not depend on $\gls{sym-b-mu}$, so $\frac{\partial \gls{sym-Gamma-c-mu}}{\partial \gls{sym-b-mu}}=0$.
Thus, the constraint limit $\gls{sym-L-mu}$ does not affect the circle center.

%%%%%%%%%%%%%%%
\paragraph{Optimal Reflection Coefficient as a Function of Constraint Limits}
Now we combine the relationship between circle radius and optimal reflection coefficient $\gls{sym-Gamma-opt}(\gls{sym-r-mu})$, with the relationship between circle radius and constraint limit $\gls{sym-r-mu}(\gls{sym-b-mu})$ derived above
to obtain a relationship between reflection coefficient magnitude and constraint limit in a given sea state, $|\gls{sym-Gamma-opt}|(\gls{sym-b-mu})$.
Substituting \Cref{eq:radius-simplified} into \Cref{eq:candidate-points} shows that
the squared optimal reflection coefficient magnitude in each sea state takes one of the following values depending on the number of active constraints:
\begin{equation}\label{eq:gamma-opt-summary}
    |\gls{sym-Gamma-opt}|^2 =
    \begin{cases}
        0
    & \text{none active} \\
        |\gls{sym-Gamma-c-mu-star}|^2
        - 2|\gls{sym-Gamma-c-mu-star}|\sqrt{\frac{\gls{sym-b-mu-star}+\gls{sym-c-mu-star}}{\gls{sym-d-mu-star}}}
        + \frac{\gls{sym-b-mu-star}+\gls{sym-c-mu-star}}{\gls{sym-d-mu-star}}
    & \text{only $\gls{sym-mu}^*$ active} \\
        \left|
            \gls{sym-Gamma-c-mu-star} + \frac{\gls{sym-Gamma-c-nu-star} - \gls{sym-Gamma-c-mu-star}}{\gls{sym-d-mu-star-nu-star}}
            \left( \gls{sym-a-mu-star-nu-star}
                \pm
                1i \sqrt{
                    \frac{\gls{sym-b-mu-star}
                    + \gls{sym-c-mu-star}}{\gls{sym-d-mu-star}}
                    -\gls{sym-a-mu-star-nu-star}^2
                    }
            \right)
        \right|^2
    & \text{$\gls{sym-mu}^*$ \& $\gls{sym-nu}^*$ active}
    \end{cases}
\end{equation}
where $\gls{sym-mu}^*,\gls{sym-nu}^*$ refer to the indices of the active constraint(s) for a given sea state,
$\gls{sym-a-mu-star-nu-star}$ and $\gls{sym-d-mu-star-nu-star}$ are real numbers defined in the discussion following \Cref{eq:circle-intersection}.
$\gls{sym-d-mu-star-nu-star}$ is independent of the constraint limits, while $\gls{sym-a-mu-star-nu-star}$ is affine in $\gls{sym-b-mu-star}$ and $\gls{sym-b-nu-star}$.
Thus, using \textrm{aff}() and \textrm{quad}() to denote affine and quadratic functions of the arguments, we have the following scale behavior:
\begin{equation}\label{eq:gamma-opt-scale}
    |\gls{sym-Gamma-opt}|^2 =
    \begin{cases}
        0
    & \text{none active} \\
        \textrm{aff} (\gls{sym-b-mu-star}) + \sqrt{\textrm{aff} (\gls{sym-b-mu-star})}
    & \text{only $\gls{sym-mu}^*$ active} \\
        \left|
            \textrm{aff} (\gls{sym-b-mu-star})+\textrm{aff} (\gls{sym-b-nu-star})
            \pm
            1i \sqrt{\textrm{quad} (\gls{sym-b-mu-star}, \gls{sym-b-nu-star})}
        \right|^2
    & \text{$\gls{sym-mu}^*$ \& $\gls{sym-nu}^*$ active} \\
    = \textrm{quad} (\gls{sym-b-mu-star}, \gls{sym-b-nu-star}) +
    (\textrm{aff} (\gls{sym-b-mu-star}) + \textrm{aff} (\gls{sym-b-nu-star}))
    \sqrt{\textrm{quad} (\gls{sym-b-mu-star}, \gls{sym-b-nu-star})} &
    \end{cases}
\end{equation}
Constraint activity indices $\gls{sym-mu}^*,\gls{sym-nu}^*$ depend on the constraint limits,
and this dependence must be accounted for before the curvature of $|\gls{sym-Gamma-opt}|^2$ can be analyzed.
From \Cref{eq:gamma-opt-quadratic}, we can define $\gls{sym-mu}^*$ and $\gls{sym-nu}^*$ more explicitly as functions of $\gls{sym-b-mu}$ by defining a penalty function $\gls{sym-f-penalty}(\gls{sym-Gamma})$ that returns $\gls{sym-Gamma}$ if $\gls{sym-Gamma}$ is feasible and $\infty$ \textrm{otherwise}.
This function can be constructed as $\gls{sym-f-penalty}(\gls{sym-Gamma}) = \gls{sym-Gamma} + \sum_{\gls{sym-mu}} \gls{sym-mathbb-I-mu}(\gls{sym-Gamma})$,
where $\gls{sym-mathbb-I-mu}(\gls{sym-Gamma})$ is an indicator function that returns $0$ if $\gls{sym-Gamma}$ satisfies the $\gls{sym-mu}$th constraint and $\infty$ \textrm{otherwise}.
The penalty function is convex in $\gls{sym-Gamma}$ if and only if all constraints are convex, i.e., if $\gls{sym-S-mu}=1 ~\forall \gls{sym-mu}$.
Then, the optimal reflection coefficient of each candidate type can be written as the minimum of the candidate points passed through the penalty function.
\begin{equation}\label{eq:gamma-opt-min-per-candidate}
\begin{aligned}
   |\gls{sym-Gamma-cl-mu-star}| &=
    \min_{\gls{sym-mu}}
    \left\{
        \left| \gls{sym-f-penalty}\left(
                                        \gls{sym-Gamma-cl-mu} (\gls{sym-b-mu})
                                  \right)
        \right|
    \right\},
    \quad
    |\gls{sym-Gamma-int-mu-star-nu-star}| =
    \min_{\gls{sym-mu},\gls{sym-nu}}
    \left\{
        \left| \gls{sym-f-penalty}\left(
                                        \gls{sym-Gamma-mu-nu}(\gls{sym-b-mu}, \gls{sym-b-nu})
                                  \right)
        \right|
    \right\}
\end{aligned}
\end{equation}
Finally, the overall optimal reflection coefficient $|\gls{sym-Gamma-opt}|$ is the minimum of all the per-candidate-type minima:
\begin{equation}\label{eq:gamma-opt-min}
\begin{aligned}
    |\gls{sym-Gamma-opt}| &=
    \min \left\{
        |\gls{sym-Gamma-cl-mu-star} (\gls{sym-b-mu}) |, 
        |\gls{sym-Gamma-int-mu-star-nu-star} (\gls{sym-b-mu}, \gls{sym-b-nu})|
    \right\}
\end{aligned}
\end{equation}

Working from the inside out, we now analyze the convexity of $|\gls{sym-Gamma-opt}|^2$ with respect to $\gls{sym-b-mu}$.
Starting with the candidates at closest points, the square root is a concave increasing function,
and a concave increasing function of an affine function is concave,
so $\sqrt{\textrm{aff}(\gls{sym-b-mu})}$ is concave.
The sum of a concave function and an affine function is concave,
so $|\gls{sym-Gamma-cl}|=\textrm{aff}(\gls{sym-b-mu}) + \sqrt{\textrm{aff}(\gls{sym-b-mu})}$ is concave.
The penalty function is convex if and only if all constraints are convex,
so $-|\gls{sym-f-penalty}(\gls{sym-Gamma-cl})|$
is concave if and only if
$\gls{sym-S-mu}=1 ~\forall \gls{sym-mu}$.
The minimum of a set of concave functions is a concave increasing function,
% or the maximum of a set of convex functions is a convex increasing function
and $|\gls{sym-Gamma}_{cl,\gls{sym-mu}^*}|$
is the minimum of a set of concave functions if and only if all constraints are concave,
so $|\gls{sym-Gamma}_{cl,\gls{sym-mu}^*}|$
is concave if and only if
$\gls{sym-S-mu}=1 ~\forall \gls{sym-mu}$.

Moving onto the candidates at intersection points, a quadratic function is concave if its leading order term is negative, and the leading order term of $\textrm{quad}(\gls{sym-b-mu}, \gls{sym-b-nu})$ is negative (arising from the subtraction of squared-affine term $\gls{sym-a-mu-star-nu-star}^2$).
%Square root is a concave function, and a concave function of a concave function is concave, so $\sqrt{\textrm{quad}(\gls{sym-b-mu}, \gls{sym-b-nu})}$ is concave.
The geometric mean of two positive values is a concave function, so the geometric mean of two positive concave functions is concave.
The term $\left(\textrm{aff}(\gls{sym-b-mu}) + \textrm{aff}(\gls{sym-b-nu})\right) \sqrt{\textrm{quad}(\gls{sym-b-mu}, \gls{sym-b-nu})}$
is the geometric mean of $(\textrm{aff}(\gls{sym-b-mu})+\textrm{aff}(\gls{sym-b-nu}))^2$ and $\textrm{quad}(\gls{sym-b-mu}, \gls{sym-b-nu})$.
The $(\textrm{aff}(\gls{sym-b-mu})+\textrm{aff}(\gls{sym-b-nu}))^2$ is positive but is convex rather than concave.
This means the curvature of $\left(\textrm{aff}(\gls{sym-b-mu}) + \textrm{aff}(\gls{sym-b-nu})\right) \sqrt{\textrm{quad}(\gls{sym-b-mu}, \gls{sym-b-nu})}$ cannot be established with the present analysis.
If the curvature of $\left(\textrm{aff}(\gls{sym-b-mu}) + \textrm{aff}(\gls{sym-b-nu})\right) \sqrt{\textrm{quad}(\gls{sym-b-mu}, \gls{sym-b-nu})}$ could be established as concave,
then the sum of two concave functions is concave, so the candidate points $|\gls{sym-Gamma-int}|$ would be concave.
If that were the case, then by the same argument as for the closest points, the minimum of feasible $|\gls{sym-Gamma-int}|$ would be concave if and only if all constraints are concave.
The optimum $|\gls{sym-Gamma-opt}|$ is the minimum of the feasible $|\gls{sym-Gamma-cl}|$ and $|\gls{sym-Gamma-int}|$, so if both of those are concave, then it is also concave.
Therefore, the criteria for the concavity of $|\gls{sym-Gamma-opt}|$ with respect to $\gls{sym-b-mu}$ is that all constraints are convex (interior of circles), and the curvature of $\left(\textrm{aff}(\gls{sym-b-mu}) + \textrm{aff}(\gls{sym-b-nu})\right) \sqrt{\textrm{quad}(\gls{sym-b-mu}, \gls{sym-b-nu})}$ is concave.
As future work, the above curvature analysis should be verified with a disciplined convex programming package such as CVXPY.

To facilitate the analysis of the next section,
we observe that although the squared optimal reflection coefficient magnitude, $|\gls{sym-Gamma-opt}|^2$,
has so far been expressed as a function only of the active constraints $(\gls{sym-mu}^*,\gls{sym-nu}^*)$,
it can be written as a sum over all constraint pairs $(\gls{sym-mu},\gls{sym-nu})$ multiplied by an appropriate Kronecker delta function:
\begin{equation}
\label{eq:gamma-opt-scale-sum}
\begin{aligned}
    |\gls{sym-Gamma-opt}|^2
    = &\sum_{\gls{sym-mu}} \sum_{\gls{sym-nu} > \gls{sym-mu}}
    \left[
        \gls{sym-delta}_{\gls{sym-mu},\gls{sym-mu}^*}(1-\gls{sym-delta}_{\gls{sym-nu},\gls{sym-nu}^*})
        \left(
                    \textrm{aff} (\gls{sym-b-mu})
            + \sqrt{\textrm{aff} (\gls{sym-b-mu})}
        \right)
     + \gls{sym-delta}_{\gls{sym-mu},\gls{sym-mu}^*}\gls{sym-delta}_{\gls{sym-nu},\gls{sym-nu}^*}
    \right.
\\
    &\left.
        \left(            \textrm{quad} (\gls{sym-b-mu}, \gls{sym-b-nu}) +
            \left(\textrm{aff} (\gls{sym-b-mu}) + \textrm{aff} (\gls{sym-b-nu})\right)
            \sqrt{\textrm{quad} (\gls{sym-b-mu}, \gls{sym-b-nu})}
        \right)
    \right]
\end{aligned}
\end{equation}
where $\gls{sym-delta}_{\gls{sym-mu},\gls{sym-mu}^*}\gls{sym-delta}_{\gls{sym-nu},\gls{sym-nu}^*}$ is non-zero when both $\gls{sym-mu}$ and $\gls{sym-nu}$ are active in that sea state,
and $\gls{sym-delta}_{\gls{sym-mu},\gls{sym-mu}^*}(1-\gls{sym-delta}_{\gls{sym-nu},\gls{sym-nu}^*})$ is non-zero when $\gls{sym-mu}$ is active and $\gls{sym-nu}$ is inactive in the sea state.
Both Kronecker delta functions are implicitly functions of the constraint limit $\gls{sym-b-mu}$, as seen in \Cref{eq:gamma-opt-min}.
% delta mu* is monotonially decreasing (negative step) with $\gls{sym-b-mu}$ when $\gls{sym-S-mu}=1$ (enlarging the mu-th circle when the mu-th constraint is feasible-inside will make more candidate points mu-feasible, so the mu-th constraint would transition from active to inactive)
% and monotonically increasing with $\gls{sym-b-mu}$ when $\gls{sym-S-mu}=-1$ (enlarging the mu-th circle when the mu-th constraint is feasible-outside will make fewer candidate points mu-feasible, so the mu-th constraint would transition from inactive to active).
% $\gls{sym-delta}_{\gls{sym-nu},\gls{sym-nu}^*}$ is monotonially decreasing with $\gls{sym-b-mu}$.

%%%%%%%%%%%%%%%
\paragraph{Average Power as a Function of Constraint Limits}
The results just obtained establish conditions for the convexity of power in a single sea state to the constraint limit in that sea state, which we now make explicit by adding a $\gls{sym-beta}$ subscript to all sea state dependent quantities.
For \gls{acr-lcoe}, we are interested in the weighted average power across sea states as a function of the constraint limits.
All constraint limits are independent of the sea state ($\gls{sym-L-mu-beta}=\gls{sym-L-mu}$), with the exception of the amplitude limits due to slamming.
\Cref{sec:appendix-slam} will show that the slamming limits are a function of $\gls{sym-dDelta-z-slam}/\gls{sym-H}$, which depends explicitly on the wave height, and $\gls{sym-theta-surf}$, which depends implicitly on both the wave period and wave height through the guessed phase $\angle\gls{sym-xi}$.
Since the \gls{acr-pto} sizing limits are of primary relevance, the following derivation assumes that the constraint limits are independent of the sea state to simplify the expression, although sea-state-dependent limits can easily be accommodated by adding a $\gls{sym-beta}$ index on $\gls{sym-b-mu}$ and carrying it through the derivation.

From \Cref{eq:power-elec-sum,eq:gamma-opt-summary}, the average power can be written as:
\begin{equation}\label{eq:power-elec-sensitivity}
    \overline{\gls{sym-P}}_{elec} = \overline{\gls{sym-P}}_{elec,0}
    - \gls{sym-eta} \sum_{\gls{sym-beta}=1}^{\gls{sym-N-sea}} \gls{sym-w-beta}
    \cdot |\gls{sym-Gamma-beta-opt}|^2
\end{equation}
where $\overline{\gls{sym-P}}_{\text{elec},0}$ is the power available if all sea states were unconstrained (impedance-matched) and
$\gls{sym-w-beta}=\gls{acr-jpd}_{\gls{sym-beta}} \cdot
    \frac{|\gls{sym-hat-V-s-th-beta}|^2}
         {8\Re (\gls{sym-Z-s-th-beta})}$ is a constant weight for each sea state.
The second term represents power lost to constraint-induced impedance mismatch and depends on the constraint limits through \Cref{eq:gamma-opt-summary}.

Substituting \Cref{eq:gamma-opt-scale-sum} into \Cref{eq:power-elec-sensitivity} requires a $\gls{sym-beta}$ subscript on the starred index to the Kronecker delta functions since constraint activity varies by sea state, as well as a $\gls{sym-beta}$ subscript on the affine and quadratic functions to capture sea-state dependence of the Th\'{e}venin parameters:
\begin{equation}
\begin{aligned}
    \overline{\gls{sym-P}}_{elec} = &\overline{\gls{sym-P}}_{elec,0}
    - \gls{sym-eta} \sum_{\gls{sym-mu}} \sum_{\gls{sym-nu} > \gls{sym-mu}}
    \sum_{\gls{sym-beta}=1}^{\gls{sym-N-sea}} \gls{sym-w-beta}
    \left[
        \gls{sym-delta}_{\gls{sym-mu}\mu^*_\beta}(1-\gls{sym-delta}_{\gls{sym-nu}\nu^*_\beta})
        \left(
            \textrm{aff}_{\gls{sym-beta}}(\gls{sym-b-mu})
            + \sqrt{\textrm{aff}_{\gls{sym-beta}}(\gls{sym-b-mu})}
        \right)
    \right.
\\
    &\left.
        + \gls{sym-delta}_{\gls{sym-mu}\mu^*_\beta}\gls{sym-delta}_{\gls{sym-nu}\nu^*_\beta}
        \left(
            \textrm{quad}_{\gls{sym-beta}}(\gls{sym-b-mu}, \gls{sym-b-nu}) +
            \left(\textrm{aff}_{\gls{sym-beta}}(\gls{sym-b-mu}) + \textrm{aff}_{\gls{sym-beta}}(\gls{sym-b-nu})\right)
            \sqrt{\textrm{quad}_{\gls{sym-beta}}(\gls{sym-b-mu}, \gls{sym-b-nu})}
        \right)
    \right]
\end{aligned}
\end{equation}
Expanding under the assumption that each constraint is active by itself in some sea states and active with another constraint in others,
we find that $\gls{sym-overline-P-elec}$ is a piecewise function of each constraint limit $\gls{sym-b-mu}$, with a quadratic leading order term as well as affine, square root of affine, and affine times square root of quadratic terms.
The two Kronecker delta terms are combined with the weighting term $\gls{sym-w-beta}$ to create revised weighting terms $\gls{sym-pi-1-mu-nu-beta}$ and $\gls{sym-pi-2-mu-nu-beta}$ that are non-zero only when the corresponding constraint (s) are active in that sea state:
\begin{equation}\label{eq:power-triple-sum}
\begin{aligned}
    \gls{sym-overline-P-elec} &= \overline{\gls{sym-P}}_{elec,0}
    - \gls{sym-eta} \sum_{\gls{sym-mu}} \sum_{\gls{sym-nu} > \gls{sym-mu}}
    \left[
        \sum_{\gls{sym-beta}=1}^{\gls{sym-N-sea}} \gls{sym-pi-1-mu-nu-beta}
        \left(
            \textrm{aff}_{\gls{sym-beta}}(\gls{sym-b-mu})
            + \sqrt{\textrm{aff}_{\gls{sym-beta}}(\gls{sym-b-mu})}
        \right)
    \right. + \\
    &\left.
     \sum_{\gls{sym-beta}=1}^{\gls{sym-N-sea}} \gls{sym-pi-2-mu-nu-beta}
        \left(
            \textrm{quad}_{\gls{sym-beta}}(\gls{sym-b-mu}, \gls{sym-b-nu}) +
            \left(\textrm{aff}_{\gls{sym-beta}}(\gls{sym-b-mu}) + \textrm{aff}_{\gls{sym-beta}}(\gls{sym-b-nu})\right)
            \sqrt{\textrm{quad}_{\gls{sym-beta}}(\gls{sym-b-mu}, \gls{sym-b-nu})}
        \right)
    \right]
\end{aligned}
\end{equation}
so the power-constraint relationship is a sum over sea states of piecewise functions of the constraint limit $\gls{sym-b-mu}$.
Recognizing that the sea state sum in \eqref{eq:power-triple-sum} essentially just scales the individual piecewise functions, the expression can be simplified by absorbing the weights into the subscripts of the affine and quadratic terms, and combining the sum of an affine and quadratic term into a single quadratic term:
\begin{equation}\label{eq:power-double-sum}
\begin{aligned}
    \gls{sym-overline-P-elec} = \overline{\gls{sym-P}}_{elec,0}
    - \gls{sym-eta} \sum_{\gls{sym-mu}\gls{sym-nu}\gls{sym-beta}}
    &\left[
        \sqrt{\textrm{aff}_{\gls{sym-mu}\gls{sym-nu}\gls{sym-beta}}(\gls{sym-b-mu})}
        + \textrm{quad}_{\gls{sym-mu}\gls{sym-nu}\gls{sym-beta}}(\gls{sym-b-mu}, \gls{sym-b-nu})
    +
        \left(\textrm{aff}_{\gls{sym-mu}\gls{sym-nu}\gls{sym-beta}}(\gls{sym-b-mu}) + \textrm{aff}_{\gls{sym-mu}\gls{sym-nu}\gls{sym-beta}}(\gls{sym-b-nu})\right)
        \sqrt{\textrm{quad}_{\gls{sym-mu}\gls{sym-nu}\gls{sym-beta}}(\gls{sym-b-mu}, \gls{sym-b-nu})}
    \right]
\end{aligned}\end{equation}
We can then use a similar analysis as before to obtain conditions for the convexity of the overall power-constraint relationship.

%The convexity of $|\gls{sym-Gamma-opt}|^2$ in $\gls{sym-b-mu}$ derived above translates directly to convexity of $\gls{sym-overline-P-elec}$ in $\gls{sym-b-mu}$ via \Cref{eq:power-elec-sensitivity}, since the weighted sum of concave functions ($-|\gls{sym-Gamma-opt}|^2$) is concave and the addition of constant $\overline{\gls{sym-P}}_{\text{elec},0}$ preserves concavity. Equivalently, $\gls{sym-overline-P-elec}$ is convex in the constraint coefficients $\gls{sym-b-mu}$ under the same conditions, namely that all constraints are convex ($\gls{sym-S-mu} = 1$) and that the curvature of $(\textrm{aff}(\gls{sym-b-mu}) + \textrm{aff}(\gls{sym-b-nu}))\sqrt{\textrm{quad}(\gls{sym-b-mu}, \gls{sym-b-nu})}$ resolves favorably (see the caveat above). This is the convexity claim referenced in \Cref{sec:multidisciplinary-insights}.

% \paragraph{Convexity of the Power-Constraint Relationship}
% To illustrate convexity, we use an intuitive geometric argument, leaving formal proof to future work.
% In a constraint that has $\gls{sym-S-mu}=1$ in a given sea state, the feasible region is the interior of a circle, so increasing the constraint limit $\gls{sym-b-mu}$ enlarges the circle radius and therefore enlarges the feasible region, which can only decrease or maintain the optimal $|\gls{sym-Gamma-beta-opt}|$ and therefore increase or maintain power.
% In a constraint with $\gls{sym-S-mu}=-1$, the feasible region is the exterior of a circle, so increasing the constraint limit $\gls{sym-b-mu}$ enlarges the circle radius and therefore shrinks the feasible region, which can only increase or maintain the optimal $|\gls{sym-Gamma-beta-opt}|$ and therefore decrease or maintain power.
% This corresonds to the \textrm{sign} of the leading order term of a given piecewise segment in the power-constraint relationship, which is negative for $\gls{sym-S-mu}=1$ and positive for $\gls{sym-S-mu}=-1$.
% Furthermore, the number of sea states in which a particular constraint is singularly or pairwise active ($\sum_{\gls{sym-beta}}\gls{sym-mathbb-I}_\gls{sym-beta}$) is monotonically non-decreasing with the constraint limit, so the leading order term in the overall power-constraint relationship (all piecewise segments together) is positive if the constraint has $\gls{sym-S-mu}=1$ at all sea states, or negative if $\gls{sym-S-mu}=-1$ at all sea states.
% This means that as long as the \textrm{sign} of $\gls{sym-S-mu}$ is consistent across sea states, the power-constraint relationship is globally convex.
% Relaxing any constraint thus always increases or maintains power, and the marginal power benefit diminishes as the constraint limit grows.

%%%%%%%%%%%%%%%
\paragraph{Extension to Dimension Sensitivity}
\label{sec:appendix-dimension-sensitivity}
In principle, the same parametric programming approach can be used to semi-analytically compute the derivatives of the constrained optimal power with respect to bulk dimensions,
though the implementation remains out of reach.
These derivatives would be highly valuable, and indeed would enable the full wave-to-wire \gls{acr-ccd} optimization problem to be solved semi-analytically, rather than the present approach which requires a numerical bulk-dimension optimization to be wrapped around the semi-analytical control optimization and simulation.
Computing bulk-dimension derivatives would require not only $\partial \gls{sym-Gamma-opt}/\partial \gls{sym-b-mu}$ as we derive here,
but also the derivatives of Th\'{e}venin parameters with respect to bulk dimensions,
because bulk dimensions affect $\gls{sym-mathbf-Q-mu}, \gls{sym-vec-a-mu}$, and the unconstrained power through the Th\'{e}venin parameters.
The sensitivity formula would become significantly more complex and less intuitive with these additional terms,
and more importantly, the Th\'{e}venin-to-dimension derivatives are not currently available in MDOcean because they depend on the hydrodynamic coefficient-to-dimension derivatives.
Application of the adjoint method has recently enabled efficient computation of hydrodynamic coefficient derivatives in the context of a numerical boundary element method solver \citep{khanal_fully_2025},
but the adjoint method has not yet been implemented for the semi-analytical \gls{acr-meem} model used here.
Implementing the adjoint for a linear solve is relatively straightforward in a \gls{acr-mdo} package like OpenMDAO,
and would unlock a fully semi-analytical differentiation workflow.


%%%%%%%%%%%%%%%%%%%%%%%%%%%%%%%%%%%%%%%%%
\FloatBarrier
\subsection{Drag Model}
\label{sec:appendix-drag}
As described in \Cref{sec:drag}, a describing function approximation is used to model the form drag.
The relative-velocity computation in \Cref{eq:drag-pressure} uses the standard finite-depth linear-wave velocity phasor evaluated at body draft,
with $\gls{sym-vec-T} = [\gls{sym-T-f-2}, \gls{sym-T-s}]$ collecting the bottom drafts of the float and spar:
\begin{equation}
    \label{eq:wave-velocity-phasor}
    \gls{sym-vec-hat-V-wave}(\gls{sym-y}) = \frac{\gls{sym-H}}{2}\,\frac{\gls{sym-g}\gls{sym-k-wavenumber}}{\gls{sym-omega}}\, \exp\{-\gls{sym-k-wavenumber}\gls{sym-vec-T}-i\gls{sym-k-wavenumber}\gls{sym-y}\}.
\end{equation}

Strip theory integrates the spatially-dependent infinitesimal drag force vector $d\gls{sym-vec-F-d}$ over the wetted surface area $\gls{sym-A-wet}$.
For a hollow cylinder of outer radius $\gls{sym-vec-R}$ and inner radius $\gls{sym-vec-R-in}$,
% derived p156-159 notebook 3/10/25
% simplified p169-173 notebook 3/12/25
\begin{equation}\label{eq:drag-force-integral}
    \gls{sym-vec-hat-F-d} = \int_{\gls{sym-A-wet}} d\gls{sym-vec-hat-F-d}(\gls{sym-y}) = -\int_{\gls{sym-y}=-\gls{sym-vec-R}}^{\gls{sym-vec-R}} \gls{sym-vec-w-y}\, \gls{sym-vec-hat-p-d}(\gls{sym-y})\, dy,
\end{equation}
where $\gls{sym-vec-w-y}$ is the local body width perpendicular to wave propagation:
\begin{equation}\label{eq:drag-width}
    \gls{sym-vec-w-y} = \begin{cases}
        2\sqrt{\gls{sym-vec-R}^2-\gls{sym-y}^2}, & \gls{sym-vec-R-in}<|\gls{sym-y}|<\gls{sym-vec-R} \\
        2\sqrt{\gls{sym-vec-R}^2-\gls{sym-y}^2} - 2\sqrt{\gls{sym-vec-R-in}^2-\gls{sym-y}^2}, & 0<|\gls{sym-y}|<\gls{sym-vec-R-in}.
    \end{cases}
\end{equation}

Plugging the describing function \Cref{eq:drag-pressure} into the drag force integral \eqref{eq:drag-force-integral} reveals the quasi-linearized drag force to be the sum of a damping term in phase with the body motion and an excitation term potentially out of phase with the body motion:
\begin{equation}\label{eq:drag-damping-excitation-2}
    \gls{sym-vec-F-d} = -\gls{sym-mathbf-B-d} \gls{sym-vec-hat-dot-xi}
    + \gls{sym-vec-gamma-d} \gls{sym-hat-zeta}
\end{equation}
The excitation term arises from the use of the relative velocity rather than direct \gls{acr-wec} velocity, which is often overlooked in \gls{acr-wec} models and has important implications for the phase of the drag force.
The damping-excitation grouping of the drag terms, as well as the dependence of the relative velocity on the direction of wave propagation $\gls{sym-y}$, improves on the approach of \citet{quartier_influence_2021} and is discussed more in \Cref{subsec:appendix-drag-others}.
The coefficients in \Cref{eq:drag-damping-excitation} include real diagonal drag damping matrix $\gls{sym-mathbf-B-d}$ and complex drag excitation vector $\gls{sym-vec-gamma-d}$ defined as follows:
\begin{equation}\label{eq:drag-coeffs}
\begin{aligned}
    \gls{sym-mathbf-B-d} &= \frac{8}{3\gls{sym-pi}}\frac{1}{2} \gls{sym-rho-w} \, \textrm{diag}\left(\gls{sym-vec-C-d} |\gls{sym-vec-hat-dot-xi}| \int_{\gls{sym-y}=-\gls{sym-vec-R}}^{\gls{sym-vec-R}} \gls{sym-vec-w-y} \gls{sym-vec-alpha-v-y} d\gls{sym-y} \right) \\
    \gls{sym-vec-gamma-d} &= \frac{8}{3\gls{sym-pi}} \frac{1}{2} \gls{sym-rho-w} \gls{sym-vec-C-d} \, |\gls{sym-vec-hat-dot-xi}| \frac{\gls{sym-g}\gls{sym-k-wavenumber}\exp(-\gls{sym-k-wavenumber}\gls{sym-vec-T})}{\gls{sym-omega}} \int_{\gls{sym-y}=-\gls{sym-vec-R}}^{\gls{sym-vec-R}} \gls{sym-vec-w-y} \gls{sym-vec-alpha-v-y} \exp(-i\gls{sym-k-wavenumber}\gls{sym-y}) d\gls{sym-y}
\end{aligned}
\end{equation}
where we use radii vectors $\gls{sym-vec-R}=[\gls{sym-D-f}/2,\gls{sym-D-d}/2]$ and $\gls{sym-vec-R-in}=[\gls{sym-D-f-in},0]$ to represent the float and spar.
$\gls{sym-vec-alpha-v-y}$ is the amplitude ratio of the relative velocity and \gls{acr-wec} velocity:
\begin{equation}\label{eq:drag-alpha-definitions}
    \gls{sym-vec-alpha-v-y} = \frac{|\gls{sym-vec-hat-V-rel}(\gls{sym-y})|}{|\gls{sym-vec-dot-xi}|} = \sqrt{1 + \gls{sym-vec-r}^{\, 2} \exp(-2\gls{sym-k-wavenumber}\gls{sym-y}) - 2 \gls{sym-vec-r} \exp(-\gls{sym-k-wavenumber}\gls{sym-y})  \sin\angle \gls{sym-vec-dot-xi}}
\end{equation}
and $\gls{sym-vec-r} = |\gls{sym-hat-V-wave}|/ |\gls{sym-vec-dot-xi}|$ is the amplitude ratio of the incident wave velocity and \gls{acr-wec} velocity.

%%%%%%%%%%%%%%%
\subsubsection{Nondimensional Drag Integrals $\gls{sym-I-B}$ and $\gls{sym-I-G}$}
The drag coefficients $\gls{sym-mathbf-B-d}$ and $\gls{sym-vec-gamma-d}$ are expressed in \Cref{eq:drag-coeffs} as dimensional integrals over the width of the \gls{acr-wec}.
Here we present an alternative form of the coefficients in terms of nondimensionalized integral functions, which we refer to as $\gls{sym-I-B}(\gls{sym-r},\gls{sym-theta},\gls{sym-kappa})$ and $\gls{sym-I-G}(\gls{sym-r},\gls{sym-theta},\gls{sym-kappa})$:
\begin{equation}\label{eq:drag-integrals-IB-IG}
\begin{aligned}
    \gls{sym-mathbf-B-d} &= \frac{8}{3\gls{sym-pi}} \gls{sym-rho-w} \gls{sym-vec-R}^2 \gls{sym-C-d} |\gls{sym-hat-dot-xi}|
    \left[
        \gls{sym-I-B}(\gls{sym-r},\angle\gls{sym-hat-dot-xi},\gls{sym-k-wavenumber}\gls{sym-vec-R}) - \gls{sym-alpha}^2 \gls{sym-I-B}(\gls{sym-r},\angle\gls{sym-hat-dot-xi},\gls{sym-alpha} \gls{sym-k-wavenumber} \gls{sym-vec-R})
    \right]
    \\
    \gls{sym-vec-gamma-d} &= \frac{8}{3\gls{sym-pi}} \gls{sym-rho-w} \gls{sym-vec-R}^2 \gls{sym-C-d} |\gls{sym-hat-dot-xi}| \frac{\gls{sym-g} \gls{sym-k-wavenumber}}{\gls{sym-omega}} \exp(-\gls{sym-k-wavenumber}\gls{sym-vec-T})
    \left[
        \gls{sym-I-G}(\gls{sym-r},\angle\gls{sym-hat-dot-xi},\gls{sym-k-wavenumber}\gls{sym-vec-R}) - \gls{sym-alpha}^2 \gls{sym-I-G}(\gls{sym-r},\angle\gls{sym-hat-dot-xi},\gls{sym-alpha} \gls{sym-k-wavenumber}\gls{sym-vec-R})
    \right]
\end{aligned}
\end{equation}
where the nondimensional integral functions $\gls{sym-I-B}$ (real) and $\gls{sym-I-G}$ (complex) are defined as:
\begin{equation}\label{eq:definition-IB-IG}
\begin{aligned}
    \gls{sym-I-B}(\gls{sym-r},\gls{sym-theta},\gls{sym-kappa})
    &= \int_{-1}^1 \sqrt{1-\gls{sym-x}^2}\sqrt{1+\gls{sym-r}^2\exp(-2\gls{sym-kappa} \gls{sym-x})
    -2\gls{sym-r}\sin\gls{sym-theta} \exp(-\gls{sym-kappa} \gls{sym-x})} dx
    \\
    \gls{sym-I-G}(\gls{sym-r},\gls{sym-theta},\gls{sym-kappa}) 
    &= \int_{-1}^1 \sqrt{1-\gls{sym-x}^2}\sqrt{1+\gls{sym-r}^2\exp(-2\gls{sym-kappa} \gls{sym-x})
    -2\gls{sym-r}\sin\gls{sym-theta} \exp(-\gls{sym-kappa} \gls{sym-x})} \exp(-i\gls{sym-kappa}\gls{sym-x}) dx
\end{aligned}
\end{equation}

The integrals must be evaluated numerically, except for special limit cases for which analytical results are available.
\Cref{fig:drag-integrals} shows the results of numerical integration for $\gls{sym-I-B}$ and $\gls{sym-I-G}$ across a range of inputs.
\begin{figure}[h!]
    \centering
    \begin{subfigure}[t]{0.32\textwidth}
        \includegraphics[\FigureOptions{aor:figs/from-matlab/drag_integral_B.pdf}]{figs/from-matlab/drag_integral_B.pdf}
        \caption{Drag damping integral $\gls{sym-I-B}$, scaled by $\frac{2}{\gls{sym-pi}}$.}
        \label{fig:drag-damping-integral}
    \end{subfigure}
    \hfill
    \begin{subfigure}[t]{0.32\textwidth}
        \includegraphics[\FigureOptions{aor:figs/from-matlab/drag_integral_G_m.pdf}]{figs/from-matlab/drag_integral_G_m.pdf}
        \caption{Magnitude of drag excitation integral $|\gls{sym-I-G}|$, scaled by $\frac{2}{\gls{sym-pi}}$.}
        \label{fig:drag-excitation-magnitude-integral}
    \end{subfigure}
    \hfill
    \begin{subfigure}[t]{0.32\textwidth}
        \includegraphics[\FigureOptions{aor:figs/from-matlab/drag_integral_G_p.pdf}]{figs/from-matlab/drag_integral_G_p.pdf}
        \caption{Phase of drag excitation integral $\angle \gls{sym-I-G}$, scaled by $\frac{1}{\gls{sym-pi}}$.}
        \label{fig:drag-excitation-phase-integral}
    \end{subfigure}
    \caption[Nondimensional drag integrals in polar coordinates]{Nondimensional drag integrals shown in polar coordinates for various values of $\gls{sym-kappa}$, where the polar radius represents $1/\gls{sym-r}$ and the polar angle represents $\gls{sym-theta}$.}
    \label{fig:drag-integrals}
\end{figure}

The limit cases in question include $\gls{sym-r} \rightarrow 0$ (the body velocity dominates the incident wave velocity) and
$\gls{sym-kappa} \rightarrow 0$ (the body diameter is small with respect to the wavelength), with the following expressions:
\begin{equation}
\begin{aligned}
    \lim_{\gls{sym-r}\rightarrow 0} \gls{sym-I-B}(\gls{sym-r},\gls{sym-theta},\gls{sym-kappa}) 
    = \frac{\gls{sym-pi}}{2},
    \qquad
    \lim_{\gls{sym-r}\rightarrow 0} \gls{sym-I-G}(\gls{sym-r},\gls{sym-theta},\gls{sym-kappa}) =
    \begin{cases}
    \displaystyle\frac{\gls{sym-pi} \gls{sym-mathrm-J-1}(\gls{sym-kappa})}{\gls{sym-kappa}} & \textrm{if } \gls{sym-kappa} \neq 0 \\
    \displaystyle\frac{\gls{sym-pi}}{2} & \textrm{if } \gls{sym-kappa} = 0
    \end{cases} \\
    \lim_{\gls{sym-kappa} \rightarrow 0} \gls{sym-I-B}(\gls{sym-r},\gls{sym-theta},\gls{sym-kappa}) 
    =  \lim_{\gls{sym-kappa} \rightarrow 0} \gls{sym-I-G}(\gls{sym-r},\gls{sym-theta},\gls{sym-kappa}) 
    = \frac{\gls{sym-pi}}{2}\sqrt{1+\gls{sym-r}^2-2\gls{sym-r}\sin\gls{sym-theta}},
\end{aligned}
\end{equation}
where $\gls{sym-mathrm-J-1}$ is the Bessel function of the first kind of order 1.
An important edge case occurs at the origin of the polar plots in \Cref{fig:drag-integrals}, when $\gls{sym-r} \rightarrow \infty$ (the incident wave velocity dominates the body velocity, so $|\gls{sym-hat-dot-xi}|\rightarrow 0$).
Here, the integrals as defined in \Cref{eq:definition-IB-IG} diverge, and the form of \Cref{eq:drag-integrals-IB-IG} does not apply.
Instead, when $\gls{sym-r}\rightarrow \infty$, we apply an alternative convention that replaces the \gls{acr-wec} speed in \Cref{eq:drag-integrals-IB-IG} with the incident wave speed and uses a primed form of the integrals, $\gls{sym-I-prime} = \gls{sym-I-drag}/\gls{sym-r}$:
\begin{equation}
\begin{aligned}
    \gls{sym-mathbf-B-d} &= \frac{8}{3\gls{sym-pi}} \gls{sym-rho-w} \gls{sym-vec-R}^2 \gls{sym-C-d} |\gls{sym-hat-V-wave}|
    \left[
        \gls{sym-I-B}'(\gls{sym-r},\angle\gls{sym-hat-dot-xi},\gls{sym-k-wavenumber}\gls{sym-vec-R}) 
        - \gls{sym-alpha}^2 \gls{sym-I-B}'(\gls{sym-r},\angle\gls{sym-hat-dot-xi},\gls{sym-alpha} \gls{sym-k-wavenumber} \gls{sym-vec-R})
    \right]
    \\
    \gls{sym-vec-gamma-d} &= \frac{8}{3\gls{sym-pi}} \gls{sym-rho-w} \gls{sym-vec-R}^2 \gls{sym-C-d} |\gls{sym-hat-V-wave}|
    \frac{\gls{sym-g}\gls{sym-k-wavenumber}}{\gls{sym-omega}} \exp(-\gls{sym-k-wavenumber}\gls{sym-vec-T})
    \left[
        \gls{sym-I-G}'(\gls{sym-r},\angle\gls{sym-hat-dot-xi},\gls{sym-k-wavenumber}\gls{sym-vec-R}) 
        - \gls{sym-alpha}^2 \gls{sym-I-G}'(\gls{sym-r},\angle\gls{sym-hat-dot-xi},\gls{sym-alpha} \gls{sym-k-wavenumber}\gls{sym-vec-R})
    \right]
\end{aligned}
\end{equation}
The relevant limits are then:
\begin{equation}
\begin{aligned}
    \lim_{\gls{sym-r}\rightarrow \infty} \gls{sym-I-B}'(\gls{sym-r},\gls{sym-theta},\gls{sym-kappa}) = 0,
    \qquad
    \lim_{\gls{sym-r}\rightarrow \infty} \gls{sym-I-G}'(\gls{sym-r},\gls{sym-theta},\gls{sym-kappa}) =
    \begin{cases}
    \displaystyle  \frac{
                        \gls{sym-pi} \gls{sym-mathrm-J-1}((1i-1)\gls{sym-kappa})
                        }
                        {
                            (1i-1)\gls{sym-kappa}
                        }
                    & \textrm{if } \gls{sym-kappa} \neq 0 \\
    \displaystyle\frac{\gls{sym-pi}}{2}
                    & \textrm{if } \gls{sym-kappa} = 0
    \end{cases}
\end{aligned}
\end{equation}

%%%%%%%%%%%%%%%
\subsubsection{Comparison to Other Drag Formulations}
\label{subsec:appendix-drag-others}
\citet{quartier_influence_2021} propose a \gls{acr-wec} drag model that also utilizes describing functions with the relative velocity but groups the drag force into damping and stiffness terms rather than damping and excitation terms.
While this representation yields identical results for a floating \gls{acr-wec} in the frequency domain, it misleadingly implies that a body fixed at the still water line would experience no drag force from incident waves.
In other words, it fails to account for the edge case of $\gls{sym-r} \rightarrow \infty$ described above, where the drag force is purely an excitation force and cannot be expressed as a damping or stiffness.
Furthermore, the drag ``damping'' term in the grouping of \citet{quartier_influence_2021} includes not just forces with a dissipative effect on the body but also the portion of nonlinear drag excitation forces that are in phase with the body velocity.
This creates the confusing possibility for negative apparent damping coefficients, leading to false concerns about violating energy conservation if erroneously interpreted as a true damping coefficient.
The damping-excitation grouping chosen in \Cref{sec:drag} avoids this confusion and guarantees positive values of the damping coefficient $\gls{sym-mathbf-B-d}$.

Another improvement over \citet{quartier_influence_2021} is that the present model considers the fact that the incident wave velocity, and thus the relative velocity, differs across the width of the \gls{acr-wec} in the direction of wave propagation $\gls{sym-y}$.
\citet{quartier_influence_2021} more simplistically assumes that the wave velocity at all points equals that at the center of the \gls{acr-wec}, which is only accurate when the incident wavelength far exceeds the \gls{acr-wec} width.
This means that \citet{quartier_influence_2021} captures only the Froude-Krylov-like component of the drag excitation, neglecting the diffraction-like component.
Early MDOcean simulations that used the \citet{quartier_influence_2021} formulation found this approximation to be unacceptable, producing unstable dynamics (negative drag damping coefficients whose absolute value exceeds the radiation damping) at high frequencies and large wave heights in the operational regime.
The strip-theory approach presented in \Cref{sec:drag} avoids this issue.
Comparisons against WEC-Sim confirm that even without strip theory, assuming that the relative velocity is perfectly in phase with the body velocity (i.e. a drag force of purely damping and in-phase excitation, with no stiffness or out-of-phase excitation) yields more accurate results than the drag formulation used in the study \citep{quartier_influence_2021}.

%%%%%%%%%%%%%%%
\subsubsection{Iterative Solution}\label{sec:appendix-drag-iterative}
With the describing function approximation, the nonlinear time-domain drag equation is now a quasi-linear frequency-domain equation, since the state-dependence of the coefficients $\gls{sym-mathbf-B-d}$ and $\gls{sym-vec-gamma-d}$ prevents true linearity.
The solution for states $|\gls{sym-vec-hat-dot-xi}|$ and $\angle \gls{sym-vec-hat-dot-xi}$ for a given controller can be obtained either through numerical iteration or analytical solution of the nonlinear equation \eqref{eq:eom}.
Iteration is chosen here, the same approach used for a frequency-domain drag simulation of floating offshore wind turbines \citep{hall_open-source_2022}.
MDOcean provides users with two options: to employ a typical nonlinear root-finding algorithm, or to simply use the equation of motion \eqref{eq:eom} to obtain subsequent iterates.
The latter is a form of fixed point iteration where $\gls{sym-vec-hat-dot-xi}_{\gls{sym-k-generic-index}} = \gls{sym-f}(\gls{sym-vec-hat-dot-xi}_{\gls{sym-k-generic-index}-1})$, where $\gls{sym-f}(\gls{sym-vec-hat-dot-xi})$ is the equation of motion \Cref{eq:eom-freq-domain} and $\gls{sym-k-generic-index}$ is the iteration index.
It is therefore guaranteed to converge as long as the dynamics are contracting at the solution, which \citet{chicone_contraction_2006} shows has the following criteria:
\begin{equation}
 \left|
    \frac{\partial}
    {\partial \gls{sym-vec-hat-dot-xi}}~\left[
    \gls{sym-f}\left(
        \gls{sym-vec-hat-dot-xi}~
    \right)\right]
\right| < 1
\end{equation}

Five to eight iterations on $|\gls{sym-vec-hat-dot-xi}|$ and $\angle \gls{sym-vec-hat-dot-xi}$ are typically required to converge all sea states to within 0.01~m and 3 degrees. % from summer research slide 44 5/30/24
%In high-frequency sea states where there is substantial phase difference between $\gls{sym-dot-X-rel}$ and $\gls{sym-dot-X}$, $\gls{sym-B-d}$ occasionally becomes negative.
At present (since MDOcean v1.0.0), the solver used for drag iteration can be simultaneously used for the numerical optimal control procedure described in \Cref{sec:optimal-control}, although since v1.2.0 by default the solver is used only for drag convergence and optimal control is done analytically.
Compared to numerical optimal control nested inside the drag solve, combining the solvers in this way has the advantages of code simplicity and reducing the number of dynamics evaluations required for a given simulation, provided that the number of iterations required to converge the coupled drag-control solve is less than the product of the number of iterations required to converge each solve individually.
However, it has the disadvantage that the error signal couples drag nonlinearities with constraint violations, which can potentially lead to slower convergence.

%%%%%%%%%%%%%%%
\subsubsection{Optimal Control Condition}
Note that in the quasi-linear formulation, the lack of complete linearity means that the complex-conjugate reactive controller is technically no longer optimal in the unconstrained case.
The optimal load damping $\gls{sym-B-l}$ is derived by setting
\begin{equation}
\begin{aligned}
0 &= \frac{\partial }{\partial \gls{sym-B-l}} \left[\gls{sym-p-avg-VI}\right]
= \frac{\partial}{\partial \gls{sym-B-l}} \left[\frac{1}{2} \gls{sym-B-l} |\gls{sym-hat-I}|^2\right]
= \frac{\partial }{\partial \gls{sym-B-l}} \left[\frac{1}{2} \gls{sym-B-l} \left|\frac{\gls{sym-hat-V-s-th}}{\gls{sym-Z-s-th}+\gls{sym-B-l}+\gls{sym-K-l}/\gls{sym-s}}\right|^2 \right]
\end{aligned}
\end{equation}
and the typical simplification $\frac{\partial \gls{sym-hat-V-s-th}}{\partial \gls{sym-B-l}}=\frac{\partial \gls{sym-Z-s-th}}{\partial \gls{sym-B-l}}=0$ must be replaced with the implicit dependence of the Th\'{e}venin equivalent parameters on $\gls{sym-B-l}$ because $\gls{sym-B-l}$ affects $\gls{sym-hat-dot-xi}$ and therefore $\gls{sym-mathbf-B-d}$ and $\gls{sym-vec-gamma-d}$ via \eqref{eq:drag-coeffs}.
In practice, however, the standard complex conjugate controller was observed to produce nearly identical power as the one incorporating the more complicated dependence, so the standard is used for simplicity.


%%%%%%%%%%%%%%%
\FloatBarrier
\subsection{Dynamic Limit Details}
\label{sec:appendix-dynamic-limit}

The operational and storm loadcases described in \Cref{sec:design-load-cases} correspond roughly with the 1-year and 50-year return periods specified in the wave energy design standard \gls{acr-iec} \gls{acr-ts} 62600-2,
and with the \gls{acr-fls} and \gls{acr-uls} in offshore wind standards.
When the \gls{acr-wec} is in survival mode, an external brake locks the \gls{acr-pto}, enforcing the float and spar to move together without loading the generator.

\ifdefined\DISSERTATION
    This section details the allowable amplitudes and forces in each case.
    Currently all sea states in the \gls{acr-jpd} with nonzero probability are considered operational, since the generator force and power limit serve to reduce loads in energetic sea states.
    Future studies could control the transition from operational to survival mode more finely, perhaps with an additional input variable representing an operational incident energy threshold.
\else
    All nonzero-probability sea states in the \gls{acr-jpd} are considered operational;
    the generator force and power limits reduce loads in the most energetic sea states.
\fi

\subsubsection{Relationship Between Peak and Fundamental Amplitudes}
\label{sec:appendix-peak-fundamental}
\Cref{tab:DLCs} summarizes the peak forces and amplitudes based on the assumptions of each design load case.

\begin{table}[htbp]
    \centering
    \caption[Forces and amplitudes for each design load case]{Forces, amplitudes, and maximum allowable amplitudes in each design load case}
\label{tab:DLCs}
\begin{tabular}{P{0.28\linewidth}P{0.3\linewidth}P{0.33\linewidth}}
    & \multicolumn{2}{c}{\textbf{Design Loadcase}}\\\cline{2-3}
    \textbf{Variable}& \textbf{1: Operational}&\textbf{2: Storm}\\ \hline
    \parbox[m]{\linewidth}{\centering \vspace{8pt}
    Peak powertrain force, $\max_{\gls{sym-t}} |\gls{sym-F-p}(\gls{sym-t})|$ \vspace{8pt} }
        & \parbox[m]{\linewidth}{\centering
        Given by \Cref{eq:peak-pto-force}}
        & \parbox[m]{\linewidth}{\centering
        Solve \Cref{eq:eom-freq-domain} with \Cref{eq:op-hydro-coeffs} coefficients for $\gls{sym-F-p}$}\\ \hline
    \parbox[m]{\linewidth}{\centering
    Peak fluid force, $\max_{\gls{sym-t}} |\gls{sym-m}\gls{sym-ddot-xi}(\gls{sym-t})-\gls{sym-F-p}(\gls{sym-t})|$ }
        & \parbox[m]{\linewidth}{\centering
        Given by \Cref{eq:peak-fluid-force} with \Cref{eq:op-hydro-coeffs} coefficients and $\gls{sym-F-p}$ from \Cref{eq:peak-pto-force}}
        & \parbox[m]{\linewidth}{\centering
        Given by \Cref{eq:peak-fluid-force} with \Cref{eq:st-hydro-coeffs} coefficients and $\gls{sym-F-p}$ from storm case}\\ \hline
    \parbox[m]{\linewidth}{\centering
    Peak float travel, $\max_{\gls{sym-t}} |\gls{sym-xi-f}(\gls{sym-t})|$ }
        & \parbox[m]{\linewidth}{\centering
        Given by \Cref{eq:eom-freq-domain} with \Cref{eq:op-hydro-coeffs} coefficients}
        & \parbox[m]{\linewidth}{\centering
        Given by \Cref{eq:eom-freq-domain} with \Cref{eq:st-hydro-coeffs} coefficients and $\gls{sym-F-p}=0$}\\ \hline
    \parbox[m]{\linewidth}{\centering
    Peak spar travel, $\max_{\gls{sym-t}} |\gls{sym-xi-s}(\gls{sym-t})|$ }
        & \parbox[m]{\linewidth}{\centering
        Given by \Cref{eq:eom-freq-domain} with \Cref{eq:op-hydro-coeffs} coefficients}
        &Equals $\gls{sym-xi-f}$ due to brake\\ \hline
    \parbox[m]{\linewidth}{\centering \vspace{8pt}
    Maximum allowable relative amplitude between float and spar, $|\gls{sym-xi-f}-\gls{sym-xi-s}|_{\text{max}}$ \vspace{8pt}}
        & \shortstack{$\min(\gls{sym-h}_{fs,\text{clear}},\gls{sym-h}_{fs,\text{up}},$\\$\gls{sym-h}_{fs,\text{down}})$}
        &N/A\\  \hline
    \parbox[m]{\linewidth}{\centering \vspace{8pt}
    Maximum allowable float amplitude, $|\gls{sym-xi-f}|_{\text{max}}$ \vspace{8pt} }
        & $\min(\gls{sym-xi-f-lin}, \gls{sym-xi-f-slam})$
        & N/A \\ \hline %$\gls{sym-X}_{\text{slam}}$\\
    \parbox[m]{\linewidth}{\centering \vspace{8pt}
    Maximum allowable spar amplitude, $|\gls{sym-xi-s}|_{\text{max}}$ }
        & $\min(\gls{sym-xi-s-lin}, \gls{sym-xi-s-slam})$
        & N/A %$\gls{sym-X}_{\text{slam}}$\\
\end{tabular}
\end{table}

The surge force for each component is calculated as \citep{newman_motions_1963}
\begin{equation}\label{eq:surge-force}
   \gls{sym-F-surge} = \frac{\gls{sym-H}\gls{sym-rho-w} \gls{sym-omega}^2 \gls{sym-A-w}}{\gls{sym-k-wavenumber}} (\exp(-\gls{sym-k-wavenumber}\gls{sym-z-top})-\exp(-\gls{sym-k-wavenumber}\gls{sym-T-draft}))
\end{equation}
using the respective waterplane areas $\gls{sym-A-w}$  and drafts $\gls{sym-T-draft}$ for the float and spar.
The top vertical height $\gls{sym-z-top}$ is set to zero for the float and $\gls{sym-T-f-2}$ for the spar.
The design surge force is the maximum surge force over all nonzero-probability sea states.
Equation~\eqref{eq:surge-force} assumes a slender cylinder with diameter much smaller than the wavelength and may be less accurate for the float due to diffraction effects.
Incorporating semi-analytical hydrodynamic coefficients for surge is a possible area of future work.
Maximum acceptable heave %and surge
forces will be evaluated with structural factors of safety in \Cref{sec:structures}.
The structures model currently lacks the ability to account for surge forces, so the surge force calculation is not used in the simulation.

\ifdefined\DISSERTATION
    The heave forces in equation \eqref{eq:eom} can be separated into fluid forces, which act on the vertically-oriented wetted surface of the body, the \gls{acr-pto} force, which acts at the connection between the float and the spar, and the inertial force, which for dynamic purposes is assumed to act at the center of mass but for structural purposes actually acts distributed over the material.
    The distributed nature of the inertial force means that the float and spar structures experience internal loads that vary not only in time but also in space along their lengths.
    Due to force balance in dynamic equilibrium, the maximum magnitude of internal heave force in space at any given time, $\gls{sym-F-heave-f}(\gls{sym-t})$, is the greater of the absolute value of the \gls{acr-pto} force and the fluid force.
    Expressing the latter as the difference between the inertial force and the \gls{acr-pto} force, this gives:
\else
    By force balance, the peak internal heave force at any cross-section of each body is the greater of the \gls{acr-pto} force and the fluid (inertial minus \gls{acr-pto}) force:
\fi
\begin{equation}\label{eq:heave-force-decomp}
\begin{aligned}
    \gls{sym-F-heave-f} (\gls{sym-t}) &=\max\!\left(|\gls{sym-F-p} (\gls{sym-t})|,~ |\gls{sym-m-f}\gls{sym-ddot-xi-f} (\gls{sym-t})-\gls{sym-F-p} (\gls{sym-t})|\right) \\
    \gls{sym-F-heave-s} (\gls{sym-t}) &=\max\!\left(|\gls{sym-F-p} (\gls{sym-t})|,~ |\gls{sym-m-s}\gls{sym-ddot-xi-s} (\gls{sym-t})-\gls{sym-F-p} (\gls{sym-t})|\right)
\end{aligned}
\end{equation}
In operational seas, $\gls{sym-F-p}(\gls{sym-t})$ is determined by the generator torque $\gls{sym-tau}(\gls{sym-t})$ cascaded through the \gls{acr-pto} dynamics matrix $\gls{sym-mathbf-a}_{\gls{sym-F-p}\gls{sym-dot-xi}\leftarrow \gls{sym-tau}\gls{sym-Omega}}$.
In storm seas, $\gls{sym-F-p}(\gls{sym-t})$ represents the brake force, and its fundamental amplitude is obtained by solving the 2-\gls{acr-dof} equation of motion for the force necessary to enforce zero relative motion between the float and spar.
\ifdefined\DISSERTATION
    Maximizing over time requires more care because nonlinearity-induced harmonics cannot be neglected for peak force analysis like they can for power production.
    The peak of $\gls{sym-tau}(\gls{sym-t})$ is lower than its fundamental amplitude due to generator force saturation at \gls{sym-tau-max}, so using \gls{sym-tau-max} rather than $|\gls{sym-hat-tau}|$ is both less conservative and more realistic.
    Assuming conservatively that the shaft torque constructively interferes with the saturated generator torque, for the operational case we use:
\else
    Saturation caps the operational generator torque at $\gls{sym-tau-max}$, which is a less conservative and more realistic peak than the fundamental amplitude $|\gls{sym-hat-tau}|$.
    Assuming constructive interference with the shaft torque,
\fi
\begin{equation}\label{eq:peak-pto-force}
\max_{\gls{sym-t}} | \gls{sym-F-p-op} (\gls{sym-t}) |
\approx \frac{1}{\gls{sym-R-gear}}
\left(
\left| \gls{sym-Z-shaft} \gls{sym-hat-Omega} \right|
                         + \gls{sym-tau-max} \right).
\end{equation}
The storm case lacks the control saturation nonlinearity and uses the fundamental amplitude $\gls{sym-hat-F-p-st}$ directly.

The fluid force nonlinearities $|\gls{sym-m}\gls{sym-ddot-xi}(\gls{sym-t})-\gls{sym-F-p}(\gls{sym-t})|$ are harder to handle, since unlike the generator torque, the peak of the drag force exceeds its fundamental.
MDOcean currently uses the fundamental amplitude for both load cases:
\begin{equation}\label{eq:peak-fluid-force}
\max_{\gls{sym-t}} \left|\gls{sym-m}\gls{sym-ddot-xi}(\gls{sym-t}) - \gls{sym-F-p}(\gls{sym-t})\right| \approx \left|-\gls{sym-m}\gls{sym-omega}^2 \gls{sym-hat-xi} - \gls{sym-hat-F-p}\right|.
\end{equation}
This approximation can either under- or overestimate the true peak depending on whether harmonics interfere constructively or destructively with the fundamental; further discussion is in \Cref{sec:appendix-peak-fluid-force}.
\ifdefined\DISSERTATION
    The fundamental amplitudes are also used to calculate the peak float and spar displacements $\max_{\gls{sym-t}} |\gls{sym-xi-f}(\gls{sym-t})|$ and $\max_{\gls{sym-t}} |\gls{sym-xi-s}(\gls{sym-t})|$ in both cases.
    Using the fundamental for the response displacements is more accurate than for the forces because the system dynamics filter out high frequency inputs.
\else
    Peak float and spar displacements likewise use fundamental amplitudes; the second-order plant filters out higher harmonics, making the approximation more accurate for displacements than for forces.
\fi

\paragraph{Conservatism of the Fundamental-Amplitude Approximation}
\label{sec:appendix-peak-fluid-force}

\Cref{eq:peak-fluid-force} approximates the peak fluid force by the magnitude of its fundamental harmonic phasor.
This approximation can either under- or overestimate the true peak depending on whether the higher-harmonic content (primarily from the drag describing function) interferes constructively or destructively with the fundamental.

The fundamental amplitudes are also used for the peak float and spar displacements $\max_{\gls{sym-t}}|\gls{sym-xi-f}(\gls{sym-t})|$ and $\max_{\gls{sym-t}}|\gls{sym-xi-s}(\gls{sym-t})|$ in both load cases.
The fundamental is more accurate for displacements than for forces because the second-order plant dynamics filter out high-frequency inputs, attenuating drag-induced harmonics in the displacement response.

A future refinement could explicitly reconstruct the time-domain drag force from the describing function and track its peak directly, at the cost of either time-domain integration or additional Fourier-component bookkeeping in the frequency-domain formulation.

\paragraph{Maximum Permissible Amplitudes}
The table also indicates the maximum permissible amplitudes.
In operational sea states, the geometric clearance between the float and spar is enforced to prevent overtravel.
The permissible heights for upward and downward motion of the float relative to the spar are:
\begin{equation}\label{eq:h-fs-up-down}
\begin{aligned}
    \gls{sym-h-fs-up} &= \gls{sym-h-s} - \gls{sym-T-s} - (\gls{sym-h-f}- \gls{sym-T-f-2}) \\
    \gls{sym-h-fs-down} &= \gls{sym-T-s} - \gls{sym-T-f-2} - \gls{sym-h-d}
\end{aligned}
\end{equation}
Additionally, a limit of $\gls{sym-h-fs-clear}$ (see definition in \Cref{fig:dims}) is required to ensure clearance of the tubular structure connecting the float and spar.

The conditions for linear wave-body interaction are enforced to maintain compatibility with the linear potential flow theory approach of \Cref{sec:hydro}.
This maximum amplitude $\gls{sym-xi-lin}$ for the float and spar respectively is
\begin{equation}
    \gls{sym-xi-f-lin} = \frac{1}{10} (\gls{sym-h}-\gls{sym-T-f-2}), \quad
    \gls{sym-xi-s-lin}  = \frac{1}{10} (\gls{sym-h}-\gls{sym-T-s})
\end{equation}
\ifdefined\DISSERTATION
    which is derived by requiring that the height of the water column below the body (the relevant dimension determining the hydrodynamic coefficients in \gls{acr-meem}) does not change by more than 10\% from its equilibrium value when the body is at its peak displacement.
\else
    derived by requiring the water column height below the body (relevant for the \gls{acr-meem} hydrodynamic coefficients) to vary by no more than 10\% from equilibrium.
\fi

In both operational and storm states, the \gls{acr-wec} surface should not transiently pierce the wave free surface, since this causes extremely high loads.
This includes slamming (the bottom surface rises out of the water and impacts on reentry) and submersion (the top of the \gls{acr-wec} becomes fully submerged).
The condition to prevent slamming and submersion in regular waves of height $\gls{sym-H}$ is derived in \Cref{sec:appendix-slam}.
These limits depend on the phase of the response, not just the magnitude, so they do not use the quadratic coefficients of \Cref{tab:qp-constraints} but can nonetheless be expressed in a quadratic form compatible with \Cref{eq:opt-problem}:
\begin{equation}\label{eq:slam-constraint-quadratic}
    \gls{sym-hat-xi}^*\gls{sym-hat-xi} 
    - \gls{sym-H} \Re(
                    \exp (i\gls{sym-k-wavenumber}\gls{sym-y})
                    \gls{sym-hat-xi}
                    )
    < \gls{sym-dDelta-z-slam}^2
    -\left(
            \frac{\gls{sym-H}}{2}
    \right)^2
\end{equation}
for relevant vertical body dimension $\gls{sym-dDelta-z-slam}$ and horizontal coordinate $\gls{sym-y}$.
The values of $\gls{sym-dDelta-z-slam}$ used for float and spar in operational and storm cases and the rationale
(avoiding slanted-surface slamming, damping plate surfacing, and related effects)
are given in \Cref{sec:appendix-slam}.


%%%%%%%%%%%%%%%%%%%%%%%%%%%%%%%%%%%%%%%%%
\subsection{Slamming and Submersion Amplitude}
\label{sec:appendix-slam}
In regular waves of height $\gls{sym-H}$, the conditions for slamming can be derived by comparing the vertical position of the bottom of the \gls{acr-wec}, $\gls{sym-xi-t}-\gls{sym-dDelta-z-slam}$, with the wave elevation $\gls{sym-zeta-y-t}$, which depends on the horizontal position coordinate $\gls{sym-y}$.
To prevent slamming, the criteria
\begin{equation}\label{eq:slamming-time-domain}
\begin{aligned}
    \gls{sym-xi-t} - \gls{sym-dDelta-z-slam} &< \gls{sym-zeta-y-t} \\
    |\gls{sym-hat-xi}|\cos(\gls{sym-omega} \gls{sym-t}+\angle \gls{sym-hat-xi}) - \gls{sym-dDelta-z-slam} &< \frac{\gls{sym-H}}{2}\cos(\gls{sym-omega} \gls{sym-t} - \gls{sym-k-wavenumber}\gls{sym-y})
\end{aligned}
\end{equation}
must be true at all times and over all positions where the body could exit the waves, $-\gls{sym-D}/2<\gls{sym-y}<\gls{sym-D}/2$.
Note that this expression assumes a sinusoidal waveshape but still applies for irregular waves if a wave-by-wave approach \citep{lin_fast_2025} is utilized.
Manipulating \Cref{eq:slamming-time-domain} to ensure the inequality is satisfied for all $\gls{sym-t}$ yields the criterion:
\begin{equation}
|\gls{sym-hat-xi}|^2 - \gls{sym-H}\cos(\gls{sym-k-wavenumber}\gls{sym-y})|\gls{sym-hat-xi}|\cos\angle\gls{sym-hat-xi} + \gls{sym-H}\sin(\gls{sym-k-wavenumber}\gls{sym-y})|\gls{sym-hat-xi}|\sin\angle\gls{sym-hat-xi} < \gls{sym-dDelta-z-slam}^2-\left(\frac{\gls{sym-H}}{2}\right)^2
\end{equation}
which is quadratic in $\gls{sym-hat-xi}$ and therefore already in a form suitable for the constrained optimal control procedure described in \Cref{sec:optimal-control,sec:appendix-qp-solution}, once the appropriate worst-case $\gls{sym-y}$ is substituted.
The remainder of this section develops an explicit expression for the allowable amplitudes under various scenarios and the corresponding worst-case $\gls{sym-y}$.
Completing the square results in:
\begin{equation}\label{eq:slamming-squared}
    \left[
        |\gls{sym-hat-xi}|- \frac{\gls{sym-H}}{2}\cos(\angle\gls{sym-hat-xi} + \gls{sym-k-wavenumber}\gls{sym-y})
    \right]^2
    <
    \gls{sym-dDelta-z-slam}^2 - (\gls{sym-H}/2)^2 \sin^2(
                                                        \angle\gls{sym-hat-xi} + \gls{sym-k-wavenumber}\gls{sym-y}
                                                        )
\end{equation}
In general, after taking the square root of \Cref{eq:slamming-squared}, both the positive and negative branches must be used.
This yields a maximum and minimum amplitude criteria to prevent slamming, $\gls{sym-xi}_{\text{min},\text{slam}}<|\gls{sym-hat-xi}|<\gls{sym-xi}_{\text{max},\text{slam}}$, with the following values:
\begin{equation}\label{eq:slamming-min-max}
\begin{aligned}
    \gls{sym-xi-max-slam} &= \frac{\gls{sym-H}}{2} \cos\gls{sym-theta-surf} + \sqrt{\gls{sym-dDelta-z-slam}^2-\left(\frac{\gls{sym-H}}{2}\sin\gls{sym-theta-surf}\right)^2 } \\
    \gls{sym-xi-min-slam} &= \frac{\gls{sym-H}}{2} \cos\gls{sym-theta-surf} - \sqrt{\gls{sym-dDelta-z-slam}^2-\left(\frac{\gls{sym-H}}{2}\sin\gls{sym-theta-surf}\right)^2 }
\end{aligned}
\end{equation}
where $\gls{sym-theta-surf}=\angle\gls{sym-hat-xi} + \gls{sym-k-wavenumber}\gls{sym-y}$ must be evaluated with the value of $\gls{sym-y}$ that creates the minimum value of $\gls{sym-xi}_{\text{max},\text{slam}}$ and the maximum value of $\gls{sym-xi}_{\text{min},\text{slam}}$ to ensure lack of slamming across the entire \gls{acr-wec} surface.
The relevant value of $\gls{sym-theta-surf}$ differs depending on the wave amplitude.

%%%%%%%%%%%%%%%
\subsubsection{Case 1: Small Wave Amplitudes}
For the case of $\gls{sym-dDelta-z-slam} >\gls{sym-H}/2$ (wave amplitudes smaller than the draft), \Cref{eq:slamming-min-max} reduces to only a maximum amplitude criteria $|\gls{sym-hat-xi}|<\gls{sym-xi}_{\text{max},\text{slam}}$ because $\gls{sym-xi}_{\text{min},\text{slam}}<0$.
This aligns with the intuition that a body fixed at the still water line would not exit the water when the wave amplitude is less than the draft.
The value of $\gls{sym-theta-surf}$ that minimizes $\gls{sym-xi-max-slam}$ over the relevant range is
\begin{equation}\label{eq:slamming-limit-small-waves}
\gls{sym-theta-surf} = \gls{sym-pi} - \max\left(
                                            0,~
                                            \frac{-\gls{sym-k-wavenumber}\gls{sym-D}}{2}
                                            + |\gls{sym-pi}-\angle \gls{sym-xi}|
                                     \right)
\end{equation}
where the symmetry of \Cref{eq:slamming-min-max} has been used to collapse possible values onto $[0,\gls{sym-pi}]$.

This expression reduces to the simple limits $\gls{sym-dDelta-z-slam} -\gls{sym-H}/2$ and $\gls{sym-dDelta-z-slam} +\gls{sym-H}/2$
for short wavelengths $(\gls{sym-theta-surf}=\gls{sym-pi})$ and long wavelengths $(\gls{sym-theta-surf}=0)$ respectively.
These limits are also the minimum and maximum values of $\gls{sym-xi-max-slam}$ for any diameter and wave condition where $\gls{sym-dDelta-z-slam} >\gls{sym-H}/2$.
For intermediate wavelengths, the variable $\gls{sym-theta-surf}$ accurately accounts for the effect of body diameter and the phase offset of body motion from the waves.
Intuitively, $\gls{sym-theta-surf}$ represents how high up the free surface the bottom of the \gls{acr-wec} is allowed to get before the edge exits the water,
with $\gls{sym-theta-surf}=0$ allowing the body to get all the way to the wave crest (top)
and $\gls{sym-theta-surf}=\gls{sym-pi}$ requiring the float to remain fully below the wave trough (bottom).

%%%%%%%%%%%%%%%
\subsubsection{Case 2: Large Wave Amplitudes}
In sufficiently large waves ($\gls{sym-H}/2>\gls{sym-dDelta-z-slam}$), \Cref{eq:slamming-limit-small-waves} does not apply.
Intuitively, when the wave amplitude exceeds the draft, a minimum amplitude of motion in-phase with the incident wave is required to prevent the wave trough from going below the body surface.
The minimum amplitude arises mathematically when $\gls{sym-xi}_{\text{min},\text{slam}}>0$, which occurs exactly for waves meeting the amplitude criterion.
The ``in-phase'' motion requirement is consistent with the observation that in this wave regime, the maximum amplitude in \Cref{eq:slamming-min-max} becomes unsatisfiable ($\gls{sym-xi}_{\text{max},\text{slam}}<0$) for \gls{acr-wec} phases that are more than 90 degrees out of phase with the incident wave ($\gls{sym-theta-surf} \in [\gls{sym-pi}/2,3\gls{sym-pi}/2]$).

For evaluating \Cref{eq:slamming-min-max}, the value of $\gls{sym-theta-surf}$ that produces both minimum $\gls{sym-xi}_{\text{max},\text{slam}}$ and maximum $\gls{sym-xi}_{\text{min},\text{slam}}$ in this regime is:
\begin{equation}
\gls{sym-theta-surf} = \angle\gls{sym-hat-xi} + \frac{\gls{sym-k-wavenumber}\gls{sym-D}}{2}\textrm{sgn}(\gls{sym-pi}-\angle\gls{sym-hat-xi})
\end{equation}
After wrapping the value onto $[0,2\gls{sym-pi}]$, it can be mapped to $[0,\gls{sym-pi}]$ without changing the value of \Cref{eq:slamming-min-max} via the identity $\cos(\gls{sym-theta-surf})=\cos(\gls{sym-pi}-|\gls{sym-pi}-\gls{sym-theta-surf}|)$.

Notably, in this wave regime it becomes possible for the right hand side of \Cref{eq:slamming-squared} to become negative for certain values of $\angle\gls{sym-hat-xi}$ and $\gls{sym-D}$.
This means that those combinations of $\angle\gls{sym-hat-xi}$ and $\gls{sym-D}$ are prohibited because slamming would occur regardless of the amplitude $|\gls{sym-hat-xi}|$.
The following requirement on $\gls{sym-D}$ ensures that there exist some values of $\angle\gls{sym-hat-xi}$ without automatic slamming:
\begin{equation}\label{eq:D-slamming}
    \gls{sym-k-wavenumber}\gls{sym-D}
    <
    2\arcsin\left(
                    \frac{\gls{sym-dDelta-z-slam}}{\gls{sym-H}/2}
            \right)
\end{equation}
which has upper bound $\gls{sym-k-wavenumber}\gls{sym-D}<\gls{sym-pi}$.
Once \Cref{eq:D-slamming} is satisfied, the set of acceptable $\angle\gls{sym-hat-xi}$ that avoid negative values for both $\gls{sym-xi}_{\text{max},\text{slam}}$ and the right hand side of \Cref{eq:slamming-squared} are:
\begin{equation}\label{eq:slamming-allowed-angles}
\begin{aligned}
    \angle\gls{sym-hat-xi} \in &\left[
                                    0,
                                    ~~
                                    \arcsin\left(
                                                \frac{\gls{sym-dDelta-z-slam}}
                                                     {\gls{sym-H}/2}
                                           \right)
                                    -\frac{
                                            \gls{sym-k-wavenumber}\gls{sym-D}
                                            }
                                            {2}
                               \right]
    \cup \left[
                2\gls{sym-pi}-\arcsin\left(
                                            \frac{\gls{sym-dDelta-z-slam}}
                                                 {\gls{sym-H}/2}
                                     \right)
                +\frac{
                        \gls{sym-k-wavenumber}\gls{sym-D}
                      }
                      {2},
                ~~2\gls{sym-pi}
        \right]
\end{aligned}
\end{equation}
At the edge-case $\angle\gls{sym-hat-xi}$ on the inner boundaries of this interval,
only a single amplitude is permitted because the values for
$\gls{sym-xi}_{\text{max},\text{slam}}$ and $\gls{sym-xi}_{\text{min},\text{slam}}$
coincide at $\sqrt{(\gls{sym-H}/2)^2-\gls{sym-dDelta-z-slam} ^2}$.
At any lower amplitude, the \gls{acr-wec} would exit the water at the wave trough,
and any higher, the \gls{acr-wec} would raise out of the water at the wave crest.
The permissible amplitude range can be widened by decreasing $\gls{sym-D}$
or moving $\angle\gls{sym-hat-xi}$ towards the closer of $0$ and $2\gls{sym-pi}$.
The widest permissible amplitude range occurs when the \gls{acr-wec} is perfectly in-phase with the incident wave
($\angle\gls{sym-hat-xi}=0$)
and is small with respect to the wavelength
($\gls{sym-k-wavenumber}\gls{sym-D}\approx0$),
which yields the limits
$\gls{sym-H}/2-\gls{sym-dDelta-z-slam} < |\gls{sym-hat-xi}| < \gls{sym-H}/2+\gls{sym-dDelta-z-slam}$.

%%%%%%%%%%%%%%%
\subsubsection{Implementation}
\Cref{fig:slam} visualizes the nondimensional maximum and minimum slamming amplitudes
$\gls{sym-xi-max-slam}/(\gls{sym-H}/2)$ and $\gls{sym-xi-min-slam}/(\gls{sym-H}/2)$
as a function of the worst-case phase angle $\gls{sym-theta-surf}$
and the draft to wave amplitude ratio $\gls{sym-dDelta-z-slam} /(\gls{sym-H}/2)$.
This reveals that besides increasing the draft, decreasing $\gls{sym-theta-surf}$ can prevent slamming,
for example by decreasing the diameter or adjusting $\angle \gls{sym-xi}$ via control.
If $\gls{sym-dDelta-z-slam} <\gls{sym-H}/2|\sin\gls{sym-theta-surf}|$,
the design is unviable and the slamming limits are undefined,
consistent with the negative right hand side of \Cref{eq:slamming-squared} discussed above.
\begin{figure}[htbp]
    \centering
    \includegraphics[\FigureOptions{aor:figs/from-matlab/slamming_amplitude.pdf}]{figs/from-matlab/slamming_amplitude.pdf}
    \caption{Nondimensional critical slamming amplitude}
    \label{fig:slam}
\end{figure}

Separately, there exists a symmetrical counterpart to the slamming condition that we will call the submersion condition.
While preventing slamming requires that the bottom of the \gls{acr-wec} remains below the free surface,
preventing submersion requires that the top of the \gls{acr-wec} remains above the free surface.
Analogous to \eqref{eq:slamming-time-domain}, we obtain:
\begin{equation}\label{eq:submersion}
    \gls{sym-xi-t} + \gls{sym-dDelta-z-sub} > \gls{sym-zeta-y-t}
\end{equation}

Conceptually, both the slamming requirement \eqref{eq:slamming-time-domain} and the submersion requirement \eqref{eq:submersion} must be applied to both the float and the spar for a total of 4 free surface constraints for each sea state.
However, the limit on $|\gls{sym-hat-xi-f}-\gls{sym-hat-xi-s}|$ described in \Cref{tab:DLCs} already ensures that the top of the spar remains above the top of the float across all dynamic conditions.
Coupled with the fact that the float diameter always exceeds the spar diameter, this means that satisfaction of the float submersion requirement implies automatic satisfaction of the spar submersion requirement.
While this logic works for submersion, it does not apply for slamming because the damping plate diameter may exceed the float diameter.
The remaining three free surface constraints are aggregated into a single constraint for each sea state, using the minimum z-dimension for each body $\gls{sym-dDelta-z-surf}=\min(\gls{sym-dDelta-z-sub}, \gls{sym-dDelta-z-slam} )$, where the values for $\gls{sym-dDelta-z-sub} $ and $\gls{sym-dDelta-z-slam}$ were given in \Cref{sec:dynamics}.
This reduces the number of free surface constraints per sea state from four to one.

For operational waves, the vertical dimension $\gls{sym-dDelta-z-surf}$ in \Cref{eq:slam-constraint-quadratic} is set as follows.
For float slamming, $\gls{sym-dDelta-z-surf} = \gls{sym-T-f-1}$ (rather than $\gls{sym-T-f-2}$):
this prevents slamming on the slanted underside of the float and additionally maintains constant waterplane area,
avoiding unmodeled nonlinearities in the hydrostatic stiffness and Froude-Krylov excitation force.
For spar slamming, $\gls{sym-dDelta-z-surf} = \gls{sym-T-s} - \gls{sym-h-d}$
prevents both slamming on the bottom of the spar and surfacing of the damping plate.
For float and spar submersion,
$\gls{sym-dDelta-z-surf}$
is set to the above-water heights
$\gls{sym-h-f} - \gls{sym-T-f-2}$
and
$\gls{sym-h-s} - \gls{sym-T-s}$
respectively.

The storm case slamming constraint is not currently applied.
A future extension should set
$\gls{sym-dDelta-z-surf} = \gls{sym-T-f-2}$
in the storm case to prevent slamming on the (now-merged) bottom surface of the float-spar system.

The amplitude limits
$\gls{sym-h-fs-clear}$,
$\gls{sym-h-fs-up}$,
$\gls{sym-h-fs-down}$,
$\gls{sym-xi-f-lin}$, and
$\gls{sym-xi-s-lin}$
are sea-state-independent upper limits, while
$\gls{sym-xi}_{\text{slam}}$
depends on the sea state and can act as both an upper and lower limit on oscillation magnitude.

Relevant model values for the nominal float design under reactive optimal control are plotted for each sea state in \Cref{fig:slamming-validation}.
Comparing the first two subplots for the phase of the float motion $\angle\gls{sym-hat-xi-f}$ and the slamming angle $\gls{sym-theta-surf}$
shows that the slamming angle roughly ``wraps'' the motion angle around a value of $\gls{sym-pi}/2$, colored white in the figure.
The third plot shows that the minimum slamming and submersion amplitude is not a concern since it is negative for all sea states and thus will never be an active constraint.
The fourth plot shows the maximum amplitude limit and the fifth depicts the actual amplitude, with amplitudes in violation of the limit indicated with an overlaid black hatching.
A significant fraction of sea states violate the limit and will require adjusted control gains computed via constrained optimal control per \Cref{sec:optimal-control,sec:appendix-qp-solution}.
Significant wave heights in excess of 5.5 meters have negative maximum slamming amplitudes (colored blue),
indicating that control adjustments in those waves must not only decrease $|\gls{sym-hat-xi-f}|$
but also adjust $\angle\gls{sym-hat-xi-f}$ to reduce $\gls{sym-theta-surf}$.
\begin{figure}[htbp]
    \centering
    \includegraphics[\FigureOptions{aor:figs/from-matlab/slamming_model_comparison_float.pdf}]{figs/from-matlab/slamming_model_comparison_float.pdf}
    \caption{Slamming Model Results for Nominal \gls{acr-rm3} Float Design in Operational Sea States}
    \label{fig:slamming-validation}
\end{figure}

%%%%%%%%%%%%%%%%%%%%%%%%%%%%%%%%%%%%%%%%%
\FloatBarrier
\subsection{Narrow-band Wave Approximation}
\label{sec:appendix-narrowband}
\Cref{sec:irregular-waves} introduced the narrow-band wave environment approximation,
in which an irregular sea state is approximated by a regular wave at its energy-equivalent period and wave height.

In addition to enabling the \gls{acr-qcqp} formulation for the linearized system to be single-frequency and thus solvable analytically,
even for the nonlinear system the regular wave assumption may be desirable to avoid the time-consuming convolution integral or state-space identification that is typically required to simulate nonlinear dynamics subject to irregular wave forcing.

In the absence of dynamic constraints, the regular wave assumption does not significantly affect power production if the ultimate design utilizes a high-order linear controller capable of impedance matching over the bandwidth of each sea state.
If the ultimate design instead utilizes a simple PI controller, the assumption slightly overestimates the power produced in a broadband wave environment.
\citet{coe_practical_2021} used a Bretschneider spectrum and found the irregular wave power generation for a point absorber with PI control to be around 92\% of the perfectly matched power, which is sufficiently close to justify the approximation for early design phases without worrying about the complexity tradeoff of controller order.

\subsection{Power Matrix Decomposition}
\label{sec:appendix-power-decomposition}

The efficiency matrices introduced in \Cref{eq:power-matrix-decomposition} are defined as follows.
The wave power density $\gls{sym-mathbf-J-H-T-wave}$ and maximum capture width  $\gls{sym-mathbf-CW-H-T-max}$ are independent of design:
\begin{equation}
\mathbf{J}^{H,T}_{wave} = \frac{\gls{sym-rho-w} \gls{sym-g}^2}{64\gls{sym-pi}} \gls{sym-T-e} \gls{sym-H-s}^2, \qquad
\mathbf{CW}^{\gls{sym-H},\gls{sym-T-period}}_{max} = \frac{\gls{sym-g} \gls{sym-T-e}^2}{4\gls{sym-pi}^2}.
\end{equation}
The maximum capture width matrix represents the radiation limit determined by the \gls{acr-wec} mode of motion
(heave in this case; for surge or pitch the value would be doubled due to the asymmetry of the outgoing wave field).

The radiation and drag efficiencies are defined relative to a drag-free simulation $\mathbf{\gls{sym-P}}_{\gls{sym-F-PTO}\gls{sym-dot-X},\text{no drag}}$ (computed with $\gls{sym-C-d-f}=\gls{sym-C-d-s}=0$, ideal \gls{acr-pto}):
\begin{equation}
\gls{sym-boldsymbol-eta-H-T-rad} = \frac{\mathbf{P}_{F\dot{X},\text{no drag}}}{\mathbf{J}^{H,T}_{wave}\,\mathbf{CW}^{\gls{sym-H},\gls{sym-T-period}}_{max}}, \qquad
\gls{sym-boldsymbol-eta-H-T-drag} = \frac{\mathbf{P}_{F\dot{X},\text{ideal}}}{\mathbf{P}_{F\dot{X},\text{no drag}}}.
\end{equation}
$\gls{sym-boldsymbol-eta-H-T-drag}$ may exceed 1 at frequencies where the drag-free system is so undamped that
the unconstrained controller from \Cref{eq:matched-load,eq:damping-control} would induce instability;
in such cases a sub-optimal controller is substituted, which lowers
$\mathbf{\gls{sym-P}}_{\gls{sym-F-PTO}\gls{sym-dot-X},\text{no drag}}$
(the denominator) rather than raising the numerator.
In cases where the unconstrained controller is stable,
$\gls{sym-boldsymbol-eta-H-T-drag} \leq 1$ is guaranteed.

The product of the first four matrices is the theoretical ideal mechanical power:
\begin{equation}
    \mathbf{P}^{H,T}_{F\dot{X},\text{ideal}} = \mathbf{J}^{H,T}_{wave}
    \,
    \mathbf{CW}^{\gls{sym-H},\gls{sym-T-period}}_{max}
    \,
    \gls{sym-boldsymbol-eta-H-T-rad}
    \,
    \gls{sym-boldsymbol-eta-H-T-drag} 
    = \frac{
                |\gls{sym-mathbf-hat-F-H-T-th}|^2
            }
            {
                8\Re(\gls{sym-mathbf-Z-H-T-i-th})
            }.
\end{equation}
The remaining two efficiency matrices capture the loss to dynamic constraints and to \gls{acr-pto} electrical inefficiency:
\begin{equation}
    \gls{sym-boldsymbol-eta-H-T-limits} = \frac{\gls{sym-mathbf-P-H-T-F-dot-X}}{\mathbf{P}^{H,T}_{F\dot{X},\text{ideal}}}, \qquad
    \gls{sym-boldsymbol-eta-H-T-elec} = \gls{sym-eta} \frac{\gls{sym-mathbf-P-H-T-IV}}{\gls{sym-mathbf-P-H-T-F-dot-X}}.
\end{equation}

Although it would be possible to formulate alternative variants such as $\gls{sym-boldsymbol-eta-H-T}_{\text{drag},\text{var}} = \gls{sym-mathbf-P-H-T-IV}/\mathbf{P}^{H,T}_{IV,\text{no drag}}$ that capture the full effect of drag on the saturated system, such matrices would represent independent sensitivities and could not be multiplied together to recover the overall power matrix.
The current decomposition preserves multiplicativity, which is the key advantage of the formulation in \Cref{eq:power-matrix-decomposition}.

%%%%%%%%%%%%%%%%%%%%%%%%%%%%%%%%%%%%%%%%%
\FloatBarrier
\subsection{Dynamic Validation}
\label{sec:appendix-dynamic-validation}
\subsubsection{WEC-Sim}
The drag describing function and \gls{acr-meem} hydrodynamic coefficients have a minor effect assuming a 1-\gls{acr-dof} system (9.7\% and 2.7\% error on the average power and maximum amplitude respectively) but a major effect on the 2-\gls{acr-dof} system (38.2\% and 28.6\% respectively).
Meanwhile, the error is extremely low in the 2-\gls{acr-dof} system enforcing the same hydrodynamic coefficients as WEC-Sim and with zero drag (0.2\% and 1.9\% in power and amplitude respectively).
This indicates that the larger error in the 2-\gls{acr-dof}-drag-\gls{acr-meem} case is not an error in the 2-\gls{acr-dof} dynamic model itself, but in the way that a 2-\gls{acr-dof} model amplifies errors in drag and hydrodynamic coefficients due to the importance of the phase of motion between each \gls{acr-dof}.
The remainder of this subsection details the four-quadrant breakdown across drag-on/drag-off and identical/different hydrodynamic coefficient configurations.

A WEC-Sim parallel multiple condition run is performed in accelerator mode for 200s with a 100s ramp time using the ode4 solver and a 10~ms timestep, with the default mass properties and body geometry of the WEC-Sim \gls{acr-rm3} tutorial.\footnote{This geometry is not identical to the nominal float design used in the Reference Model Report \citep{RM3} and the rest of this paper, with a shorter spar draft, thicker damping plate, wider float frustum diameter, and assumed deep-water depth regime.
For consistency, MDOcean also uses this alternate geometry for WEC-Sim comparisons only.}
Each of the 210 sea states are simulated individually as the equivalent regular wave.
Body to body interactions (nonzero off-diagonal hydrodynamic coefficients) and state-space computations are turned on.
%\Cref{tab:dynamic-validation} and
\Cref{fig:dynamic-validation} compares the results.
The $\gls{sym-C-d}=0$ simulation with identical hydrodynamic coefficients (top left of \Cref{fig:dynamic-validation}) is conducted to verify correct implementation of the linear dynamics and consistency of sign conventions.
The validation obtains an error of less than $0.2\%$ for wave periods below 12 seconds.
This aligns with the perfect matching that is expected in the absence of drag non-linearities, and we attribute the tiny error that remains to the discrete-time nature of WEC-Sim's numerical integration and post-processing.
For wave periods above 12 seconds, MDOcean significantly overpredicts the mechanical power.
An additional validation for this case with the spar held fixed (not shown) maintained the excellent $<0.2\%$ accuracy even for $\gls{sym-T-period}>12 \gls{sym-s}$, so we attribute the error to resonant instabilities that occur in the multi-body hydrodynamic interactions in the highly underdamped drag-free dynamics.
Turning on drag ($\gls{sym-C-d}=1$, top right of \Cref{fig:dynamic-validation}) substantially resolves this inaccuracy, with power errors of $<5\%$ for the most common moderate sea states and error exceeding $40\%$ at the most energetic sea states.
While determining the exact source of the error requires further investigation, it is deemed low-impact enough to justify the model's use in an early-stage design setting.

Moving to the bottom left corner of \Cref{fig:dynamic-validation}, we examine the effect of the error in MDOcean's \gls{acr-meem} hydrodynamic coefficients on its dynamic power production.
The error closely aligns with the error in the identical hydro coefficient case, including the large errors at longer wave periods, confirming that the small error in the \gls{acr-meem} hydro coefficients has minimal effect on the overall dynamic response.
Finally, the bottom right corner of \Cref{fig:dynamic-validation} shows the combined effect of drag and hydro coefficient errors, which is the best representation of MDOcean's total dynamic error in the absence of further model improvements.
Power is underpredicted by up to 40\% at sea states with moderate (10-12s) periods and low wave heights, and overpredicted at both shorter and longer wave periods.
This results in a beneficial positive and negative error cancellation that explains the low error in weighted average power described in \Cref{sec:dynamic-validation}.
Meanwhile, MDOcean underpredicts the peak power by around 25\%.
Error scales more closely with wave period than wave height for periods under 12 seconds, while the opposite is true for periods above 12 seconds.

% \begin{table}[htbp]
%     \centering
%     \begin{tabular}{|c|c|c|} \hline
%          Power Percent Error: Average, Maximum&  $\gls{sym-C-d}=0$& $\gls{sym-C-d}=1$\\ \hline
%          WAMIT&  1.90\%, 1.96\%& 5.21\%, -70\%\\ \hline
%          MEEM&  -6.30\%, -14\%& -2.79\%, -80\%\\ \hline
%     \end{tabular}
%     \caption{Dynamic validation showing error in a variety of modeling assumptions}
%     \label{tab:dynamic-validation}
% \end{table}
\begin{landscape}
\begingroup
\begin{figure}[htbp]
\centering
\begin{tabular}{c m{1em} | M{.47\linewidth} | M{.47\linewidth} }
  && \multicolumn{2}{c}{Drag Coefficient}\\
         && $\gls{sym-C-d}=0$& $\gls{sym-C-d}=1$\\ \hline
    \multirow{2}{*}[-3em]{\rot{Hydro Coeffs}} &\rot{Identical}&
    \includegraphics[\FigureOptions{aor:figs/from-matlab/wecsim_wcsm_multi_drag_off_meem_off__power_mech_unsat.pdf}]{figs/from-matlab/wecsim_wcsm_multi_drag_off_meem_off__power_mech_unsat.pdf}
  & \includegraphics[\FigureOptions{aor:figs/from-matlab/wecsim_wcsm_multi_drag_on_meem_off__power_mech_unsat.pdf}]{figs/from-matlab/wecsim_wcsm_multi_drag_on_meem_off__power_mech_unsat.pdf}
\\ &\rot{Different} &
    \includegraphics[\FigureOptions{aor:figs/from-matlab/wecsim_wcsm_multi_drag_off_meem_on__power_mech_unsat.pdf}]{figs/from-matlab/wecsim_wcsm_multi_drag_off_meem_on__power_mech_unsat.pdf}
  & \includegraphics[\FigureOptions{aor:figs/from-matlab/wecsim_wcsm_multi_drag_on_meem_on__power_mech_unsat.pdf}]{figs/from-matlab/wecsim_wcsm_multi_drag_on_meem_on__power_mech_unsat.pdf} \\
\end{tabular}
\caption[Dynamic validation error in unsaturated mechanical power]{Dynamic validation showing error in unsaturated mechanical power generation for a variety of modeling assumptions}
\label{fig:dynamic-validation}
\fillandplacepagenumber
\end{figure}
\endgroup
\end{landscape}

Besides mechanical power, a variety of other signals including amplitudes, phases, and forces are compared between the models.
While most signals show good agreement between the models, interestingly, large errors are observed in the drag forces, although these errors do not cause significant deviations in the overall system response.
\Cref{fig:drag-validation} compares the drag force on the float and spar in each model using identical hydro coefficients and $\gls{sym-C-d}=1$.

\begin{landscape}
\begingroup
\begin{figure}[htbp]
\centering
\begin{tabular}{c m{1em} | M{.47\linewidth} | M{.47\linewidth} }
  && \multicolumn{2}{c}{Drag Force}\\
         && Magnitude& Phase\\ \hline
    \multirow{2}{*}[-3em]{\rot{Body}} &\rot{Float}&
    \includegraphics[\FigureOptions{aor:figs/from-matlab/wecsim_wcsm_multi_drag_on_meem_off__float_drag_force_fund.pdf}]{figs/from-matlab/wecsim_wcsm_multi_drag_on_meem_off__float_drag_force_fund.pdf}
  & \includegraphics[\FigureOptions{aor:figs/from-matlab/wecsim_wcsm_multi_drag_on_meem_off__float_drag_force_phase.pdf}]{figs/from-matlab/wecsim_wcsm_multi_drag_on_meem_off__float_drag_force_phase.pdf}
\\ &\rot{Spar} &
    \includegraphics[\FigureOptions{aor:figs/from-matlab/wecsim_wcsm_multi_drag_on_meem_off__spar_drag_force_fund.pdf}]{figs/from-matlab/wecsim_wcsm_multi_drag_on_meem_off__spar_drag_force_fund.pdf}
  & \includegraphics[\FigureOptions{aor:figs/from-matlab/wecsim_wcsm_multi_drag_on_meem_off__spar_drag_force_phase.pdf}]{figs/from-matlab/wecsim_wcsm_multi_drag_on_meem_off__spar_drag_force_phase.pdf} \\
\end{tabular}
\caption[Drag force comparison across modeling assumptions]{Drag force comparison showing large errors in drag force, despite low errors in amplitude and power}
\label{fig:drag-validation}
\fillandplacepagenumber
\end{figure}
\endgroup
\end{landscape}
The reason for this discrepancy remains unclear, since the expected error of the describing function approximation due to higher harmonics would increase with frequency, whereas the observed error is worse at mid-range frequencies.
Evidently, amplitude and power results are quite sensitive to the presence of drag, but not very sensitive to the exact magnitude and phase of the drag force.
The WEC-Sim simulation uses drag-only Morison elements with a maximum width of one tenth the incident wavelength to ensure accuracy of strip theory, and uses a Fourier transform to extract the fundamental amplitude and phase of the nonlinear drag waveform.
Future work should investigate the source of this error and whether it can be reduced by adjusting the drag model.

We also validate the describing function approximation for the drag force.
The general application of the describing function technique requires that the fundamental frequency $\gls{sym-omega}$ contains the vast majority of the energy, resulting in body displacements of the form $\gls{sym-xi}(\gls{sym-t}) \sim \cos(\gls{sym-omega} \gls{sym-t}+\gls{sym-phi})$.
\Cref{fig:wecsim-thd} shows the total harmonic distortion of the float displacement for each sea state in the $\gls{sym-C-d}=1$ case, which peaks at 1\% in the worst-case sea state and is below 0.5\% for the majority of sea states.
This confirms that higher harmonics have a negligible impact on the system and that the describing function method is a valid approach for modeling drag in this context.
The particular describing function relation used for drag assumes that the drag force has the form
$\gls{sym-F-d-t} \sim |\cos(\gls{sym-omega} \gls{sym-t}+\gls{sym-phi})|
                       \cos(\gls{sym-omega} \gls{sym-t}+\gls{sym-phi})$.
\Cref{fig:wecsim-drag-waveform} shows the assumed and actual drag force waveform for four representative sea states corresponding to the four corners of the \gls{acr-jpd}.
The match is excellent in all four cases, confirming that the drag force is indeed well approximated by this form, even in high-amplitude, high-frequency conditions where the nonlinearity is expected to be strongest.

\begin{figure}[htbp]
    \centering
    \includegraphics[\FigureOptions{aor:figs/from-matlab/position_THD_contour.pdf}]{figs/from-matlab/position_THD_contour.pdf}
    \caption[Total harmonic distortion of float displacement with drag]{Total Harmonic Distortion of Float Displacement in WEC-Sim Simulations with Drag}
    \label{fig:wecsim-thd}
\end{figure}
\begin{figure}[htbp]
    \centering
    \includegraphics[\FigureOptions{aor:figs/from-matlab/drag_force_desc_fcn.pdf}]{figs/from-matlab/drag_force_desc_fcn.pdf}
    \caption{Comparison of Assumed and Actual Drag Force Signal Shape}
    \label{fig:wecsim-drag-waveform}
\end{figure}

\subsubsection{Reference Model Report}
While the original reference model report \citep{RM3} does not document the intermediate values used in dynamics calculations such as response amplitude or hydrodynamic coefficients,
it does record the mechanical power for each sea state.
\Cref{fig:report-power-validation} compares these values against those predicted by WEC-Sim and MDOcean.
\begin{figure}[htbp]
    \centering
    \includegraphics[\FigureOptions{aor:figs/from-matlab/wecsim_rpt_multi_drag_on_meem_on__power_mech_unsat.pdf}]{figs/from-matlab/wecsim_rpt_multi_drag_on_meem_on__power_mech_unsat.pdf}
    \caption{Comparison of mechanical power for \gls{acr-rm3} report \cite{RM3} and MDOcean}
    \label{fig:report-power-validation}
\end{figure}
Note that the data is independent of \Cref{fig:dynamic-validation} due to the disparate geometry between the reference model report and provided WEC-Sim hydrodynamic coefficients,
and that the power multiplier mentioned in \Cref{sec:dynamic-validation} has not been applied in the figure.
We observe a rather poor match, with MDOcean underestimating power in the low-period sea states and overestimating in the high-period sea states.
The discrepancy is of minimal concern for the reasons described in \Cref{sec:dynamic-validation}.
Nonetheless, researchers and practitioners should be aware of the inconsistency and exercise caution before using the reference model report's power values for technology benchmarking or model validation.


%%%%%%%%%%%%%%%%%%%%%%%%%%%%%%%%%%%%%%%%%
%%%%%%%%%%%%%%%%%%%%%%%%%%%%%%%%%%%%%%%%%
\section{Economics Module Details}
\label{sec:appendix-econ}
%%%%%%%%%%%%%%%%%%%%%%%%%%%%%%%%%%%%%%%%%
\subsection{\gls{acr-rm3} \gls{acr-cbs}}
\label{sec:appendix-econ-cbs}

The MDOcean economic model is calibrated against the \gls{acr-rm3} \gls{acr-cbs} of \citet{neary_reference_2014}, reproduced in \Cref{tab:CBS}.
Percentages are shown for both single-device ($\gls{sym-N-WEC}=1$) and array ($\gls{sym-N-WEC}=100$) cases to illustrate the economies-of-scale trend captured by the power-law scaling \Cref{eq:cost-power-law}.

\begin{table}[htbp]
    \centering

    \caption{\gls{acr-cbs}}
    \label{tab:CBS}
    \expanded{
        \noexpand\tabular{\TableColumnSpecs{aor:cbs}}
    }
         \gls{acr-cbs} Category&  Nominal \% for $\gls{sym-N-WEC}=1$& Nominal \% for $\gls{sym-N-WEC}=100$&Scales with design?\\\hline
         1.1 - Development&  26\%&  3\%&No\\
         1.2 - Infrastructure&  6\%&  4\%&No\\
         1.3 - Mooring/Foundation&  3\%&  12\%&No\\
         1.4 - Structural&  17\%&  46\%&$\gls{sym-V-struct}$\\
         1.5 - \gls{acr-pto}&  4\%&  11\%&$\gls{sym-F-max}, \gls{sym-P-max}$\\
         1.6 - Integration \& Profit Margin&  2\%&  6\%&No\\
         1.7 - Installation&  34\%&  10\%&No\\
         1.8 - Decommissioning&  9\%&  9\%&No\\
         1.9 - Contingency&  9\%&  9\%&No\\
 \textbf{Total 1.1-1.9 (\gls{acr-capex})}& 93.3\%& 97.6\%&-\\\hline
 2.1 - Insurance& 1.3\%& 0.4\%&No\\
 2.2 - Environmental/Regulatory& 4.1\%& 0.3\%&No\\
 2.3 - Marine Operations& 0.2\%& 0.3\%&No\\
 2.4 - Shoreside Operations& 0.8\%& 0.2\%&No\\
 2.5 - Replacement Parts& 0.3\%& 1.0\%&No\\
 2.6 - Consumables& $<$0.1\%& 0.2\%&No\\
 \textbf{Total 2.1-2.6 (\gls{acr-opex})}& 6.7\%& 2.4\%&-\\
    \endtabular
    \end{table}

%%%%%%%%%%%%%%%%%%%%%%%%%%%%%%%%%%%%%%%%%
\subsection{Power-Law Scaling with Number of Devices}
\label{sec:appendix-econ-power-law}

The per-\gls{acr-wec} unit costs and prices in \Cref{eq:unit-cost}
($\gls{sym-C}_{\text{fixed}}$,
 $\gls{sym-C}_{pto,\text{constant}}$, 
 $\gls{sym-p-P}$,
 $\gls{sym-p-F}$,
 $\gls{sym-p-s}$)
each decrease with the number of devices $\gls{sym-N-WEC}$ according to
\begin{equation}\label{eq:cost-power-law}
    \gls{sym-C} = \gls{sym-C-infty}+ (\gls{sym-C-1} - \gls{sym-C-infty})
                    \cdot (\gls{sym-N-WEC})^{-\gls{sym-beta-expo}},
\end{equation}
where $\gls{sym-C-infty}$ is the asymptotic large-scale unit cost,
$\gls{sym-C-1}$ is the single-unit cost,
and $\gls{sym-beta-expo}$ is the rate of decrease.

The values of $\gls{sym-C-infty}$, $\gls{sym-C-1}$, and $\gls{sym-beta-expo}$ for each \gls{acr-cbs} category
were obtained by curve-fitting \Cref{eq:cost-power-law} to the \gls{acr-cbs} estimates
for $\gls{sym-N-WEC} \in \{1, 10, 50, 100\}$ \citep{neary_reference_2014},
and are tabulated in \Cref{tab:econ-model-values}.

\begin{table}[h]
    \centering
    \caption{Cost model values for each \gls{acr-cbs} category. Cost values are in 2012 USD.}
    \label{tab:econ-model-values}
    \expanded{
        \noexpand\tabular{\TableColumnSpecs{aor:econ-model-values}}
    }
         \gls{acr-cbs} Category& Expo. $\gls{sym-beta-expo}$&Unit cost at scale $\gls{sym-C-infty}$&Single-unit cost \newline $\gls{sym-C-1}$\\\hline
         \textbf{1.1-1.3, 1.6-1.9} - Design-Indep. \gls{acr-capex}& 0.741& 1.24 \$M &13.92 \$M\\
         \textbf{1.4} - Structural \gls{acr-capex} & 0.481& 2.387 \$/kg&4.294 \$/kg\\
         \textbf{1.5} - \gls{acr-pto} \gls{acr-capex}& 0.206& Constant: 92.59 \$k Power: 0.4454 \$/W Force: 0.0086 \$/N &Constant: 93.64 \$k Power: 1.355 \$/W Force: 0.0204 \$/N\\
 \textbf{2.1-2.6} - \gls{acr-opex}&0.557& 0&1.193 \$M\\
    \endtabular
    \end{table}

%%%%%%%%%%%%%%%%%%%%%%%%%%%%%%%%%%%%%%%%%
\subsection{Linear \gls{acr-pto} Cost Scaling}
\label{sec:appendix-econ-pto-scaling}

Linear \gls{acr-pto} cost scaling is common for \gls{acr-wec} techno-economic analysis.
The \gls{acr-rm3} cost numbers \citep{RM3} which MDOcean uses as a baseline are themselves scaled from estimates for a smaller device: the reference model project applies linear scaling with peak power to an unpublished 2011 cost estimate by ReVision Consulting for a 100~kW peak hydraulic \gls{acr-pto} (except for the riser cable and control system, whose cost does not scale with power).
The \gls{acr-nrel} \gls{acr-sam} assumes that hydraulic \gls{acr-pto} cost scales slightly less than linearly with power (exponent 0.91), and one study uses regression to obtain separate linear models for the cost of commercial induction and permanent magnet generators as a function of torque capability \citep{nakhai_techno-economic_2022}.
%the cheaper of the induction and permanent magnet machines is chosen with a torque threshold:
% \begin{equation}
%     \gls{sym-C}_{gen} = \begin{cases}
%     80.86 \gls{sym-tau} + 1032, & 0~\leq\gls{sym-tau} <47~\text{Nm, ~~~~induction generator} \\
%     9.151 \gls{sym-tau} + 4386, & 47\leq\gls{sym-tau} \leq 4000~\text{Nm, permanent magnet generator}
%     \end{cases}
% \end{equation}

%%%%%%%%%%%%%%%%%%%%%%%%%%%%%%%%%%%%%%%%%
\subsection{Economic Validation}
\label{sec:appendix-econ-validation}
\Cref{fig:econ-nwec-validate} shows the scaling behavior of economic outputs (\gls{acr-capex}, \gls{acr-opex}, and \gls{acr-lcoe}) against the number of \glspl{acr-wec} in the farm, comparing MDOcean results against the \citet{RM3} reference values.
\begin{figure}[htbp]
    \centering
    \includegraphics[\FigureOptions{aor:figs/from-matlab/cost_vs_N_WEC.pdf}]{figs/from-matlab/cost_vs_N_WEC.pdf}
    \caption{Validation for cost scaling with number of \glspl{acr-wec}}
    \label{fig:econ-nwec-validate}
\end{figure}

%%%%%%%%%%%%%%%%%%%%%%%%%%%%%%%%%%%%%%%%%
%%%%%%%%%%%%%%%%%%%%%%%%%%%%%%%%%%%%%%%%%
\section{Structures Module Details}
\label{sec:appendix-structures}
Inputs to the structures module include forces, bulk and structural dimensions, and material constants.
It outputs a factor of safety for each limit case.
First, it obtains equivalent section properties based on the plate and stiffener dimensions.
Second, it relates applied loads to stresses and deflection using analytical solutions to structural boundary value problems obtained from the well-known Roark's handbook \citep{young_roarks_2001} and the references therein.
Finally, it utilizes design standards from organizations like Det Norske Veritas, the American Bureau of Shipping, and the American Iron and Steel Institute to develop limit-state expressions for each possible failure mode of each major structural element under each design load case.
For consistency with the reference model report \citep{RM3}, all structural elements are assumed to be welded, although future work should consider the use of pinned joints for certain interfaces to reduce reaction moments and enhance structural efficiency.

%%%%%%%%%%%%%%%%%%%%%%%%%%%%%%%%%%%%%%%%%
\subsection{Equivalent Section Properties for Stiffened Plates}
\label{sec:appendix-equivalent-thickness}
Stiffeners are structural elements, frequently with I or T profiles, that can be welded or riveted to flat plates to provide additional load-carrying capacity and prevent buckling.
Stiffened plates are common in marine and aerospace structural design because they are an efficient way to carry spatially distributed fluid loads.
To analyze these compound sections, equivalent properties can be established that describe the overall bending stiffness of the combined section as if it were a simple uniform section, while still utilizing the maximum distance from the neutral axis of the irregular section to accurately determine stress.
\Cref{fig:stiffened-plate} illustrates the concept of equivalent thickness for a stiffened plate.
\begin{figure}[hbtp]
    \centering
    \includegraphics[\FigureOptions{aor:figs/from-matlab/equivalent_stiffness.pdf}]{figs/from-matlab/equivalent_stiffness.pdf}
    \caption{Equivalent thickness for a stiffened plate}
    \label{fig:stiffened-plate}
\end{figure}

Specifically, the equivalent thickness of the plate $\gls{sym-t-eq}$ is derived by equating the stiffened section's second moment of area per unit width with that of a uniform section:
\begin{equation}
    \int \gls{sym-z}^2dz = \frac{\gls{sym-t-eq}^3}{12}
\end{equation}
Next, deflection of the stiffened plate is equal to the deflection of the equivalent unstiffened plate with flexural rigidity $\gls{sym-D-eq}$:
\begin{equation}
    \gls{sym-D-eq} = \frac{\gls{sym-E}~  \gls{sym-t-eq}^3}{12(1-\gls{sym-nu-poisson}^2)}
\end{equation}

The maximum stress $\gls{sym-sigma}$ of the stiffened plate is then:
\begin{equation}\label{eq:plate-stress}
    \gls{sym-sigma}= 12\gls{sym-M} \gls{sym-z-max}/\gls{sym-t-eq}^3
\end{equation}
where $\gls{sym-M}$ is the maximum moment per unit length and $z_{\text{max}}$ is the maximum distance from the neutral axis of the stiffened section.
This parallels the stress for an unstiffened plate of thickness $\gls{sym-t}$, which is $\gls{sym-sigma}=6\gls{sym-M}/\gls{sym-t}^2 = 12 \gls{sym-M}(\gls{sym-t}/2)/\gls{sym-t}^3$.

Note that this equivalent-thickness method of accounting for stiffeners assumes that as a whole, the stiffener-plate system deflects like a single plate element, rather than as a set of multiple plate elements separated by stiffeners.
This is a reasonable assumption when the stiffeners are small and densely spaced with respect to the plate, but breaks down if the stiffeners come to dominate the system.
The so-called effective breadth ratio $\gls{sym-rho}$ is used to quantify the validity of this assumption, where the equivalent-thickness method requires $\gls{sym-rho}=1$:
\begin{equation}
   \gls{sym-rho} = \gls{sym-rho-epsilon} \gls{sym-rho-psi}
\quad
\gls{sym-rho-epsilon} = \begin{cases}
        1, & \gls{sym-epsilon-slenderness}\leq 0.673 \\
        \frac{1-0.22/\gls{sym-epsilon-slenderness}}{\gls{sym-epsilon-slenderness}}, & \gls{sym-epsilon-slenderness} > 0.673
    \end{cases} \quad
    \gls{sym-rho-psi} = \begin{cases}
        1, & \gls{sym-psi}> -0.236 \\
        \frac{1}{2}+\frac{1}{3-\gls{sym-psi}}, & \gls{sym-psi} \leq -0.236
    \end{cases}
\end{equation}
The effective breadth ratio depends on the slenderness factor $\gls{sym-epsilon-slenderness}$ and load distribution factor $\gls{sym-psi}$, found as follows:
\begin{equation}
\begin{aligned}
    \gls{sym-epsilon-slenderness} &= \frac{1.052}{\sqrt{\gls{sym-k-buckle}}} \frac{\gls{sym-w}}{\gls{sym-thick}} \sqrt{\frac{\gls{sym-f-1}}{\gls{sym-E}}} \\
    \gls{sym-psi}&=\frac{\gls{sym-f-2}}{\gls{sym-f-1}}
    \end{aligned}
\end{equation}
with $\gls{sym-t}$ referring to the plate unstiffened thickness, $\gls{sym-w}$ as the distance between stiffeners, $\gls{sym-f-1}$ as the maximum compression stress along the width, and $\gls{sym-f-2}$ as the maximum tension stress along the width (or minimum compression stress, if there is no tension).
By convention, compression stress is positive and tension is negative, so $\gls{sym-psi}$ ranges from $-\infty$ to $+1$.
The plate buckling coefficient $\gls{sym-k-buckle}$ is expressed as:
\begin{equation}\label{eq:k-buckle}
  \gls{sym-k-buckle} = 4+2\left(1-\gls{sym-psi}\right)^3+2\left(1-\gls{sym-psi}\right)
\end{equation}
This effective breadth procedure follows the American Iron and Steel Institute design manual \citep{american_iron_and_steel_institute_cold-formed_1991}.
Rather than enforcing $\gls{sym-epsilon-slenderness}<0.673$, corresponding with $\gls{sym-rho-epsilon}=1$, MDOcean more flexibly requires $\gls{sym-epsilon-slenderness}<0.809$, corresponding to $\gls{sym-rho-epsilon}=0.9$ and thereby capping the error due to insufficient slenderness at around 10\%.
MDOcean also uses $\gls{sym-k-buckle}=4$, the minimum value that equation \eqref{eq:k-buckle} can take on, conservatively maximizing the slenderness factor $\gls{sym-epsilon-slenderness}$.

Broadening the design space to allow for more dominant stiffeners with higher slenderness factors would require modeling the stiffener-plate system not as a plate with the stiffeners absorbed into an equivalent thickness, but as individual stiffeners with the plate absorbed into the effective breadth.
That model has other complexities such as the shear lag phenomenon so is left to future extensions \citep{wierzbicki_lecture_2013,american_iron_and_steel_institute_cold-formed_1991}.

%\hl{Describe how this informs the lower design variable bounds on the structural stiffnesses.}

%%%%%%%%%%%%%%%%%%%%%%%%%%%%%%%%%%%%%%%%%
\subsection{Float}
The float is composed of 12 watertight stiffened shells in the shape of trapezoidal prisms.
Each shell consists of a top and a bottom trapezoidal plate, and an inner, an outer, and two side rectangular plates.
All edges are welded.
Rather than model the deflections of each edge and apply compatibility, for simplicity the edges of each plate are conservatively modeled as fixed.
The top and bottom plates are the only ones with external heave loads, arising from the float-spar tubular connection and the hydrodynamic pressure respectively.

The trapezoidal plates are isosceles trapezoids with smaller base $\gls{sym-b-1} = \gls{sym-pi} \gls{sym-D-f-in}/12$,
larger base $\gls{sym-b-2}=\gls{sym-pi} \gls{sym-D-f}/12$,
and perpendicular height $\gls{sym-h-0}=(\gls{sym-D-f}-\gls{sym-D-f-in})/2$ (see \Cref{fig:trapezoid}).
The references consulted do not contain structural results for trapezoidal plates, so geometric intuition is used to scale available solutions.
For example, \citet{young_roarks_2001} compile models for rectangular plates under perpendicular loading.
They show that the maximum bending stress $\gls{sym-sigma}$ scales with the square of the shorter side length, $x_{\text{min}}^2$. % in rectangular plates, with the square of the altitude in triangular plates, and with the square of the radius in circular sectors. 
This makes sense because for fixed-edge plates with distributed loads,
it is the square of the minimum distance from an edge to the point of highest deflection,
$(x_{\text{min}}/2)^2$, that geometrically sets the maximum curvature, so $\gls{sym-sigma}\sim x_{\text{min}}^2$.
For trapezoids, one then expects approximate rectangles ($\gls{sym-b-2}/\gls{sym-b-1}\approx1$) to have $x_{\text{min}}=\min(\gls{sym-b-1},\gls{sym-h-0})$,
short trapezoids ($\gls{sym-h-0}<\gls{sym-b-1})$ to have $x_{\text{min}}=\gls{sym-h-0}$,
tall trapezoids ($\gls{sym-h-0}\gg \gls{sym-b-2}$) to have $x_{\text{min}}=\gls{sym-b-1}$,
and intermediate trapezoids to have an $x_{\text{min}}$ that depends on the trapezoid slope.
The nominal \gls{acr-rm3} design falls into this intermediate case,
with $\gls{sym-b-1}\approx1.7$~m,
$\gls{sym-b-2}\approx5.2$~m,
and $\gls{sym-h-0}\approx6.8$~m.
These dimensions are shown in \Cref{fig:trapezoid}.

\begin{figure}[htbp]
    \centering
    \includegraphics[\FigureOptions{aor:figs/from-matlab/trapezoid.pdf}]{figs/from-matlab/trapezoid.pdf}
    \caption[Trapezoidal float plate with inscribed circle]{Trapezoidal float plate with inscribed circle used to determine $x_{\text{min}}$}
    \label{fig:trapezoid}
\end{figure}

To approximately capture the dimensions of a trapezoid, which effectively set the deflection curvature and therefore stress, $x_{\text{min}}$ is set to the diameter of the largest circle that can be inscribed in the trapezoid tangent to the shorter base $\gls{sym-b-1}$.
This inscribed circle scheme satisfies all limit cases mentioned above and can be calculated as
\begin{equation}
    \gls{sym-x-min} =
    \begin{cases}
        \gls{sym-b-1} \left( \gls{sym-m-trap} + \sqrt{1+\gls{sym-m-trap}^2}\right), & \gls{sym-h-0} \geq \sqrt{\gls{sym-b-1}\gls{sym-b-2}}
        \\
        \gls{sym-h-0}, & \gls{sym-h-0} < \sqrt{\gls{sym-b-1}\gls{sym-b-2}}
    \end{cases}
\end{equation}
where $\gls{sym-m-trap}=(\gls{sym-b-2}-\gls{sym-b-1})/(2\gls{sym-h-0})$ is the slope of the trapezoid.
The expression is continuous at $\gls{sym-h-0}=\sqrt{\gls{sym-b-1}\gls{sym-b-2}}$, the point where the inscribed circle becomes tangent to base $\gls{sym-b-2}$ in addition to $\gls{sym-b-1}$.
Continuity of the model across the design space avoids problems in the gradient-based optimization due to undefined gradients.
The inscribed circle model gives $\gls{sym-x-min}\approx2.2$~m for the nominal float design.
In addition to the shorter side length $\gls{sym-x-min}$, stress in a fixed-edge rectangular plate also depends weakly on the longer side length of the plate $\gls{sym-x-max}$, which for the trapezoid is taken as
\begin{equation}
    \gls{sym-x-max} =
    \begin{cases}
        \gls{sym-h-0},& \gls{sym-h-0} \geq \sqrt{\gls{sym-b-1}\gls{sym-b-2}}
        \\
        \frac{\gls{sym-b-1}+\gls{sym-b-2}}{2} + \gls{sym-h}(1-\sqrt{1+\gls{sym-m-trap}^2}), & \gls{sym-h-0} < \sqrt{\gls{sym-b-1} \gls{sym-b-2}}
    \end{cases}
\end{equation}
using a similar inscribed circle scheme to maintain continuity.

\subsubsection{Float Bottom Plate}
Under a distributed pressure $\gls{sym-q-pressure}$, the bottom plate maximum moment per unit length is then:
\begin{equation}\label{eq:bottom-plate-moment}
\gls{sym-M-bot}=\gls{sym-beta-plate-0} \gls{sym-q-pressure} \gls{sym-x-min}^2
\end{equation}
for $\gls{sym-beta-plate-0}$ as a function of $x_{\text{max}}/x_{\text{min}}$ that \citet{timoshenko_theory_1959} tabulates, and the bending stress is found by plugging this moment into equation~\eqref{eq:plate-stress}.
While the simulation utilizes the real $\gls{sym-beta-plate-0}$, constant $\gls{sym-beta-plate-0}$ is assumed for the scaling laws of \Cref{eq:struct-scaling} because $\gls{sym-beta-plate-0}$ varies only weakly with $x_{\text{max}}/x_{\text{min}}$.

\subsubsection{Float Top Plate}
The top float plate has equivalent thickness $t_{f,\text{top},eq}$ and the same dimensions and edge fixity as the bottom plate.
Unlike the uniform pressure on the bottom plate, the top plate is subject to the welded tubular float-spar attachment force and moment, modeled as a load distributed over a small circle of radius $\gls{sym-r-0}$, where $\gls{sym-r-0}$ is assumed to be at least half the thickness.
Under these conditions, either the stress at the center $\gls{sym-sigma-cent}$ or along the longer edge $\gls{sym-sigma-edge}$ may dominate:
\begin{equation}\label{eq:top-plate-stress}
    \begin{aligned}
    \gls{sym-sigma-edge}&= \frac{3\gls{sym-W}}{2\gls{sym-pi} \gls{sym-t-f-top-eq}^2}\left( (1+\gls{sym-nu-poisson})\ln\frac{2x_{\text{min}}}{\gls{sym-pi} \gls{sym-r-0}}+\gls{sym-beta-plate-1}\right)\\
        \gls{sym-sigma-cent}&=-\gls{sym-beta-plate-2} \gls{sym-W}/\gls{sym-t-f-top-eq}^2
    \end{aligned}
\end{equation}
using $\gls{sym-beta-plate-1}$ and $\gls{sym-beta-plate-2}$ tabulated as a function of $\gls{sym-x-max}/\gls{sym-x-min}$ by \citet{young_roarks_2001}.
The bending moment per unit length is then found as $\gls{sym-M-top}=\gls{sym-sigma-edge} \gls{sym-t-f-top-eq}^2/6$.
However, unlike \Cref{eq:bottom-plate-moment}, it is found that \Cref{eq:top-plate-stress} predicts stresses that do not match well with the finite element results for the nominal \gls{acr-rm3} design \citep{RM3,previsic_reference_2011}.
It is unknown whether this is an implementation or modeling issue, and remains as future work to address.
In the meantime, the top plate thickness is set by scaling the bottom plate thickness, as described in \Cref{tab:struct-calc-vs-scale}.

\subsubsection{Other Float Plates}
The other float plates besides the top and bottom ones experience only the edge reaction moments that the top and bottom plates apply at their interface, with no additional external loads in heave.
Assuming that the maximum allowable stress is the same in all float plates, the required side plate thickness can be found by simply scaling the input float bottom plate thickness using the float thickness ratios in the nominal \gls{acr-rm3} design.
This prevents the need for a separate structural analysis of the side plates.

%%%%%%%%%%%%%%%%%%%%%%%%%%%%%%%%%%%%%%%%%
\subsection{Damping Plate}
For an annular plate with free outer radius $\gls{sym-a-outer-damping-plate}$ and fixed inner radius $\gls{sym-b-inner-damping-plate}$, the nondimensional plate deflection $\gls{sym-overline-delta-plate}$ and bending moment $\gls{sym-overline-M-r}$ can be calculated as a function of plate aspect ratio $\gls{sym-b-inner-damping-plate}/\gls{sym-a-outer-damping-plate}$ for various load cases.
In this case, the plate is loaded by a distributed load (subscript $dis$) equal to $F_{\text{heave}}$, and a concentrated load (subscript $con$) at four equidistant points along the edge, equal to the tubular support reaction force \gls{sym-F-tube}.
Formulas for the concentrated loading are given by \citet{boedo_corrected_1998}, and formulas for distributed loading are given by \citet[case 2L, table 11.2]{young_roarks_2001}.
The 24 radial plate stiffeners are taken into account with the procedure described in \Cref{sec:appendix-equivalent-thickness}, producing equivalent flexural rigidity $\gls{sym-D-eq}$ and equivalent section modulus $\gls{sym-S-eq}$.
The tube force \gls{sym-F-tube} is statically indeterminate and is solved with compatibility by equating the plate and tubular support displacements, $\delta_{\text{plate}}=\delta_{\text{tube}}$, where the tube displacement is expressed in terms of its bending stiffness $\gls{sym-K-tube}$.
This results in the following expression for the radial plate bending stress $\gls{sym-sigma-r}$:
\begin{equation}
    \gls{sym-sigma-r} =  \frac{\gls{sym-F-heave}}{\gls{sym-S-eq}}
    \left(
            \gls{sym-overline-M-r}_{,dis} + \gls{sym-overline-M-r}_{,con} \frac{
                                                                                \gls{sym-overline-delta-plate}_{dis}
                                                                                }
                                                                                {
                                                                                \frac{\gls{sym-D-eq}}
                                                                                     {\gls{sym-a-outer-damping-plate}^2 \gls{sym-K-tube}}
                                                                                - \gls{sym-overline-delta-plate}_{con}
                                                                                }
    \right)
\end{equation}

The critical buckling stress $\gls{sym-sigma-buckle}$ is found with the \citet{american_bureau_of_shipping_requirements_2022} methodology.
The ultimate stress is then $\sqrt{\gls{sym-sigma-Y} \gls{sym-sigma-buckle}}$  \citep{wierzbicki_lecture_2013-1} and the endurance limit is taken as half of ultimate.
The factor of safety is then the ratio of the maximum (ultimate or endurance, depending on the design load case) stress to the radial plate bending stress.
The procedure is summarized in \Cref{fig:damping-plate-flowchart}.
\begin{figure}[htbp]
    \centering
    \includegraphics[\FigureOptions{aor:figs/from-matlab/damping_plate_flowchart.pdf}]{figs/from-matlab/damping_plate_flowchart.pdf}
    \caption{Calculation process flowchart for the damping plate structural assessment}
    \label{fig:damping-plate-flowchart}
\end{figure}
%\hl{Todo: typeset equations for all the inputs with leftmost boxes: Seq, Deq, Ktube, the nondim variables, and sigma buckle.

%Todo: describe the plots of plate stress/deflection over space, maybe add dots in the 4 corners showing the point of load application}
\Cref{fig:damping-plate-struct} shows the spatial distribution of radial bending moment and deflection across the plate radius for the nominal design, computed using the analytical solutions described above.
\begin{figure}[htbp]
    \centering
    \begin{subfigure}[t]{0.48\linewidth}
        \includegraphics[\FigureOptions{aor:figs/from-matlab/damping_plate_moment.pdf}]{figs/from-matlab/damping_plate_moment.pdf}
        \caption{Damping plate moment}
        \label{fig:damping-plate-moment}
    \end{subfigure}
    \hfill
    \begin{subfigure}[t]{0.48\linewidth}
        \includegraphics[\FigureOptions{aor:figs/from-matlab/damping_plate_deflection.pdf}]{figs/from-matlab/damping_plate_deflection.pdf}
        \caption{Damping plate deflection}
        \label{fig:damping-plate-deflection}
    \end{subfigure}
    \caption[Damping plate structural analysis: bending moment and deflection]{Damping plate structural analysis: radial bending moment \ref{fig:damping-plate-moment} and deflection \ref{fig:damping-plate-deflection} over the plate radius.}
    \label{fig:damping-plate-struct}
\end{figure}
In \Cref{fig:damping-plate-maxs}, peak stress and deflection are plotted as a function of damping plate aspect ratio, normalized by their values in the nominal design.

%%%%%%%%%%%%%%%%%%%%%%%%%%%%%%%%%%%%%%%%%
\subsection{Column}
As stated in \Cref{sec:structures}, the column's small slenderness ratio means that both axial compression and global Euler buckling must be taken into account.
Meanwhile, the hydrostatic pressure at the bottom creates a substantial compressive hoop stress.

The Euler critical buckling force $\gls{sym-F-crit}$ is:
\begin{equation}
    \gls{sym-F-crit} = \frac{\gls{sym-pi}^2 \gls{sym-E} \gls{sym-I-zz}}{(\gls{sym-K-end} \gls{sym-L})^2}
\end{equation}
where $\gls{sym-E}$ is the Young's modulus, $\gls{sym-I-zz}$ is the second moment of area, $\gls{sym-K-end}=2$ is the end condition
(fixed-free since the hydrostatic rotational stiffness provides a restoring moment at the surface),
and $\gls{sym-L}$ is the beam length, set to $\gls{sym-h-s}$ the full spar height.
The spar factor of safety $FOS$ to prevent yield and global buckling under the combined loading is \citep{american_bureau_of_shipping_requirements_2022}:
\begin{equation}
    FOS = \frac{\gls{sym-sigma}_{\gls{sym-s},\text{max}}}{\gls{sym-sigma-s}}
    = \frac{
                \gls{sym-sigma-Y} \frac{
                                        \gls{sym-zeta-s} + \sqrt{
                                                                \gls{sym-zeta-s}^2
                                                                + 4\gls{sym-omega-s}
                                                                }
                                        }
                                        {2}
            }
            {
                \frac{\gls{sym-F-heave}}{\gls{sym-A}} + \gls{sym-q-pressure}
            }
\end{equation}
where $\gls{sym-zeta-s}$ and $\gls{sym-omega-s}$ are defined as follows:
\begin{equation}
\begin{aligned}
     \gls{sym-zeta-s} &= 1 - \gls{sym-P-r}(1 - \gls{sym-P-r})
                    ~\frac{
                            \gls{sym-sigma-Y}
                          }
                          {
                            \gls{sym-F-crit}/\gls{sym-A}
                          }
                    - \frac{
                            \gls{sym-sigma-theta}
                            }
                            {
                            \gls{sym-sigma-Y}
                            },
     \qquad
    \gls{sym-omega-s} = \frac{\gls{sym-sigma-theta}}{2\gls{sym-sigma-Y}} 
                        \left(
                            1 - \frac{
                                        \gls{sym-sigma-theta}
                                    }
                                    {
                                        2\gls{sym-sigma-Y}
                                    }
                        \right)
\end{aligned}
\end{equation}
and the hoop stress $\gls{sym-sigma-theta}$ is
\begin{equation}
     \gls{sym-sigma-theta} = \frac{
                                    \gls{sym-q-pressure}\gls{sym-D-s}
                                  }
                                  {
                                    2~\gls{sym-t-s-r}
                                  }
\end{equation}
$\gls{sym-P-r}$ is a constant and $\gls{sym-q-pressure}$ is the distributed hydrostatic pressure at the bottom of the spar.
These equations assume that the spar is compact,
meaning that local buckling of the tube as a plate element is not a concern \citep{american_bureau_of_shipping_requirements_2022}.
%These stresses are combined into the von Mises stress $\sigma_{vm}$ as defined in equation \eqref{von}.
% \begin{equation}
% \begin{split}
%     \gls{sym-sigma}_{vm}^2 = \frac{1}{2} \left( (\gls{sym-sigma}_{11} - \gls{sym-sigma}_{22})^2 + (\gls{sym-sigma}_{22} - \gls{sym-sigma}_{33})^2 + (\gls{sym-sigma}_{33} - \gls{sym-sigma}_{11})^2 \right) \\ + 3 (\gls{sym-sigma}_{12}^2 + \gls{sym-sigma}_{23}^2 + \gls{sym-sigma}_{31}^2)
%     \label{von}
% \end{split}
% \end{equation}

% Three yield factors of safety $FOS_y$ (for the float, vertical column, and damping plate) and one buckling factor of safety $FOS_b$ (for the vertical column) are computed for each of two forces (the powertrain force $F_p$ and the hydrodynamic excitation force $F_h$). This gives a total of eight FOS:
% \begin{equation}
% \begin{aligned}
%     FOS_{\gls{sym-y}~{i,\gls{sym-j}}} = \frac{\gls{sym-sigma-Y}}{\gls{sym-sigma}_{vm,ij}} ~~&~~
%     FOS_{\gls{sym-b},\gls{sym-j}} = \frac{\gls{sym-F-crit}}{\gls{sym-F}_{\gls{sym-j}}}} \\
%     i = 1..3, ~& \gls{sym-j} = 1..2
% \end{aligned}
% \end{equation}

%%%%%%%%%%%%%%%%%%%%%%%%%%%%%%%%%%%%%%%%%
%%%%%%%%%%%%%%%%%%%%%%%%%%%%%%%%%%%%%%%%%
\section{Parameters}
\label{sec:appendix-parameters}
\Cref{tab:parameters} lists the constant parameters assumed in this work, including their descriptions, values, and units.

\begin{table}[ht]
\centering
\caption{Selected Parameters}
\label{tab:parameters}
%\setlength{\tabrowsep}{} % default value:
{\renewcommand{\arraystretch}{1.4}

\begin{tabular}{cP{0.35\linewidth}P{0.3\linewidth}c}
\textbf{Parameter} & \textbf{Description} & \textbf{Value} & \textbf{Units}  \\ \hline
$\gls{sym-rho-w}$ & Seawater density& 1000 & kg/$m^3$  \\ 
$\gls{sym-g}$& Acceleration of gravity& 9.8 & m/$s^2$  \\ 
$\gls{sym-h}$& Water depth & 100 & m  \\ 
\gls{acr-jpd}& Wave joint probability distribution & read from file \citep{janzou_sam_2022}  &\%  \\ 
$\gls{sym-H-s}$ & Significant wave height \citep{janzou_sam_2022}& 0.25 : 0.5 : 6.75  & m  \\ 
$H_{s,\text{struct}}$ & 100 year wave height \citep{berg_extreme_2011}& [13.4, 18.8, 24.2, 30.1, 24.2, 18.8, 13.4] & m  \\ 
$\gls{sym-T-e}$& Wave energy period \citep{janzou_sam_2022}& 4.5 : 1 : 18.5  & s  \\ 
$T_{\text{struct}}$ & 100 year wave peak period \citep{berg_extreme_2011}& [5.57 8.76 12.18 17.26 21.09 24.92 31.70] & s  \\ 
$\gls{sym-sigma-Y}$ & Material yield strength& 248 & MPa  \\ 
$\gls{sym-rho-M}$ & Material density& 7850 & kg/$m^3$  \\ 
$\gls{sym-E}$& Material Young's modulus& 200 & GPa  \\ 
$\text{cost}_{\gls{sym-M}}$ & Material cost& 1.89 & 2012\$/kg \\ 
%$t_{ft}$ & \begin{tabular}[c]{@{}c@{}}Surface Float Top Plate\vspace{-1.5 mm} \\ Thickness\end{tabular} & 0.013 & m  \\ 
%$t_{fr}$ & \begin{tabular}[c]{@{}c@{}}Float Radial Wall\vspace{-1.5 mm} \\ Thickness\end{tabular} & 0.011 & m  \\ 
%$t_{fc}$ & \begin{tabular}[c]{@{}c@{}} Float Circumferential\vspace{-1.5 mm} \\ Gusset Thickness\end{tabular} & 0.011 & m  \\ 
%$t_{fb}$ & Float Bottom Thickness & 0.014 & m  \\ 
%$t_{sr}$ & Vertical Column Thickness & 0.025 & m  \\ 
%$B_{\text{min}}$ & Minimum Buoyancy Ratio & 1 & -  \\ 
$FOS_{\text{min}}$ & Minimum factor of safety& 1.5 & -  \\ 
$D_{d_{\text{min}}}$ & Minimum damping plate diameter & 30 & m  \\ 
$\gls{acr-fcr}$ & Fixed charge rate& 10.8 & \%  \\ 
$\gls{sym-N-WEC}$ & Number of \glspl{acr-wec} in array& 100 & -  \\ 
$\gls{sym-D-d}/{\gls{sym-D-s}}$ & Normalized damping plate diameter & 5 & -  \\ 
$\gls{sym-T-s}/{\gls{sym-D-s}}$ & Normalized spar draft& 5.83 & -  \\ 
$\gls{sym-h-d}/{\gls{sym-D-s}}$ & Normalized damping plate thickness & 0.004 & -  \\ 
$\gls{sym-T-f-2}/{\gls{sym-h-f}}$ & Float submergence ratio & 0.615 & -  \\ 
$F_{\text{heave},\text{mult}}$ & Heave force multiplier & 1.65 & - \\ 
\end{tabular}
}
 \end{table}

\Cref{tab:hyperparams} documents the hyper-parameters used to control convergence and other numerical aspects of the simulation.
\begin{table}[htbp]
    \centering
    \caption{Hyper-parameters used to control numerics}
    \label{tab:hyperparams}
    \begin{tabular}{ccc} \toprule
          Variable&  Value& Unit\\ \midrule
          $\gls{sym-xi-tol}$&  0.01 & m \\
          $\gls{sym-angle-xi-tol}$& 3 & degrees\\
          Maximum drag iterations& 100 & -\\
          \gls{sym-N-harmonics} & 10 & - \\
          \gls{sym-b-max} & 700.5 & - \\ \bottomrule
    \end{tabular}
\end{table}


\clearpage


\ifdefined\JOURNALTWOCOLUMN
    \twocolumn%
\fi


\printcredits

\bibliographystyle{cas-model2-names}
\bibliography{references,zotero-meem-refs}

\end{document}
